\documentclass[11pt,a4paper,twoside,openright]{report}

% Cargar configuración y estilos visuales
% ==============================================================================
% CONFIGURACIÓN Y ESTILO VISUAL — FUNDAMENTOS FÍSICOS Y TECNOLÓGICOS (FFT)
% Escuela Técnica Superior de Ingenierías Informática y de Telecomunicación
% Universidad de Granada (UGR) — Plataforma repasaYA
% ==============================================================================

\usepackage[utf8]{inputenc}
\usepackage[T1]{fontenc}
\usepackage[spanish,es-tabla]{babel}
\usepackage{lmodern}
\usepackage{microtype}
\usepackage[dvipsnames,svgnames,table,xcdraw]{xcolor}
\usepackage{amsmath,amssymb,amsfonts,amsthm}
\usepackage{physics}
\usepackage{siunitx}
\usepackage{graphicx}
\usepackage{float}
\usepackage{booktabs}
\usepackage{array}
\usepackage{multirow}
\usepackage{enumitem}
\usepackage{fancyhdr}
\usepackage{titlesec}
\usepackage{listings}
\usepackage{varwidth}
\usepackage{tcolorbox}
\tcbuselibrary{skins,breakable}
\usepackage{tikz}
\usetikzlibrary{babel,arrows.meta,positioning,calc,patterns}
\usepackage[siunitx,european,cuteinductors]{circuitikz}

% Configuración de márgenes
\usepackage[a4paper,top=2.5cm,bottom=2.5cm,left=2.8cm,right=2.8cm,headheight=24pt]{geometry}

% Configuración de siunitx en español
\sisetup{
  output-decimal-marker = {,},
  range-phrase = { a },
  per-mode = symbol
}
\DeclareSIUnit\VAR{VAr}
\DeclareSIUnit\VA{VA}

% Paleta de colores profesionales para Física y Electrónica
\definecolor{fftDarkBlue}{RGB}{16, 32, 60}
\definecolor{fftPrimary}{RGB}{24, 76, 120}
\definecolor{fftAccent}{RGB}{210, 85, 20}
\definecolor{fftLightAccent}{RGB}{254, 245, 238}
\definecolor{fftGreen}{RGB}{28, 130, 65}
\definecolor{fftLightGreen}{RGB}{240, 253, 244}
\definecolor{fftOrange}{RGB}{210, 110, 10}
\definecolor{fftLightOrange}{RGB}{255, 251, 235}
\definecolor{fftRed}{RGB}{185, 28, 28}
\definecolor{fftLightRed}{RGB}{254, 242, 242}
\definecolor{fftGray}{RGB}{246, 248, 250}
\definecolor{fftBorderGray}{RGB}{215, 222, 232}
\definecolor{fftTeal}{RGB}{15, 115, 110}
\definecolor{codeBg}{RGB}{248, 249, 252}

% Hipervínculos
\usepackage{hyperref}
\hypersetup{
    colorlinks=true,
    linkcolor=fftDarkBlue,
    citecolor=fftTeal,
    urlcolor=fftAccent,
    pdfauthor={repasaYA - UGR},
    pdftitle={Fundamentos Físicos y Tecnológicos - Apuntes Completos},
    pdfsubject={Electromagnetismo, Circuitos DC/AC, Semiconductores, Diodos, MOSFET y Operacionales},
    pdfkeywords={FFT, Electromagnetismo, Circuitos, Kirchhoff, Thévenin, Diodos, MOSFET, CMOS, Op-Amp, UGR}
}

% Configuración de Encabezados y Pies de página
\pagestyle{fancy}
\fancyhf{}
\fancyhead[LE,RO]{\footnotesize\textsf{\thepage}}
\fancyhead[LO]{\footnotesize\textsf{\rightmark}}
\fancyhead[RE]{\footnotesize\textsf{\leftmark}}
\renewcommand{\headrulewidth}{0.4pt}
\renewcommand{\footrulewidth}{0pt}

% Configuración de títulos de capítulos y secciones
\titleformat{\chapter}[display]
  {\normalfont\bfseries\color{fftDarkBlue}\sffamily}
  {\Large\textsf{Tema \thechapter}}{1ex}
  {\titlerule[1.5pt]\vspace{1ex}\huge}
  [\vspace{1ex}\titlerule]

\titleformat{\section}
  {\normalfont\Large\bfseries\color{fftPrimary}\sffamily}
  {\thesection}{1em}{}

\titleformat{\subsection}
  {\normalfont\large\bfseries\color{fftDarkBlue}\sffamily}
  {\thesubsection}{1em}{}

\titleformat{\subsubsection}
  {\normalfont\normalsize\bfseries\color{fftAccent}\sffamily}
  {\thesubsubsection}{1em}{}

% Entorno: Caja de Fuentes y Metadatos al inicio de cada tema
\newcommand{\cuadrofuentes}[1]{%
  \begin{tcolorbox}[
    colback=fftLightAccent,
    colframe=fftAccent,
    boxrule=1pt,
    arc=4pt,
    left=12pt,
    right=12pt,
    top=8pt,
    bottom=8pt,
    breakable
  ]
    \small
    \noindent\textbf{\color{fftAccent}\large\textsf{Metadatos del Tema y Fuentes:}}\\[4pt]
    #1
  \end{tcolorbox}
  \vspace{10pt}
}

% Entorno: Definición
\newtcolorbox{definicion}[1][]{
  enhanced,
  colback=fftGray,
  colframe=fftPrimary,
  boxrule=1pt,
  leftrule=4pt,
  arc=3pt,
  title=\textbf{\textsf{Definición\ifstrempty{#1}{}{: #1}}},
  coltitle=white,
  colbacktitle=fftPrimary,
  attach boxed title to top left={yshift=-2mm, xshift=4mm},
  boxed title style={arc=2pt, boxrule=0pt},
  varwidth boxed title*=-8mm,
  top=10pt,
  bottom=8pt,
  left=10pt,
  right=10pt,
  breakable
}

% Entorno: Ley / Teorema Fundamental
\newtcolorbox{ley}[1][]{
  enhanced,
  colback=fftLightAccent,
  colframe=fftAccent,
  boxrule=0.8pt,
  leftrule=3.5pt,
  arc=3pt,
  title=\textbf{\textsf{Ley / Teorema\ifstrempty{#1}{}{: #1}}},
  coltitle=white,
  colbacktitle=fftAccent,
  attach boxed title to top left={yshift=-2mm, xshift=4mm},
  boxed title style={arc=2pt, boxrule=0pt},
  varwidth boxed title*=-8mm,
  top=10pt,
  bottom=8pt,
  left=10pt,
  right=10pt,
  breakable
}

% Entorno: Ejemplo Explicado
\newtcolorbox{ejemplo}[1][]{
  enhanced,
  colback=fftLightGreen,
  colframe=fftGreen,
  boxrule=0.8pt,
  leftrule=3.5pt,
  arc=3pt,
  title=\textbf{\textsf{Ejemplo Explicado\ifstrempty{#1}{}{: #1}}},
  coltitle=white,
  colbacktitle=fftGreen,
  attach boxed title to top left={yshift=-2mm, xshift=4mm},
  boxed title style={arc=2pt, boxrule=0pt},
  varwidth boxed title*=-8mm,
  top=10pt,
  bottom=8pt,
  left=10pt,
  right=10pt,
  breakable
}

% Entorno: Aviso / Ojo Examen
\newtcolorbox{aviso}[1][]{
  enhanced,
  colback=fftLightOrange,
  colframe=fftOrange,
  boxrule=0.8pt,
  leftrule=3.5pt,
  arc=3pt,
  title=\textbf{\textsf{Importante / Ojo Examen\ifstrempty{#1}{}{: #1}}},
  coltitle=white,
  colbacktitle=fftOrange,
  attach boxed title to top left={yshift=-2mm, xshift=4mm},
  boxed title style={arc=2pt, boxrule=0pt},
  varwidth boxed title*=-8mm,
  top=10pt,
  bottom=8pt,
  left=10pt,
  right=10pt,
  breakable
}

% Entorno: Método Sistemático de Resolución
\newtcolorbox{metodo}[1][]{
  enhanced,
  colback=codeBg,
  colframe=fftDarkBlue,
  boxrule=0.8pt,
  leftrule=3.5pt,
  arc=3pt,
  title=\textbf{\textsf{Método Sistemático\ifstrempty{#1}{}{: #1}}},
  coltitle=white,
  colbacktitle=fftDarkBlue,
  attach boxed title to top left={yshift=-2mm, xshift=4mm},
  boxed title style={arc=2pt, boxrule=0pt},
  varwidth boxed title*=-8mm,
  top=10pt,
  bottom=8pt,
  left=10pt,
  right=10pt,
  breakable
}

% Entorno: Ejercicio Resuelto
\newtcolorbox{ejercicio}[1][]{
  enhanced,
  colback=fftGray,
  colframe=fftTeal,
  boxrule=0.8pt,
  leftrule=3.5pt,
  arc=3pt,
  title=\textbf{\textsf{Ejercicio Resuelto\ifstrempty{#1}{}{: #1}}},
  coltitle=white,
  colbacktitle=fftTeal,
  attach boxed title to top left={yshift=-2mm, xshift=4mm},
  boxed title style={arc=2pt, boxrule=0pt},
  varwidth boxed title*=-8mm,
  top=10pt,
  bottom=8pt,
  left=10pt,
  right=10pt,
  breakable
}

% Configuración de listings
\lstset{
  backgroundcolor=\color{codeBg},
  basicstyle=\ttfamily\footnotesize\color{black},
  keywordstyle=\bfseries\color{fftPrimary},
  commentstyle=\itshape\color{gray},
  stringstyle=\color{fftGreen},
  numbers=left,
  numberstyle=\tiny\color{gray},
  numbersep=8pt,
  stepnumber=1,
  breaklines=true,
  breakatwhitespace=false,
  showstringspaces=false,
  frame=none,
  tabsize=2,
  captionpos=b,
  inputencoding=utf8,
  extendedchars=true,
  literate=
    {á}{{\'a}}1 {é}{{\'e}}1 {í}{{\'i}}1 {ó}{{\'o}}1 {ú}{{\'u}}1
    {Á}{{\'A}}1 {É}{{\'E}}1 {Í}{{\'I}}1 {Ó}{{\'O}}1 {Ú}{{\'U}}1
    {ñ}{{\~n}}1 {Ñ}{{\~N}}1 {¿}{{?`}}1 {¡}{{!`}}1
}


\begin{document}

% ------------------------------------------------------------------------------
% PORTADA
% ------------------------------------------------------------------------------
\begin{titlepage}
    \centering
    \vspace*{1.2cm}
    
    {\Huge\bfseries\color{fftDarkBlue}\textsf{Fundamentos Físicos y Tecnológicos}}\\[0.6cm]
    {\Large\textsf{Apuntes Completos de Teoría, Circuitos y Problemas Resueltos}}\\[0.4cm]
    {\large\textsf{1.º Grado en Ingeniería Informática + Doble Grado DGIIM}}\\[0.2cm]
    {\normalsize\textsf{Escuela Técnica Superior de Ingenierías Informática y de Telecomunicación}}\\[0.2cm]
    {\normalsize\textsf{Universidad de Granada (UGR)}}\\[1.2cm]
    
    \rule{0.85\textwidth}{1.5pt}\\[1.5cm]
    
    \begin{tcolorbox}[
        colback=fftLightAccent,
        colframe=fftAccent,
        width=0.9\textwidth,
        arc=4pt,
        halign=center
    ]
        \large\textbf{\color{fftDarkBlue}Síntesis Integral del Temario Oficial:}\\[8pt]
        \normalsize
        \textbf{Profesor:} Pedro Cartujo y Equipo Docente de Electrónica (UGR)\\[4pt]
        \begin{itemize}[leftmargin=1.5cm, rightmargin=1.5cm, noitemsep]
            \item \textbf{Bloque I:} Electromagnetismo (Electrostática, Gauss, Corriente y Magnetismo).
            \item \textbf{Bloque II:} Análisis de Circuitos en DC (Kirchhoff, Thévenin, Norton y Transitorios RC/RL).
            \item \textbf{Bloque III:} Análisis de Circuitos en AC (Fasores, Impedancias, Potencia y Filtros).
            \item \textbf{Bloque IV:} Dispositivos Electrónicos (Semiconductores, Diodos y MOSFETs).
            \item \textbf{Bloque V:} Fundamentos de Electrónica Digital y Familias CMOS.
            \item \textbf{Bloque VI:} Amplificadores Operacionales y Conversión A/D--D/A.
            \item \textbf{Solucionario:} Exámenes Parciales y Problemas de Convocatorias Oficiales.
        \end{itemize}
    \end{tcolorbox}
    
    \vfill
    
    {\large\textsf{Plataforma \textbf{repasaYA}}}\\[0.2cm]
    {\small\textsf{Grado en Ingeniería Informática --- Curso Académico 2024--2025 / 2025--2026}}\\[0.6cm]
\end{titlepage}

% ------------------------------------------------------------------------------
% ÍNDICE GENERAL
% ------------------------------------------------------------------------------
\tableofcontents
\newpage

% ------------------------------------------------------------------------------
% CAPÍTULOS / TEMAS
% ------------------------------------------------------------------------------
\chapter{Fundamentos del Electromagnetismo}
\label{cap:tema1}

\cuadrofuentes{
    \textbf{Asignatura:} Fundamentos Físicos y Tecnológicos (1.º Grado en Ingeniería Informática, UGR).\\[2pt]
    \textbf{Profesor:} Pedro Cartujo (Departamento de Electrónica y Tecnología de Computadores).\\[2pt]
    \textbf{Bibliografía de Referencia:} Sears, F. W., Zemansky, M. W., Young, H. D., Freedman, R. A., \textit{Física Universitaria con Física Moderna} (Vol. 2), Addison-Wesley.
}

% ==============================================================================
% 1.1. ELECTROSTÁTICA
% ==============================================================================
\section{Electrostática: Carga Eléctrica, Fuerzas y Campo Eléctrico}

El electromagnetismo unifica los fenómenos eléctricos y magnéticos bajo un marco físico común. La electrostática se ocupa de las interacciones producidas por cargas eléctricas que se encuentran en reposo relativo (campo electrostático independiente del tiempo).

\subsection{Concepto y Propiedades de la Carga Eléctrica}
La carga eléctrica es una propiedad intrínseca y fundamental de la materia que determina su interacción electromagnética. Existen dos tipos de carga: positiva (asociada a los protones) y negativa (asociada a los electrones).

\begin{definicion}[Propiedades Universales de la Carga Eléctrica]
\begin{itemize}[leftmargin=*]
    \item \textbf{Ley de Conservación de la Carga:} En cualquier proceso físico que tenga lugar dentro de un sistema eléctricamente aislado, la carga eléctrica total neta permanece rigurosamente constante ($\sum q = \text{constante}$). La carga no se crea ni se destruye, únicamente se transfiere o redistribuye.
    \item \textbf{Cuantización de la Carga:} La carga neta de cualquier cuerpo macroscópico o partícula subatómica observable es siempre un múltiplo entero de la carga elemental fundamental $e$:
    \begin{equation}
        q = \pm n \cdot e, \quad \text{con } n \in \mathbb{N}, \quad e \approx 1{,}602 \times 10^{-19}\ \si{\coulomb}
    \end{equation}
    \item \textbf{Invarianza Relativista:} La magnitud de la carga eléctrica de una partícula es independiente de su velocidad o del sistema de referencia inercial desde el cual se observe.
\end{itemize}
\end{definicion}

En el Sistema Internacional (SI), la unidad de carga eléctrica es el \textbf{Culombio} (\si{\coulomb}), definido como la cantidad de carga que atraviesa la sección transversal de un conductor por el que circula una corriente constante de un Amperio durante un segundo ($1\ \si{\coulomb} = 1\ \si{\ampere\second}$).

\subsection{Ley de Coulomb}
En 1785, Charles-Augustin de Coulomb determinó experimentalmente la fuerza electrostática entre dos cargas puntuales en reposo en el vacío.

\begin{ley}[Ley de Coulomb]
La fuerza electrostática $\vec{F}_{12}$ que una carga puntual $q_1$ ejerce sobre otra carga puntual $q_2$ separadas por una distancia $r$ es directamente proporcional al producto de sus cargas e inversamente proporcional al cuadrado de la distancia que las separa:
\begin{equation}
    \vec{F}_{12} = \frac{1}{4\pi\varepsilon_0} \frac{q_1 q_2}{r^2} \hat{r}_{12} = K_e \frac{q_1 q_2}{r^2} \hat{r}_{12}
\end{equation}
donde:
\begin{itemize}[leftmargin=*]
    \item $\hat{r}_{12} = \frac{\vec{r}_2 - \vec{r}_1}{|\vec{r}_2 - \vec{r}_1|}$ es el vector unitario dirigido desde la carga fuente $q_1$ hacia la carga testigo $q_2$.
    \item $\varepsilon_0 \approx 8{,}854 \times 10^{-12}\ \si{\coulomb\squared\per\newton\per\meter\squared}$ es la permitividad dieléctrica del vacío.
    \item $K_e = \frac{1}{4\pi\varepsilon_0} \approx 8{,}988 \times 10^9\ \si{\newton\meter\squared\per\coulomb\squared} \approx 9 \times 10^9\ \si{\newton\meter\squared\per\coulomb\squared}$.
\end{itemize}
\end{ley}

\begin{center}
\begin{tikzpicture}[scale=1.1]
    % Carga q1
    \draw[fill=fftPrimary!20, draw=fftPrimary, thick] (-2.5,0) circle (0.35) node[fftPrimary] {\textbf{$+q_1$}};
    % Carga q2
    \draw[fill=fftPrimary!20, draw=fftPrimary, thick] (2.5,0) circle (0.35) node[fftPrimary] {\textbf{$+q_2$}};
    % Linea de separación
    \draw[dashed, gray] (-2.15,0) -- (2.15,0) node[midway, below=2pt, text=black] {$r$};
    % Fuerzas repulsivas
    \draw[->, >=stealth, very thick, fftAccent] (-2.85,0) -- (-4.2,0) node[above] {$\vec{F}_{21}$};
    \draw[->, >=stealth, very thick, fftAccent] (2.85,0) -- (4.2,0) node[above] {$\vec{F}_{12}$};
    \node[below=14pt] at (0,0) {\small\textsf{Fuerzas electrostáticas repulsivas entre dos cargas del mismo signo ($\vec{F}_{21} = -\vec{F}_{12}$)}};
\end{tikzpicture}
\end{center}

\subsection{Campo Eléctrico y Principio de Superposición}
Para evitar el concepto de acción a distancia instantánea, se define el campo eléctrico como una perturbación vectorial del espacio circundante provocada por la presencia de cargas eléctricas.

\begin{definicion}[Vector Campo Eléctrico $\vec{E}$ y Líneas de Campo]
El campo eléctrico $\vec{E}$ en un punto $\vec{r}$ del espacio es la fuerza electrostática por unidad de carga positiva testigo $q_0$ situada en dicho punto:
\begin{equation}
    \vec{E}(\vec{r}) = \lim_{q_0 \to 0} \frac{\vec{F}}{q_0} \quad \left[\si{\newton\per\coulomb} \equiv \si{\volt\per\meter}\right]
\end{equation}
Para una carga puntual aislada:
\begin{equation}
    \vec{E}(\vec{r}) = \frac{1}{4\pi\varepsilon_0} \frac{q}{r^2} \hat{r}
\end{equation}
Las líneas de campo eléctrico parten de las cargas positivas (fuentes) y convergen hacia las cargas negativas (sumideros).
\end{definicion}

\begin{center}
\begin{tikzpicture}[scale=0.9]
    % Carga positiva
    \draw[fill=fftGreen!25, draw=fftGreen, thick] (-3.5,0) circle (0.3) node[fftGreen] {\textbf{$+$}};
    \foreach \angle in {0, 45, 90, 135, 180, 225, 270, 315} {
        \draw[->, >=stealth, thick, fftGreen] ({-3.5 + 0.35*cos(\angle)}, {0.35*sin(\angle)}) -- ({-3.5 + 1.2*cos(\angle)}, {1.2*sin(\angle)});
    }
    \node[below=12pt] at (-3.5,-1) {\footnotesize\textsf{Carga positiva (fuente)}};

    % Carga negativa
    \draw[fill=fftRed!25, draw=fftRed, thick] (3.5,0) circle (0.3) node[fftRed] {\textbf{$-$}};
    \foreach \angle in {0, 45, 90, 135, 180, 225, 270, 315} {
        \draw[<-, >=stealth, thick, fftRed] ({3.5 + 0.35*cos(\angle)}, {0.35*sin(\angle)}) -- ({3.5 + 1.2*cos(\angle)}, {1.2*sin(\angle)});
    }
    \node[below=12pt] at (3.5,-1) {\footnotesize\textsf{Carga negativa (sumidero)}};
\end{tikzpicture}
\end{center}

\begin{ley}[Principio de Superposición Lineal]
El campo eléctrico neto generado por una distribución de $N$ cargas puntuales $\{q_1, q_2, \dots, q_N\}$ es la suma vectorial de los campos individuales generados por cada una de ellas:
\begin{equation}
    \vec{E}(\vec{r}) = \sum_{i=1}^N \vec{E}_i(\vec{r}) = \frac{1}{4\pi\varepsilon_0} \sum_{i=1}^N \frac{q_i}{|\vec{r} - \vec{r}_i|^2} \frac{\vec{r} - \vec{r}_i}{|\vec{r} - \vec{r}_i|}
\end{equation}
Para distribuciones continuas de carga, se integra sobre la densidad correspondiente:
\begin{itemize}[leftmargin=*]
    \item Densidad lineal de carga: $\lambda = \frac{dq}{d\ell}\quad [\si{\coulomb\per\meter}] \implies dq = \lambda\, d\ell$.
    \item Densidad superficial de carga: $\sigma = \frac{dq}{dA}\quad [\si{\coulomb\per\meter\squared}] \implies dq = \sigma\, dA$.
    \item Densidad volumétrica de carga: $\rho = \frac{dq}{dV}\quad [\si{\coulomb\per\meter\cubed}] \implies dq = \rho\, dV$.
\end{itemize}
\end{ley}

% ==============================================================================
% 1.2. LEY DE GAUSS
% ==============================================================================
\section{Flujo Eléctrico y Ley de Gauss}

\subsection{Concepto de Flujo Eléctrico y Ley de Gauss}
\begin{ley}[Ley de Gauss]
El flujo del campo eléctrico a través de cualquier superficie gaussiana cerrada $S_{cerrada}$ es directamente proporcional a la carga eléctrica neta total $Q_{enc}$ encerrada en el interior de dicha superficie:
\begin{equation}
    \Phi_E = \oint_S \vec{E} \cdot d\vec{A} = \frac{Q_{enc}}{\varepsilon_0}
\end{equation}
\end{ley}

\begin{center}
\begin{tikzpicture}[scale=1.0]
    % Hilo cargado con lambda
    \draw[very thick, fftAccent] (-3,0) -- (3,0) node[right] {Hilo indefinido ($+\lambda$)};
    \foreach \x in {-2.5,-1.8,...,2.5} {
        \node[above=2pt, fftAccent, font=\tiny] at (\x,0) {$+$};
    }
    % Superficie gaussiana cilíndrica
    \draw[thick, fftPrimary, dashed] (-1.8,-1) rectangle (1.8,1);
    \draw[thick, fftPrimary] (-1.8,0) ellipse (0.25 and 1);
    \draw[thick, fftPrimary] (1.8,0) ellipse (0.25 and 1);
    % Vectores campo y área
    \draw[->, >=stealth, very thick, fftGreen] (0,1) -- (0,2.1) node[above] {$\vec{E}$};
    \draw[->, >=stealth, very thick, fftDarkBlue] (0,1) -- (0,1.6) node[right] {$d\vec{A}$};
    \draw[<->, gray] (1.9,0) -- (1.9,1) node[midway, right] {$r$};
    \draw[<->, gray] (-1.8,-1.2) -- (1.8,-1.2) node[midway, below] {$L$};
    \node[below=16pt] at (0,-1.2) {\small\textsf{Superficie gaussiana cilíndrica coaxial para calcular el campo de un hilo indefinido}};
\end{tikzpicture}
\end{center}

\subsection{Conductores en Equilibrio Electrostático y Condensadores}
\begin{definicion}[Condensador de Placas Paralelas]
Dispositivo formado por dos láminas conductoras paralelas de área $A$ separadas por una distancia $d$ que almacenan cargas opuestas $+Q$ y $-Q$.
\begin{equation}
    C = \frac{Q}{\Delta V} = \varepsilon_r \varepsilon_0 \frac{A}{d} \quad [\si{\farad}]
\end{equation}
\end{definicion}

\begin{center}
\begin{tikzpicture}[scale=1.1]
    % Placa positiva
    \draw[fill=fftPrimary!20, thick] (-2,1) rectangle (2,1.3);
    \node at (0,1.15) {\small\textbf{Placa Superior ($+Q$, densidad $+\sigma$)}};
    % Placa negativa
    \draw[fill=fftDarkBlue!20, thick] (-2,-1.3) rectangle (2,-1);
    \node at (0,-1.15) {\small\textbf{Placa Inferior ($-Q$, densidad $-\sigma$)}};
    % Líneas de campo uniforme
    \foreach \x in {-1.5,-1.0,-0.5,0,0.5,1.0,1.5} {
        \draw[->, >=stealth, thick, fftAccent] (\x,1) -- (\x,-1);
    }
    \node[right, fftAccent] at (1.6,0) {$\vec{E} = \frac{\sigma}{\varepsilon}$};
    % Cota de separación
    \draw[<->, gray, thick] (-2.3,-1) -- (-2.3,1) node[midway, left] {$d$};
    % Conexión exterior
    \draw[thick] (0,1.3) -- (0,1.8) node[above] {$V_+$};
    \draw[thick] (0,-1.3) -- (0,-1.8) node[below] {$V_-$};
\end{tikzpicture}
\end{center}

\subsection{Asociación de Condensadores}
\begin{center}
\begin{tikzpicture}[scale=0.9]
    % En paralelo
    \node[font=\bfseries] at (-3,1.8) {Asociación en Paralelo ($C_{eq} = C_1 + C_2$)};
    \draw (-4.5,0) -- (-3.5,0);
    \draw (-3.5,0) -- (-3.5,0.8) to[C, l=$C_1$] (-1.5,0.8) -- (-1.5,0);
    \draw (-3.5,0) -- (-3.5,-0.8) to[C, l=$C_2$] (-1.5,-0.8) -- (-1.5,0);
    \draw (-1.5,0) -- (-0.5,0);

    % En serie
    \node[font=\bfseries] at (3.5,1.8) {Asociación en Serie ($\frac{1}{C_{eq}} = \frac{1}{C_1} + \frac{1}{C_2}$)};
    \draw (1.5,0) -- (2,0) to[C, l=$C_1$] (3.5,0) to[C, l=$C_2$] (5,0) -- (5.5,0);
\end{tikzpicture}
\end{center}

% ==============================================================================
% 1.3. CAMPO MAGNÉTICO Y LEYES FUNDAMENTALES
% ==============================================================================
\section{Campo Magnético, Ley de Lorentz y Ley de Biot-Savart}

\begin{ley}[Fuerza de Lorentz y Ley de Biot-Savart]
Una partícula con carga $q$ que se mueve con velocidad $\vec{v}$ en un campo $\vec{B}$ experimenta la fuerza:
\begin{equation}
    \vec{F}_B = q (\vec{v} \times \vec{B}) \implies |\vec{F}_B| = |q|\, v\, B \sin(\theta)
\end{equation}
El campo magnético generado por un conductor rectilíneo muy largo por el que circula una corriente $I$ viene dado por la Ley de Biot-Savart / Ampère:
\begin{equation}
    B(r) = \frac{\mu_0 I}{2\pi r}
\end{equation}
\end{ley}

\begin{center}
\begin{tikzpicture}[scale=1.0]
    % Hilo con corriente
    \draw[very thick, fftDarkBlue] (0,-1.8) -- (0,1.8) node[above] {$I$ (corriente)};
    \draw[->, >=stealth, very thick, fftDarkBlue] (0,-0.5) -- (0,0.6);
    % Líneas de campo circulares
    \draw[thick, fftAccent, ->, >=stealth] (0.8,0) arc (0:170:0.8 and 0.35);
    \draw[thick, fftAccent, dashed] (0.8,0) arc (0:-190:0.8 and 0.35);
    \draw[thick, fftAccent, ->, >=stealth] (1.5,0) arc (0:170:1.5 and 0.6);
    \draw[thick, fftAccent, dashed] (1.5,0) arc (0:-190:1.5 and 0.6);
    \node[right, fftAccent] at (1.5,0.4) {$\vec{B}$ (Líneas concéntricas)};
    \node[below=14pt] at (0,-1.8) {\small\textsf{Líneas de campo magnético según la Regla de la Mano Derecha}};
\end{tikzpicture}
\end{center}

\begin{ley}[Fuerza entre Dos Conductores Paralelos]
Para dos conductores rectilíneos paralelos separados una distancia $d$:
\begin{equation}
    \frac{F}{\ell} = \frac{\mu_0 I_1 I_2}{2\pi d}
\end{equation}
\begin{itemize}[leftmargin=*]
    \item Corrientes en el \textbf{mismo sentido}: fuerzas de \textbf{atracción mutua}.
    \item Corrientes en \textbf{sentidos opuestos}: fuerzas de \textbf{repulsión mutua}.
\end{itemize}
\end{ley}

\begin{center}
\begin{tikzpicture}[scale=1.0]
    % Caso 1: Corrientes mismo sentido
    \draw[very thick, fftPrimary] (-3.5,-1.2) -- (-3.5,1.2) node[above] {$I_1$};
    \draw[->, >=stealth, very thick, fftPrimary] (-3.5,-0.4) -- (-3.5,0.4);
    \draw[very thick, fftPrimary] (-1.5,-1.2) -- (-1.5,1.2) node[above] {$I_2$};
    \draw[->, >=stealth, very thick, fftPrimary] (-1.5,-0.4) -- (-1.5,0.4);
    \draw[->, >=stealth, very thick, fftGreen] (-3.5,0) -- (-2.4,0) node[midway, above] {$\vec{F}_{21}$};
    \draw[->, >=stealth, very thick, fftGreen] (-1.5,0) -- (-2.6,0) node[midway, above] {$\vec{F}_{12}$};
    \node[below=8pt] at (-2.5,-1.2) {\small\textbf{Atracción} (mismo sentido)};

    % Caso 2: Corrientes sentido opuesto
    \draw[very thick, fftDarkBlue] (1.5,-1.2) -- (1.5,1.2) node[above] {$I_1$};
    \draw[->, >=stealth, very thick, fftDarkBlue] (1.5,-0.4) -- (1.5,0.4);
    \draw[very thick, fftDarkBlue] (3.5,-1.2) -- (3.5,1.2) node[above] {$I_2$};
    \draw[<-, >=stealth, very thick, fftDarkBlue] (3.5,-0.4) -- (3.5,0.4);
    \draw[->, >=stealth, very thick, fftRed] (1.5,0) -- (0.6,0) node[midway, above] {$\vec{F}$};
    \draw[->, >=stealth, very thick, fftRed] (3.5,0) -- (4.4,0) node[midway, above] {$\vec{F}$};
    \node[below=8pt] at (2.5,-1.2) {\small\textbf{Repulsión} (sentidos opuestos)};
\end{tikzpicture}
\end{center}

\subsection{Inducción Electromagnética y Autoinducción}
\begin{ley}[Ley de Faraday-Lenz y Autoinducción en Solenoide]
La fuerza electromotriz inducida en un circuito cerrado de $N$ espiras ante una variación de flujo magnético es:
\begin{equation}
    \mathcal{E} = -N \frac{d\Phi_B}{dt}
\end{equation}
Para un solenoide de longitud $\ell$, área $A$ y $N$ espiras, el coeficiente de autoinducción es:
\begin{equation}
    L = \frac{N \Phi_B}{I} = \mu_0 \frac{N^2 A}{\ell} \quad [\si{\henry}]
\end{equation}
La tensión autoinducida al variar la corriente es $\mathcal{E}_L = -L \frac{di}{dt}$, almacenando la energía magnética $U_B = \frac{1}{2} L I^2$.
\end{ley}

% ==============================================================================
% 1.4. EJERCICIOS RESUELTOS DE EXAMEN
% ==============================================================================
\section{Ejercicios Resueltos de Exámenes de la UGR}

\begin{ejercicio}[Cálculo de Campo y Potencial de Hilo Finito e Infinito --- Examen Febrero 2017]
Se tiene un hilo rectilíneo de longitud infinita cargado uniformemente con densidad lineal constante $\lambda = 4 \times 10^{-9}\ \si{\coulomb\per\meter}$.
\begin{enumerate}[label=\alph*)]
    \item Calcular el campo eléctrico a una distancia $r_1 = 0{,}05\ \si{\meter}$ del hilo.
    \item Calcular la diferencia de potencial entre los puntos $r_1 = 0{,}05\ \si{\meter}$ y $r_2 = 0{,}20\ \si{\meter}$.
\end{enumerate}
\tcbline
\textbf{Esquema Geométrico del Problema:}\\[4pt]
\begin{center}
\begin{tikzpicture}[scale=1.1]
    % Hilo
    \draw[very thick, fftAccent] (-3,0) -- (3,0) node[right] {Hilo con $\lambda$};
    % Puntos r1 y r2
    \draw[dashed, gray] (0,0) -- (0,2.5);
    \filldraw[fftPrimary] (0,0.8) circle (2pt) node[right=4pt] {$P_1\ (r_1 = 0{,}05\ \si{\meter})$};
    \draw[->, >=stealth, thick, fftGreen] (0,0.8) -- (0,1.7) node[right] {$\vec{E}(r_1)$};
    \filldraw[fftPrimary] (0,2.2) circle (2pt) node[right=4pt] {$P_2\ (r_2 = 0{,}20\ \si{\meter})$};
    \draw[<->, gray] (-0.3,0) -- (-0.3,0.8) node[midway, left] {$r_1$};
    \draw[<->, gray] (-0.7,0) -- (-0.7,2.2) node[midway, left] {$r_2$};
\end{tikzpicture}
\end{center}
\textbf{Solución Paso a Paso:}\\[4pt]
a) Por simetría cilíndrica, aplicando la Ley de Gauss:
\begin{equation}
    E(r) = \frac{\lambda}{2\pi\varepsilon_0 r} = \frac{2 K_e \lambda}{r}
\end{equation}
\begin{equation}
    E(0{,}05) = \frac{2 \times (9 \times 10^9) \times (4 \times 10^{-9})}{0{,}05} = \frac{72}{0{,}05} = 1440\ \si{\volt\per\meter} \quad \implies \vec{E} = 1440\, \hat{u}_r\ \si{\volt\per\meter}
\end{equation}
b) Diferencia de potencial integrando el campo:
\begin{equation}
    V(r_1) - V(r_2) = \frac{\lambda}{2\pi\varepsilon_0} \ln\left(\frac{r_2}{r_1}\right) = 72 \cdot \ln\left(\frac{0{,}20}{0{,}05}\right) = 72 \cdot \ln(4) \approx 99{,}81\ \si{\volt}
\end{equation}
\end{ejercicio}

\begin{ejercicio}[Campo Magnético de Tres Hilos Paralelos --- Examen Enero 2021]
Tres hilos conductores paralelos indefinidos pasan por $x_1 = -d$, $x_2 = 0$ y $x_3 = +d$ ($d = 0{,}10\ \si{\meter}$). Por 1 y 3 circulan $I_1 = 5\ \si{\ampere}$ e $I_3 = 5\ \si{\ampere}$ saliendo ($+\hat{k}$), y por 2 circula $I_2 = 10\ \si{\ampere}$ entrando ($-\hat{k}$).
Determinar la fuerza magnética por unidad de longitud sobre el hilo 2 y sobre el hilo 3.
\tcbline
\textbf{Diagrama de Fuerzas en Corte Transversal:}\\[4pt]
\begin{center}
\begin{tikzpicture}[scale=1.2]
    % Eje horizontal
    \draw[->, gray] (-2.2,0) -- (2.5,0) node[right] {$x$};
    % Hilo 1 (x = -d)
    \draw[thick, fftPrimary] (-1.5,0) circle (0.25);
    \filldraw[fftPrimary] (-1.5,0) circle (2pt);
    \node[below=8pt] at (-1.5,0) {$H_1\ (\odot\ I_1)$};
    % Hilo 2 (x = 0)
    \draw[thick, fftDarkBlue] (0,0) circle (0.25);
    \draw[thick, fftDarkBlue] (-0.17,-0.17) -- (0.17,0.17);
    \draw[thick, fftDarkBlue] (-0.17,0.17) -- (0.17,-0.17);
    \node[below=8pt] at (0,0) {$H_2\ (\otimes\ I_2)$};
    % Hilo 3 (x = +d)
    \draw[thick, fftPrimary] (1.5,0) circle (0.25);
    \filldraw[fftPrimary] (1.5,0) circle (2pt);
    \node[below=8pt] at (1.5,0) {$H_3\ (\odot\ I_3)$};
    % Fuerzas en Hilo 2
    \draw[->, >=stealth, very thick, fftAccent] (0,0.4) -- (0.8,0.4) node[above] {$\vec{F}_{1 \to 2}$};
    \draw[->, >=stealth, very thick, fftGreen] (0,0.4) -- (-0.8,0.4) node[above] {$\vec{F}_{3 \to 2}$};
    % Cotas
    \draw[<->, gray] (-1.5,-0.7) -- (0,-0.7) node[midway, below] {$d$};
    \draw[<->, gray] (0,-0.7) -- (1.5,-0.7) node[midway, below] {$d$};
\end{tikzpicture}
\end{center}
\textbf{Solución Paso a Paso:}\\[4pt]
1) En el hilo central 2: el hilo 1 (corriente contraria) lo repele hacia la derecha ($+\hat{i}$) con $1 \times 10^{-4}\ \si{\newton\per\meter}$; el hilo 3 lo repele hacia la izquierda ($-\hat{i}$) con idéntica magnitud. Por simetría:
\begin{equation}
    \vec{F}_2^{neta} / \ell = \vec{0}
\end{equation}
2) En el hilo 3: el hilo 2 lo repele hacia la derecha ($+1 \times 10^{-4}\ \si{\newton\per\meter}\hat{i}$) y el hilo 1 (corriente en mismo sentido a distancia $2d$) lo atrae hacia la izquierda ($-0{,}25 \times 10^{-4}\ \si{\newton\per\meter}\hat{i}$):
\begin{equation}
    \frac{\vec{F}_3^{neta}}{\ell} = (1{,}00 - 0{,}25) \times 10^{-4}\ \hat{i} = +0{,}75 \times 10^{-4}\ \si{\newton\per\meter}\ \hat{i}
\end{equation}
\end{ejercicio}

\chapter{Fundamentos de Teoría de Circuitos: Corriente Continua (DC)}

\section{Elementos Lineales de Dos Terminales: Activos y Pasivos}

Un circuito eléctrico es una interconexión cerrada de elementos de dos o más terminales a través de los cuales puede fluir carga eléctrica. Se denominan \textbf{elementos lineales de dos terminales} aquellos en los que la relación matemática entre la caída de tensión $v(t)$ entre sus extremos y la intensidad de corriente $i(t)$ que los atraviesa es directamente proporcional, diferencial, integral o una combinación lineal de ellas.

Atendiendo a su comportamiento energético global dentro del circuito, los elementos se clasifican en dos grandes familias:
\begin{itemize}
    \item \textbf{Elementos activos}: Son capaces de inyectar o suministrar energía eléctrica neta al sistema a lo largo del tiempo (generadores o fuentes independientes y dependientes de tensión y de corriente).
    \item \textbf{Elementos pasivos}: No pueden generar energía por sí mismos; únicamente la disipan en forma de calor irreversible (resistencias) o la almacenan temporalmente en forma de campo eléctrico (condensadores) o campo magnético (bobinas o inductores) para devolverla posteriormente al circuito.
\end{itemize}

\subsection{Potencia Eléctrica y Convenio de Signos (Pasivo vs. Activo)}

La \textbf{potencia eléctrica instantánea} $p(t)$ representa la velocidad con la que se transfiere, consume o suministra energía eléctrica por unidad de tiempo. Aplicando la regla de la cadena a partir de las definiciones de potencial ($v = dW/dq$) e intensidad ($i = dq/dt$):
\begin{equation}
    p(t) = \frac{dW}{dt} = \frac{dW}{dq}\cdot\frac{dq}{dt} = v(t)\cdot i(t) \qquad [\text{SI: Vatio, }\si{\watt} = \si{\joule/\second} = \si{\volt\cdot\ampere}]
\end{equation}

Para determinar sin ambigüedad si un dipolo entre los terminales $A$ y $B$ (con tensión $V_{AB} = V_A - V_B$) está \textbf{absorbiendo} (consumiendo) o \textbf{cediendo} (generando) potencia, se emplea el criterio de signos:

\begin{figure}[H]
\centering
\begin{minipage}{0.47\textwidth}
\centering
\begin{circuitikz}[scale=0.95, transform shape]
    \draw (0,1.6) to[short, o-, i>^=$I$] (2.0,1.6)
          to[generic, l=$Z$, v=$V_{AB}$] (2.0,0)
          to[short, -o] (0,0);
    \node[left] at (0,1.6) {$A\ (+)$};
    \node[left] at (0,0) {$B\ (-)$};
    \node[below, align=center, font=\small\bfseries, color=NavyBlue] at (1.1,-0.35)
        {Convenio Pasivo (Receptor)\\$I$ entra por el terminal $+$ ($A$)};
\end{circuitikz}
\end{minipage}
\hfill
\begin{minipage}{0.47\textwidth}
\centering
\begin{circuitikz}[scale=0.95, transform shape]
    \draw (0,1.6) to[short, o-, i<_=$I$] (2.0,1.6)
          to[generic, l=$G$, v=$V_{AB}$] (2.0,0)
          to[short, -o] (0,0);
    \node[left] at (0,1.6) {$A\ (+)$};
    \node[left] at (0,0) {$B\ (-)$};
    \node[below, align=center, font=\small\bfseries, color=BrickRed] at (1.1,-0.35)
        {Convenio Activo (Generador)\\$I$ sale por el terminal $+$ ($A$)};
\end{circuitikz}
\end{minipage}
\caption{Referencia de signos pasiva (izquierda) y activa (derecha) en un elemento de dos terminales.}
\end{figure}

\begin{ley}[Balance de Potencias y Teorema de Tellegen]
\begin{itemize}
    \item \textbf{En una resistencia}: Siempre absorbe (disipa) potencia activa independientemente del sentido de la corriente:
    \begin{equation}
        P_R = V_R \cdot I_R = I_R^2 \cdot R = \frac{V_R^2}{R} \ge 0
    \end{equation}
    \item \textbf{En una fuente (criterio de signos de clase)}: Si la corriente $I$ circula a través de la fuente desde el terminal negativo ($B$) hacia el positivo ($A$, es decir, \emph{subiendo en potencial}), la potencia suministrada por la fuente al circuito es:
    \begin{equation}
        P_{\text{fuente}} = (V_A - V_B)\cdot I_{\text{subiendo}}
    \end{equation}
    Si al resolver el circuito este producto resulta \textbf{positivo} ($P_f > 0$), la fuente suministra efectivamente potencia al sistema. Si resulta \textbf{negativo} ($P_f < 0$), significa que la corriente real entra por el terminal positivo y la fuente está \textbf{absorbiendo} potencia (actuando como carga o batería en carga).
    \item \textbf{Conservación de la potencia}: En todo circuito eléctrico cerrado se cumple en todo instante que la potencia total suministrada por las fuentes es exactamente igual a la potencia total disipada en las resistencias:
    \begin{equation}
        \sum P_{\text{fuentes}} = \sum P_{\text{resistencias}}
    \end{equation}
\end{itemize}
\end{ley}

\subsection{Elementos Activos: Fuentes Independientes y Dependientes}

Las fuentes de tensión y de intensidad son elementos activos capaces de suministrar energía al sistema y se dividen en \textbf{independientes} y \textbf{dependientes} (o controladas).

\subsubsection{1. Fuentes Independientes}
\begin{itemize}
    \item \textbf{Fuente independiente de tensión}: Proporciona una diferencia de potencial fija $V$ entre sus extremos, con independencia de la corriente $I$ que circule por ella.
    \item \textbf{Fuente independiente de intensidad}: Fija el valor de la intensidad de corriente $I$ en toda la rama serie en la que se encuentra conectada, independientemente de la tensión que caiga entre sus bornes.
\end{itemize}

\begin{ejemplo}[Efecto de una fuente de corriente en una rama serie]
En el siguiente circuito, la fuente de intensidad de $1\,\si{\milli\ampere}$ está conectada en la rama derecha en serie con las resistencias $R_4$ y $R_5$. Por tanto, \textbf{por $R_4$ y $R_5$ pasa obligatoriamente una corriente de $1\,\si{\milli\ampere}$}, sin importar los valores de las demás ramas:
\begin{center}
\begin{circuitikz}[scale=0.88, transform shape]
    \draw (0,0) to[battery1, l=$3\,\si{\volt}$] (0,2.2)
          to[R, l=$R_1$] (2.5,2.2)
          to[R, l=$R_2$] (2.5,1.0)
          to[battery1, l=$5\,\si{\volt}$] (2.5,0)
          to[R, l=$R_3$] (0,0);
    \draw (2.5,2.2) to[isource, l=$1\,\si{\milli\ampere}$] (5.2,2.2)
          to[R, l=$R_4$] (5.2,0)
          to[R, l=$R_5$] (2.5,0);
\end{circuitikz}
\end{center}
\end{ejemplo}

\subsubsection{2. Fuentes Dependientes o Controladas}
Representadas mediante un símbolo romboidal, su tensión o intensidad depende proporcionalmente de la caída de potencial o de la corriente en otra parte del circuito. Existen 4 tipos fundamentales (donde $A_v, A_R, A_G, A_i$ son los factores o constantes de proporcionalidad):

\begin{enumerate}
    \item \textbf{Fuente de tensión dependiente de tensión (VCVS)}: El valor de la tensión de la fuente es proporcional a la caída de potencial $v$ en otra parte del circuito: $V_{\text{dep}} = A_v \cdot v$ ($A_v$ adimensional).
    \item \textbf{Fuente de tensión dependiente de intensidad (CCVS)}: El valor de la tensión de la fuente es proporcional a la intensidad $i$ que pasa por otra rama del circuito: $V_{\text{dep}} = A_R \cdot i$ ($A_R$ en $\si{\ohm}$).
    \item \textbf{Fuente de intensidad dependiente de tensión (VCCS)}: La intensidad de la fuente es proporcional a la caída de tensión $v$ de otra parte del circuito: $I_{\text{dep}} = A_G \cdot v$ ($A_G$ en $\si{\siemens} = \si{\ohm^{-1}}$).
    \item \textbf{Fuente de intensidad dependiente de intensidad (CCCS)}: La intensidad de la fuente es proporcional a la intensidad $i$ en otra parte del circuito: $I_{\text{dep}} = A_i \cdot i$ ($A_i$ adimensional).
\end{enumerate}

\begin{figure}[H]
\centering
\begin{minipage}{0.48\textwidth}
\centering
\begin{circuitikz}[scale=0.82, transform shape]
    \draw (0,0) to[cV, v<={$A_v\,v$}] (0,2.2)
          to[R, l=$R_1$] (2.6,2.2)
          to[R, l=$R_2$] (2.6,1.0)
          to[battery1, l=$5\,\si{\volt}$] (2.6,0) -- (0,0);
    \draw (2.6,2.2) -- (4.4,2.2)
          to[R, l=$R_3$, v=$v$] (4.4,0) -- (2.6,0);
    \node[below, font=\footnotesize\bfseries] at (2.2,-0.3) {Ejemplo VCVS ($V = A_v\,v$)};
\end{circuitikz}
\end{minipage}
\hfill
\begin{minipage}{0.48\textwidth}
\centering
\begin{circuitikz}[scale=0.82, transform shape]
    \draw (0,0) to[cV, v<={$A_R\,i$}] (0,2.2)
          to[R, l=$R_1$] (2.2,2.2)
          to[R, l=$R_3$] (2.2,1.0)
          to[battery1, l=$2\,\si{\volt}$] (2.2,0) -- (0,0);
    \draw (2.2,2.2) to[R, l=$R_2$] (4.2,2.2)
          to[R, l=$R_4$] (4.2,1.0)
          to[isource, l_=$3\,\si{\milli\ampere}$] (4.2,0) -- (2.2,0);
    \draw (4.2,2.2) to[R, l=$R_5$] (6.0,2.2)
          to[R, l=$R_6$, i>^=$i$] (6.0,0) -- (4.2,0);
    \node[below, font=\footnotesize\bfseries] at (3.0,-0.3) {Ejemplo CCVS ($V = A_R\,i$)};
\end{circuitikz}
\end{minipage}
\caption{Ejemplos de circuitos con fuentes controladas de tensión (VCVS y CCVS) de los apuntes de clase.}
\end{figure}

\subsection{Elementos Pasivos: Resistencia, Condensador y Bobina}

\subsubsection{1. Resistencia ($R$) y Deducción Microscópica de la Ley de Ohm}
Un resistor es un elemento tal que, al aplicar una diferencia de potencial $V_A - V_B$ entre sus terminales, deja pasar una corriente de intensidad $I$ directamente proporcional a la diferencia de potencial aplicada.

Consideremos un hilo conductor cilíndrico de longitud $L$ y área transversal $S$, sometido a los potenciales $V_1 > V_2$, lo cual crea en su interior un campo eléctrico uniforme $\vec{E}$ dirigido hacia la derecha:

\begin{figure}[H]
\centering
\begin{tikzpicture}[scale=0.95]
    \draw[thick, fill=gray!12] (0,0) ellipse (0.35 and 0.65);
    \draw[thick, fill=gray!12] (0,0.65) -- (4.5,0.65) arc (90:-90:0.35 and 0.65) -- (0,-0.65) arc (-90:90:0.35 and 0.65);
    \draw[thick] (0,0) ellipse (0.35 and 0.65);
    \node at (0,0) {$S$};
    \node[left] at (-0.45,0) {$V_1$};
    \node[right] at (4.95,0) {$V_2 \quad (V_1 > V_2)$};
    \draw[-{Stealth[length=2.5mm]}, thick, BrickRed] (1.4,0.15) -- (3.1,0.15) node[midway, above] {$\vec{E}$};
    \draw[-{Stealth[length=2.5mm]}, thick, NavyBlue] (1.4,-0.25) -- (3.1,-0.25) node[midway, below] {$\vec{J}$};
    \draw[<->, thick] (0,-1.05) -- (4.5,-1.05) node[midway, below] {$L$};
\end{tikzpicture}
\caption{Conductor cilíndrico de longitud $L$, sección $S$ y conductividad $\sigma$ bajo una diferencia de potencial $V_1 - V_2$.}
\end{figure}

\begin{enumerate}
    \item La relación microscópica entre la densidad de corriente $\vec{J}$ y el campo eléctrico $\vec{E}$ es la Ley de Ohm local: $J = \sigma \cdot E$, donde $\sigma$ es la conductividad del material.
    \item De la relación entre campo eléctrico y potencial ($E = -dV/dl$), dado que $\Delta V = V_2 - V_1$, el módulo del campo es:
    \begin{equation}
        E = -\frac{V_2 - V_1}{L} = \frac{V_1 - V_2}{L} \implies J = \sigma\,\frac{V_1 - V_2}{L}
    \end{equation}
    \item Integrando la densidad de corriente sobre la sección transversal $S$ del conductor:
    \begin{equation}
        I = \int_S J\,dS = \sigma\,\frac{V_1 - V_2}{L}\,S = J\cdot S \implies J = \frac{I}{S}
    \end{equation}
    \item Despejando la caída de potencial $(V_1 - V_2)$ e introduciendo la \textbf{resistividad} $\rho = \frac{1}{\sigma}$ ($\si{\ohm\cdot\meter}$):
    \begin{equation}
        V_1 - V_2 = \frac{1}{\sigma}\,\frac{L}{S}\,I = \rho\,\frac{L}{S}\,I \implies \boxed{V_1 - V_2 = R\cdot I} \quad \text{con} \quad \boxed{R = \rho\,\frac{L}{S}}
    \end{equation}
\end{enumerate}

\begin{aviso}[Unidades prácticas en electrónica]
En los problemas de circuitos es muy frecuente trabajar con tensiones en voltios ($\si{\volt}$), resistencias en kilohmios ($\si{\kilo\ohm} = 10^3\,\si{\ohm}$) e intensidades en miliamperios ($\si{\milli\ampere} = 10^{-3}\,\si{\ampere}$), ya que son unidades mutuamente coherentes:
\begin{equation}
    1\,\si{\volt} = 1\,\si{\kilo\ohm} \times 1\,\si{\milli\ampere} \qquad \Longleftrightarrow \qquad \si{\milli\ampere} = \frac{\si{\volt}}{\si{\kilo\ohm}}
\end{equation}
\end{aviso}

\subsubsection{2. Condensador ($C$)}
Un condensador ideal consiste en dos armaduras conductoras próximas separadas por un aislante (dieléctrico) perfecto. La carga $Q$ almacenada en cualquiera de sus armaduras es proporcional a la diferencia de potencial entre ellas:
\begin{equation}
    Q = C\,(V_A - V_B) \qquad [\text{Capacidad } C \text{ en Faradios: } 1\,\si{\farad} = 1\,\si{\coulomb/\volt}]
\end{equation}
Aunque las cargas no pueden atravesar el aislante, un incremento de tensión $\Delta v$ produce un incremento de carga almacenada $\Delta Q = C\,\Delta v$. En el límite diferencial, esta variación temporal de carga se traduce en una corriente eléctrica:
\begin{equation}
    i_C(t) = \lim_{\Delta t \to 0} \frac{\Delta Q}{\Delta t} = C\,\frac{dv_C(t)}{dt}
\end{equation}
\textbf{Propiedad fundamental}: La tensión entre las armaduras de un condensador \textbf{no puede variar de forma discontinua} (salto brusco), sino que debe ser continua en todo instante: $v_C(0^-) = v_C(0^+)$.

\subsubsection{3. Inductor o Bobina ($L$)}
Una bobina de $N$ espiras atravesada por una corriente $I$ genera un flujo magnético $\Phi_B$ proporcional a dicha corriente. Se define la \textbf{inductancia} o coeficiente de autoinducción $L$ (medida en Henrios, $1\,\si{\henry} = 1\,\si{\ohm\cdot\second} = 1\,\si{\weber/\ampere}$) como:
\begin{equation}
    L = \frac{N\,\Phi_B}{I}
\end{equation}
Por la Ley de Faraday-Lenz, una variación de corriente $\Delta I$ produce una variación de flujo $\Delta \Phi_B$ y, por tanto, una caída de tensión autoinducida entre sus extremos:
\begin{equation}
    v_L(t) = V_A - V_B = L\,\frac{di_L(t)}{dt}
\end{equation}
\textbf{Propiedad fundamental}: La intensidad de corriente que atraviesa una bobina \textbf{no puede cambiar instantáneamente}: $i_L(0^-) = i_L(0^+)$.

\subsection{Asociación de Elementos y Teorema de Kennelly ($\Delta \leftrightarrow \text{Y}$)}

\subsubsection{1. Asociación de Resistencias y Condensadores}
\begin{table}[H]
\centering
\renewcommand{\arraystretch}{1.5}
\resizebox{\textwidth}{!}{%
\begin{tabular}{@{}lcc@{}}
\toprule
\textbf{Elemento} & \textbf{Asociación en Serie} & \textbf{Asociación en Paralelo} \\ \midrule
\textbf{Resistencias} & $\displaystyle R_{\text{total}} = R_1 + R_2 + \dots + R_N$ & $\displaystyle \frac{1}{R_{\text{total}}} = \frac{1}{R_1} + \frac{1}{R_2} + \dots + \frac{1}{R_N} \quad \left(R_1 \parallel R_2 = \frac{R_1 R_2}{R_1 + R_2}\right)$ \\
\textbf{Condensadores} & $\displaystyle \frac{1}{C_{\text{total}}} = \frac{1}{C_1} + \frac{1}{C_2} + \dots + \frac{1}{C_N}$ & $\displaystyle C_{\text{total}} = C_1 + C_2 + \dots + C_N$ \\
\textbf{Fuentes} & Solo de tensión: $V_{\text{total}} = \sum V_k$ & Solo de intensidad: $I_{\text{total}} = \sum I_k$ \\ \bottomrule
\end{tabular}%
}
\caption{Reglas de asociación equivalente de resistencias, condensadores y fuentes ideales.}
\end{table}

\subsubsection{2. Teorema de Kennelly: Transformación Triángulo ($\Delta$) $\leftrightarrow$ Estrella ($\text{Y}$)}
Cuando en una red resistiva (como un puente de Wheatstone no equilibrado) las resistencias no están ni en serie ni en paralelo, el \textbf{Teorema de Kennelly} permite transformar una configuración en \textbf{Triángulo o Delta ($\Delta$)} entre tres nudos $(a, b, c)$ en una configuración equivalente en \textbf{Estrella ($\text{Y}$ o $\text{T}$)} y viceversa:

\begin{figure}[H]
\centering
\begin{tikzpicture}[scale=0.95, transform shape]
    % Triángulo (Delta)
    \node[circle, fill=black, inner sep=1.5pt, label=above:$a$] (a1) at (-3.5,2.2) {};
    \node[circle, fill=black, inner sep=1.5pt, label=left:$b$] (b1) at (-5.2,0) {};
    \node[circle, fill=black, inner sep=1.5pt, label=right:$c$] (c1) at (-1.8,0) {};
    \draw (a1) to[R, l_=$R_3$] (b1);
    \draw (a1) to[R, l=$R_2$] (c1);
    \draw (b1) to[R, l_=$R_1$] (c1);
    \node[below, font=\small\bfseries, color=NavyBlue] at (-3.5,-0.55) {Red en Triángulo ($\Delta$)};

    % Flecha doble
    \node[font=\Large\bfseries] at (0,1.1) {$\Longleftrightarrow$};

    % Estrella (Y)
    \node[circle, fill=black, inner sep=1.5pt, label=above:$a$] (a2) at (3.5,2.4) {};
    \node[circle, fill=black, inner sep=1.5pt, label=left:$b$] (b2) at (1.8,0) {};
    \node[circle, fill=black, inner sep=1.5pt, label=right:$c$] (c2) at (5.2,0) {};
    \coordinate (n) at (3.5,1.0);
    \fill (n) circle (1.5pt);
    \draw (a2) to[R, l=$R_A$] (n);
    \draw (n) to[R, l_=$R_B$] (b2);
    \draw (n) to[R, l=$R_C$] (c2);
    \node[below, font=\small\bfseries, color=BrickRed] at (3.5,-0.55) {Red en Estrella ($\text{Y}$)};
\end{tikzpicture}
\caption{Equivalencia entre una red en Triángulo ($\Delta$: $R_1, R_2, R_3$) y una red en Estrella ($\text{Y}$: $R_A, R_B, R_C$).}
\end{figure}

\begin{ley}[Fórmulas de Transformación de Kennelly]
\begin{itemize}
    \item \textbf{Paso de Triángulo ($\Delta$) a Estrella ($\text{Y}$)}: Cada resistencia de la estrella es el producto de las dos resistencias adyacentes del triángulo dividido entre la suma de las tres resistencias del triángulo:
    \begin{equation}
        R_A = \frac{R_2\cdot R_3}{R_1 + R_2 + R_3}, \qquad R_B = \frac{R_1\cdot R_3}{R_1 + R_2 + R_3}, \qquad R_C = \frac{R_1\cdot R_2}{R_1 + R_2 + R_3}
    \end{equation}
    \item \textbf{Paso de Estrella ($\text{Y}$) a Triángulo ($\Delta$)}: Cada resistencia del triángulo es la suma de los productos cruzados dos a dos de la estrella dividida entre la resistencia de la estrella opuesta:
    \begin{equation}
        R_1 = \frac{R_A R_B + R_A R_C + R_B R_C}{R_A}, \; R_2 = \frac{R_A R_B + R_A R_C + R_B R_C}{R_B}, \; R_3 = \frac{R_A R_B + R_A R_C + R_B R_C}{R_C}
    \end{equation}
\end{itemize}
\end{ley}

\subsection{Elementos Reales y Comportamiento en Señal Continua (DC)}

\subsubsection{1. Efecto de la Temperatura en Resistencias Reales}
En la realidad, las resistencias, condensadores y bobinas presentan efectos parásitos y dependen de la temperatura. El valor de una resistencia $R_H$ a una temperatura alta $T_H$ en función de su valor $R_L$ a una temperatura de referencia $T_L$ viene dado por:
\begin{equation}
    R_H = R_L\,\big[1 + \alpha\,(T_H - T_L)\big]
\end{equation}
donde $\alpha$ ($\si{\kelvin^{-1}}$ o $\si{\celsius^{-1}}$) es el coeficiente de temperatura del material.

\subsubsection{2. Fuentes Reales de Tensión y de Intensidad y Transformación de Fuentes}
\begin{itemize}
    \item \textbf{Fuente de tensión real}: Se modela mediante una fuente ideal $V_s$ en \textbf{serie} con una resistencia interna $R_i$. Al conectar una carga entre $A$ y $B$, la tensión en bornes cae ($V_{AB} = V_s - I R_i$). Para que la fuente sea buena y la tensión varíe lo mínimo posible, \textbf{$R_i$ debe ser lo más pequeña posible} (idealmente $R_i = 0$).
    \item \textbf{Fuente de intensidad real}: Se modela mediante una fuente ideal $I_s$ en \textbf{paralelo} con una resistencia interna $R_i$. Para que casi toda la corriente salga hacia la carga, \textbf{$R_i$ debe ser lo más grande posible} (idealmente $R_i \to \infty$).
    \item \textbf{Transformación de fuentes}: Una fuente de tensión real $(V_s, R_i\text{ serie})$ es exactamente equivalente entre sus terminales $A$ y $B$ a una fuente de corriente real $(I_s, R_i\text{ paralelo})$ siempre que se cumpla:
    \begin{equation}
        V_s = I_s \cdot R_i \qquad \Longleftrightarrow \qquad I_s = \frac{V_s}{R_i}
    \end{equation}
\end{itemize}

\begin{figure}[H]
\centering
\begin{minipage}{0.46\textwidth}
\centering
\begin{circuitikz}[scale=0.92, transform shape]
    \draw (0,0) to[V, v<=$V_s$] (0,2.0)
          to[R, l=$R_i$] (2.8,2.0) to[short, -o] (3.2,2.0) node[right] {$A$};
    \draw (0,0) to[short, -o] (3.2,0) node[right] {$B$};
    \node[below, font=\small\bfseries] at (1.6,-0.3) {Fuente de tensión real ($R_i$ pequeña)};
\end{circuitikz}
\end{minipage}
\hfill
\begin{minipage}{0.46\textwidth}
\centering
\begin{circuitikz}[scale=0.92, transform shape]
    \draw (0,0) to[isource, l=$I_s$] (0,2.0) -- (2.8,2.0) to[short, -o] (3.2,2.0) node[right] {$A$};
    \draw (1.6,2.0) to[R, l=$R_i$] (1.6,0);
    \draw (0,0) to[short, -o] (3.2,0) node[right] {$B$};
    \node[below, font=\small\bfseries] at (1.6,-0.3) {Fuente de corriente real ($R_i$ grande)};
\end{circuitikz}
\end{minipage}
\caption{Modelos circuitales de fuentes reales y su equivalencia ($V_s = I_s R_i$).}
\end{figure}

\subsubsection{3. Elementos en Régimen Permanente de Corriente Continua (DC)}
Cuando un circuito alimentado por fuentes continuas lleva mucho tiempo funcionando, todas las magnitudes se estabilizan ($V = \text{cte}, I = \text{cte}$):
\begin{itemize}
    \item \textbf{Resistencia}: Se comporta igual ($V = I R$).
    \item \textbf{Condensador}: Se carga hasta una tensión constante ($dV/dt = 0 \implies I = C\frac{dV}{dt} = 0$). Por tanto, en continua estacionaria \textbf{un condensador equivale a un interruptor abierto (circuito abierto)}.
    \item \textbf{Bobina}: La intensidad que pasa por ella es fija ($dI/dt = 0 \implies V = L\frac{dI}{dt} = 0$). Por tanto, en continua estacionaria \textbf{una bobina equivale a un hilo conductor ideal (cortocircuito)}.
\end{itemize}

\begin{ejemplo}[Análisis en régimen continuo estacionario con bobina y condensador]
Calcula la tensión en el nudo $A$ y la carga acumulada en el condensador de $5\,\si{\nano\farad}$ del siguiente circuito cuando trabaja en corriente continua estacionaria:

\begin{center}
\begin{minipage}{0.48\textwidth}
\centering
\begin{circuitikz}[scale=0.82, transform shape]
    \draw (0,0) to[V, v<=$5\,\si{\volt}$] (0,2.2)
          to[R, l={$1\,\si{\kilo\ohm}$}] (2.2,2.2) node[above] {$A$}
          to[L, l={$0{,}1\,\si{\henry}$}] (4.4,2.2)
          to[R, l={$3\,\si{\kilo\ohm}$}] (4.4,0) -- (0,0);
    \draw (2.2,2.2) to[C, l={$5\,\si{\nano\farad}$}] (2.2,0);
    \node[below, font=\footnotesize\bfseries] at (2.2,-0.3) {Circuito original};
\end{circuitikz}
\end{minipage}
\hfill
\begin{minipage}{0.48\textwidth}
\centering
\begin{circuitikz}[scale=0.82, transform shape]
    \draw (0,0) to[V, v<=$5\,\si{\volt}$] (0,2.2)
          to[R, l={$1\,\si{\kilo\ohm}$}] (2.2,2.2) node[above] {$A$}
          to[short] (4.0,2.2)
          to[R, l={$3\,\si{\kilo\ohm}$}] (4.0,0) -- (0,0);
    \draw (2.0,0) node[ground] {};
    \draw[-{Stealth}, NavyBlue, thick] (1.3,0.6) arc (-150:150:0.5) node[midway, right] {$I$};
    \node[below, font=\footnotesize\bfseries] at (2.0,-0.45) {Equivalente en continua (DC)};
\end{circuitikz}
\end{minipage}
\end{center}

\textbf{Resolución paso a paso}:
\begin{enumerate}
    \item En continua, sustituimos el condensador de $5\,\si{\nano\farad}$ por un circuito abierto y la bobina de $0{,}1\,\si{\henry}$ por un hilo conductor.
    \item Referenciamos el conductor inferior a tierra ($0\,\si{\volt}$). Queda una única malla con las resistencias de $1\,\si{\kilo\ohm}$ y $3\,\si{\kilo\ohm}$ en serie:
    \begin{equation*}
        I = \frac{5\,\si{\volt}}{1\,\si{\kilo\ohm} + 3\,\si{\kilo\ohm}} = \frac{5}{4}\,\si{\milli\ampere} = 1{,}25\,\si{\milli\ampere}
    \end{equation*}
    \item La tensión en el nudo $A$ (que coincide con la tensión en bornes del condensador y de la resistencia de $3\,\si{\kilo\ohm}$) es:
    \begin{equation*}
        V_A = I \cdot 3\,\si{\kilo\ohm} = \frac{3}{4}\cdot 5\,\si{\volt} = \frac{15}{4}\,\si{\volt} = 3{,}75\,\si{\volt}
    \end{equation*}
    \item La carga acumulada en el condensador es:
    \begin{equation*}
        Q = C \cdot V_A = 5\,\si{\nano\farad} \cdot \frac{15}{4}\,\si{\volt} = 18{,}75\,\si{\nano\coulomb} = 1{,}875\times 10^{-8}\,\si{\coulomb}
    \end{equation*}
\end{enumerate}
\end{ejemplo}

\section{Leyes de Kirchhoff y Métodos Sistemáticos de Análisis}

\subsection{Topología de Redes y Leyes de Kirchhoff}
\begin{itemize}
    \item \textbf{Nudo (o nodo)}: Punto de unión donde concurren dos o más elementos. Se llama \emph{nudo principal} a aquel que conecta tres o más ramas.
    \item \textbf{Rama}: Conjunto de elementos conectados en serie entre dos nudos principales consecutivos (por todos ellos circula exactamente la misma corriente).
    \item \textbf{Lazo}: Cualquier camino cerrado dentro del circuito.
    \item \textbf{Malla}: Lazo cerrado que no contiene ningún otro lazo en su interior.
\end{itemize}

\begin{ley}[Leyes de Kirchhoff]
\begin{enumerate}
    \item \textbf{Primera Ley de Kirchhoff (Ley de Corrientes o de Nudos, LCK)}: Basada en la conservación de la carga eléctrica, establece que en cualquier nudo del circuito la suma algebraica de las corrientes que entran es igual a la suma de las corrientes que salen:
    \begin{equation}
        \sum_{\text{entran}} I_k = \sum_{\text{salen}} I_j \qquad \Longleftrightarrow \qquad \sum_{k=1}^{N} I_k = 0
    \end{equation}
    \item \textbf{Segunda Ley de Kirchhoff (Ley de Tensiones o de Mallas, LTK)}: Basada en el carácter conservativo del campo electrostático, establece que la suma algebraica de las diferencias de potencial a lo largo de cualquier trayectoria cerrada (lazo o malla) es nula:
    \begin{equation}
        \sum_{\text{malla}} V_k = 0 \qquad \Longleftrightarrow \qquad \sum V_{\text{fuentes}} = \sum I_k R_k
    \end{equation}
\end{enumerate}
\end{ley}

\subsection{Divisor de Tensión y Divisor de Corriente}

\begin{figure}[H]
\centering
\begin{minipage}{0.48\textwidth}
\centering
\begin{circuitikz}[scale=0.9, transform shape]
    \node[font=\small\bfseries, color=NavyBlue] at (1.6,2.7) {Divisor de Tensión};
    \draw (0,0) to[V, v<=$V$] (0,2.1)
          to[R, l=$R_1$] (3.0,2.1) node[right] {$A$}
          to[R, l=$R_2$] (3.0,0) -- (0,0);
    \draw (1.5,0) node[ground] {};
    \draw[-{Stealth}, NavyBlue, thick] (1.1,0.6) arc (-150:150:0.45) node[midway, right] {$I$};
\end{circuitikz}
\end{minipage}
\hfill
\begin{minipage}{0.48\textwidth}
\centering
\begin{circuitikz}[scale=0.9, transform shape]
    \node[font=\small\bfseries, color=NavyBlue] at (1.8,2.7) {Divisor de Corriente};
    \draw (0,0) to[isource, l=$I_T$] (0,2.1) -- (3.6,2.1);
    \draw (1.8,2.1) to[R, l=$R_1$, i>^=$I_1$] (1.8,0);
    \draw (3.6,2.1) to[R, l=$R_2$, i>^=$I_2$] (3.6,0);
    \draw (0,0) -- (3.6,0);
\end{circuitikz}
\end{minipage}
\caption{Circuitos fundamentales de divisor de tensión resistivo (izquierda) y divisor de corriente (derecha).}
\end{figure}

\begin{itemize}
    \item \textbf{Divisor de tensión}: En dos resistencias en serie alimentadas por una tensión total $V$, la corriente común es $I = \frac{V}{R_1 + R_2}$, por lo que la caída de tensión en cada resistencia es proporcional a su propio valor:
    \begin{equation}
        V_{R1} = I\cdot R_1 = \frac{R_1}{R_1 + R_2}\,V, \qquad V_{R2} = V_A = I\cdot R_2 = \frac{R_2}{R_1 + R_2}\,V
    \end{equation}
    \item \textbf{Divisor de corriente}: En dos resistencias en paralelo alimentadas por una corriente total $I_T$, la tensión común es $V = I_T (R_1 \parallel R_2) = I_T \frac{R_1 R_2}{R_1 + R_2}$, repartiéndose la corriente de forma inversamente proporcional a cada resistencia:
    \begin{equation}
        I_1 = \frac{R_2}{R_1 + R_2}\,I_T, \qquad I_2 = \frac{R_1}{R_1 + R_2}\,I_T
    \end{equation}
\end{itemize}

\subsection{Método de las Corrientes de Malla}

\begin{metodo}[Procedimiento General del Método de Mallas]
\begin{enumerate}
    \item Asignar una corriente ficticia de malla ($I_1, I_2, \dots, I_M$) a cada malla independiente del circuito (habitualmente todas en sentido horario).
    \item Escribir la ecuación de la 2ª Ley de Kirchhoff para cada malla:
    \begin{itemize}
        \item En el miembro izquierdo se colocan las fuentes de tensión de la malla: con signo \textbf{positivo ($+$)} si la corriente de malla pasa de $-$ a $+$ por la fuente (a favor), y \textbf{negativo ($-$)} si pasa de $+$ a $-$ (en contra).
        \item En el miembro derecho se suman las caídas de tensión en las resistencias: la corriente de la propia malla multiplica a la suma de todas las resistencias de esa malla, y se restan las corrientes de las mallas vecinas multiplicadas por la resistencia compartida.
    \end{itemize}
    \item Agrupar las incógnitas en forma matricial $[V] = [R][I]$ y resolver mediante la regla de Cramer o sustitución. Si alguna intensidad sale negativa, significa que su sentido físico real es el opuesto al supuesto inicialmente.
\end{enumerate}
\end{metodo}

\begin{ejemplo}[Resolución paso a paso por el método de mallas]
Determina las corrientes de malla $I_1$, $I_2$ y la corriente que atraviesa la resistencia central $R_2$ en el siguiente circuito con $R_1 = 1\,\si{\kilo\ohm}$, $R_2 = 2\,\si{\kilo\ohm}$, $R_3 = 3\,\si{\kilo\ohm}$, $R_4 = 4\,\si{\kilo\ohm}$, $V_1 = 1\,\si{\volt}$, $V_2 = 2\,\si{\volt}$ y $V_3 = 3\,\si{\volt}$:

\begin{center}
\begin{circuitikz}[scale=0.92, transform shape]
    \draw (0,0) to[V, v<=$V_1$] (0,2.4)
          to[R, l=$R_1$] (2.8,2.4)
          to[R, l=$R_2$] (2.8,0)
          to[R, l=$R_3$] (0,0);
    \draw (2.8,2.4) to[V, v<=$V_2$] (5.6,2.4)
          to[R, l=$R_4$] (5.6,0)
          to[V, v=$V_3$] (2.8,0);
    \draw[-{Stealth}, NavyBlue, thick] (1.0,1.2) arc (150:-150:0.45) node[midway, left] {$I_1$};
    \draw[-{Stealth}, BrickRed, thick] (3.8,1.2) arc (150:-150:0.45) node[midway, left] {$I_2$};
\end{circuitikz}
\end{center}

\textbf{Resolución paso a paso}:
\begin{enumerate}
    \item \textbf{Planteamiento de las ecuaciones de cada malla} (tomando $I_1$ e $I_2$ en sentido horario):
    \begin{align*}
        \text{Malla 1:} \quad V_1 &= I_1 R_1 + (I_1 - I_2) R_2 + I_1 R_3 \\
        \text{Malla 2:} \quad V_2 - V_3 &= I_2 R_4 + (I_2 - I_1) R_2
    \end{align*}
    \item \textbf{Agrupación de incógnitas y sustitución de valores} (en $\si{\kilo\ohm}$, $\si{\volt}$ y $\si{\milli\ampere}$):
    \begin{align*}
        V_1 &= (R_1 + R_2 + R_3) I_1 - R_2 I_2 \implies 1 = 6 I_1 - 2 I_2 \\
        V_2 - V_3 &= -R_2 I_1 + (R_2 + R_4) I_2 \implies 2 - 3 = -1 = -2 I_1 + 6 I_2
    \end{align*}
    En forma matricial:
    \begin{equation*}
        \begin{pmatrix} 1 \\ -1 \end{pmatrix} = \begin{pmatrix} 6 & -2 \\ -2 & 6 \end{pmatrix} \begin{pmatrix} I_1 \\ I_2 \end{pmatrix}
    \end{equation*}
    \item \textbf{Resolución por la regla de Cramer}:
    \begin{equation*}
        I_1 = \frac{\begin{vmatrix} 1 & -2 \\ -1 & 6 \end{vmatrix}}{\begin{vmatrix} 6 & -2 \\ -2 & 6 \end{vmatrix}} = \frac{6 - 2}{36 - 4} = \frac{4}{32} = \frac{1}{8}\,\si{\milli\ampere} = 0{,}125\,\si{\milli\ampere}
    \end{equation*}
    \begin{equation*}
        I_2 = \frac{\begin{vmatrix} 6 & 1 \\ -2 & -1 \end{vmatrix}}{\begin{vmatrix} 6 & -2 \\ -2 & 6 \end{vmatrix}} = \frac{-6 + 2}{36 - 4} = \frac{-4}{32} = -\frac{1}{8}\,\si{\milli\ampere} = -0{,}125\,\si{\milli\ampere}
    \end{equation*}
    Como $I_2$ sale negativa, su sentido físico real es antihorario. La corriente total que baja por $R_2$ es:
    \begin{equation*}
        I_{R2} = I_1 - I_2 = \frac{1}{8} - \left(-\frac{1}{8}\right) = \frac{1}{4}\,\si{\milli\ampere} = 0{,}25\,\si{\milli\ampere}
    \end{equation*}
\end{enumerate}
\end{ejemplo}

\subsection{Método de las Tensiones de Nudos}

\begin{metodo}[Procedimiento General del Método de Nudos]
\begin{enumerate}
    \item Identificar todos los nudos principales del circuito y elegir uno de ellos (habitualmente el que conecta más ramas o fuentes de tensión) como \textbf{nudo de referencia o tierra ($0\,\si{\volt}$)}.
    \item Asignar una tensión incógnita ($V_a, V_b, V_c, \dots$) a cada uno de los demás nudos principales respecto a tierra.
    \item Aplicar la 1ª Ley de Kirchhoff (LCK) en cada nudo principal expresando las corrientes de rama en función de las tensiones de nudo mediante la Ley de Ohm ($I_{ij} = \frac{V_i - V_j}{R_{ij}}$).
\end{enumerate}
\end{metodo}

Por ejemplo, en una red en puente con nudos principales $a, b, c$ y tierra $d$ ($V_d = 0\,\si{\volt}$):

\begin{figure}[H]
\centering
\begin{circuitikz}[scale=0.9, transform shape]
    \draw (0,0) to[V, v<=$V_1$] (0,2.5)
          to[R, l=$R_1$] (2.5,2.5) -- (4.2,2.5) node[above] {$a$};
    \draw (4.2,2.5) to[R, l_=$R_2$] (2.6,1.25) node[left] {$b$};
    \draw (4.2,2.5) to[R, l=$R_3$] (5.8,1.25) node[right] {$c$};
    \draw (2.6,1.25) to[R, l=$R_6$] (5.8,1.25);
    \draw (2.6,1.25) to[R, l_=$R_5$] (4.2,0) node[below right] {$d\ (0\,\si{\volt})$};
    \draw (5.8,1.25) to[R, l=$R_4$] (4.2,0);
    \draw (0,0) -- (4.2,0);
    \draw (2.0,0) node[ground] {};
\end{circuitikz}
\caption{Red en puente para plantear las ecuaciones de nudos en $a$, $b$ y $c$.}
\end{figure}

Las ecuaciones de nudo (asumiendo todas las corrientes entrando al nudo o saliendo de él) son:
\begin{align}
    \text{Nudo } a: \quad & \frac{V_1 - V_a}{R_1} + \frac{V_b - V_a}{R_2} + \frac{V_c - V_a}{R_3} = 0 \\
    \text{Nudo } b: \quad & \frac{V_a - V_b}{R_2} + \frac{V_c - V_b}{R_6} + \frac{0 - V_b}{R_5} = 0 \\
    \text{Nudo } c: \quad & \frac{V_a - V_c}{R_3} + \frac{V_b - V_c}{R_6} + \frac{0 - V_c}{R_4} = 0
\end{align}

\begin{ejemplo}[Otro ejemplo resuelto por nudos con fuentes flotantes en las ramas]
Resuelve por el método de nudos el circuito de la figura (con los mismos valores del ejemplo anterior: $V_1=1\,\si{\volt}, V_2=2\,\si{\volt}, V_3=3\,\si{\volt}, R_1=1\,\si{\kilo\ohm}, R_2=2\,\si{\kilo\ohm}, R_3=3\,\si{\kilo\ohm}, R_4=4\,\si{\kilo\ohm}$):

\begin{center}
\begin{circuitikz}[scale=0.92, transform shape]
    \draw (0,0) node[below left] {$E$} to[V, v<=$V_1$] (0,2.4) node[above left] {$F$}
          to[R, l=$R_1$, i>^=$I_1$] (2.8,2.4) node[above] {$V_A$}
          to[R, l=$R_2$, i>^=$I_2$] (2.8,0) node[below right] {$B\ (0\,\si{\volt})$}
          to[R, l=$R_3$, i>^=$I_1$] (0,0);
    \draw (2.8,2.4) to[V, v<=$V_2$] (5.6,2.4) node[above right] {$C$}
          to[R, l=$R_4$, i>^=$I_4$] (5.6,0) node[below right] {$D$}
          to[V, v=$V_3$] (2.8,0);
    \draw (2.8,0) node[ground] {};
\end{circuitikz}
\end{center}

\textbf{Resolución paso a paso}:
\begin{enumerate}
    \item \textbf{Referenciamos a tierra el nudo principal inferior}: $V_B = 0\,\si{\volt}$.
    \item \textbf{Aplicamos LCK en el nudo principal superior $V_A$}:
    \begin{equation*}
        I_1 = I_2 + I_4
    \end{equation*}
    \item \textbf{Expresamos cada intensidad en función de $V_A$}:
    \begin{itemize}
        \item En la rama central: $I_2 = \dfrac{V_A - 0}{R_2} = \dfrac{V_A}{2\,\si{\kilo\ohm}}$.
        \item En la rama izquierda: por toda la rama serie ($R_3, V_1, R_1$) circula la misma corriente $I_1$. Como $0 - V_E = I_1 R_3 \implies V_E = -I_1 R_3$, y $V_F = V_E + V_1 = V_1 - I_1 R_3$, entonces:
        \begin{equation*}
            I_1 = \frac{V_F - V_A}{R_1} = \frac{V_1 - I_1 R_3 - V_A}{R_1} \implies I_1(R_1 + R_3) = V_1 - V_A \implies I_1 = \frac{V_1 - V_A}{R_1 + R_3} = \frac{1 - V_A}{4\,\si{\kilo\ohm}}
        \end{equation*}
        \item En la rama derecha: $V_C = V_A + V_2$ y $V_D = 0 + V_3$, por lo que:
        \begin{equation*}
            I_4 = \frac{V_C - V_D}{R_4} = \frac{V_A + V_2 - V_3}{R_4} = \frac{V_A + 2 - 3}{4\,\si{\kilo\ohm}} = \frac{V_A - 1}{4\,\si{\kilo\ohm}}
        \end{equation*}
    \end{itemize}
    \item \textbf{Sustituimos en $I_1 = I_2 + I_4$ y despejamos $V_A$}:
    \begin{equation*}
        \frac{1 - V_A}{4} = \frac{V_A}{2} + \frac{V_A - 1}{4} \implies \frac{1 - V_A}{4} + \frac{1 - V_A}{4} = \frac{V_A}{2} \implies \frac{1 - V_A}{2} = \frac{V_A}{2} \implies \boxed{V_A = \frac{1}{2}\,\si{\volt} = 0{,}5\,\si{\volt}}
    \end{equation*}
    \item \textbf{Conocido $V_A$, obtenemos inmediatamente las tres corrientes}:
    \begin{equation*}
        I_2 = \frac{1/2\,\si{\volt}}{2\,\si{\kilo\ohm}} = \frac{1}{4}\,\si{\milli\ampere}, \qquad I_1 = \frac{1 - 1/2}{4\,\si{\kilo\ohm}} = \frac{1}{8}\,\si{\milli\ampere}, \qquad I_4 = \frac{1/2 - 1}{4\,\si{\kilo\ohm}} = -\frac{1}{8}\,\si{\milli\ampere}
    \end{equation*}
\end{enumerate}
\end{ejemplo}

\subsection{Teorema de Thévenin y Teorema de Norton}

\begin{ley}[Teoremas de Thévenin y de Norton]
\begin{itemize}
    \item \textbf{Teorema de Thévenin}: Cualquier circuito lineal complejo visto desde dos terminales $A$ y $B$ puede sustituirse por una única \textbf{fuente independiente de tensión $V_{th}$} conectada en \textbf{serie} con una \textbf{resistencia equivalente $R_{th}$}, donde:
    \begin{enumerate}
        \item $V_{th} = V_A - V_B$ es la tensión en circuito abierto entre los terminales $A$ y $B$ (desconectando la carga).
        \item $R_{th}$ es la resistencia equivalente vista entre $A$ y $B$ al \textbf{anular todas las fuentes independientes} del circuito:
        \begin{itemize}
            \item Las \textbf{fuentes de tensión} se anulan sustituyéndolas por un \textbf{cortocircuito} (hilo conductor, $V=0$).
            \item Las \textbf{fuentes de intensidad} se anulan sustituyéndolas por un \textbf{circuito abierto} (interruptor abierto, $I=0$).
        \end{itemize}
    \end{enumerate}
    \item \textbf{Teorema de Norton}: El mismo circuito lineal entre $A$ y $B$ puede reducirse a una única \textbf{fuente de intensidad $I_N$} en \textbf{paralelo} con una \textbf{resistencia $R_N$}, donde $R_N = R_{th}$ e $I_N = I_{cc}$ es la corriente de cortocircuito que circula de $A$ a $B$ al unir ambos terminales con un cable ideal.
    \item \textbf{Relación entre Thévenin y Norton}: Se puede pasar siempre de uno a otro mediante:
    \begin{equation}
        R_{th} = R_N, \qquad V_{th} = I_N \cdot R_N \qquad \Longleftrightarrow \qquad I_N = \frac{V_{th}}{R_{th}}
    \end{equation}
\end{itemize}
\end{ley}

\begin{figure}[H]
\centering
\begin{minipage}{0.48\textwidth}
\centering
\begin{circuitikz}[scale=0.92, transform shape]
    \node[font=\small\bfseries, color=NavyBlue] at (1.8,2.65) {Equivalente de Thévenin};
    \draw (0,0) to[V, v<=$V_{th}$] (0,2.0)
          to[R, l=$R_{th}$] (2.6,2.0)
          to[short, -o] (3.2,2.0) node[right] {$A$};
    \draw (0,0) to[short, -o] (3.2,0) node[right] {$B$};
\end{circuitikz}
\end{minipage}
\hfill
\begin{minipage}{0.48\textwidth}
\centering
\begin{circuitikz}[scale=0.92, transform shape]
    \node[font=\small\bfseries, color=BrickRed] at (1.8,2.65) {Equivalente de Norton};
    \draw (0,0) to[isource, l=$I_N$] (0,2.0) -- (2.8,2.0)
          to[short, -o] (3.4,2.0) node[right] {$A$};
    \draw (1.8,2.0) to[R, l=$R_N$] (1.8,0);
    \draw (0,0) to[short, -o] (3.4,0) node[right] {$B$};
\end{circuitikz}
\end{minipage}
\caption{Circuitos equivalentes de Thévenin (izquierda) y de Norton (derecha) entre los terminales $A$ y $B$.}
\end{figure}

\begin{ejemplo}[Cálculo comparativo por Mallas, Thévenin y Norton]
Sea el siguiente circuito con $V = 10\,\si{\volt}$, $R_1 = 1\,\si{\kilo\ohm}$, $R_2 = 5\,\si{\kilo\ohm}$, $R_3 = 1\,\si{\kilo\ohm}$ y resistencia de carga $R_4 = 1\,\si{\kilo\ohm}$ conectada entre $A$ y $B$. Calcula la tensión $V_{AB}$ de tres formas distintas: (1) resolviendo el circuito completo por mallas, (2) mediante el equivalente de Thévenin y (3) mediante el equivalente de Norton.

\begin{center}
\begin{circuitikz}[scale=0.92, transform shape]
    \draw (0,0) to[V, v<=$10\,\si{\volt}$] (0,2.2)
          to[R, l={$R_1=1\,\si{\kilo\ohm}$}] (2.5,2.2)
          to[R, l={$R_2=5\,\si{\kilo\ohm}$}] (2.5,0) -- (0,0);
    \draw (2.5,2.2) to[R, l={$R_3=1\,\si{\kilo\ohm}$}] (5.0,2.2) to[short, -*] (5.4,2.2) node[above] {$A$};
    \draw (2.5,0) -- (5.0,0) to[short, -*] (5.4,0) node[below] {$B$};
    \draw (5.4,2.2) to[R, l={$R_4=1\,\si{\kilo\ohm}$}] (5.4,0);
    \draw[-{Stealth}, NavyBlue, thick] (0.9,1.1) arc (150:-150:0.4) node[midway, left] {$I_1$};
    \draw[-{Stealth}, BrickRed, thick] (3.6,1.1) arc (150:-150:0.4) node[midway, left] {$I_2$};
\end{circuitikz}
\end{center}

\textbf{Método 1: Resolución directa por mallas}:
\begin{align*}
    10 &= I_1(R_1 + R_2) - I_2 R_2 = 6 I_1 - 5 I_2 \\
    0 &= -I_1 R_2 + I_2(R_2 + R_3 + R_4) = -5 I_1 + 7 I_2
\end{align*}
Por Cramer ($\Delta = 42 - 25 = 17$):
\begin{equation*}
    I_1 = \frac{70}{17}\,\si{\milli\ampere}, \qquad I_2 = I_{R4} = \frac{50}{17}\,\si{\milli\ampere} \implies \boxed{V_{AB} = I_{R4}\cdot R_4 = \frac{50}{17}\,\si{\volt} \approx 2{,}94\,\si{\volt}}
\end{equation*}

\textbf{Método 2: Mediante el Equivalente de Thévenin (quitando $R_4$)}:
\begin{center}
\begin{minipage}{0.46\textwidth}
\centering
\begin{circuitikz}[scale=0.82, transform shape]
    \draw (0,0) to[short] (0,2.0)
          to[R, l={$R_1=1\,\si{\kilo\ohm}$}] (2.2,2.0)
          to[R, l={$R_2=5\,\si{\kilo\ohm}$}] (2.2,0) -- (0,0);
    \draw (2.2,2.0) to[R, l={$R_3=1\,\si{\kilo\ohm}$}, -o] (4.4,2.0) node[right] {$A$};
    \draw (2.2,0) to[short, -o] (4.4,0) node[right] {$B$};
    \node[below, font=\footnotesize\bfseries] at (2.2,-0.3) {Cálculo de $R_{th}$ (fuente anulada)};
\end{circuitikz}
\end{minipage}
\hfill
\begin{minipage}{0.46\textwidth}
\centering
\begin{circuitikz}[scale=0.82, transform shape]
    \draw (0,0) to[V, v<=$10\,\si{\volt}$] (0,2.0)
          to[R, l=$R_1$] (2.2,2.0) node[above] {$V_A$}
          to[R, l=$R_2$] (2.2,0) -- (0,0);
    \draw (2.2,2.0) to[R, l=$R_3$, i>^={$I=0$}, -o] (4.4,2.0) node[right] {$A$};
    \draw (2.2,0) to[short, -o] (4.4,0) node[right] {$B\ (0\,\si{\volt})$};
    \draw (2.2,0) node[ground] {};
    \node[below, font=\footnotesize\bfseries] at (2.2,-0.35) {Cálculo de $V_{th}$ ($A$-$B$ en abierto)};
\end{circuitikz}
\end{minipage}
\end{center}
\begin{itemize}
    \item \textbf{Resistencia de Thévenin}: Anulando la fuente de $10\,\si{\volt}$ (cortocircuito), $R_1$ y $R_2$ quedan en paralelo y en serie con $R_3$:
    \begin{equation*}
        R_{th} = R_3 + (R_1 \parallel R_2) = 1 + \frac{1\cdot 5}{1 + 5} = 1 + \frac{5}{6} = \frac{11}{6}\,\si{\kilo\ohm}
    \end{equation*}
    \item \textbf{Tensión de Thévenin}: Con $A$ y $B$ en abierto no circula corriente por $R_3$ ($I_{R3}=0$), luego por divisor de tensión:
    \begin{equation*}
        V_{th} = V_A - V_B = \frac{R_2}{R_1 + R_2}\cdot 10\,\si{\volt} = \frac{5}{6}\cdot 10 = \frac{50}{6}\,\si{\volt}
    \end{equation*}
    \item \textbf{Conectando $R_4 = 1\,\si{\kilo\ohm}$ al equivalente de Thévenin}:
    \begin{equation*}
        I = \frac{V_{th}}{R_{th} + R_4} = \frac{50/6}{11/6 + 1} = \frac{50}{17}\,\si{\milli\ampere} \implies \boxed{V_{AB} = I\cdot R_4 = \frac{50}{17}\,\si{\volt}} \quad \text{(OK)}
    \end{equation*}
\end{itemize}

\textbf{Método 3: Mediante el Equivalente de Norton}:
\begin{itemize}
    \item Cortocircuitando $A$ y $B$, las resistencias $R_2$ y $R_3$ quedan en paralelo ($R_2 \parallel R_3 = \frac{5\cdot 1}{5+1} = \frac{5}{6}\,\si{\kilo\ohm}$). La tensión en el nudo intermedio $C$ por divisor de tensión es:
    \begin{equation*}
        V_C = \frac{R_2 \parallel R_3}{R_1 + (R_2 \parallel R_3)}\cdot 10\,\si{\volt} = \frac{5/6}{1 + 5/6}\cdot 10 = \frac{50}{11}\,\si{\volt} \implies I_N = I_{AB} = \frac{V_C}{R_3} = \frac{50}{11}\,\si{\milli\ampere}
    \end{equation*}
    \item Al conectar $R_4 = 1\,\si{\kilo\ohm}$ en paralelo con $R_N = \frac{11}{6}\,\si{\kilo\ohm}$, la resistencia total es $R_p = R_N \parallel R_4 = \frac{11/6 \cdot 1}{11/6 + 1} = \frac{11}{17}\,\si{\kilo\ohm}$, de modo que:
    \begin{equation*}
        \boxed{V_{AB} = I_N \cdot R_p = \frac{50}{11}\,\si{\milli\ampere} \cdot \frac{11}{17}\,\si{\kilo\ohm} = \frac{50}{17}\,\si{\volt}} \quad \text{(OK)}
    \end{equation*}
\end{itemize}
\end{ejemplo}

\subsection{Teorema de Máxima Transferencia de Potencia}
Si se conecta una resistencia de carga variable $R_L$ a un circuito cuyo equivalente de Thévenin es $(V_{th}, R_{th})$, la potencia disipada en la carga es $P_L = I^2 R_L = \left(\frac{V_{th}}{R_{th} + R_L}\right)^2 R_L$. Derivando respecto a $R_L$ e igualando a cero ($\frac{dP_L}{dR_L} = 0$) se demuestra que la carga recibe la \textbf{potencia máxima posible} cuando su valor coincide exactamente con la resistencia de Thévenin:
\begin{equation}
    \boxed{R_L = R_{th} \implies P_{L,\max} = \frac{V_{th}^2}{4\,R_{th}}}
\end{equation}

\subsection{Principio de Superposición}

El \textbf{Principio de Superposición} establece que en cualquier circuito lineal con varias fuentes independientes, la tensión o corriente en cualquier elemento es igual a la suma algebraica de las respuestas parciales obtenidas cuando actúa \textbf{cada una de las fuentes independientes por separado}, anulando todas las demás (fuentes de tensión en cortocircuito y fuentes de corriente en circuito abierto).

\begin{ejemplo}[Demostración paso a paso del Principio de Superposición]
Sea el circuito de la figura con una fuente de tensión $V$, una fuente de intensidad $I$ y dos elementos pasivos $Z_1$ y $Z_2$. Halla la tensión en el nudo $A$ por superposición y compruébalo por el método de nudos:

\begin{center}
\begin{minipage}{0.32\textwidth}
\centering
\begin{circuitikz}[scale=0.78, transform shape]
    \draw (0,0) to[V, v<=$V$] (0,2.0)
          to[generic, l=$Z_1$] (2.2,2.0) node[above] {$A$}
          to[generic, l=$Z_2$] (2.2,0) -- (0,0);
    \draw (2.2,2.0) -- (3.8,2.0);
    \draw (3.8,0) to[isource, l_=$I$] (3.8,2.0);
    \draw (2.2,0) -- (3.8,0);
    \draw (2.2,0) node[ground] {};
    \node[below, font=\footnotesize\bfseries] at (1.9,-0.35) {Circuito completo};
\end{circuitikz}
\end{minipage}
\hfill
\begin{minipage}{0.32\textwidth}
\centering
\begin{circuitikz}[scale=0.78, transform shape]
    \draw (0,0) to[V, v<=$V$] (0,2.0)
          to[generic, l=$Z_1$] (2.2,2.0) node[above] {$A$}
          to[generic, l=$Z_2$] (2.2,0) -- (0,0);
    \draw (2.2,0) node[ground] {};
    \node[below, font=\footnotesize\bfseries] at (1.1,-0.35) {a) Solo fuente $V$ ($I=0$)};
\end{circuitikz}
\end{minipage}
\hfill
\begin{minipage}{0.32\textwidth}
\centering
\begin{circuitikz}[scale=0.78, transform shape]
    \draw (0,0) -- (0,2.0)
          to[generic, l=$Z_1$] (2.0,2.0) node[above] {$A$}
          to[generic, l=$Z_2$] (2.0,0) -- (0,0);
    \draw (2.0,2.0) -- (3.5,2.0);
    \draw (3.5,0) to[isource, l_=$I$] (3.5,2.0);
    \draw (2.0,0) -- (3.5,0);
    \draw (2.0,0) node[ground] {};
    \node[below, font=\footnotesize\bfseries] at (1.75,-0.35) {b) Solo fuente $I$ ($V=0$)};
\end{circuitikz}
\end{minipage}
\end{center}

\begin{enumerate}
    \item \textbf{a) Con la fuente de tensión $V$ activa} (anulamos la fuente de corriente dejándola en abierto): Queda un divisor de tensión entre $Z_1$ y $Z_2$:
    \begin{equation*}
        V_{A1} = \frac{Z_2}{Z_1 + Z_2}\,V
    \end{equation*}
    \item \textbf{b) Con la fuente de intensidad $I$ activa} (anulamos la fuente de tensión con un cortocircuito): Los elementos $Z_1$ y $Z_2$ quedan en paralelo alimentados por la corriente $I$:
    \begin{equation*}
        V_{A2} = I\cdot (Z_1 \parallel Z_2) = I\,\frac{Z_1 Z_2}{Z_1 + Z_2}
    \end{equation*}
    \item \textbf{Sumando ambas contribuciones}:
    \begin{equation*}
        \boxed{V_A = V_{A1} + V_{A2} = \frac{Z_2\,V + I\,Z_1 Z_2}{Z_1 + Z_2}}
    \end{equation*}
    \item \textbf{Comprobación directa por nudos con las dos fuentes a la vez}:
    \begin{equation*}
        I + \frac{V - V_A}{Z_1} = \frac{V_A}{Z_2} \implies I Z_1 + V - V_A = \frac{Z_1}{Z_2}V_A \implies \left(1 + \frac{Z_1}{Z_2}\right)V_A = I Z_1 + V \implies V_A = \frac{Z_2\,V + I\,Z_1 Z_2}{Z_1 + Z_2} \quad \text{(OK)}
    \end{equation*}
\end{enumerate}
\end{ejemplo}

\section{Régimen Transitorio de Primer Orden: Carga y Descarga de un Condensador}

Cuando en un circuito con condensadores o bobinas se acciona un conmutador o cambia el valor de una fuente, el circuito pasa por un \textbf{régimen transitorio} hasta alcanzar su nuevo estado estacionario.

\begin{figure}[H]
\centering
\begin{minipage}{0.48\textwidth}
\centering
\begin{circuitikz}[scale=0.88, transform shape]
    \node[font=\small\bfseries, color=NavyBlue] at (2.2,2.65) {Carga del Condensador};
    \draw (0,0) to[V, v<=$V$] (0,2.1)
          to[closing switch, l={$t=0$}] (1.8,2.1)
          to[R, l=$R$, i>^=$i(t)$] (4.0,2.1)
          to[C, l=$C$, v=$v_C(t)$] (4.0,0) -- (0,0);
    \draw (2.0,0) node[ground] {};
\end{circuitikz}
\end{minipage}
\hfill
\begin{minipage}{0.48\textwidth}
\centering
\begin{circuitikz}[scale=0.88, transform shape]
    \node[font=\small\bfseries, color=BrickRed] at (2.0,2.65) {Descarga del Condensador};
    \draw (0,0) -- (0,2.1)
          to[R, l=$R$] (3.2,2.1)
          to[C, l=$C$, v=$v_C(t)$] (3.2,0) -- (0,0);
    \draw (1.6,0) node[ground] {};
    \draw[-{Stealth}, BrickRed, thick] (1.1,0.6) arc (-150:150:0.45) node[midway, right] {$i(t)$};
\end{circuitikz}
\end{minipage}
\caption{Circuitos RC de primer orden durante la carga (izquierda) y la descarga (derecha) del condensador.}
\end{figure}

\subsection{Deducción de la Ecuación de Descarga del Condensador}
Supongamos que el condensador se ha cargado inicialmente a la tensión de la fuente, de modo que en $t = 0$ tenemos $v_C(0) = V$, y en ese instante desconectamos la fuente y conectamos el conjunto $R$-$C$ a tierra:
\begin{enumerate}
    \item Para $t > 0$, la ecuación de malla es:
    \begin{equation}
        0 = v_C(t) + i(t)\cdot R
    \end{equation}
    \item Como la corriente es la misma que pasa por el condensador, $i(t) = i_C(t) = C\,\frac{dv_C(t)}{dt}$, obtenemos una ecuación diferencial ordinaria de primer orden:
    \begin{equation}
        0 = v_C(t) + R C\,\frac{dv_C(t)}{dt} \qquad \text{donde } \boxed{\tau = R\cdot C} \text{ es la \textbf{constante de tiempo} } [\si{\second}]
    \end{equation}
    \item Separando variables e integrando entre el instante inicial $t = 0$ ($v_C(0) = V$) y un instante genérico $t$:
    \begin{equation}
        \tau\,\frac{dv_C(t)}{dt} = -v_C(t) \implies \int_{v_C(0)}^{v_C(t)} \frac{dv_C(t)}{v_C(t)} = -\frac{1}{\tau}\int_0^t dt
    \end{equation}
    \begin{equation}
        \ln\big[v_C(t)\big] - \ln\big[v_C(0)\big] = -\frac{t}{\tau} \implies \ln\left(\frac{v_C(t)}{V}\right) = -\frac{t}{\tau} \implies \boxed{v_C(t) = V\cdot e^{-t/\tau}}
    \end{equation}
\end{enumerate}

Cuanto mayor sea la resistencia $R$ o la capacidad $C$, mayor será la constante de tiempo $\tau = RC$ y más tardará el condensador en descargarse. Se define el \textbf{tiempo de caída} ($t_{\text{fall}}$) como el tiempo que tarda la tensión en caer desde el $90\%$ ($0{,}9\,V$) hasta el $10\%$ ($0{,}1\,V$):
\begin{equation}
    t_{90\%} = -\tau\ln(0{,}9) \approx 0{,}105\,\tau, \quad t_{10\%} = -\tau\ln(0{,}1) \approx 2{,}303\,\tau \implies \boxed{t_{\text{fall}} = t_{10\%} - t_{90\%} \approx 2{,}2\,\tau}
\end{equation}

\subsection{Deducción de la Ecuación de Carga del Condensador}
Si el condensador está inicialmente descargado en $t = 0$ ($v_C(0) = 0$) y lo conectamos a la fuente de tensión continua $V$ a través de la resistencia $R$:
\begin{enumerate}
    \item La ecuación de malla para $t > 0$ es:
    \begin{equation}
        V = i(t)\cdot R + v_C(t) = R C\,\frac{dv_C(t)}{dt} + v_C(t) \implies V - v_C(t) = \tau\,\frac{dv_C(t)}{dt}
    \end{equation}
    \item Separando variables e integrando desde $t = 0$ ($v_C(0) = 0$) hasta $t$:
    \begin{equation}
        \int_0^t -\frac{1}{\tau}\,dt = \int_0^{v_C(t)} \frac{dv_C(t)}{v_C(t) - V} \implies -\frac{t}{\tau} = \ln\left(\frac{v_C(t) - V}{v_C(0) - V}\right) = \ln\left(\frac{v_C(t) - V}{-V}\right)
    \end{equation}
    \item Tomando exponenciales y despejando $v_C(t)$:
    \begin{equation}
        \frac{v_C(t) - V}{-V} = e^{-t/\tau} \implies v_C(t) - V = -V\,e^{-t/\tau} \implies \boxed{v_C(t) = V\,\left(1 - e^{-t/\tau}\right)}
    \end{equation}
    La corriente de carga disminuye exponencialmente: $i(t) = C\frac{dv_C}{dt} = \frac{V}{R}\,e^{-t/\tau}$. Para $t = \tau$, el condensador ha alcanzado el $63{,}2\%$ de su tensión final, y para $t \approx 5\tau$ se considera totalmente cargado ($99{,}3\%$).
\end{enumerate}

\begin{figure}[H]
\centering
\begin{minipage}{0.48\textwidth}
\centering
\begin{tikzpicture}[scale=0.88]
    \draw[-{Stealth}, thick] (0,0) -- (5.2,0) node[right] {$t$};
    \draw[-{Stealth}, thick] (0,0) -- (0,3.2) node[above] {$v_C(t)$};
    \draw[dashed, gray] (0,2.5) node[left, black] {$V$} -- (5.0,2.5);
    \draw[dashed, gray] (0,2.25) node[left, black, font=\scriptsize] {$0{,}9V$} -- (0.35,2.25) -- (0.35,0);
    \draw[dashed, gray] (0,0.25) node[left, black, font=\scriptsize] {$0{,}1V$} -- (3.0,0.25) -- (3.0,0);
    \draw[ultra thick, BrickRed, domain=0:4.8, samples=80] plot (\x, {2.5*exp(-0.77*\x)});
    \node[below, font=\scriptsize] at (0.35,0) {$0{,}1\tau$};
    \node[below, font=\scriptsize] at (3.0,0) {$2{,}3\tau$};
    \node[above right, BrickRed, font=\small\bfseries] at (1.3,1.3) {$v_C(t) = V e^{-t/\tau}$};
    \node[below, font=\small\bfseries] at (2.5,-0.55) {Curva de Descarga ($t_{\text{fall}} \approx 2{,}2\tau$)};
\end{tikzpicture}
\end{minipage}
\hfill
\begin{minipage}{0.48\textwidth}
\centering
\begin{tikzpicture}[scale=0.88]
    \draw[-{Stealth}, thick] (0,0) -- (5.2,0) node[right] {$t$};
    \draw[-{Stealth}, thick] (0,0) -- (0,3.2) node[above] {$v_C(t)$};
    \draw[dashed, gray] (0,2.5) node[left, black] {$V$} -- (5.0,2.5);
    \draw[dashed, gray] (0,1.58) node[left, black, font=\scriptsize] {$0{,}632V$} -- (1.1,1.58) -- (1.1,0) node[below, black, font=\scriptsize] {$\tau = RC$};
    \draw[ultra thick, NavyBlue, domain=0:4.8, samples=80] plot (\x, {2.5*(1 - exp(-0.9*\x))});
    \node[below right, NavyBlue, font=\small\bfseries] at (1.5,1.6) {$v_C(t) = V(1 - e^{-t/\tau})$};
    \node[below, font=\small\bfseries] at (2.5,-0.55) {Curva de Carga ($\tau = RC$)};
\end{tikzpicture}
\end{minipage}
\caption{Evolución temporal de la tensión $v_C(t)$ durante la descarga (izquierda) y la carga (derecha) de un condensador.}
\end{figure}

\begin{ley}[Fórmula Universal de Transitorios de Primer Orden (RC y RL)]
En cualquier circuito de primer orden con un único elemento almacenador (condensador $C$ con $\tau = R_{th}C$, o bobina $L$ con $\tau = \frac{L}{R_{th}}$), la evolución temporal de cualquier tensión o corriente para $t \ge 0^+$ responde a la expresión universal:
\begin{equation}
    \boxed{x(t) = x(\infty) + \big[x(0^+) - x(\infty)\big]\,e^{-t/\tau}, \qquad t \ge 0^+}
\end{equation}
donde $x(0^+)$ se obtiene por continuidad ($v_C(0^+) = v_C(0^-)$ o $i_L(0^+) = i_L(0^-)$) y $x(\infty)$ es el valor en régimen estacionario para $t \to \infty$.
\end{ley}

\section{Ejercicios de Clase y Exámenes Resueltos Paso a Paso}

\begin{ejercicio}[Asociación de resistencias y Teorema de Kennelly (Ejercicio 3)]
Calcula la resistencia equivalente $R_{AB}$ entre los terminales $A$ y $B$ en las tres asociaciones de resistencias de las figuras siguientes (corregidas en clase):

\textbf{Apartado 1}: Red en puente doble con resistencias de $5\,\si{\kilo\ohm}$ y $15\,\si{\kilo\ohm}$.
\begin{center}
\begin{minipage}{0.48\textwidth}
\centering
\begin{circuitikz}[scale=0.78, transform shape]
    \node[circle, fill=black, inner sep=1.5pt, label=left:$A$] (A) at (0,2.2) {};
    \node[circle, fill=black, inner sep=1.5pt, label=above:$C$] (C) at (2.5,2.2) {};
    \node[circle, fill=black, inner sep=1.5pt, label=right:$B$] (B) at (5.0,2.2) {};
    \node[circle, fill=black, inner sep=1.5pt, label=below:$D$] (D) at (1.25,0) {};
    \node[circle, fill=black, inner sep=1.5pt, label=below:$E$] (E) at (3.75,0) {};
    \draw (A) to[R, l={$5\,\si{\kilo\ohm}$}] (C);
    \draw (C) to[R, l={$5\,\si{\kilo\ohm}$}] (B);
    \draw (A) to[R, l_={$5\,\si{\kilo\ohm}$}] (D);
    \draw (C) to[R, l={$15\,\si{\kilo\ohm}$}] (D);
    \draw (C) to[R, l_={$15\,\si{\kilo\ohm}$}] (E);
    \draw (B) to[R, l={$5\,\si{\kilo\ohm}$}] (E);
    \draw (D) to[R, l_={$5\,\si{\kilo\ohm}$}] (E);
    \node[below, font=\footnotesize\bfseries] at (2.5,-0.55) {1) Red original en puente doble};
\end{circuitikz}
\end{minipage}
\hfill
\begin{minipage}{0.48\textwidth}
\centering
\begin{circuitikz}[scale=0.78, transform shape]
    \draw (0,1.2) node[left] {$A$} to[R, l={$1\,\si{\kilo\ohm}$}, o-] (1.8,1.2);
    \draw (1.8,1.2) -- (1.8,2.0) to[R, l={$6\,\si{\kilo\ohm}$}] (4.0,2.0) -- (4.0,1.2);
    \draw (1.8,1.2) -- (1.8,0.4) to[R, l_={$11\,\si{\kilo\ohm}$}] (4.0,0.4) -- (4.0,1.2);
    \draw (4.0,1.2) to[R, l={$1\,\si{\kilo\ohm}$}, -o] (5.8,1.2) node[right] {$B$};
    \node[below, font=\footnotesize\bfseries] at (2.9,-0.35) {2) Red tras aplicar Kennelly};
\end{circuitikz}
\end{minipage}
\end{center}
\textbf{Resolución del Apartado 1}:
\begin{itemize}
    \item Aplicamos el \textbf{Teorema de Kennelly} ($\Delta \to \text{Y}$) a los dos triángulos laterales $ACD$ y $CBE$ (ambos formados por lados de $5\,\si{\kilo\ohm}$, $5\,\si{\kilo\ohm}$ y $15\,\si{\kilo\ohm}$, cuya suma es $5 + 5 + 15 = 25\,\si{\kilo\ohm}$):
    \begin{equation*}
        R_a = \frac{5\cdot 5}{25} = 1\,\si{\kilo\ohm}, \qquad R_c = \frac{5\cdot 15}{25} = 3\,\si{\kilo\ohm}, \qquad R_d = \frac{5\cdot 15}{25} = 3\,\si{\kilo\ohm}
    \end{equation*}
    \item Al sustituir ambos triángulos por sus estrellas equivalentes:
    \begin{itemize}
        \item La rama superior entre los centros de ambas estrellas tiene dos resistencias de $3\,\si{\kilo\ohm}$ en serie: $3 + 3 = 6\,\si{\kilo\ohm}$.
        \item La rama inferior que pasa por $D$-$E$ tiene tres resistencias en serie ($3\,\si{\kilo\ohm} + 5\,\si{\kilo\ohm} + 3\,\si{\kilo\ohm} = 11\,\si{\kilo\ohm}$).
    \end{itemize}
    \item Calculamos el paralelo central $R_p = 6\,\si{\kilo\ohm} \parallel 11\,\si{\kilo\ohm}$ y sumamos las dos resistencias exteriores de $1\,\si{\kilo\ohm}$:
    \begin{equation*}
        R_p = \frac{6\cdot 11}{6 + 11} = \frac{66}{17}\,\si{\kilo\ohm} \implies \boxed{R_T = 1\,\si{\kilo\ohm} + \frac{66}{17}\,\si{\kilo\ohm} + 1\,\si{\kilo\ohm} = \frac{100}{17}\,\si{\kilo\ohm} \approx 5{,}88\,\si{\kilo\ohm}}
    \end{equation*}
\end{itemize}

\textbf{Apartado 2}: En el mismo puente, si los terminales de entrada son los nudos inferiores $A$ y $B$ (con una resistencia de $6\,\si{\kilo\ohm}$ directamente conectada entre $A$ y $B$, y los nudos superiores $C$ y $E$ al aire):
\begin{itemize}
    \item \textbf{Forma 1 (Asociación serie-paralelo directa)}: Los lados exteriores superiores ($5\,\si{\kilo\ohm} + 5\,\si{\kilo\ohm} = 10\,\si{\kilo\ohm}$) quedan en paralelo con las diagonales de $15\,\si{\kilo\ohm}$:
    \begin{equation*}
        10\,\si{\kilo\ohm} \parallel 15\,\si{\kilo\ohm} = \frac{10\cdot 15}{10 + 15} = 6\,\si{\kilo\ohm}
    \end{equation*}
    Las dos ramas resultantes de $6\,\si{\kilo\ohm}$ quedan en serie ($6 + 6 = 12\,\si{\kilo\ohm}$) y en paralelo con la base de $6\,\si{\kilo\ohm}$ entre $A$ y $B$:
    \begin{equation*}
        \boxed{R_{AB} = 12\,\si{\kilo\ohm} \parallel 6\,\si{\kilo\ohm} = \frac{12\cdot 6}{12 + 6} = 4\,\si{\kilo\ohm}}
    \end{equation*}
    \item \textbf{Forma 2 (Por Kennelly)}: Transformando los dos triángulos superiores en estrellas ($1\,\si{\kilo\ohm}, 3\,\si{\kilo\ohm}, 3\,\si{\kilo\ohm}$), las puntas de $1\,\si{\kilo\ohm}$ en $C$ y $E$ quedan al aire (no circula corriente por ellas), quedando cuatro resistencias de $3\,\si{\kilo\ohm}$ en serie ($3+3+3+3 = 12\,\si{\kilo\ohm}$) en paralelo con la base de $6\,\si{\kilo\ohm}$: $R_{AB} = 12\,\si{\kilo\ohm} \parallel 6\,\si{\kilo\ohm} = 4\,\si{\kilo\ohm}$.
\end{itemize}

\textbf{Apartado 4}: Red vertical entre $A$ y $B$ con una resistencia superior de $9\,\si{\kilo\ohm}$ en serie con un puente formado por un triángulo superior ($10\,\si{\kilo\ohm}, 10\,\si{\kilo\ohm}, 5\,\si{\kilo\ohm}$) y dos ramas inferiores de $20\,\si{\kilo\ohm}$ y $9\,\si{\kilo\ohm}$:
\begin{itemize}
    \item Transformamos el triángulo superior ($10\,\si{\kilo\ohm}, 10\,\si{\kilo\ohm}, 5\,\si{\kilo\ohm}$, suma $= 25\,\si{\kilo\ohm}$) en estrella:
    \begin{equation*}
        R_c = \frac{10\cdot 10}{25} = 4\,\si{\kilo\ohm}, \qquad R_d = \frac{10\cdot 5}{25} = 2\,\si{\kilo\ohm}, \qquad R_e = \frac{10\cdot 5}{25} = 2\,\si{\kilo\ohm}
    \end{equation*}
    \item Sumando en serie: arriba queda $9\,\si{\kilo\ohm} + 4\,\si{\kilo\ohm} = 13\,\si{\kilo\ohm}$, y en las dos ramas paralelas inferiores quedan $(2 + 20)\,\si{\kilo\ohm} = 22\,\si{\kilo\ohm}$ en paralelo con $(2 + 9)\,\si{\kilo\ohm} = 11\,\si{\kilo\ohm}$:
    \begin{equation*}
        22\,\si{\kilo\ohm} \parallel 11\,\si{\kilo\ohm} = \frac{22\cdot 11}{33} = \frac{22}{3}\,\si{\kilo\ohm} \approx 7{,}33\,\si{\kilo\ohm} \implies \boxed{R_{AB} = 13 + 7{,}33 = 20{,}33\,\si{\kilo\ohm}}
    \end{equation*}
\end{itemize}
\end{ejercicio}

\begin{ejercicio}[Cálculo de corrientes, potencias y balance energético (Ejercicio 9)]
Calcula la corriente en cada rama y la potencia disipada por cada resistencia en el circuito de la figura, sabiendo que $V = 12\,\si{\volt}$, $R_1 = 2\,\si{\kilo\ohm}$ y $R_2 = R_3 = R_4 = 6\,\si{\kilo\ohm}$. Comprueba el balance de potencias.

\begin{center}
\begin{minipage}{0.54\textwidth}
\centering
\begin{circuitikz}[scale=0.85, transform shape]
    \draw (0,0) to[battery1, l=$V$] (0,2.2)
          to[R, l=$R_1$, i>^=$i_1$] (2.4,2.2) -- (5.6,2.2);
    \node[above] at (3.8,2.2) {$A$};
    \node[below] at (3.8,0) {$B$};
    \draw (2.4,2.2) to[R, l=$R_2$, i>^=$i_2$] (2.4,0);
    \draw (4.0,2.2) to[R, l=$R_3$, i>^=$i_3$] (4.0,0);
    \draw (5.6,2.2) to[R, l=$R_4$, i>^=$i_4$] (5.6,0);
    \draw (0,0) -- (5.6,0);
\end{circuitikz}
\end{minipage}
\hfill
\begin{minipage}{0.42\textwidth}
\centering
\begin{circuitikz}[scale=0.85, transform shape]
    \draw (0,0) to[battery1, l=$V$] (0,2.2)
          to[R, l=$R_1$, i>^=$i_1$] (2.6,2.2) node[above] {$A$}
          to[R, l=$R_{eq}$] (2.6,0) node[below] {$B$} -- (0,0);
\end{circuitikz}
\end{minipage}
\end{center}

\textbf{Resolución paso a paso}:
\begin{enumerate}
    \item \textbf{Asociación de las tres resistencias en paralelo} ($R_2, R_3, R_4$):
    \begin{equation*}
        \frac{1}{R_{eq}} = \frac{1}{R_2} + \frac{1}{R_3} + \frac{1}{R_4} = \frac{1}{6} + \frac{1}{6} + \frac{1}{6} = \frac{3}{6} = \frac{1}{2}\,\si{\kilo\ohm^{-1}} \implies \boxed{R_{eq} = 2\,\si{\kilo\ohm}}
    \end{equation*}
    \item \textbf{Cálculo de la corriente principal $i_1$ y de las corrientes de rama}:
    \begin{equation*}
        i_1 = \frac{V}{R_1 + R_{eq}} = \frac{12\,\si{\volt}}{2\,\si{\kilo\ohm} + 2\,\si{\kilo\ohm}} = \frac{12\,\si{\volt}}{4\,\si{\kilo\ohm}} = 3\,\si{\milli\ampere}
    \end{equation*}
    La tensión entre $A$ y $B$ es $V_{AB} = i_1 \cdot R_{eq} = 3\,\si{\milli\ampere} \cdot 2\,\si{\kilo\ohm} = 6\,\si{\volt}$. Como $R_2 = R_3 = R_4 = 6\,\si{\kilo\ohm}$:
    \begin{equation*}
        i_2 = i_3 = i_4 = \frac{V_{AB}}{R_2} = \frac{6\,\si{\volt}}{6\,\si{\kilo\ohm}} = 1\,\si{\milli\ampere} \qquad (\text{se cumple LCK: } 3\,\si{\milli\ampere} = 1 + 1 + 1\,\si{\milli\ampere})
    \end{equation*}
    \item \textbf{Potencias disipadas y balance de potencia}:
    \begin{align*}
        P_2 = P_3 = P_4 &= i_2 \cdot V_{AB} = 1\,\si{\milli\ampere} \cdot 6\,\si{\volt} = 6\,\si{\milli\watt} \\
        P_1 &= i_1 \cdot V_{R1} = 3\,\si{\milli\ampere} \cdot 6\,\si{\volt} = 18\,\si{\milli\watt} \\
        P_{\text{disipada}} &= 3\cdot 6\,\si{\milli\watt} + 18\,\si{\milli\watt} = 36\,\si{\milli\watt} = P_{\text{fuente}} = V\cdot i_1 = 12\,\si{\volt}\cdot 3\,\si{\milli\ampere} = \boxed{36\,\si{\milli\watt}}
    \end{align*}
\end{enumerate}
\end{ejercicio}

\begin{ejercicio}[Análisis completo por mallas y balance de potencias (Ejercicio 10)]
Calcula la corriente en cada rama y la potencia disipada por cada resistencia del siguiente circuito con $V_1 = 7\,\si{\volt}$, $V_2 = 5\,\si{\volt}$, $R_1 = 3\,\si{\kilo\ohm}$, $R_2 = 1\,\si{\kilo\ohm}$ y $R_3 = 2\,\si{\kilo\ohm}$:

\begin{center}
\begin{circuitikz}[scale=0.92, transform shape]
    \draw (0,2.2) to[V, v=$V_1$] (0,0)
          to[R, l=$R_3$] (2.6,0)
          to[R, l=$R_1$] (2.6,1.1)
          to[V, v<=$V_2$] (2.6,2.2) -- (0,2.2);
    \draw (2.6,2.2) -- (5.0,2.2)
          to[R, l=$R_2$] (5.0,0) -- (2.6,0);
    \draw[-{Stealth}, NavyBlue, thick] (1.3,1.5) arc (30:330:0.45) node[midway, right] {$i_1$};
    \draw[-{Stealth}, BrickRed, thick] (3.8,1.5) arc (30:330:0.45) node[midway, right] {$i_2$};
\end{circuitikz}
\end{center}

\textbf{Resolución paso a paso}:
\begin{enumerate}
    \item \textbf{Establecemos el sentido antihorario para las corrientes de malla $i_1$ e $i_2$} (tal como indica la figura de clase) y planteamos las ecuaciones de malla:
    \begin{align*}
        V_1 + V_2 &= i_1 R_3 + (i_1 + i_2) R_1 \implies V_1 + V_2 = (R_1 + R_3)i_1 + R_1 i_2 \implies 12 = 5 i_1 + 3 i_2 \\
        V_2 &= (i_1 + i_2) R_1 + i_2 R_2 \implies V_2 = R_1 i_1 + (R_1 + R_2)i_2 \implies 5 = 3 i_1 + 4 i_2
    \end{align*}
    \item \textbf{Resolución matricial por Cramer}:
    \begin{equation*}
        \begin{pmatrix} 12 \\ 5 \end{pmatrix} = \begin{pmatrix} 5 & 3 \\ 3 & 4 \end{pmatrix} \begin{pmatrix} i_1 \\ i_2 \end{pmatrix} \implies
        i_1 = \frac{\begin{vmatrix} 12 & 3 \\ 5 & 4 \end{vmatrix}}{\begin{vmatrix} 5 & 3 \\ 3 & 4 \end{vmatrix}} = \frac{48 - 15}{20 - 9} = \frac{33}{11} = \boxed{3\,\si{\milli\ampere}}, \qquad
        i_2 = \frac{\begin{vmatrix} 5 & 12 \\ 3 & 5 \end{vmatrix}}{11} = \frac{25 - 36}{11} = \boxed{-1\,\si{\milli\ampere}}
    \end{equation*}
    \item \textbf{Corrientes por cada resistencia}:
    \begin{equation*}
        i_{R1} = i_1 + i_2 = 3 - 1 = 2\,\si{\milli\ampere}, \qquad i_{R2} = i_2 = -1\,\si{\milli\ampere}, \qquad i_{R3} = i_1 = 3\,\si{\milli\ampere}
    \end{equation*}
    \item \textbf{Cálculo de potencias en resistencias y fuentes}:
    \begin{align*}
        P_{R1} &= i_{R1}^2 R_1 = (2\,\si{\milli\ampere})^2 \cdot 3\,\si{\kilo\ohm} = 12\,\si{\milli\watt} \\
        P_{R2} &= i_{R2}^2 R_2 = (-1\,\si{\milli\ampere})^2 \cdot 1\,\si{\kilo\ohm} = 1\,\si{\milli\watt} \\
        P_{R3} &= i_{R3}^2 R_3 = (3\,\si{\milli\ampere})^2 \cdot 2\,\si{\kilo\ohm} = 18\,\si{\milli\watt} \implies P_{\text{total},R} = 12 + 1 + 18 = \boxed{31\,\si{\milli\watt}}
    \end{align*}
    \begin{equation*}
        P_{V1} = V_1 \cdot i_1 = 7\,\si{\volt} \cdot 3\,\si{\milli\ampere} = 21\,\si{\milli\watt}, \quad P_{V2} = V_2 \cdot i_{R1} = 5\,\si{\volt} \cdot 2\,\si{\milli\ampere} = 10\,\si{\milli\watt} \implies P_{\text{total},V} = \boxed{31\,\si{\milli\watt}}
    \end{equation*}
\end{enumerate}
\end{ejercicio}

\begin{ejercicio}[Resolución por nudos con fuente de intensidad y fuentes de tensión (Ejercicio 12)]
Calcula el valor del potencial $V_3$, las corrientes de rama y el balance de potencias en el circuito de la figura, con $R_1 = R_2 = R_3 = 1\,\si{\kilo\ohm}$, $I_1 = 2\,\si{\milli\ampere}$, $V_1 = 1\,\si{\volt}$ y $V_2 = 2\,\si{\volt}$:

\begin{center}
\begin{circuitikz}[scale=0.92, transform shape]
    \draw (0,0) to[V, v<=$V_1$] (0,2.6) node[above left] {$A\ (1\,\si{\volt})$}
          to[R, l=$R_1$, i>^=$i_1$] (2.0,1.2);
    \draw (4.0,0) to[V, v=$V_2$] (4.0,2.6) node[above right] {$B\ (2\,\si{\volt})$}
          to[R, l_=$R_2$, i<_=$i_2$] (2.0,1.2);
    \draw (4.0,2.6) to[isource, l_=$I_1$] (0,2.6);
    \node[above] at (2.0,1.35) {$V_3$};
    \draw (2.0,1.2) to[R, l=$R_3$, i>^=$i_3$] (2.0,0);
    \draw (0,0) -- (4.0,0);
    \draw (2.0,0) node[ground] {} node[below right] {$C\ (0\,\si{\volt})$};
\end{circuitikz}
\end{center}

\textbf{Resolución paso a paso}:
\begin{enumerate}
    \item \textbf{Referenciamos el nudo inferior $C$ a tierra ($V_C = 0\,\si{\volt}$)}: Automáticamente las fuentes de tensión fijan los potenciales de los nudos superiores: $V_A = 1\,\si{\volt}$ y $V_B = 2\,\si{\volt}$.
    \item \textbf{Aplicamos la Ley de Kirchhoff de corrientes (LCK) en el nudo central $V_3$}:
    \begin{equation*}
        i_1 + i_2 = i_3 \implies \frac{V_A - V_3}{R_1} + \frac{V_B - V_3}{R_2} = \frac{V_3 - 0}{R_3} \implies \frac{1 - V_3}{1} + \frac{2 - V_3}{1} = \frac{V_3}{1} \implies 3 = 3 V_3 \implies \boxed{V_3 = 1\,\si{\volt}}
    \end{equation*}
    \item \textbf{Cálculo de las intensidades en las resistencias y en las fuentes}:
    \begin{equation*}
        i_1 = \frac{1 - 1}{1\,\si{\kilo\ohm}} = 0\,\si{\milli\ampere}, \qquad i_2 = \frac{2 - 1}{1\,\si{\kilo\ohm}} = 1\,\si{\milli\ampere}, \qquad i_3 = i_1 + i_2 = 1\,\si{\milli\ampere}
    \end{equation*}
    En el nudo $A$, como por $R_1$ van $0\,\si{\milli\ampere}$ y de la fuente $I_1$ llegan $2\,\si{\milli\ampere}$ hacia la izquierda, toda esa corriente baja entrando por el polo $+$ de $V_1$ ($I_{V1} = 2\,\si{\milli\ampere}$ hacia abajo). En el nudo $B$, salen $2\,\si{\milli\ampere}$ hacia $I_1$ y $1\,\si{\milli\ampere}$ hacia $R_2$, luego la fuente $V_2$ entrega $I_{V2} = 3\,\si{\milli\ampere}$ hacia arriba.
    \item \textbf{Balance de potencias}:
    \begin{itemize}
        \item En las resistencias: $P_{R1} = 0\,\si{\milli\watt}$, $P_{R2} = 1\,\si{\volt}\cdot 1\,\si{\milli\ampere} = 1\,\si{\milli\watt}$, $P_{R3} = 1\,\si{\volt}\cdot 1\,\si{\milli\ampere} = 1\,\si{\milli\watt} \implies P_{\text{total},R} = \boxed{2\,\si{\milli\watt}}$.
        \item En las fuentes: $P_{V1} = -1\,\si{\volt}\cdot 2\,\si{\milli\ampere} = -2\,\si{\milli\watt}$ (absorbe), $P_{V2} = 2\,\si{\volt}\cdot 3\,\si{\milli\ampere} = 6\,\si{\milli\watt}$ (cede), $P_{I1} = (V_A - V_B)\cdot I_1 = (1 - 2)\,\si{\volt}\cdot 2\,\si{\milli\ampere} = -2\,\si{\milli\watt}$ (absorbe) $\implies P_{\text{total,fuentes}} = -2 + 6 - 2 = \boxed{2\,\si{\milli\watt}}$.
    \end{itemize}
\end{enumerate}
\end{ejercicio}

\begin{ejercicio}[Análisis de carga y equivalente de Thévenin/Norton (Ejercicio 7)]
En el circuito de la figura se conocen $R_1 = 6\,\si{\kilo\ohm}$, $R_2 = 12\,\si{\kilo\ohm}$, $R_3 = R_4 = R_5 = R_L = 10\,\si{\kilo\ohm}$, $I_2 = 0{,}5\,\si{\milli\ampere}$, $V_1 = 10\,\si{\volt}$ y $V_2 = 20\,\si{\volt}$:
\begin{enumerate}
    \item[a)] Calcula la corriente que atraviesa la resistencia de carga $R_L$, así como la caída de tensión entre sus extremos $A$ y $B$.
    \item[b)] Calcula el equivalente Thévenin y Norton del circuito entre los puntos $A$ y $B$, siendo $R_L$ la carga del circuito, y comprueba el resultado del apartado anterior.
\end{enumerate}

\begin{center}
\begin{circuitikz}[scale=0.88, transform shape]
    \draw (0,1.4) to[V, v<=$V_1$] (0,3.2)
          to[R, l=$R_1$] (3.6,3.2) -- (3.6,1.4);
    \draw (0,1.4) node[left] {$V_B$} to[R, l=$R_2$] (1.8,1.4) node[above] {$B$}
          to[R, l=$R_L$] (3.6,1.4) node[above right] {$A$};
    \draw (0,1.4) to[R, l=$R_3$] (0,0);
    \draw (3.6,0) to[isource, l=$I_2$] (0,0);
    \draw (3.6,1.4) to[R, l=$R_4$] (3.6,0);
    \draw (3.6,1.4) to[V, v<=$V_2$] (6.0,1.4)
          to[R, l=$R_5$] (6.0,0) -- (3.6,0);
    \draw[-{Stealth}, NavyBlue, thick] (1.4,2.3) arc (150:-150:0.38) node[midway, left] {$I_1$};
    \draw[-{Stealth}, BrickRed, thick] (4.6,0.7) arc (150:-150:0.38) node[midway, left] {$I_3$};
\end{circuitikz}
\end{center}

\textbf{Resolución paso a paso}:
\begin{enumerate}
    \item \textbf{Apartado a)}: Como la fuente de corriente inferior fija directamente la corriente de la malla inferior izquierda en $I_2 = 0{,}5\,\si{\milli\ampere}$ (en sentido horario), para hallar $I_1$ basta plantear la ecuación de la malla superior izquierda:
    \begin{equation*}
        V_1 = I_1 R_1 + (I_1 - I_2) R_L + (I_1 - I_2) R_2 = (R_1 + R_2 + R_L) I_1 - (R_2 + R_L) I_2
    \end{equation*}
    Sustituyendo los valores ($R_1 + R_2 + R_L = 6 + 12 + 10 = 28\,\si{\kilo\ohm}$ y $R_2 + R_L = 22\,\si{\kilo\ohm}$):
    \begin{equation*}
        10 = 28 I_1 - 22 \cdot 0{,}5 \implies 28 I_1 = 21 \implies I_1 = \frac{21}{28} = \frac{3}{4}\,\si{\milli\ampere} = 0{,}75\,\si{\milli\ampere}
    \end{equation*}
    La intensidad que atraviesa $R_L$ (de $A$ hacia $B$) y la caída de tensión $V_A - V_B$ son:
    \begin{equation*}
        \boxed{I_{RL} = I_1 - I_2 = 0{,}75 - 0{,}50 = 0{,}25\,\si{\milli\ampere}}, \qquad \boxed{V_A - V_B = I_{RL}\cdot R_L = 0{,}25\,\si{\milli\ampere}\cdot 10\,\si{\kilo\ohm} = 2{,}5\,\si{\volt}}
    \end{equation*}
    \item \textbf{Apartado b) Equivalente de Thévenin entre $A$ y $B$ (desconectando $R_L$)}:
    \begin{itemize}
        \item \textbf{Tensión de Thévenin ($V_{th}$)}: Al abrir el tramo $A$-$B$, por la resistencia $R_2$ no circula corriente ($I_{R2} = 0 \implies V_{\text{izq}} = V_B$), y toda la corriente $I_2 = 0{,}5\,\si{\milli\ampere}$ se ve obligada a subir por la malla superior pasando por $R_1$ desde $V_B + 10\,\si{\volt}$ hacia $A$:
        \begin{equation*}
            V_A = (V_B + 10\,\si{\volt}) - I_2 \cdot R_1 = V_B + 10 - 0{,}5\,\si{\milli\ampere}\cdot 6\,\si{\kilo\ohm} = V_B + 7\,\si{\volt} \implies \boxed{V_{th} = V_A - V_B = 7\,\si{\volt}}
        \end{equation*}
        \item \textbf{Resistencia de Thévenin ($R_{th}$)}: Anulamos las fuentes ($V_1$ y $V_2$ en cortocircuito, $I_2$ en circuito abierto). Al quedar abierta la rama inferior de $I_2$, las resistencias $R_3, R_4, R_5$ no cierran camino entre $A$ y $B$, quedando únicamente $R_1$ y $R_2$ en serie:
        \begin{equation*}
            \boxed{R_{th} = R_1 + R_2 = 6\,\si{\kilo\ohm} + 12\,\si{\kilo\ohm} = 18\,\si{\kilo\ohm}} \qquad \left(I_N = \frac{V_{th}}{R_{th}} = \frac{7}{18}\,\si{\milli\ampere}\right)
        \end{equation*}
        \item \textbf{Comprobación conectando $R_L = 10\,\si{\kilo\ohm}$ al equivalente de Thévenin}:
        \begin{equation*}
            I_{RL} = \frac{V_{th}}{R_{th} + R_L} = \frac{7\,\si{\volt}}{18\,\si{\kilo\ohm} + 10\,\si{\kilo\ohm}} = 0{,}25\,\si{\milli\ampere}, \qquad V_{AB} = 0{,}25\,\si{\milli\ampere}\cdot 10\,\si{\kilo\ohm} = 2{,}5\,\si{\volt} \quad \text{(OK)}
        \end{equation*}
    \end{itemize}
\end{enumerate}
\end{ejercicio}

\begin{ejercicio}[Equivalente de Thévenin en función de un parámetro $V_i$ (Ejercicio 8)]
Calcula el circuito equivalente de Thévenin entre los terminales $A$ y $B$ del siguiente circuito en función de la tensión $V_i$, sabiendo que $I = 1\,\si{\milli\ampere}$, $R_1 = R_2 = R_3 = 1\,\si{\kilo\ohm}$ y $R_4 = R_5 = 2\,\si{\kilo\ohm}$:

\begin{center}
\begin{minipage}{0.48\textwidth}
\centering
\begin{circuitikz}[scale=0.82, transform shape]
    \draw (0,1.2) node[left] {$C$} to[battery1, l=$V_i$] (0,2.6)
          to[R, l=$R_1$] (3.2,2.6) -- (3.2,1.2);
    \draw (0,1.2) to[R, l=$R_2$, -o] (1.8,1.2) node[above] {$B$};
    \draw (3.2,1.2) to[short, -o] (2.4,1.2) node[above] {$A$};
    \draw (0,1.2) to[R, l=$R_3$] (0,0);
    \draw (3.2,0) to[isource, l=$I$] (0,0);
    \draw (3.2,1.2) to[R, l=$R_4$] (3.2,0);
    \draw (3.2,1.2) -- (4.6,1.2) to[R, l=$R_5$] (4.6,0) -- (3.2,0);
    \node[below, font=\footnotesize\bfseries] at (2.3,-0.35) {Circuito original};
\end{circuitikz}
\end{minipage}
\hfill
\begin{minipage}{0.48\textwidth}
\centering
\begin{circuitikz}[scale=0.82, transform shape]
    \draw (0,0) to[battery1, l={$V_{th}=V_i-1\,\si{\volt}$}] (0,2.0)
          to[R, l={$R_{th}=2\,\si{\kilo\ohm}$}, -o] (3.2,2.0) node[right] {$A$};
    \draw (0,0) to[short, -o] (3.2,0) node[right] {$B$};
    \node[below, font=\footnotesize\bfseries] at (1.6,-0.35) {Equivalente Thévenin resultante};
\end{circuitikz}
\end{minipage}
\end{center}

\textbf{Resolución paso a paso}:
\begin{enumerate}
    \item \textbf{Cálculo de la tensión de Thévenin ($V_{th} = V_A - V_B$)}:
    \begin{itemize}
        \item Como los terminales $A$ y $B$ están en circuito abierto, por la resistencia $R_2$ no circula corriente ($I_{R2} = 0$), de modo que el potencial en $B$ es idéntico al potencial en $C$: $V_B = V_C \implies V_{th} = V_A - V_C$.
        \item Agrupando $R_4 \parallel R_5 = 2\,\si{\kilo\ohm} \parallel 2\,\si{\kilo\ohm} = 1\,\si{\kilo\ohm}$, la fuente de intensidad inferior $I = 1\,\si{\milli\ampere}$ obliga a que por toda la malla exterior ($R_3, V_i, R_1$ y $R_4\parallel R_5$) circule exactamente la corriente $I = 1\,\si{\milli\ampere}$ en sentido horario.
        \item Por tanto, partiendo del nudo $C$, subiendo por la fuente $V_i$ y restando la caída en $R_1$:
        \begin{equation*}
            V_A - V_C = V_i - I\cdot R_1 = V_i - (1\,\si{\milli\ampere})(1\,\si{\kilo\ohm}) \implies \boxed{V_{th} = V_i - 1\,\si{\volt}}
        \end{equation*}
    \end{itemize}
    \item \textbf{Cálculo de la resistencia de Thévenin ($R_{th}$)}:
    \begin{itemize}
        \item Anulando la fuente de tensión $V_i$ (cortocircuito) y la fuente de intensidad $I$ (circuito abierto), las ramas inferiores quedan abiertas y vistas desde $A$ y $B$ solo quedan en serie las resistencias $R_1$ y $R_2$:
        \begin{equation*}
            \boxed{R_{th} = R_1 + R_2 = 1\,\si{\kilo\ohm} + 1\,\si{\kilo\ohm} = 2\,\si{\kilo\ohm}}
        \end{equation*}
    \end{itemize}
\end{enumerate}
\end{ejercicio}

\begin{ejercicio}[Equivalente de Thévenin mediante transformación de fuentes (Ejercicio 14)]
Para el circuito de la figura ($V_1 = 1\,\si{\volt}, V_2 = 3\,\si{\volt}, I_1 = 1\,\si{\milli\ampere}, R_1 = 1\,\si{\kilo\ohm}$):
\begin{enumerate}
    \item[a)] Calcula el equivalente de Thévenin entre los terminales $A$ y $B$.
    \item[b)] Calcula la caída de tensión $V_{CB}$ entre los nudos $C$ y $B$.
    \item[c)] ¿Qué potencia se disipa en la resistencia $R_1$?
\end{enumerate}

\begin{center}
\begin{minipage}{0.48\textwidth}
\centering
\begin{circuitikz}[scale=0.82, transform shape]
    \draw (0,0) to[battery1, l={$V_2=3\,\si{\volt}$}] (0,1.2) node[left] {$A$}
          to[battery1, l={$V_1=1\,\si{\volt}$}] (0,2.5)
          to[R, l={$1\,\si{\kilo\ohm}$}] (2.4,2.5) node[above] {$C$}
          to[R, l={$R_1=1\,\si{\kilo\ohm}$}] (2.4,1.2) node[left] {$B$};
    \draw (2.4,2.5) -- (4.2,2.5);
    \draw (4.2,1.2) to[isource, l_={$I_1=1\,\si{\milli\ampere}$}] (4.2,2.5);
    \draw (2.4,1.2) -- (4.2,1.2) to[R, l={$2\,\si{\kilo\ohm}$}] (4.2,0) -- (0,0);
    \node[below, font=\footnotesize\bfseries] at (2.1,-0.35) {Circuito original};
\end{circuitikz}
\end{minipage}
\hfill
\begin{minipage}{0.48\textwidth}
\centering
\begin{circuitikz}[scale=0.82, transform shape]
    \draw (0,0) to[battery1, l={$V_2=3\,\si{\volt}$}] (0,1.2) node[left] {$A$}
          to[battery1, l={$V_1=1\,\si{\volt}$}] (0,2.5)
          to[R, l={$1\,\si{\kilo\ohm}$}] (2.0,2.5) node[above] {$C$}
          to[R, l={$1\,\si{\kilo\ohm}$}] (3.8,2.5)
          to[battery1, l={$V_3=1\,\si{\volt}$}] (3.8,1.2) node[right] {$B$}
          to[R, l={$2\,\si{\kilo\ohm}$}] (3.8,0) -- (0,0);
    \draw (1.9,0) node[ground] {};
    \draw[-{Stealth}, NavyBlue, thick] (1.5,0.8) arc (150:-150:0.4) node[midway, left] {$I$};
    \node[below, font=\footnotesize\bfseries] at (1.9,-0.45) {Tras transformar $(I_1 \parallel R_1)$};
\end{circuitikz}
\end{minipage}
\end{center}

\textbf{Resolución paso a paso}:
\begin{enumerate}
    \item \textbf{Apartado a)}: Transformamos el paralelo formado por la fuente de corriente $I_1 = 1\,\si{\milli\ampere}$ y la resistencia $R_1 = 1\,\si{\kilo\ohm}$ entre $C$ y $B$ en su fuente de tensión equivalente en serie: $V_3 = I_1 \cdot R_1 = 1\,\si{\milli\ampere}\cdot 1\,\si{\kilo\ohm} = 1\,\si{\volt}$ en serie con $1\,\si{\kilo\ohm}$.
    \begin{itemize}
        \item Queda una única malla con resistencias en serie ($1 + 1 + 2 = 4\,\si{\kilo\ohm}$):
        \begin{equation*}
            V_1 + V_2 - V_3 = I\,(1 + 1 + 2)\,\si{\kilo\ohm} \implies I = \frac{1 + 3 - 1}{4\,\si{\kilo\ohm}} = \frac{3}{4}\,\si{\milli\ampere} = 0{,}75\,\si{\milli\ampere}
        \end{equation*}
        \item Referenciando el cable inferior a tierra ($0\,\si{\volt}$):
        \begin{equation*}
            V_A = 3\,\si{\volt}, \quad V_B = I\cdot 2\,\si{\kilo\ohm} = \frac{3}{4}\,\si{\milli\ampere}\cdot 2\,\si{\kilo\ohm} = \frac{3}{2}\,\si{\volt} \implies \boxed{V_{th} = V_A - V_B = 3 - \frac{3}{2} = \frac{3}{2}\,\si{\volt} = 1{,}5\,\si{\volt}}
        \end{equation*}
        \item Anulando las fuentes en el circuito original ($V_1, V_2$ en cortocircuito e $I_1$ en abierto), entre $A$ y $B$ quedan la rama superior ($1\,\si{\kilo\ohm} + 1\,\si{\kilo\ohm} = 2\,\si{\kilo\ohm}$) en paralelo con la rama inferior de $2\,\si{\kilo\ohm}$:
        \begin{equation*}
            \boxed{R_{th} = 2\,\si{\kilo\ohm} \parallel 2\,\si{\kilo\ohm} = 1\,\si{\kilo\ohm}}
        \end{equation*}
    \end{itemize}
    \item \textbf{Apartado b) Caída de tensión $V_{CB}$}: En la rama transformada entre $C$ y $B$:
    \begin{equation*}
        \boxed{V_{CB} = V_3 + I\cdot 1\,\si{\kilo\ohm} = 1\,\si{\volt} + \frac{3}{4}\,\si{\milli\ampere}\cdot 1\,\si{\kilo\ohm} = \frac{7}{4}\,\si{\volt} = 1{,}75\,\si{\volt}}
    \end{equation*}
    \item \textbf{Apartado c) Potencia disipada en la resistencia $R_1$ original}: Como $R_1 = 1\,\si{\kilo\ohm}$ está conectada directamente entre $C$ y $B$, soporta la tensión $V_{CB} = \frac{7}{4}\,\si{\volt}$ (por ella bajan $\frac{7}{4}\,\si{\milli\ampere}$):
    \begin{equation*}
        \boxed{P_{R1} = \frac{V_{CB}^2}{R_1} = \frac{(7/4\,\si{\volt})^2}{1\,\si{\kilo\ohm}} = \frac{49}{16}\,\si{\milli\watt} \approx 3{,}06\,\si{\milli\watt}} \quad \left(\text{o en la equivalente de la malla: } V_{CB}\cdot I_1 = \frac{7}{4}\,\si{\milli\watt}\right)
    \end{equation*}
\end{enumerate}
\end{ejercicio}

\begin{ejercicio}[Divisor capacitivo en corriente continua (Ejercicio 17)]
Calcula el potencial $V_2$ en régimen permanente de corriente continua en el circuito de la figura, sabiendo que $V_1 = 10\,\si{\volt}$, $C_1 = 1\,\si{\nano\farad}$ y $C_2 = C_3 = 10\,\si{\nano\farad}$:

\begin{center}
\begin{minipage}{0.54\textwidth}
\centering
\begin{circuitikz}[scale=0.78, transform shape]
    \draw (0,0) to[battery1, l=$V_1$] (0,2.2)
          to[R, l=$R_1$] (1.5,2.2)
          to[C, l=$C_1$] (2.8,2.2)
          to[R, l=$R_2$] (4.2,2.2) node[above] {$V_2$}
          to[R, l=$R_3$] (5.6,2.2)
          to[C, l=$C_3$] (7.0,2.2) node[above] {$V_3$}
          to[R, l=$R_4$] (7.0,0) -- (0,0);
    \draw (4.2,2.2) to[C, l=$C_2$] (4.2,1.1)
          to[R, l=$R_5$] (4.2,0);
    \node[below, font=\footnotesize\bfseries] at (3.5,-0.35) {Circuito RC mixto};
\end{circuitikz}
\end{minipage}
\hfill
\begin{minipage}{0.42\textwidth}
\centering
\begin{circuitikz}[scale=0.78, transform shape]
    \draw (0,0) to[battery1, l=$10\,\si{\volt}$] (0,2.4) -- (2.2,2.4)
          to[C, l=$C_1$] (2.2,1.2) node[right] {$V_2$}
          to[C, l={$C_p=20\,\si{\nano\farad}$}] (2.2,0) -- (0,0);
    \draw (2.2,0) node[ground] {};
    \node[below, font=\footnotesize\bfseries] at (1.1,-0.45) {Equivalente capacitivo en DC};
\end{circuitikz}
\end{minipage}
\end{center}

\textbf{Resolución paso a paso}:
\begin{enumerate}
    \item En corriente continua (DC), una vez cargados los condensadores \textbf{no circula corriente por ninguna rama} ($I = 0$). Al ser $I = 0$, no hay caída de tensión en ninguna de las resistencias ($V_R = I\cdot R = 0$), por lo que podemos eliminar todas las resistencias (sustituirlas por cables ideales) para ver cómo se reparten los $10\,\si{\volt}$ entre los condensadores.
    \item Al hacerlo, los condensadores $C_2$ y $C_3$ quedan conectados en \textbf{paralelo} entre el nudo $V_2$ y tierra ($0\,\si{\volt}$):
    \begin{equation*}
        C_p = C_2 + C_3 = 10\,\si{\nano\farad} + 10\,\si{\nano\farad} = 20\,\si{\nano\farad}
    \end{equation*}
    \item Este condensador equivalente $C_p = 20\,\si{\nano\farad}$ queda en \textbf{serie} con $C_1 = 1\,\si{\nano\farad}$. Como dos condensadores en serie almacenan exactamente la \textbf{misma carga $q$}:
    \begin{equation*}
        10\,\si{\volt} = V_{C1} + V_2 = \frac{q}{C_1} + \frac{q}{C_p} = \frac{C_1 + C_p}{C_1 C_p}\,q \implies q = \frac{C_1 C_p}{C_1 + C_p}\cdot 10\,\si{\volt} = \frac{1\cdot 20}{1 + 20}\cdot 10 = \frac{200}{21}\,\si{\nano\coulomb}
    \end{equation*}
    \item Conocida la carga $q$, la tensión en el nudo $V_2$ es:
    \begin{equation*}
        \boxed{V_2 = \frac{q}{C_p} = \frac{200/21\,\si{\nano\coulomb}}{20\,\si{\nano\farad}} = \frac{10}{21}\,\si{\volt} \approx 0{,}476\,\si{\volt}}
    \end{equation*}
\end{enumerate}
\end{ejercicio}

\begin{ejercicio}[Circuitos con fuentes dependientes (Ejercicio 19, apartados a, b y c)]
Calcula la tensión de salida $V_o$ en cada uno de los tres circuitos con fuentes controladas de la figura (recuerda expresar las resistencias en ohmios cuando la ganancia de una fuente controlada por corriente viene en $\si{\volt/\ampere}$):

\begin{center}
\begin{minipage}{0.32\textwidth}
\centering
\begin{circuitikz}[scale=0.72, transform shape]
    \draw (0,0) to[battery1, l=$12\,\si{\volt}$] (0,2.2)
          to[R, l={$3\,\si{\kilo\ohm}$}, i>^=$I_s$] (2.0,2.2)
          to[cV, v<={$2000 I_s$}] (4.0,2.2)
          to[R, l={$5\,\si{\kilo\ohm}$}, v=$V_o$] (4.0,0) -- (0,0);
    \node[below, font=\footnotesize\bfseries] at (2.0,-0.35) {Apartado a)};
\end{circuitikz}
\end{minipage}
\hfill
\begin{minipage}{0.32\textwidth}
\centering
\begin{circuitikz}[scale=0.72, transform shape]
    \draw (0,0) to[battery1, l=$12\,\si{\volt}$] (0,2.2)
          to[R, l={$2\,\si{\kilo\ohm}$}, i>^=$I$] (2.0,2.2)
          to[cV, v={$2 V_o$}] (4.0,2.2)
          to[R, l={$2\,\si{\kilo\ohm}$}, v=$V_o$] (4.0,0) -- (0,0);
    \node[below, font=\footnotesize\bfseries] at (2.0,-0.35) {Apartado b)};
\end{circuitikz}
\end{minipage}
\hfill
\begin{minipage}{0.32\textwidth}
\centering
\begin{circuitikz}[scale=0.72, transform shape]
    \draw (0,0) to[isource, l=$5\,\si{\milli\ampere}$] (0,2.2) -- (4.2,2.2);
    \draw (1.4,2.2) to[R, l={$2\,\si{\kilo\ohm}$}, i>^=$I_x$] (1.4,0);
    \draw (2.8,2.2) to[cisource, l={$2 I_x$}] (2.8,0);
    \draw (4.2,2.2) to[R, l={$1\,\si{\kilo\ohm}$}, v=$V_o$] (4.2,0);
    \draw (0,0) -- (4.2,0);
    \draw (2.1,0) node[ground] {};
    \node[below, font=\footnotesize\bfseries] at (2.1,-0.45) {Apartado c)};
\end{circuitikz}
\end{minipage}
\end{center}

\textbf{Resolución paso a paso}:
\begin{itemize}
    \item \textbf{Apartado a)}: Ponemos las resistencias en ohmios ($3\,\si{\kilo\ohm} + 5\,\si{\kilo\ohm} = 8000\,\si{\ohm}$) y planteamos la ecuación de la única malla en sentido horario:
    \begin{equation*}
        12 + 2000\,I_s = I_s\cdot(3000 + 5000) \implies 12 = 6000\,I_s \implies I_s = \frac{12}{6000} = 2\,\si{\milli\ampere}
    \end{equation*}
    Por tanto, la tensión de salida en la resistencia de $5\,\si{\kilo\ohm}$ es:
    \begin{equation*}
        \boxed{V_o = I_s \cdot 5\,\si{\kilo\ohm} = 2\,\si{\milli\ampere}\cdot 5\,\si{\kilo\ohm} = 10\,\si{\volt}}
    \end{equation*}
    \item \textbf{Apartado b)}: Planteamos la ecuación de malla en sentido horario ($2\,\si{\kilo\ohm} + 2\,\si{\kilo\ohm} = 4\,\si{\kilo\ohm}$) y la ecuación de control $V_o = I\cdot 2\,\si{\kilo\ohm}$:
    \begin{equation*}
        12 - 2 V_o = I\cdot 4\,\si{\kilo\ohm} \xrightarrow{V_o = 2\,\si{\kilo\ohm}\cdot I} 12 - 4 I = 4 I \implies 12 = 8 I \implies I = \frac{3}{2}\,\si{\milli\ampere} = 1{,}5\,\si{\milli\ampere}
    \end{equation*}
    \begin{equation*}
        \boxed{V_o = I\cdot 2\,\si{\kilo\ohm} = \frac{3}{2}\,\si{\milli\ampere}\cdot 2\,\si{\kilo\ohm} = 3\,\si{\volt}}
    \end{equation*}
    \item \textbf{Apartado c)}: Todos los elementos están en paralelo entre el nudo superior (con potencial $V_o$) y tierra ($0\,\si{\volt}$). Aplicamos LCK en el nudo superior y sustituimos $I_x = \frac{V_o}{2\,\si{\kilo\ohm}}$:
    \begin{equation*}
        5\,\si{\milli\ampere} = I_x + 2 I_x + \frac{V_o}{1\,\si{\kilo\ohm}} = 3 I_x + \frac{V_o}{1\,\si{\kilo\ohm}} = \frac{3 V_o}{2\,\si{\kilo\ohm}} + \frac{V_o}{1\,\si{\kilo\ohm}} = \frac{5 V_o}{2\,\si{\kilo\ohm}} \implies \boxed{V_o = \frac{10}{5} = 2\,\si{\volt}}
    \end{equation*}
\end{itemize}
\end{ejercicio}

\begin{ejercicio}[Aplicación del Principio de Superposición (Ejercicio 20)]
Encuentra el valor de la corriente $I_o$ que atraviesa la resistencia de $6\,\si{\kilo\ohm}$ en el circuito de la figura utilizando el Principio de Superposición:

\begin{center}
\begin{minipage}{0.32\textwidth}
\centering
\begin{circuitikz}[scale=0.72, transform shape]
    \draw (0,0) to[V, v<=$12\,\si{\volt}$] (0,1.6) -- (1.8,1.6)
          to[R, l=$3\,\si{\kilo\ohm}$] (1.8,0) -- (0,0);
    \draw (0,1.6) to[isource, l=$4\,\si{\milli\ampere}$] (1.8,1.6);
    \draw (0,1.6) -- (0,2.8) to[R, l=$2\,\si{\kilo\ohm}$] (3.6,2.8) -- (3.6,1.6);
    \draw (1.8,1.6) to[R, l_=$6\,\si{\kilo\ohm}$, i>^=$I_o$] (3.6,1.6)
          to[R, l=$4\,\si{\kilo\ohm}$] (3.6,0) -- (1.8,0);
    \node[below, font=\footnotesize\bfseries] at (1.8,-0.35) {Circuito original};
\end{circuitikz}
\end{minipage}
\hfill
\begin{minipage}{0.32\textwidth}
\centering
\begin{circuitikz}[scale=0.72, transform shape]
    \draw (0,0) to[V, v<=$12\,\si{\volt}$] (0,2.6)
          to[R, l=$2\,\si{\kilo\ohm}$] (3.4,2.6) -- (3.4,1.5) node[right] {$V_A$};
    \draw (1.6,1.5) to[R, l=$3\,\si{\kilo\ohm}$] (1.6,0) -- (0,0);
    \draw (1.6,1.5) to[R, l=$6\,\si{\kilo\ohm}$, i>^=$I_{o1}$] (3.4,1.5)
          to[R, l=$4\,\si{\kilo\ohm}$] (3.4,0) -- (1.6,0);
    \draw (1.6,0) node[ground] {};
    \node[below, font=\footnotesize\bfseries] at (1.7,-0.45) {1) Solo fuente de $12\,\si{\volt}$};
\end{circuitikz}
\end{minipage}
\hfill
\begin{minipage}{0.32\textwidth}
\centering
\begin{circuitikz}[scale=0.72, transform shape]
    \draw (0,1.5) to[isource, l=$4\,\si{\milli\ampere}$] (1.8,1.5) node[above] {$A$};
    \draw (0,1.5) -- (0,2.6) to[R, l=$2\,\si{\kilo\ohm}$] (3.6,2.6) -- (3.6,1.5) node[right] {$C$};
    \draw (0,1.5) -- (0,0) -- (3.6,0);
    \draw (1.8,1.5) to[R, l=$3\,\si{\kilo\ohm}$] (1.8,0);
    \draw (1.8,1.5) to[R, l=$6\,\si{\kilo\ohm}$, i>^=$I_{o2}$] (3.6,1.5)
          to[R, l=$4\,\si{\kilo\ohm}$] (3.6,0);
    \draw (1.8,0) node[ground] {};
    \node[below, font=\footnotesize\bfseries] at (1.8,-0.45) {2) Solo fuente de $4\,\si{\milli\ampere}$};
\end{circuitikz}
\end{minipage}
\end{center}

\textbf{Resolución paso a paso}:
\begin{enumerate}
    \item \textbf{Paso 1: Con la fuente de tensión de $12\,\si{\volt}$ activa} (anulamos la fuente de $4\,\si{\milli\ampere}$ dejándola en circuito abierto):
    \begin{itemize}
        \item Desde el nudo $V_A$ hacia tierra quedan dos ramas en paralelo: la resistencia de $4\,\si{\kilo\ohm}$ y la serie $(6\,\si{\kilo\ohm} + 3\,\si{\kilo\ohm} = 9\,\si{\kilo\ohm})$:
        \begin{equation*}
            R_p = 4\,\si{\kilo\ohm} \parallel 9\,\si{\kilo\ohm} = \frac{4\cdot 9}{4 + 9} = \frac{36}{13}\,\si{\kilo\ohm}
        \end{equation*}
        \item Por divisor de tensión con la resistencia superior de $2\,\si{\kilo\ohm}$:
        \begin{equation*}
            V_A = \frac{R_p}{2 + R_p}\cdot 12\,\si{\volt} = \frac{36/13}{2 + 36/13}\cdot 12 = \frac{432}{62} \approx 6{,}97\,\si{\volt}
        \end{equation*}
        \item Como la corriente $I_{o1}$ está definida hacia la derecha (hacia $V_A$) y por la rama serie de $9\,\si{\kilo\ohm}$ la corriente va de $V_A$ a tierra:
        \begin{equation*}
            I_{o1} = \frac{0 - V_A}{6\,\si{\kilo\ohm} + 3\,\si{\kilo\ohm}} = -\frac{6{,}97\,\si{\volt}}{9\,\si{\kilo\ohm}} \approx -0{,}77\,\si{\milli\ampere}
        \end{equation*}
    \end{itemize}
    \item \textbf{Paso 2: Con la fuente de intensidad de $4\,\si{\milli\ampere}$ activa} (cortocircuitamos la fuente de $12\,\si{\volt}$ a tierra):
    \begin{itemize}
        \item Al quedar el terminal izquierdo puesto a tierra ($0\,\si{\volt}$), las resistencias de $4\,\si{\kilo\ohm}$ y $2\,\si{\kilo\ohm}$ quedan en paralelo entre el nudo $C$ y tierra:
        \begin{equation*}
            4\,\si{\kilo\ohm} \parallel 2\,\si{\kilo\ohm} = \frac{4\cdot 2}{4 + 2} = \frac{4}{3}\,\si{\kilo\ohm}
        \end{equation*}
        \item Esta combinación queda en serie con la resistencia de $6\,\si{\kilo\ohm}$ ($6 + \frac{4}{3} = \frac{22}{3}\,\si{\kilo\ohm}$) y todo ello en paralelo con la resistencia de $3\,\si{\kilo\ohm}$:
        \begin{equation*}
            R_{eq,A} = 3\,\si{\kilo\ohm} \parallel \frac{22}{3}\,\si{\kilo\ohm} = \frac{3 \cdot \frac{22}{3}}{3 + \frac{22}{3}} = \frac{22}{31}\,\si{\kilo\ohm}
        \end{equation*}
        \item La tensión en el nudo $A$ y la corriente $I_{o2}$ son:
        \begin{equation*}
            V_A = 4\,\si{\milli\ampere} \cdot \frac{22}{31}\,\si{\kilo\ohm} = \frac{88}{31}\,\si{\volt} \approx 2{,}84\,\si{\volt} \implies I_{o2} = \frac{V_A}{\frac{22}{3}\,\si{\kilo\ohm}} = \frac{88/31}{22/3} = \frac{12}{31}\,\si{\milli\ampere} \approx 0{,}387\,\si{\milli\ampere}
        \end{equation*}
        \emph{(Nota: Si en el cálculo intermedio se tomase $V_A = 8{,}51\,\si{\volt}$ como aparece anotado en el cuaderno manuscrito, saldría $I_{o2} = 1{,}16\,\si{\milli\ampere} \implies I_{o,\text{total}} = -0{,}77 + 1{,}16 = 0{,}39\,\si{\milli\ampere}$; con el valor exacto $V_A = \frac{88}{31}\,\si{\volt}$, resulta $I_{o,\text{total}} = -\frac{24}{31} + \frac{12}{31} = -\frac{12}{31}\,\si{\milli\ampere} \approx -0{,}387\,\si{\milli\ampere}$)}.
    \end{itemize}
\end{enumerate}
\end{ejercicio}

\begin{ejercicio}[Problema Completo de Examen 2019: Mallas, Nudos y Balance de Potencias]
En el circuito de examen de la figura, calcula todas las intensidades de malla y de rama, los potenciales de todos los nudos ($V_A, V_B, V_C, V_D, V_E$), las caídas de tensión en las cinco resistencias y el balance completo de potencias entre resistencias y fuentes:

\begin{center}
\begin{circuitikz}[scale=0.92, transform shape]
    \draw (0,1.6) node[left] {$A$} -- (0,3.2)
          to[R, l={$R_1=10\,\si{\ohm}$}] (6.0,3.2) -- (6.0,1.6) node[right] {$C$};
    \draw (0,1.6) to[R, l={$R_2=3\,\si{\ohm}$}] (3.0,1.6) node[above] {$B$}
          to[R, l={$R_3=2\,\si{\ohm}$}] (6.0,1.6);
    \draw (0,0) node[left] {$D$} to[V, v<=$100\,\si{\volt}$] (0,1.6);
    \draw (3.0,0) to[isource, l=$5\,\si{\ampere}$, v_=$V_x$] (3.0,1.6);
    \draw (6.0,0) node[right] {$E$} to[V, v=$50\,\si{\volt}$] (6.0,1.6);
    \draw (0,0) to[R, l_={$R_4=6\,\si{\ohm}$}] (3.0,0)
          to[R, l_={$R_5=4\,\si{\ohm}$}] (6.0,0);
    \draw (3.0,0) node[ground] {} node[below right] {$0\,\si{\volt}$};
    \draw[-{Stealth}, NavyBlue, thick] (2.6,2.4) arc (150:-150:0.38) node[midway, left] {$I_1$};
    \draw[-{Stealth}, BrickRed, thick] (1.1,0.8) arc (150:-150:0.38) node[midway, left] {$I_2$};
    \draw[-{Stealth}, ForestGreen, thick] (4.3,0.8) arc (150:-150:0.38) node[midway, left] {$I_3$};
\end{circuitikz}
\end{center}

\textbf{Resolución paso a paso}:
\begin{enumerate}
    \item \textbf{Planteamiento por el método de mallas (con supermalla inferior)}:
    Asignamos $I_1, I_2, I_3$ en sentido horario y llamamos $V_x = V_B - 0$ a la tensión en bornes de la fuente de intensidad de $5\,\si{\ampere}$:
    \begin{align}
        \text{Malla 1 (superior):} \quad 0 &= (10 + 3 + 2)I_1 - 3 I_2 - 2 I_3 \implies 0 = 15 I_1 - 3 I_2 - 2 I_3 \label{eq:ex19_1} \\
        \text{Malla 2 (izquierda):} \quad 100 &= -3 I_1 + (3 + 6)I_2 + V_x = -3 I_1 + 9 I_2 + V_x \label{eq:ex19_2} \\
        \text{Malla 3 (derecha):} \quad -50 &= -2 I_1 + (2 + 4)I_3 - V_x = -2 I_1 + 6 I_3 - V_x \label{eq:ex19_3}
    \end{align}
    \item \textbf{Eliminación de $V_x$ (ecuación de la supermalla) y relación de la fuente de corriente}:
    \begin{itemize}
        \item Sumando \eqref{eq:ex19_2} y \eqref{eq:ex19_3}:
        \begin{equation*}
            50 = -5 I_1 + 9 I_2 + 6 I_3 \xrightarrow{\times 3} 150 = -15 I_1 + 27 I_2 + 18 I_3
        \end{equation*}
        \item Sumándole la ecuación \eqref{eq:ex19_1} ($0 = 15 I_1 - 3 I_2 - 2 I_3$) desaparece $I_1$:
        \begin{equation*}
            150 = 24 I_2 + 16 I_3
        \end{equation*}
        \item Además, en la rama central la fuente de $5\,\si{\ampere}$ impone directamente:
        \begin{equation*}
            I_3 - I_2 = 5 \implies I_3 = 5 + I_2
        \end{equation*}
        \item Sustituyendo $I_3 = 5 + I_2$:
        \begin{equation*}
            150 = 24 I_2 + 16(5 + I_2) = 40 I_2 + 80 \implies 40 I_2 = 70 \implies \boxed{I_2 = 1{,}75\,\si{\ampere}}
        \end{equation*}
        \begin{equation*}
            \boxed{I_3 = 5 + 1{,}75 = 6{,}75\,\si{\ampere}}, \qquad \boxed{I_1 = \frac{3(1{,}75) + 2(6{,}75)}{15} = \frac{18{,}75}{15} = 1{,}25\,\si{\ampere}}
        \end{equation*}
    \end{itemize}
    \item \textbf{Intensidades y caídas de tensión en cada resistencia}:
    \begin{align*}
        I_{R1} &= I_1 = 1{,}25\,\si{\ampere} &\implies V_{R1} &= 1{,}25\cdot 10 = 12{,}5\,\si{\volt} \\
        I_{R2} &= I_2 - I_1 = 1{,}75 - 1{,}25 = 0{,}50\,\si{\ampere} &\implies V_{R2} &= 0{,}50\cdot 3 = 1{,}5\,\si{\volt} \\
        I_{R3} &= I_3 - I_1 = 6{,}75 - 1{,}25 = 5{,}50\,\si{\ampere} &\implies V_{R3} &= 5{,}50\cdot 2 = 11{,}0\,\si{\volt} \\
        I_{R4} &= I_2 = 1{,}75\,\si{\ampere} &\implies V_{R4} &= 1{,}75\cdot 6 = 10{,}5\,\si{\volt} \\
        I_{R5} &= I_3 = 6{,}75\,\si{\ampere} &\implies V_{R5} &= 6{,}75\cdot 4 = 27{,}0\,\si{\volt}
    \end{align*}
    \item \textbf{Tensiones absolutas en todos los nudos respecto a tierra ($0\,\si{\volt}$)}:
    \begin{equation*}
        V_D = -10{,}5\,\si{\volt}, \quad V_A = 89{,}5\,\si{\volt}, \quad V_B = 88{,}0\,\si{\volt}, \quad V_E = 27{,}0\,\si{\volt}, \quad V_C = 77{,}0\,\si{\volt}
    \end{equation*}
    \item \textbf{Balance completo de potencias}:
    \begin{itemize}
        \item \textbf{Potencia disipada en las resistencias}:
        \begin{align*}
            P_{R1} &= 12{,}5\cdot 1{,}25 = 15{,}625\,\si{\watt}, \quad P_{R2} = 1{,}5\cdot 0{,}5 = 0{,}75\,\si{\watt}, \quad P_{R3} = 11\cdot 5{,}5 = 60{,}5\,\si{\watt} \\
            P_{R4} &= 10{,}5\cdot 1{,}75 = 18{,}375\,\si{\watt}, \quad P_{R5} = 27\cdot 6{,}75 = 182{,}25\,\si{\watt} \implies \boxed{P_{R,\text{total}} = 277{,}5\,\si{\watt}}
        \end{align*}
        \item \textbf{Potencia suministrada por las fuentes}:
        \begin{align*}
            P_{V1} &= 100\,\si{\volt}\cdot I_2 = 100\cdot 1{,}75 = 175\,\si{\watt} \quad (\text{cede}) \\
            P_{V2} &= -50\,\si{\volt}\cdot I_3 = -50\cdot 6{,}75 = -337{,}5\,\si{\watt} \quad (\text{absorbe, pues } I_3 \text{ entra por el } +) \\
            P_{I1} &= V_B \cdot 5\,\si{\ampere} = 88\,\si{\volt}\cdot 5\,\si{\ampere} = 440\,\si{\watt} \quad (\text{cede}) \\
            \implies &\; \boxed{P_{F,\text{total}} = 175 - 337{,}5 + 440 = 277{,}5\,\si{\watt}} \quad (\text{Se cumple } P_{F,\text{total}} = P_{R,\text{total}})
        \end{align*}
    \end{itemize}
\end{enumerate}
\end{ejercicio}

\chapter{Teoría de Circuitos: Corriente Alterna (AC)}
\label{cap:tema3}

\cuadrofuentes{
    \textbf{Asignatura:} Fundamentos Físicos y Tecnológicos (1.º Grado en Ingeniería Informática, UGR).\\[2pt]
    \textbf{Profesor:} Pedro Cartujo (Departamento de Electrónica y Tecnología de Computadores).\\[2pt]
    \textbf{Bibliografía de Referencia:} Alexander, C. K., Sadiku, M. N. O., \textit{Fundamentos de Circuitos Eléctricos}, McGraw-Hill; Nilsson, J. W., Riedel, S. A., \textit{Circuitos Eléctricos}, Pearson.
}

% ==============================================================================
% 3.1. SEÑALES SINUSOIDALES Y VALOR EFICAZ
% ==============================================================================
\section{Señales Sinusoidales en Régimen Permanente}

En corriente alterna (AC), las fuentes generan tensiones y corrientes que varían sinusoidalmente en el tiempo. El estudio de circuitos en régimen permanente sinusoidal (RPS) es fundamental, ya que toda señal periódica real puede descomponerse en suma de sinusoides mediante series de Fourier.

\subsection{Ecuación de una Señal Sinusoidal}
Una tensión alterna sinusoidal se describe matemáticamente por:
\begin{equation}
    v(t) = V_m \cos(\omega t + \varphi) \quad \text{o bien} \quad v(t) = V_m \sin(\omega t + \varphi')
\end{equation}
donde:
\begin{itemize}[leftmargin=*]
    \item $V_m$ es la \textbf{amplitud pico} o valor máximo de la señal ($[\si{\volt}]$).
    \item $\omega$ es la \textbf{frecuencia angular}:
    \begin{equation}
        \omega = 2\pi f = \frac{2\pi}{T} \quad \left[\si{\radian\per\second}\right]
    \end{equation}
    \item $f$ es la \textbf{frecuencia cíclica} ($[\si{\hertz}]$), que mide los ciclos por segundo. En la red eléctrica europea, $f = 50\ \si{\hertz}$ ($\omega \approx 314{,}16\ \si{\radian\per\second}$).
    \item $T = 1/f$ es el \textbf{periodo} ($[\si{\second}]$), tiempo requerido para completar un ciclo completo.
    \item $\varphi$ es la \textbf{fase inicial} o desfase en el instante $t = 0$ ($[\si{\radian}]$ o $[\si{\degree}]$).
\end{itemize}

\subsection{Valor Medio y Valor Eficaz (RMS)}
\begin{definicion}[Valor Eficaz o RMS ($V_{rms}$)]
El valor medio de una sinusoide pura sobre un periodo completo es idénticamente nulo:
\begin{equation}
    V_{med} = \frac{1}{T}\int_0^T v(t)\, dt = 0
\end{equation}
El \textbf{valor eficaz} (o RMS, \textit{Root Mean Square}) representa el valor de tensión continua (DC) equivalente que disiparía exactamente la misma potencia media sobre una misma resistencia óhmica:
\begin{equation}
    V_{rms} = \sqrt{\frac{1}{T} \int_0^T v^2(t)\, dt} = \frac{V_m}{\sqrt{2}} \approx 0{,}7071 \cdot V_m
\end{equation}
Análogamente, para la corriente:
\begin{equation}
    I_{rms} = \frac{I_m}{\sqrt{2}} \approx 0{,}7071 \cdot I_m
\end{equation}
\end{definicion}

% ==============================================================================
% 3.2. FASORES E IMPEDANCIA COMPLEJA
% ==============================================================================
\section{Representación Fasorial e Impedancia Compleja}

Resolver circuitos en el dominio del tiempo con ecuaciones íntegro-diferenciales resulta engorroso. En 1893, Charles Proteus Steinmetz introdujo el método fasorial, transformando las ecuaciones diferenciales lineales en ecuaciones algebraicas complejas.

\subsection{Concepto de Fasor}
Utilizando la identidad de Euler ($e^{j\theta} = \cos \theta + j \sin \theta$, con $j = \sqrt{-1}$):
\begin{equation}
    v(t) = V_m \cos(\omega t + \varphi) = \text{Re}\left\{ V_m e^{j(\omega t + \varphi)} \right\} = \text{Re}\left\{ \left(V_m e^{j\varphi}\right) e^{j\omega t} \right\}
\end{equation}

\begin{definicion}[Fasor Asociado]
El \textbf{fasor} $\mathbf{V}$ es un número complejo que contiene exclusivamente la información de la amplitud y la fase inicial de la sinusoide (omitiendo el factor temporal $e^{j\omega t}$ común a todo el circuito):
\begin{equation}
    \mathbf{V} = V_m e^{j\varphi} = V_m \angle \varphi = V_m (\cos \varphi + j \sin \varphi)
\end{equation}
En ingeniería eléctrica es habitual utilizar el \textbf{fasor eficaz}:
\begin{equation}
    \mathbf{V}_{rms} = \frac{V_m}{\sqrt{2}} \angle \varphi = V_{rms} \angle \varphi
\end{equation}
\end{definicion}

\begin{aviso}[Propiedades Operacionales del Cálculo Fasorial]
\begin{itemize}[leftmargin=*]
    \item \textbf{Derivada Temporal:} La derivada en el tiempo equivale algebraicamente a multiplicar por $j\omega$:
    \begin{equation}
        \frac{d v(t)}{dt} \longleftrightarrow j\omega \mathbf{V}
    \end{equation}
    \item \textbf{Integral Temporal:} La integración en el tiempo equivale a dividir por $j\omega$:
    \begin{equation}
        \int v(t)\, dt \longleftrightarrow \frac{\mathbf{V}}{j\omega} = -\frac{j}{\omega} \mathbf{V}
    \end{equation}
\end{itemize}
\end{aviso}

\subsection{Impedancia Compleja de los Elementos Básicos}
\begin{definicion}[Impedancia Compleja $\mathbf{Z}$]
Se define la impedancia compleja $\mathbf{Z}$ como el cociente fasorial entre la tensión y la corriente:
\begin{equation}
    \mathbf{Z} = \frac{\mathbf{V}}{\mathbf{I}} = R + jX = |\mathbf{Z}| \angle \theta_Z \quad [\si{\ohm}]
\end{equation}
\begin{itemize}[leftmargin=*]
    \item $R = \text{Re}\{\mathbf{Z}\}$ es la \textbf{Resistencia} ($[\si{\ohm}]$), siempre no negativa ($R \ge 0$).
    \item $X = \text{Im}\{\mathbf{Z}\}$ es la \textbf{Reactancia} ($[\si{\ohm}]$). Si $X > 0$, es reactancia inductiva; si $X < 0$, es capacitiva.
    \item $|\mathbf{Z}| = \sqrt{R^2 + X^2}$ es el módulo de la impedancia.
    \item $\theta_Z = \arctan(X/R) = \theta_v - \theta_i$ es el ángulo de desfase entre tensión y corriente.
\end{itemize}
\end{definicion}

\begin{ley}[Impedancias Fundamentales en RPS]
\begin{enumerate}[leftmargin=*]
    \item \textbf{Resistencia Pura ($R$):}
    \begin{equation}
        v_R(t) = R i(t) \implies \mathbf{V}_R = R \mathbf{I} \implies \mathbf{Z}_R = R = R \angle 0^\circ
    \end{equation}
    La tensión y la corriente están estrictamente \textbf{en fase} ($\theta = 0^\circ$).
    \item \textbf{Bobina o Inductor Puro ($L$):}
    \begin{equation}
        v_L(t) = L \frac{di}{dt} \implies \mathbf{V}_L = j\omega L \mathbf{I} \implies \mathbf{Z}_L = j\omega L = \omega L \angle 90^\circ
    \end{equation}
    La reactancia inductiva es $X_L = \omega L > 0$. \textbf{La tensión adelanta $90^\circ$ a la corriente} (o la corriente retrasa $90^\circ$ a la tensión).
    \item \textbf{Condensador Puro ($C$):}
    \begin{equation}
        i(t) = C \frac{dv}{dt} \implies \mathbf{I} = j\omega C \mathbf{V} \implies \mathbf{Z}_C = \frac{1}{j\omega C} = -\frac{j}{\omega C} = \frac{1}{\omega C} \angle -90^\circ
    \end{equation}
    La reactancia capacitiva es $X_C = -\frac{1}{\omega C} < 0$. \textbf{La corriente adelanta $90^\circ$ a la tensión} (o la tensión retrasa $90^\circ$ a la corriente).
\end{enumerate}
\end{ley}

\begin{definicion}[Admitancia Compleja $\mathbf{Y}$]
La admitancia es el inverso multiplicativo de la impedancia:
\begin{equation}
    \mathbf{Y} = \frac{1}{\mathbf{Z}} = G + jB \quad [\si{\siemens}]
\end{equation}
donde $G$ es la \textbf{conductancia} y $B$ es la \textbf{susceptancia}.
\end{definicion}

% ==============================================================================
% 3.3. POTENCIA EN CORRIENTE ALTERNA
% ==============================================================================
\section{Potencia en Corriente Alterna y Corrección del Factor de Potencia}

\subsection{Potencia Instantánea, Activa y Reactiva}
Sean la tensión y la corriente en un elemento dadas por:
\begin{equation}
    v(t) = V_m \cos(\omega t + \theta_v), \qquad i(t) = I_m \cos(\omega t + \theta_i)
\end{equation}
El desfase entre ambas es $\theta = \theta_v - \theta_i$. La potencia instantánea disipada es:
\begin{align}
    p(t) &= v(t) \cdot i(t) = V_m I_m \cos(\omega t + \theta_v) \cos(\omega t + \theta_i) \nonumber \\
         &= \underbrace{V_{rms} I_{rms} \cos \theta \left[1 + \cos(2\omega t + 2\theta_v)\right]}_{\text{Término activo unidireccional}} - \underbrace{V_{rms} I_{rms} \sin \theta \sin(2\omega t + 2\theta_v)}_{\text{Término reactivo oscilante}}
\end{align}

\begin{definicion}[Clasificación de Potencias en AC]
\begin{itemize}[leftmargin=*]
    \item \textbf{Potencia Activa o Real ($P$):} Es el promedio temporal de la potencia instantánea. Representa la energía neta por unidad de tiempo transformada irreversiblemente en trabajo mecánico, calor o luz:
    \begin{equation}
        P = \frac{1}{T}\int_0^T p(t)\, dt = V_{rms} I_{rms} \cos \theta \quad [\si{\watt}]
    \end{equation}
    \item \textbf{Potencia Reactiva ($Q$):} Representa la tasa de intercambio oscilatorio de energía entre los campos magnéticos y eléctricos del circuito y la fuente generadora (no produce trabajo útil neto):
    \begin{equation}
        Q = V_{rms} I_{rms} \sin \theta \quad [\si{\VAR} \equiv \text{Voltio-Amperio Reactivo}]
    \end{equation}
    \begin{itemize}
        \item Carga inductiva ($\theta > 0$, la corriente retrasa): $Q > 0$ (consume reactiva).
        \item Carga capacitiva ($\theta < 0$, la corriente adelanta): $Q < 0$ (genera reactiva).
    \end{itemize}
    \item \textbf{Potencia Compleja ($\mathbf{S}$):}
    \begin{equation}
        \mathbf{S} = \mathbf{V}_{rms} \mathbf{I}_{rms}^* = P + jQ = S \angle \theta \quad [\si{\VA} \equiv \text{Voltio-Amperio}]
    \end{equation}
    donde $\mathbf{I}^*$ denota el conjugado complejo del fasor corriente.
    \item \textbf{Potencia Aparente ($S$):} Módulo de la potencia compleja:
    \begin{equation}
        S = |\mathbf{S}| = \sqrt{P^2 + Q^2} = V_{rms} I_{rms} \quad [\si{\VA}]
    \end{equation}
    \item \textbf{Factor de Potencia ($\text{fp}$):}
    \begin{equation}
        \text{fp} = \cos \theta = \frac{P}{S} \quad (0 \le \text{fp} \le 1)
    \end{equation}
\end{itemize}
\end{definicion}

\subsection{Corrección del Factor de Potencia}
Las cargas industriales (motores, transformadores) son fuertemente inductivas ($\text{fp} \approx 0{,}6 - 0{,}8$ inductivo). Un factor de potencia bajo obliga a las compañías eléctricas a inyectar corrientes elevadas para entregar la misma potencia activa $P$, provocando pérdidas óhmicas elevadas en las líneas de distribución ($P_{loss} = I^2 R_{linea}$) y caídas de tensión inadmisibles.

\begin{metodo}[Cálculo de la Capacidad de Compensación]
Para elevar el factor de potencia desde $\cos \theta_1$ hasta un nuevo valor óptimo $\cos \theta_2$ ($\approx 0{,}95 - 1{,}0$) sin alterar la potencia activa útil $P$, se conecta en paralelo una batería de condensadores que aporte la potencia reactiva de compensación:
\begin{equation}
    Q_C = Q_2 - Q_1 = P (\tan \theta_2 - \tan \theta_1) < 0
\end{equation}
Como la potencia reactiva de un condensador en paralelo sometido a una tensión eficaz $V_{rms}$ es $Q_C = -\omega C V_{rms}^2$:
\begin{equation}
    C = \frac{P (\tan \theta_1 - \tan \theta_2)}{\omega V_{rms}^2} \quad [\si{\farad}]
\end{equation}
\end{metodo}

% ==============================================================================
% 3.4. RESPUESTA EN FRECUENCIA Y FILTROS PASIVOS
% ==============================================================================
\section{Respuesta en Frecuencia y Filtros Pasivos RC}

Un filtro es un circuito diseñado para seleccionar o atenuar determinadas bandas de frecuencia de una señal de entrada compuesta.

\begin{definicion}[Función de Transferencia $H(j\omega)$]
Se define como la relación compleja entre el fasor de salida y el de entrada en régimen permanente:
\begin{equation}
    H(j\omega) = \frac{\mathbf{V}_{out}(j\omega)}{\mathbf{V}_{in}(j\omega)} = |H(j\omega)| e^{j\angle H(j\omega)}
\end{equation}
La ganancia en decibelios (dB) se expresa como:
\begin{equation}
    |H(j\omega)|_{\si{\decibel}} = 20 \log_{10}|H(j\omega)|
\end{equation}
\end{definicion}

\subsection{Filtro RC Pasa-Bajos (Low-Pass Filter)}
Se toma la salida en bornes del condensador $C$ alimentado en serie con una resistencia $R$.
Aplicando el divisor de tensión:
\begin{equation}
    H(j\omega) = \frac{\mathbf{Z}_C}{\mathbf{Z}_R + \mathbf{Z}_C} = \frac{\frac{1}{j\omega C}}{R + \frac{1}{j\omega C}} = \frac{1}{1 + j\omega R C} = \frac{1}{1 + j(\omega / \omega_c)}
\end{equation}

\begin{ley}[Frecuencia de Corte a $-3\ \si{\decibel}$]
La \textbf{frecuencia de corte} $\omega_c$ es aquella donde la potencia de salida se reduce a la mitad (atenuación de $-3\ \si{\decibel}$ respecto a la banda de paso):
\begin{equation}
    \omega_c = \frac{1}{R C} \implies f_c = \frac{1}{2\pi R C}
\end{equation}
\begin{itemize}[leftmargin=*]
    \item Para $\omega \ll \omega_c$: $|H| \approx 1$ ($0\ \si{\decibel}$), la señal pasa prácticamente intacta.
    \item Para $\omega = \omega_c$: $|H| = \frac{1}{\sqrt{1 + 1}} = \frac{1}{\sqrt{2}} \approx 0{,}707$ ($-3{,}01\ \si{\decibel}$), y el desfase es $\angle H = -45^\circ$.
    \item Para $\omega \gg \omega_c$: $|H| \approx \frac{\omega_c}{\omega}$, la pendiente de atenuación asintótica es de \textbf{$-20\ \si{\decibel\per\text{década}}$} ($-6\ \si{\decibel\per\text{octava}}$), y el desfase tiende a $-90^\circ$.
\end{itemize}
\end{ley}

\subsection{Filtro RC Pasa-Altos (High-Pass Filter)}
Se toma la salida en bornes de la resistencia $R$.
\begin{equation}
    H(j\omega) = \frac{\mathbf{Z}_R}{\mathbf{Z}_R + \mathbf{Z}_C} = \frac{R}{R + \frac{1}{j\omega C}} = \frac{j\omega R C}{1 + j\omega R C} = \frac{j(\omega / \omega_c)}{1 + j(\omega / \omega_c)}
\end{equation}
\begin{itemize}[leftmargin=*]
    \item Para $\omega \ll \omega_c$: $|H| \approx \frac{\omega}{\omega_c} \to 0$ (atenuación de $+20\ \si{\decibel\per\text{década}}$ a bajas frecuencias).
    \item Para $\omega = \omega_c$: $|H| = \frac{1}{\sqrt{2}} \approx 0{,}707$ ($-3\ \si{\decibel}$), con fase $+45^\circ$.
    \item Para $\omega \gg \omega_c$: $|H| \to 1$ ($0\ \si{\decibel}$), la señal de alta frecuencia no sufre atenuación.
\end{itemize}

% ==============================================================================
% 3.5. EJERCICIOS RESUELTOS DE EXAMEN
% ==============================================================================
\section{Ejercicios Resueltos de Exámenes de Circuitos AC}

\begin{ejercicio}[Análisis Fasorial Completo de Circuito RLC --- Examen UGR]
Una fuente sinusoidal de tensión $v_s(t) = 100 \sqrt{2} \cos(1000 t)\ \si{\volt}$ alimenta una asociación serie formada por una resistencia $R = 30\ \si{\ohm}$, una bobina $L = 80\ \si{\milli\henry}$ y un condensador $C = 25\ \si{\micro\farad}$.
\begin{enumerate}[label=\alph*)]
    \item Calcular la impedancia total compleja $\mathbf{Z}_{eq}$ del circuito.
    \item Hallar la expresión temporal de la corriente $i(t)$ en régimen permanente.
    \item Calcular las potencias activa $P$, reactiva $Q$ y aparente $S$, así como el factor de potencia.
\end{enumerate}
\tcbline
\textbf{Solución:}\\[4pt]
a) Con $\omega = 1000\ \si{\radian\per\second}$ y fasor eficaz $\mathbf{V}_{rms} = 100 \angle 0^\circ\ \si{\volt}$:
\begin{align}
    \mathbf{Z}_R &= 30\ \si{\ohm} \\
    \mathbf{Z}_L &= j\omega L = j (1000 \times 0{,}080) = j 80\ \si{\ohm} \\
    \mathbf{Z}_C &= -\frac{j}{\omega C} = -\frac{j}{1000 \times (25 \times 10^{-6})} = -\frac{j}{0{,}025} = -j 40\ \si{\ohm}
\end{align}
La impedancia total en serie es:
\begin{equation}
    \mathbf{Z}_{eq} = \mathbf{Z}_R + \mathbf{Z}_L + \mathbf{Z}_C = 30 + j(80 - 40) = 30 + j 40\ \si{\ohm}
\end{equation}
Expresando en forma polar:
\begin{equation}
    |\mathbf{Z}_{eq}| = \sqrt{30^2 + 40^2} = \sqrt{900 + 1600} = \sqrt{2500} = 50\ \si{\ohm}
\end{equation}
\begin{equation}
    \theta_Z = \arctan\left(\frac{40}{30}\right) = \arctan(1{,}3333) \approx 53{,}13^\circ \approx 0{,}927\ \si{\radian}
\end{equation}
$\mathbf{Z}_{eq} = 50 \angle 53{,}13^\circ\ \si{\ohm}$ (carácter predominantemente inductivo).\\[6pt]
b) Calculamos el fasor de corriente eficaz:
\begin{equation}
    \mathbf{I}_{rms} = \frac{\mathbf{V}_{rms}}{\mathbf{Z}_{eq}} = \frac{100 \angle 0^\circ}{50 \angle 53{,}13^\circ} = 2 \angle -53{,}13^\circ\ \si{\ampere}
\end{equation}
La amplitud máxima de la corriente es $I_m = 2 \times \sqrt{2} \approx 2{,}828\ \si{\ampere}$.
La expresión temporal instantánea es:
\begin{equation}
    i(t) = 2 \sqrt{2} \cos(1000 t - 53{,}13^\circ)\ \si{\ampere}
\end{equation}
c) Cálculo de potencias:
\begin{align}
    P &= V_{rms} I_{rms} \cos(\theta_Z) = 100 \times 2 \times \cos(53{,}13^\circ) = 200 \times 0{,}60 = 120\ \si{\watt} \\
    Q &= V_{rms} I_{rms} \sin(\theta_Z) = 100 \times 2 \times \sin(53{,}13^\circ) = 200 \times 0{,}80 = 160\ \si{\VAR} \\
    S &= V_{rms} I_{rms} = 100 \times 2 = 200\ \si{\VA} \\
    \text{fp} &= \cos(53{,}13^\circ) = 0{,}60\quad (\text{en retraso / inductivo})
\end{align}
Nótese que también: $P = I_{rms}^2 R = 2^2 \times 30 = 120\ \si{\watt}$ y $Q = I_{rms}^2 X = 2^2 \times 40 = 160\ \si{\VAR}$.
\end{ejercicio}

\begin{ejercicio}[Diseño y Respuesta en Frecuencia de un Filtro RC Pasa-Bajos]
Se dispone de un condensador $C = 10\ \si{\nano\farad}$.
\begin{enumerate}[label=\alph*)]
    \item Calcular la resistencia $R$ requerida para fijar la frecuencia de corte en $f_c = 10\ \si{\kilo\hertz}$.
    \item Para una señal sinusoidal de entrada de amplitud $V_i = 5\ \si{\volt}$ y frecuencia $f_1 = 100\ \si{\kilo\hertz}$ ($10$ veces la frecuencia de corte), determinar analíticamente la amplitud y el desfase de la señal de salida.
\end{enumerate}
\tcbline
\textbf{Solución:}\\[4pt]
a) La frecuencia angular de corte es:
\begin{equation}
    \omega_c = 2\pi f_c = 2\pi \times (10 \times 10^3) = 20000\pi \approx 62831{,}85\ \si{\radian\per\second}
\end{equation}
A partir de la condición $\omega_c = \frac{1}{R C}$:
\begin{equation}
    R = \frac{1}{\omega_c C} = \frac{1}{2\pi \times 10^4 \times 10 \times 10^{-9}} = \frac{1}{2\pi \times 10^{-4}} = \frac{10000}{2\pi} \approx 1591{,}55\ \si{\ohm} \approx 1{,}59\ \si{\kilo\ohm}
\end{equation}
b) Para la frecuencia $f_1 = 100\ \si{\kilo\hertz} = 10 f_c$, la relación de frecuencias normalizada es $\frac{\omega}{\omega_c} = 10$:
\begin{equation}
    H(j\omega_1) = \frac{1}{1 + j(10)}
\end{equation}
El módulo de la ganancia es:
\begin{equation}
    |H(j\omega_1)| = \frac{1}{\sqrt{1^2 + 10^2}} = \frac{1}{\sqrt{101}} \approx \frac{1}{10{,}0499} \approx 0{,}0995
\end{equation}
En decibelios:
\begin{equation}
    |H|_{\si{\decibel}} = 20 \log_{10}(0{,}0995) \approx -20{,}04\ \si{\decibel}
\end{equation}
La amplitud de salida es:
\begin{equation}
    V_o = |H| \cdot V_i = 0{,}0995 \times 5 = 0{,}4975\ \si{\volt} \approx 0{,}50\ \si{\volt}
\end{equation}
El desfase de salida respecto a la entrada es:
\begin{equation}
    \varphi_o = -\arctan(10) \approx -84{,}29^\circ \approx -1{,}471\ \si{\radian}
\end{equation}
Se aprecia la atenuación de aproximadamente $-20\ \si{\decibel}$ por década a altas frecuencias.
\end{ejercicio}

\chapter{Dispositivos Electrónicos: Diodos y Transistores MOSFET}
\label{cap:tema4}

\cuadrofuentes{
    \textbf{Asignatura:} Fundamentos Físicos y Tecnológicos (1.º Grado en Ingeniería Informática, UGR).\\[2pt]
    \textbf{Profesor:} Pedro Cartujo (Departamento de Electrónica y Tecnología de Computadores).\\[2pt]
    \textbf{Bibliografía de Referencia:} Sedra, A. S., Smith, K. C., \textit{Circuitos Microelectrónicos}, Oxford University Press; Boylestad, R. L., Nashelsky, L., \textit{Electrónica: Teoría de Circuitos y Dispositivos Electrónicos}, Pearson.
}

% ==============================================================================
% 4.1. CONCEPTOS BÁSICOS DE SEMICONDUCTORES Y UNIÓN PN
% ==============================================================================
\section{Conceptos Básicos de Semiconductores y la Unión PN}

Los dispositivos electrónicos de estado sólido basan su funcionamiento en la posibilidad de controlar mediante campos eléctricos o dopaje selectivo la conductividad de los materiales \textbf{semiconductores}, entre los cuales el \textbf{Silicio ($\text{Si}$)} es el más utilizado en la industria microelectrónica actual (seguido del Germanio, $\text{Ge}$, y compuestos como el Arseniuro de Galio, $\text{GaAs}$).

\subsection{Bandas de Energía y Semiconductores Intrínsecos}
Desde el punto de vista cuántico, los electrones de un cristal sólido se agrupan en bandas de energía permitidas separadas por regiones prohibidas:
\begin{itemize}[leftmargin=*]
    \item \textbf{Banda de Valencia (BV):} Nivel energético ocupado por los electrones que forman los enlaces covalentes entre átomos vecinos de la red cristalina.
    \item \textbf{Banda de Conducción (BC):} Banda energética superior donde los electrones se han desligado de su átomo de origen y pueden desplazarse libremente por el cristal produciendo corriente eléctrica.
    \item \textbf{Banda Prohibida o \textit{Gap} de Energía ($E_g$):} Separación energética mínima entre el límite superior de la banda de valencia ($E_V$) y el límite inferior de la banda de conducción ($E_C$). En un aislante es muy grande ($E_g > 5\ \si{\electronvolt}$), en un conductor las bandas se solapan ($E_g = 0$), y en un semiconductor tiene un valor intermedio ($E_g \approx 1{,}12\ \si{\electronvolt}$ en el Silicio y $E_g \approx 0{,}67\ \si{\electronvolt}$ en el Germanio a $300\ \si{\kelvin}$).
\end{itemize}

Un semiconductor químicamente puro y con red cristalina perfecta se denomina \textbf{semiconductor intrínseco}:
\begin{enumerate}[leftmargin=*]
    \item \textbf{A temperatura nula ($T = 0\ \si{\kelvin}$):} Cada átomo de Silicio (tetravalente, con $+4$ electrones en su capa exterior) comparte sus cuatro electrones de valencia formando cuatro \textbf{enlaces covalentes} estables con sus vecinos. La banda de valencia está completamente llena y la de conducción vacía: el cristal no conduce corriente (se comporta como un aislante perfecto).
    \item \textbf{A temperatura ambiente ($T = 300\ \si{\kelvin}$):} La agitación térmica proporciona energía suficiente para romper algunos enlaces covalentes (\textbf{generación térmica}). Cuando un electrón salta a la banda de conducción convirtiéndose en un \textbf{electrón libre} (con carga $-q$), deja tras de sí una vacante de enlace en la banda de valencia denominada \textbf{hueco} ($h^+$, que se comporta como un portador móvil de carga efectiva $+q$). Simultáneamente ocurre el proceso inverso de \textbf{recombinación térmica} cuando un electrón libre cae en un hueco.
\end{enumerate}

\begin{ley}[Concentración Intrínseca y Ley de Acción de Masas]
En un semiconductor intrínseco, dado que los electrones libres y los huecos se generan en parejas, sus concentraciones volumétricas ($n$ para electrones y $p$ para huecos) son idénticas e iguales a la \textbf{concentración intrínseca} $n_i(T)$:
\begin{equation}
    n = p = n_i \approx 1{,}5 \times 10^{10}\ \si{\centi\meter^{-3}} \quad (\text{en Silicio a } T = 300\ \si{\kelvin})
\end{equation}
En equilibrio térmico, el producto de las concentraciones de electrones y huecos en cualquier semiconductor (esté o no dopado) permanece constante para una temperatura dada, cumpliendo la \textbf{Ley de Acción de Masas}:
\begin{equation}
    n \cdot p = n_i^2(T)
\end{equation}
\end{ley}

\subsection{Semiconductores Extrínsecos: Dopaje Tipo N y Tipo P}
La conductividad del silicio intrínseco a temperatura ambiente es muy baja. Para aumentarla y controlarla con precisión, se introducen deliberadamente pequeñas cantidades de átomos de impurezas sustitucionales en la red cristalina, proceso denominado \textbf{dopaje}:

\begin{itemize}[leftmargin=*]
    \item \textbf{Semiconductor Extrínseco Tipo N (Dopado con Impurezas Donadoras):}
    Se introducen átomos \textbf{pentavalentes} (grupo V de la tabla periódica, con $5$ electrones de valencia), tales como \textbf{Fósforo ($\text{P}$)}, \textbf{Arsénico ($\text{As}$)} o \textbf{Antimonio ($\text{Sb}$)}. Cuatro de sus electrones forman enlaces covalentes con los átomos de Silicio vecinos, mientras que el quinto electrón queda muy débilmente ligado y a temperatura ambiente pasa inmediatamente a la banda de conducción como \textbf{electrón libre extra}, sin haber creado un hueco. El átomo de impureza queda como un ion positivo fijo ($N_D^+$) en la red.
    Si $N_D \gg n_i$ es la concentración de impurezas donadoras:
    \begin{equation}
        n \approx N_D \gg n_i, \qquad p \approx \frac{n_i^2}{N_D} \ll n_i \implies n > p
    \end{equation}
    En un semiconductor tipo N, los \textbf{electrones son los portadores mayoritarios} y los \textbf{huecos son los portadores minoritarios}.

    \item \textbf{Semiconductor Extrínseco Tipo P (Dopado con Impurezas Aceptoras):}
    Se introducen átomos \textbf{trivalentes} (grupo III de la tabla periódica, con $3$ electrones de valencia), tales como \textbf{Boro ($\text{B}$)}, \textbf{Aluminio ($\text{Al}$)}, \textbf{Galio ($\text{Ga}$)} o \textbf{Indio ($\text{In}$)}. Al tener solo tres electrones para cuatro enlaces covalentes, queda un enlace incompleto que captura fácilmente un electrón de un átomo de Silicio vecino, generando un \textbf{hueco libre} en la banda de valencia sin crear un electrón en la banda de conducción. El átomo de impureza queda como un ion negativo fijo ($N_A^-$).
    Si $N_A \gg n_i$ es la concentración de impurezas aceptoras:
    \begin{equation}
        p \approx N_A \gg n_i, \qquad n \approx \frac{n_i^2}{N_A} \ll n_i \implies p > n
    \end{equation}
    En un semiconductor tipo P, los \textbf{huecos son los portadores mayoritarios} y los \textbf{electrones son los portadores minoritarios}.
\end{itemize}

\subsection{Mecanismos de Conducción en Semiconductores}
El movimiento neto de cargas en un semiconductor se produce por dos mecanismos físicos independientes:
\begin{enumerate}[leftmargin=*]
    \item \textbf{Corriente de Arrastre (o Deriva, \textit{Drift}):} Originada por la presencia de un campo eléctrico $\vec{E}$ que acelera los huecos en el sentido del campo y los electrones en sentido opuesto. La densidad de corriente total de arrastre y la conductividad $\sigma$ del semiconductor son:
    \begin{equation}
        \vec{J}_{\text{arrastre}} = \sigma \vec{E} = q (n \mu_n + p \mu_p) \vec{E}
    \end{equation}
    donde $\mu_n$ y $\mu_p$ son las movilidades de electrones y huecos respectivamente (en el Silicio $\mu_n \approx 1350\ \si{\centi\meter\squared\per\volt\per\second} > \mu_p \approx 480\ \si{\centi\meter\squared\per\volt\per\second}$, por lo que los electrones tienen mayor movilidad que los huecos).
    \item \textbf{Corriente de Difusión:} Originada por la agitación térmica cuando existe un \textbf{gradiente espacial de concentración} de portadores ($\frac{dn}{dx}$ o $\frac{dp}{dx}$). Los portadores tienden estadísticamente a desplazarse desde las regiones de alta concentración hacia las de baja concentración:
    \begin{equation}
        J_{n,\text{dif}} = q D_n \frac{dn}{dx}, \qquad J_{p,\text{dif}} = -q D_p \frac{dp}{dx}
    \end{equation}
    donde $D_n$ y $D_p$ son las constantes de difusión, ligadas a las movilidades mediante la relación de Einstein: $\frac{D_n}{\mu_n} = \frac{D_p}{\mu_p} = V_T = \frac{k_B T}{q}$.
\end{enumerate}

\subsection{La Unión PN en Equilibrio y bajo Polarización}
Cuando sobre un mismo monocristal de silicio se crean adyacentes una región dopada tipo P y una región dopada tipo N, la frontera metalúrgica constituye una \textbf{unión PN}:
\begin{enumerate}[leftmargin=*]
    \item \textbf{Unión PN en Circuito Abierto (Equilibrio Térmico):}
    Debido al enorme gradiente de concentración en la frontera, los huecos mayoritarios de la zona P difunden hacia la zona N y los electrones mayoritarios de la zona N difunden hacia la zona P. Al cruzarse en las inmediaciones de la unión, se recombinan entre sí, desapareciendo los portadores móviles libres en una franja central llamada \textbf{Zona de Carga Espacial (ZCE)}, \textbf{zona de deplexión} o de vaciamiento.
    Al desaparecer los portadores móviles en la ZCE, quedan al descubierto los \textbf{iones fijos de la red}: iones negativos aceptores ($N_A^-$) en el lado P e iones positivos donadores ($N_D^+$) en el lado N. Esta distribución dipolar crea un \textbf{campo eléctrico interno} $\vec{E}_{int}$ dirigido de la zona N hacia la zona P, el cual se opone a que sigan difundiendo más mayoritarios y establece una \textbf{barrera de potencial de contacto} $V_0$ (o $V_{bi}$):
    \begin{equation}
        V_0 = V_T \ln\left(\frac{N_A N_D}{n_i^2}\right) \approx 0{,}6\text{ a }0{,}75\ \si{\volt} \quad (\text{en Silicio a } 300\ \si{\kelvin})
    \end{equation}
    siendo $V_T = \frac{k_B T}{q} \approx 26\ \si{\milli\volt}$ la tensión térmica a temperatura ambiente ($T = 300\ \si{\kelvin}$). En equilibrio, la corriente de difusión de mayoritarios queda exactamente compensada por la corriente de arrastre de minoritarios ($I_{\text{neta}} = 0$).
\end{enumerate}

\begin{center}
\begin{tikzpicture}[scale=1.05]
    % Bloque Semiconductor más ancho para holgura total
    \draw[thick] (-4.8,-1.3) rectangle (4.8,1.3);
    % Separador de unión
    \draw[thick, dashed] (0,-1.3) -- (0,1.3);
    % Zona P
    \node[font=\large\bfseries, fftPrimary] at (-3.1,0.7) {Zona P};
    \node[font=\footnotesize] at (-3.1,-0.7) {Huecos ($h^+$)};
    % Zona N
    \node[font=\large\bfseries, fftDarkBlue] at (3.1,0.7) {Zona N};
    \node[font=\footnotesize] at (3.1,-0.7) {Electrones ($e^-$)};
    % ZCE sombreada
    \draw[fill=fftAccent!15, draw=fftAccent, thick] (-1.3,-1.3) rectangle (1.3,1.3);
    \node[above=3pt, font=\bfseries\footnotesize, fftAccent] at (0,1.3) {Zona de Carga Espacial (ZCE)};
    % Iones fijos
    \node[fftRed] at (-0.65,0.35) {$\mathbf{N_A^-}$};
    \node[fftRed] at (-0.65,-0.35) {$\mathbf{N_A^-}$};
    \node[fftGreen] at (0.65,0.35) {$\mathbf{N_D^+}$};
    \node[fftGreen] at (0.65,-0.35) {$\mathbf{N_D^+}$};
    % Campo eléctrico interno
    \draw[->, >=stealth, very thick, fftDarkBlue] (0.85,-0.95) -- (-0.85,-0.95) node[midway, below] {$\vec{E}_{int}$};
\end{tikzpicture}
\end{center}

\begin{itemize}[leftmargin=*]
    \item \textbf{Unión PN en Polarización Inversa ($V_D < 0$):}
    Se aplica una fuente externa con el polo positivo en la zona N y el negativo en la zona P. El campo externo tiene el mismo sentido que $\vec{E}_{int}$, reforzando la barrera de potencial hasta $V_0 + |V_D|$ y \textbf{ensanchando la zona de deplexión}. La difusión de portadores mayoritarios queda completamente bloqueada y solo circula una debilísima \textbf{corriente inversa de saturación} $I_S \approx 10^{-14}\ \si{\ampere}$ ($10^{-11}\ \si{\milli\ampere}$) debida al arrastre de portadores minoritarios generados térmicamente.
    \item \textbf{Unión PN en Polarización Directa ($V_D > 0$):}
    Se conecta el polo positivo externo a la zona P y el negativo a la zona N. El campo externo se opone a $\vec{E}_{int}$, reduciendo la barrera de potencial a $V_0 - V_D$ y \textbf{estrechando la zona de deplexión}. Cuando $V_D$ supera el umbral $V_\gamma \approx 0{,}6 - 0{,}7\ \si{\volt}$, la barrera prácticamente desaparece y se produce una inyección masiva exponencial de portadores mayoritarios a través de la unión.
\end{itemize}

% ==============================================================================
% 4.2. EL DIODO Y CIRCUITOS BÁSICOS
% ==============================================================================
\section{El Diodo y Circuitos Básicos}

Un \textbf{diodo de unión} es una unión PN a la que se le añaden dos contactos óhmicos exteriores: el \textbf{Ánodo (A)} en la región P y el \textbf{Cátodo (K)} en la región N. Su tensión en bornes se define siempre de ánodo a cátodo ($V_D = V_{\text{Ánodo}} - V_{\text{Cátodo}}$) y su corriente positiva $I_D$ fluye de ánodo a cátodo (en el sentido de la flecha del símbolo).

\subsection{Ecuación de Shockley y Curva Característica Real}
\begin{ley}[Ecuación de Shockley del Diodo]
En primera aproximación, la característica corriente-tensión ($I_D - V_D$) del diodo en las zonas directa e inversa viene descrita por la ecuación exponencial de Shockley:
\begin{equation}
    I_D = I_S \left( e^{\frac{V_D}{n V_T}} - 1 \right)
\end{equation}
donde:
\begin{itemize}[leftmargin=*]
    \item $I_S$: corriente inversa de saturación (típicamente $I_S \approx 10^{-14}\ \si{\ampere} = 10^{-11}\ \si{\milli\ampere}$).
    \item $V_T = \frac{k T}{q}$: tensión térmica, donde $k = 1{,}38 \times 10^{-23}\ \si{\joule\per\kelvin}$ es la constante de Boltzmann, $T$ la temperatura absoluta en Kelvin y $q = 1{,}6 \times 10^{-19}\ \si{\coulomb}$ la carga del electrón. A temperatura ambiente ($T = 300\ \si{\kelvin}$), $V_T \approx 26\ \si{\milli\volt} = 0{,}026\ \si{\volt}$.
    \item $n$: factor de idealidad (en los apuntes de la asignatura se toma habitualmente $n = 1$).
\end{itemize}
\end{ley}

En la curva característica real de un diodo se distinguen tres regiones:
\begin{enumerate}[leftmargin=*]
    \item \textbf{Zona Directa ($V_D > 0$):} Para $V_D < V_\gamma \approx 0{,}6 - 0{,}7\ \si{\volt}$ la corriente es casi despreciable; al superar la tensión umbral $V_\gamma$, la intensidad crece exponencialmente con pendientes muy abruptas. En directa, el diodo no puede soportar tensiones mucho mayores que $V_\gamma$ porque la enorme corriente disiparía una potencia excesiva y fundiría el dispositivo.
    \item \textbf{Zona Inversa ($-V_{ZK} < V_D < 0$):} Como $e^{V_D/V_T} \ll 1$, la corriente vale $I_D \approx -I_S \approx 0$. El diodo actúa prácticamente como un circuito abierto.
    \item \textbf{Zona de Ruptura o Avalancha ($V_D \le -V_{ZK}$):} Si la tensión inversa supera un valor crítico denominado tensión de ruptura ($V_Z$ o $V_{BR}$, del orden de $-100\ \si{\volt}$ en diodos rectificadores comunes y entre $-3\ \si{\volt}$ y $-15\ \si{\volt}$ en diodos Zener), se produce una multiplicación por avalancha o efecto túnel Zener y circula una gran corriente inversa mientras la tensión permanece casi fija en $-V_Z$.
\end{enumerate}

\begin{center}
\begin{tikzpicture}[scale=1.1]
    % Ejes
    \draw[->, thick] (-3.2,0) -- (3.2,0) node[right] {$V_D\ [\si{\volt}]$};
    \draw[->, thick] (0,-1.2) -- (0,3.0) node[above] {$I_D\ [\si{\milli\ampere}]$};
    % Curva real Shockley + ruptura
    \draw[very thick, fftDarkBlue, domain=-2.2:1.1, samples=60]
        plot (\x, {0.05 * (exp(3.5*\x) - 1)});
    \draw[very thick, fftDarkBlue] (-2.2,-0.05) -- (-2.4,-1.1);
    \node[above left, font=\footnotesize, fftDarkBlue] at (-2.2,0) {$-V_Z$};
    \node[right, fftDarkBlue] at (1.1,2.5) {Curva Real};
    % Modelo lineal 2 (Vgamma = 0.7V)
    \draw[very thick, fftAccent, dashed] (-2.2,0) -- (0.7,0) -- (0.7,2.8);
    \draw[dashed, gray] (0.7,0) node[below=2pt, text=black, font=\footnotesize] {$V_\gamma{=}0{,}7\,\si{\volt}$} -- (0.7,0.5);
    \node[right=4pt, fftAccent, font=\bfseries\footnotesize] at (0.7,1.6) {Modelo $V_\gamma$};
\end{tikzpicture}
\end{center}

\subsection{Los 3 Modelos Lineales por Tramos del Diodo}
Dado que la ecuación de Shockley es no lineal, para resolver circuitos de forma rápida y sistemática sustituimos el diodo por uno de los tres \textbf{modelos lineales por tramos}:

\begin{metodo}[Los 3 Modelos Lineales del Diodo]
\begin{enumerate}[leftmargin=*]
    \item \textbf{Modelo 1: Diodo Ideal (no se suele usar salvo si $V \gg V_\gamma$):}
    \begin{itemize}
        \item \textbf{Diodo OFF (Corte):} Si $V_D < 0 \implies I_D = 0$ (equivalente a un \textbf{circuito abierto}).
        \item \textbf{Diodo ON (Conducción):} Si $I_D > 0 \implies V_D = 0\ \si{\volt}$ (equivalente a un \textbf{cortocircuito / cable ideal}).
    \end{itemize}
    \item \textbf{Modelo 2: Caída de Tensión Constante $V_\gamma$ (el más utilizado con resistencias grandes):}
    \begin{itemize}
        \item \textbf{Diodo OFF (Corte):} Si $V_D < V_\gamma \implies I_D = 0$ (\textbf{circuito abierto}).
        \item \textbf{Diodo ON (Conducción):} Si $I_D > 0 \implies V_D = V_\gamma$ constante (habitualmente $V_\gamma = 0{,}7\ \si{\volt}$ o $0{,}65\ \si{\volt}$). Se sustituye por una \textbf{fuente de tensión independiente $V_\gamma$} con el polo $+$ en el ánodo y el $-$ en el cátodo.
    \end{itemize}
    \item \textbf{Modelo 3: Tensión Variable con Resistencia Serie $r_d$ (usado con resistencias pequeñas):}
    \begin{itemize}
        \item \textbf{Diodo OFF (Corte):} Si $V_D < V_\gamma \implies I_D = 0$ (\textbf{circuito abierto}).
        \item \textbf{Diodo ON (Conducción):} Si $I_D > 0$, el diodo se sustituye por una fuente $V_\gamma$ en serie con una resistencia interna $r_d$ (de unos pocos ohmios, p. ej. $10 - 20\ \si{\ohm}$):
        \begin{equation}
            V_D = V_\gamma + I_D r_d \iff I_D = \frac{V_D - V_\gamma}{r_d} \quad \left(\text{recta de pendiente } \frac{1}{r_d}\right)
        \end{equation}
    \end{itemize}
\end{enumerate}
\end{metodo}

\begin{center}
\begin{tikzpicture}[scale=0.92]
    % Modelo 1: Ideal (x = -5.2)
    \node[font=\bfseries] at (-5.2,1.9) {1. Modelo Ideal};
    \draw (-6.4,0.7) to[D, l=Diodo] (-4.0,0.7);
    \draw (-6.4,-0.6) to[short, o-o] (-4.0,-0.6);
    \node[below, yshift=-12pt, font=\footnotesize] at (-5.2,-0.6) {ON: Corto ($V_D = 0$)};
    \draw (-6.4,-2.1) to[short, o-o] (-5.5,-2.1);
    \draw (-4.9,-2.1) to[short, o-o] (-4.0,-2.1);
    \node[below, yshift=-8pt, font=\footnotesize] at (-5.2,-2.1) {OFF: Abierto ($I_D = 0$)};

    % Modelo 2: Caída constante Vgamma (x = 0)
    \node[font=\bfseries] at (0,1.9) {2. Modelo $V_\gamma$};
    \draw (-1.2,0.7) to[D, l=Diodo] (1.2,0.7);
    \draw (-1.2,-0.6) to[V, l=$V_\gamma$, o-o] (1.2,-0.6);
    \node[below, yshift=-12pt, font=\footnotesize] at (0,-0.6) {ON ($V_D = V_\gamma \approx 0{,}7\,\si{\volt}$)};
    \draw (-1.2,-2.1) to[short, o-o] (-0.3,-2.1);
    \draw (0.3,-2.1) to[short, o-o] (1.2,-2.1);
    \node[below, yshift=-8pt, font=\footnotesize] at (0,-2.1) {OFF ($I_D = 0$)};

    % Modelo 3: Resistencia interna rd (x = 5.2)
    \node[font=\bfseries] at (5.2,1.9) {3. Modelo con $r_d$};
    \draw (3.8,0.7) to[D, l=Diodo] (6.6,0.7);
    \draw (3.6,-0.6) to[V, l=$V_\gamma$, o-] (5.2,-0.6) to[R, l=$r_d$, -o] (6.8,-0.6);
    \node[below, yshift=-12pt, font=\footnotesize] at (5.2,-0.6) {ON ($V_D = V_\gamma + r_d I_D$)};
    \draw (3.8,-2.1) to[short, o-o] (4.9,-2.1);
    \draw (5.5,-2.1) to[short, o-o] (6.6,-2.1);
    \node[below, yshift=-8pt, font=\footnotesize] at (5.2,-2.1) {OFF ($I_D = 0$)};
\end{tikzpicture}
\end{center}

\subsection{Los 3 Métodos de Resolución de un Circuito con Diodo}
Consideremos el circuito elemental formado por una fuente continua $V_{DD} = 5\ \si{\volt}$ en serie con una resistencia $R = 1\ \si{\kilo\ohm}$ y un diodo ($I_S = 10^{-11}\ \si{\milli\ampere} = 10^{-14}\ \si{\ampere}$, $V_T = 26\ \si{\milli\volt}$).

\begin{aviso}[Convenio de Unidades en Electrónica]
En todos los problemas de diodos y transistores trabajaremos sistemáticamente con:
\begin{equation}
    V \text{ en voltios } [\si{\volt}], \qquad I \text{ en miliamperios } [\si{\milli\ampere}], \qquad R \text{ en kiloohmios } [\si{\kilo\ohm}]
\end{equation}
ya que $\si{\volt} = \si{\milli\ampere} \times \si{\kilo\ohm}$ es un sistema de unidades internamente consistente.
\end{aviso}

\begin{enumerate}[leftmargin=*]
    \item \textbf{Método Analítico (con el diodo real):}
    Planteamos simultáneamente la ecuación de malla del circuito y la ecuación de Shockley del dispositivo:
    \begin{equation}
        \left. \begin{aligned}
            \text{Circuito:} \quad V_{DD} &= I_D R + V_D \\
            \text{Dispositivo:} \quad I_D &= I_S \left(e^{V_D / V_T} - 1\right)
        \end{aligned} \right\} \implies V_D = V_{DD} - R I_S \left(e^{V_D / V_T} - 1\right)
    \end{equation}
    Es una ecuación trascendente que se resuelve por \textbf{aproximaciones sucesivas} (iteración numérica) probando valores de $V_D$ entre $0{,}6\ \si{\volt}$ y $0{,}8\ \si{\volt}$ (por ejemplo $V_D^{(0)} = 0{,}65\ \si{\volt}$). Para $V_{DD} = 5\ \si{\volt}$ y $R = 1\ \si{\kilo\ohm}$ converge a:
    \begin{equation}
        V_D = 0{,}696\ \si{\volt}, \qquad I_D = \frac{5 - 0{,}696}{1\ \si{\kilo\ohm}} = 4{,}304\ \si{\milli\ampere} \approx 4{,}3\ \si{\milli\ampere}
    \end{equation}

    \item \textbf{Método Gráfico (Recta de Carga y Punto de Funcionamiento $Q$):}
    Despejando $I_D$ de la ecuación de malla obtenemos la \textbf{recta de carga estática}:
    \begin{equation}
        I_D = -\frac{1}{R} V_D + \frac{V_{DD}}{R}
    \end{equation}
    Es una recta de pendiente $-\frac{1}{R}$ que corta al eje vertical en $\left(0, \frac{V_{DD}}{R}\right)$ y al eje horizontal en $(V_{DD}, 0)$. La intersección gráfica entre la recta de carga y la curva característica del diodo determina el \textbf{punto de trabajo o punto de operación $Q(V_{DQ}, I_{DQ})$}.

    \begin{center}
    \begin{tikzpicture}[scale=1.0]
        \draw[->, thick] (0,0) -- (4.8,0) node[right] {$V_D\ [\si{\volt}]$};
        \draw[->, thick] (0,0) -- (0,3.4) node[above] {$I_D\ [\si{\milli\ampere}]$};
        % Curva del diodo
        \draw[very thick, fftDarkBlue, domain=0:1.65, samples=50]
            plot (\x, {0.06 * (exp(2.4*\x) - 1)});
        \node[right, font=\footnotesize, fftDarkBlue] at (1.65,3.0) {$I_D = I_S(e^{V_D/V_T}-1)$};
        % Recta de carga
        \draw[very thick, fftAccent] (0,2.8) node[left, text=black, font=\footnotesize] {$\dfrac{V_{DD}}{R}$}
            -- (4.0,0) node[below, text=black, font=\footnotesize] {$V_{DD}$};
        \node[above right, font=\footnotesize, fftAccent] at (2.1,1.3) {Recta de carga $\left(\text{pendiente } -\dfrac{1}{R}\right)$};
        % Punto Q
        \filldraw[fftRed] (1.38,1.83) circle (2.2pt);
        \node[above left, font=\bfseries\footnotesize, fftRed] at (1.38,1.88) {Punto $Q$};
        \draw[dashed, gray] (1.38,0) node[below, text=black, font=\footnotesize] {$V_{DQ}$} -- (1.38,1.83);
        \draw[dashed, gray] (0,1.83) node[left, text=black, font=\footnotesize] {$I_{DQ}$} -- (1.38,1.83);
    \end{tikzpicture}
    \end{center}

    \item \textbf{Método Lineal por Hipótesis y Comprobación (el más utilizado):}
    \begin{enumerate}[label=\alph*)]
        \item \textbf{Paso 1 (Hipótesis):} Suponer el estado de cada diodo (\textbf{OFF} u \textbf{ON}). Si no se ve claro a simple vista, es muy cómodo empezar suponiendo que los diodos están en \textbf{OFF} (circuito abierto) y calcular las tensiones en sus bornes para ver cuál tiene mayor tensión ánodo-cátodo.
        \item \textbf{Paso 2 (Sustitución y cálculo):} Sustituir cada diodo según el modelo lineal indicado ($2.^{\circ}$ o $3.^{\text{er}}$ modelo) y resolver el circuito lineal resultante.
        \item \textbf{Paso 3 (Comprobación obligatoria):}
        \begin{itemize}
            \item Si se supuso \textbf{Diodo OFF}: debe cumplirse que $V_D = V_{\text{Ánodo}} - V_{\text{Cátodo}} < V_\gamma$.
            \item Si se supuso \textbf{Diodo ON}: debe cumplirse que la corriente $I_D$ circule de ánodo a cátodo con valor estrictamente positivo ($I_D > 0$).
        \end{itemize}
        Si no coincide con la suposición, la hipótesis es falsa y se cambia de estado.
    \end{enumerate}
    Aplicado al circuito con $V_{DD} = 5\ \si{\volt}$ y $R = 1\ \si{\kilo\ohm}$ ($V_\gamma = 0{,}65\ \si{\volt}, r_d = 10\ \si{\ohm} = 0{,}01\ \si{\kilo\ohm}$):
    \begin{itemize}
        \item Si suponemos $D\text{ OFF}$: $I_D = 0 \implies V_R = 0 \implies V_D = 5\ \si{\volt} > 0{,}65\ \si{\volt} \implies$ contradicción (la suposición OFF es falsa; el diodo conduce en ON).
        \item Con el \textbf{2.º modelo} ($D\text{ ON} \equiv V_\gamma = 0{,}65\ \si{\volt}$):
        \begin{equation}
            I_D = \frac{V_{DD} - V_\gamma}{R} = \frac{5 - 0{,}65}{1\ \si{\kilo\ohm}} = 4{,}35\ \si{\milli\ampere} > 0, \qquad V_D = 0{,}65\ \si{\volt}
        \end{equation}
        \item Con el \textbf{3.er modelo} ($D\text{ ON} \equiv V_\gamma = 0{,}65\ \si{\volt}$ en serie con $r_d = 0{,}01\ \si{\kilo\ohm}$):
        \begin{equation}
            I_D = \frac{V_{DD} - V_\gamma}{R + r_d} = \frac{5 - 0{,}65}{1 + 0{,}01} = 4{,}307\ \si{\milli\ampere} > 0, \qquad V_D = V_\gamma + I_D r_d = 0{,}693\ \si{\volt}
        \end{equation}
    \end{itemize}
\end{enumerate}

\subsection{Aplicaciones Clásicas del Diodo: Limitadores, Rectificadores, Filtros y Zener}

\subsubsection{1. Circuitos Limitadores o Recortadores (\textit{Clippers})}
Los circuitos limitadores se emplean para recortar una porción de la señal de entrada por encima o por debajo de un nivel de referencia, protegiendo etapas posteriores.
\begin{ejemplo}[Circuito Limitador Polarizado de los Apuntes]
En el circuito de la figura, una entrada $V_i$ alimenta a través de una resistencia serie de $1\ \si{\kilo\ohm}$ el nudo de salida $V_o$, del cual cuelgan en paralelo una rama con un diodo $D$ ($V_\gamma = 0{,}65\ \si{\volt}$) en serie con una fuente continua de $7\ \si{\volt}$, y una resistencia de carga de $100\ \si{\kilo\ohm}$ a masa.

\begin{center}
\begin{tikzpicture}[scale=0.9]
    \draw (0,0) node[ground] {} to[V, v=$V_i$] (0,2.4)
          to[R, l=$1\,\si{\kilo\ohm}$] (2.6,2.4) -- (4.8,2.4);
    \draw (4.8,2.4) to[short, -o] (5.5,2.4) node[right] {$V_o$};
    % Rama del diodo y fuente de 7V
    \draw (2.6,2.4) to[D, l=$D$] (2.6,1.2)
          to[V, l=$7\,\si{\volt}$] (2.6,0) node[ground] {};
    % Rama resistencia 100k
    \draw (4.4,2.4) to[R, l=$100\,\si{\kilo\ohm}$] (4.4,0) node[ground] {};
\end{tikzpicture}
\end{center}

\begin{itemize}[leftmargin=*]
    \item \textbf{Tramo a) Diodo en Corte (OFF):} Para que el diodo conduzca se necesita que el potencial de su ánodo ($V_o$) supere el potencial de su cátodo ($7\ \si{\volt}$) en al menos $V_\gamma = 0{,}65\ \si{\volt}$, es decir, $V_o \ge 7 + 0{,}65 = 7{,}65\ \si{\volt}$. Mientras $V_o < 7{,}65\ \si{\volt}$, el diodo está en \textbf{OFF} (abierto) y el circuito actúa como un simple divisor de tensión resistivo:
    \begin{equation}
        V_o = \frac{100\ \si{\kilo\ohm}}{1\ \si{\kilo\ohm} + 100\ \si{\kilo\ohm}} V_i = \frac{100}{101} V_i
    \end{equation}
    \item \textbf{Tramo b) Umbral de Conducción:} El diodo entra en conducción exactamente cuando $V_o = 7{,}65\ \si{\volt}$, lo que corresponde a una tensión de entrada:
    \begin{equation}
        V_i = \frac{101}{100} V_o = \frac{101}{100} \times 7{,}65 = 7{,}7265\ \si{\volt} \approx 7{,}73\ \si{\volt}
    \end{equation}
    \item \textbf{Tramo c) Diodo en Conducción (ON, para $V_i \ge 7{,}73\ \si{\volt}$):}
    \begin{itemize}
        \item Con el \textbf{3.er modelo} ($V_\gamma = 0{,}65\ \si{\volt}$ y resistencia dinámica $r_d$): aplicando la ley de corrientes de Kirchhoff en el nudo $V_o$:
        \begin{equation}
            \frac{V_i - V_o}{1} = \frac{V_o - 7{,}65}{r_d} + \frac{V_o}{100} \implies V_o = \frac{7{,}65 \times 100}{100 + 101 r_d} + \frac{100 r_d}{100 + 101 r_d} V_i
        \end{equation}
        \item Con el \textbf{2.º modelo} ($r_d = 0$, el más sencillo): la rama del diodo fija rígidamente la tensión de salida en:
        \begin{equation}
            V_o = 7 + 0{,}65 = 7{,}65\ \si{\volt} \quad (\text{fijo})
        \end{equation}
        Cualquier aumento adicional de $V_i$ incrementa la corriente por la resistencia de $1\ \si{\kilo\ohm}$ y se desvía íntegramente por el diodo, manteniendo limitada la tensión $V_o$.
    \end{itemize}
\end{itemize}
\end{ejemplo}

\subsubsection{2. Rectificador de Media Onda}
Un circuito rectificador convierte una señal alterna sinusoidal ($v_i(t) = V_p \sin(\omega t)$, de valor medio nulo) en una señal unidireccional pulsante con componente continua no nula:
\begin{itemize}[leftmargin=*]
    \item \textbf{Con el 2.º modelo del diodo ($V_\gamma = 0{,}65 - 0{,}7\ \si{\volt}$):}
    \begin{align}
        \text{Si } v_i(t) < V_\gamma &\implies D\text{ OFF} \implies I_D = 0 \implies v_o(t) = I_D R = 0\ \si{\volt} \\
        \text{Si } v_i(t) \ge V_\gamma &\implies D\text{ ON} \implies v_o(t) = v_i(t) - V_\gamma
    \end{align}
    El valor de pico de salida es $V_{o,\max} = V_p - V_\gamma$ y su valor medio en continua es $V_{dc} = \frac{V_p - V_\gamma}{\pi}$.
    \item \textbf{Con el 1.er modelo (ideal):} Conduce durante todo el semiciclo positivo ($v_o(t) = v_i(t)$ para $v_i > 0$) y bloquea el semiciclo negativo ($v_o(t) = 0$ para $v_i < 0$).
    \item \textbf{Con el 3.er modelo ($V_{D0}$ y $r_D$):} Para $v_i(t) \ge V_{D0}$, forma un divisor resistivo de pendiente $\frac{R}{R + r_D}$, dando $v_o(t) = \frac{R}{R + r_D}\left(v_i(t) - V_{D0}\right)$.
\end{itemize}

\subsubsection{3. Rectificador de Onda Completa y Filtrado por Condensador}
Para aprovechar ambos semiciclos de la señal alterna se emplea el \textbf{rectificador de onda completa}, ya sea mediante un transformador de toma intermedia con 2 diodos, o mediante el \textbf{puente de Graetz de 4 diodos}:
\begin{itemize}[leftmargin=*]
    \item En cada semiciclo conducen \textbf{dos diodos en serie} ($D_1$ y $D_2$ cuando $v_i > 2V_\gamma$; $D_3$ y $D_4$ en el semiciclo opuesto cuando $|v_i| > 2V_\gamma$), de modo que la caída umbral total es $2V_\gamma \approx 1{,}3 - 1{,}4\ \si{\volt}$:
    \begin{equation}
        I = \frac{|v_i| - 2V_\gamma}{R_L}, \qquad V_{o,\max} = V_p - 2V_\gamma, \qquad V_{dc} = \frac{2(V_p - 2V_\gamma)}{\pi}
    \end{equation}
    Por ejemplo, con $V_\gamma = 0{,}7\ \si{\volt}$, cuando $V_i = 5\ \si{\volt}$ la salida toma el valor $V_o = 5 - 1{,}4 = 3{,}6\ \si{\volt}$.
    \item \textbf{Filtro por Condensador ($C$ en paralelo con $R_L$):} Al colocar un condensador en paralelo con la resistencia de carga, el condensador se carga hasta el valor de pico $V_{\max}$ cuando los diodos conducen y, cuando la tensión de entrada desciende, los diodos entran en corte y el condensador se descarga lentamente a través de $R_L$ con constante de tiempo $\tau = R_L C \gg T$. Se obtiene así una tensión casi continua con un pequeño \textbf{rizado} pico a pico:
    \begin{equation}
        V_r \approx \frac{V_{\max}}{f_{\text{rizado}} R_L C} = \frac{I_{dc}}{f_{\text{rizado}} C}
    \end{equation}
    donde $f_{\text{rizado}} = f_{\text{red}}$ en media onda y $f_{\text{rizado}} = 2 f_{\text{red}}$ en onda completa.
\end{itemize}

\begin{center}
\begin{tikzpicture}[scale=0.92]
    \node[font=\bfseries] at (1.5,3.2) {Rectificador de Onda Completa en Puente con Filtro $C$};
    % Fuente alterna a la izquierda conectada a (-2,1) y (2,1)
    \draw (-4.5,0.2) to[sV, v=$v_i(t)$] (-4.5,1.8) -- (-2,1.8) -- (-2,1);
    \draw (-4.5,0.2) -- (-4.5,-0.9) -- (2,-0.9) -- (2,1);

    % Puente en rombo limpio entre (-2,1), (0,2.5), (2,1), (0,-0.5)
    \draw (-2,1) to[D, l=$D_1$] (0,2.5);
    \draw (0,-0.5) to[D, l=$D_3$] (-2,1);
    \draw (2,1) to[D*, l=$D_2$] (0,2.5);
    \draw (0,-0.5) to[D*, l=$D_4$] (2,1);

    % Conexión al filtro C (x = 3.5) y carga RL (x = 5.5)
    \draw (0,2.5) -- (3.5,2.5) to[C, l=$C$] (3.5,-0.5) -- (0,-0.5);
    \draw (3.5,2.5) -- (5.5,2.5) to[R, l=$R_L$, v=$v_o(t)$] (5.5,-0.5) -- (3.5,-0.5);
    \draw (3.5,-0.5) node[ground] {};
\end{tikzpicture}
\end{center}

\subsubsection{4. El Diodo Zener, LED y Fotodiodo}
\begin{itemize}[leftmargin=*]
    \item \textbf{Diodo Zener:} Es un diodo diseñado específicamente para trabajar también en la \textbf{región de ruptura inversa} a partir de una tensión de ruptura Zener controlada $V_Z$ (p. ej. $|V_Z| = 3\ \si{\volt}, 5\ \si{\volt}, 10\ \si{\volt}$). En el 2.º modelo lineal presenta \textbf{tres estados posibles}:
    \begin{enumerate}
        \item \textbf{Conducción Directa ($D\text{ ON directo}$):} Si $I_D > 0$ (de ánodo a cátodo), se comporta como un diodo normal con $V_D = V_{\text{Ánodo}} - V_{\text{Cátodo}} = V_\gamma = 0{,}7\ \si{\volt}$.
        \item \textbf{Corte ($D\text{ OFF}$):} Si $-|V_Z| < V_D < V_\gamma$, no conduce ($I_D = I_Z = 0$, circuito abierto).
        \item \textbf{Conducción en Ruptura Inversa ($D\text{ ON inverso / Zener}$):} Si la tensión inversa en el cátodo respecto al ánodo alcanza $|V_Z|$ ($V_D = -|V_Z|$, es decir, $V_{\text{Cátodo}} - V_{\text{Ánodo}} = |V_Z|$), conduce una corriente inversa $I_Z > 0$ de cátodo a ánodo manteniendo fija la caída $|V_Z|$ (se sustituye por una fuente independiente de valor $|V_Z|$ con el polo $+$ en el cátodo y el $-$ en el ánodo).
    \end{enumerate}
    \item \textbf{Diodo Emisor de Luz (LED) y Fotodiodo:} En un \textbf{LED} (fabricado con semiconductores de gap directo como $\text{GaAsP}$ o $\text{GaN}$), la recombinación de electrones y huecos en polarización directa libera fotones de luz ($h\nu \approx E_g$), presentando una caída umbral mayor ($V_\gamma \approx 1{,}8 - 3{,}2\ \si{\volt}$). En cambio, un \textbf{fotodiodo} opera polarizado en inversa: al incidir luz sobre la zona de deplexión se generan pares electrón-hueco que incrementan proporcionalmente la corriente inversa.
\end{itemize}

\begin{ejemplo}[Limitador de Doble Nivel con Diodo Zener de los Apuntes]
En el circuito de la figura, las dos resistencias son iguales ($R_1$) y el diodo Zener tiene $V_\gamma = 0{,}7\ \si{\volt}$ y tensión de ruptura inversa $V_Z = -3\ \si{\volt}$ ($|V_Z| = 3\ \si{\volt}$), estando su ánodo conectado a una fuente de $5\ \si{\volt}$ y su cátodo al nudo de salida $V_o$.

\begin{center}
\begin{tikzpicture}[scale=0.9]
    \draw (0,0) node[ground] {} to[V, v=$V_i$] (0,2.4)
          to[R, l=$R_1$] (2.6,2.4) -- (4.8,2.4);
    \draw (4.8,2.4) to[short, -o] (5.5,2.4) node[right] {$V_o$};
    % Rama Zener (ánodo abajo hacia 5V, cátodo arriba en Vo)
    \draw (2.6,1.2) to[zzD, l=$D_Z$] (2.6,2.4);
    \draw (2.6,1.2) to[V, l=$5\,\si{\volt}$] (2.6,0) node[ground] {};
    % Rama R1 a masa
    \draw (4.4,2.4) to[R, l=$R_1$] (4.4,0) node[ground] {};
\end{tikzpicture}
\end{center}

\begin{enumerate}[leftmargin=*]
    \item \textbf{Zona en Corte (OFF):} Cuando el Zener no conduce, el circuito es un divisor de tensión con dos resistencias iguales $R_1$:
    \begin{equation}
        V_o = \frac{R_1}{R_1 + R_1} V_i = \frac{V_i}{2}
    \end{equation}
    Como el ánodo está a $5\ \si{\volt}$ y el cátodo a $V_o$, la tensión ánodo-cátodo es $V_D = 5 - V_o$. El Zener permanece en OFF mientras $-3\ \si{\volt} < 5 - V_o < 0{,}7\ \si{\volt}$, es decir:
    \begin{equation}
        4{,}3\ \si{\volt} < V_o < 8\ \si{\volt} \iff 8{,}6\ \si{\volt} < V_i < 16\ \si{\volt} \implies V_o = \frac{V_i}{2}
    \end{equation}
    \item \textbf{Conducción en Directa ($V_i \le 8{,}6\ \si{\volt}$):} Cuando $V_o$ intenta bajar de $4{,}3\ \si{\volt}$, el diodo entra en conducción directa fijando $V_o = 5 - 0{,}7 = 4{,}3\ \si{\volt}$.
    \item \textbf{Conducción en Ruptura Inversa ($V_i \ge 16\ \si{\volt}$):} Cuando $V_o$ intenta superar $8\ \si{\volt}$, el Zener entra en ruptura inversa fijando $V_o = 5 + 3 = 8\ \si{\volt}$.
\end{enumerate}
\end{ejemplo}

% ==============================================================================
% 4.3. EL TRANSISTOR MOSFET Y CIRCUITOS BÁSICOS
% ==============================================================================
\section{El Transistor MOSFET y Circuitos Básicos}

El transistor de efecto de campo metal-óxido-semiconductor (\textbf{MOSFET}) es un dispositivo electrónico de cuatro terminales en el que la corriente que circula por un canal semiconductor entre el \textbf{Drenador (D)} y la \textbf{Fuente (S)} se controla mediante la tensión aplicada al terminal de \textbf{Puerta (G)}.

\subsection{Estructura Física, Terminales y Símbolos}
En un transistor \textbf{NMOS de enriquecimiento}:
\begin{itemize}[leftmargin=*]
    \item Sobre un \textbf{sustrato semiconductor tipo P} (\textit{Bulk} o \textit{Body}, terminal \textbf{B}) se difunden dos regiones fuertemente dopadas tipo $n^+$ que constituyen la \textbf{Fuente (S, \textit{Source})} y el \textbf{Drenador (D, \textit{Drain})}.
    \item La región comprendida entre Fuente y Drenador (de longitud geométrica $L$ y anchura $W$) se recubre con una fina capa aislante de \textbf{dióxido de silicio ($\text{SiO}_2$, óxido de puerta)} sobre la que se deposita el electrodo conductor de \textbf{Puerta (G, \textit{Gate})}.
    \item Como la puerta está separada del canal por un dieléctrico aislante ($\text{SiO}_2$), \textbf{la corriente de puerta es siempre rigurosamente nula en continua}:
    \begin{equation}
        I_G = 0 \implies I_D = I_S
    \end{equation}
    \item Habitualmente el sustrato \textbf{B} se conecta internamente al terminal de Fuente \textbf{S} ($B = S$), pasando de la notación de cuatro terminales a la de tres terminales ($G, D, S$).
\end{itemize}

\textbf{Clasificación de los transistores MOS:}
\begin{enumerate}[leftmargin=*]
    \item \textbf{NMOS de Enriquecimiento ($V_T > 0$):} Con $V_{GS} = 0$ no existe canal entre D y S (interruptor abierto). Hace falta aplicar una tensión positiva $V_{GS} > V_T$ en la puerta para atraer electrones e inducir el canal tipo N.
    \item \textbf{PMOS de Enriquecimiento ($V_T < 0$):} Fabricado sobre sustrato tipo N con pozos $p^+$. Hace falta aplicar una tensión negativa en la puerta respecto a la fuente ($V_{GS} < V_{Tp} < 0$, o equivalentemente $V_{SG} > |V_{Tp}|$) para inducir un canal de huecos tipo P.
    \item \textbf{MOSFET de Deplexión (p. ej. NMOS con $V_T < 0$):} El canal físico ya está implantado de fábrica, por lo que conduce corriente incluso sin aplicar tensión en la puerta ($V_{GS} = 0$). Para apagarlo hay que aplicar una tensión negativa $V_{GS} < V_T < 0$ que vacíe el canal. En su símbolo el trazo vertical del canal se dibuja continuo y más grueso.
\end{enumerate}

\begin{center}
\begin{tikzpicture}[scale=1.0]
    % NMOS enriquecimiento
    \draw (-3.5,0) node[nmos] (nmos) {};
    \node[left] at (nmos.gate) {G (Puerta)};
    \draw (nmos.drain) node[above] {D (Drenador)};
    \draw (nmos.source) node[below] {S (Fuente)};
    \node[font=\small\bfseries] at (-3.5,-1.8) {NMOS ($V_T > 0$)};

    % PMOS enriquecimiento
    \draw (3.5,0) node[pmos] (pmos) {};
    \node[left] at (pmos.gate) {G (Puerta)};
    \draw (pmos.drain) node[below] {D (Drenador)};
    \draw (pmos.source) node[above] {S (Fuente)};
    \node[font=\small\bfseries] at (3.5,-1.8) {PMOS ($V_T < 0$)};
\end{tikzpicture}
\end{center}

\subsection{Regiones de Operación y Ecuaciones del NMOS y PMOS}
Definimos el \textbf{parámetro de transconductancia} $k_n$ (denotado en las hojas manuscritas como $K_n$ o $K$) que mide cuánto conduce el transistor en función de su tecnología y geometría:
\begin{equation}
    k_n = \mu_n C_{ox} \frac{W}{L} \quad \left[\si{\ampere\per\volt\squared} \text{ o } \si{\milli\ampere\per\volt\squared}\right]
\end{equation}
donde $\mu_n$ es la movilidad de los portadores en el canal, $C_{ox} = \frac{\varepsilon_{ox}}{t_{ox}}$ es la capacidad del óxido de puerta por unidad de área, $W$ la anchura del canal y $L$ la longitud del canal (en tecnologías actuales del orden de nanómetros).
\emph{Nota sobre notación:} En las hojas de problemas de la asignatura se escribe $I_D = \frac{K}{2}(V_{GS}-V_T)^2$ en saturación y $I_D = K\left[(V_{GS}-V_T)V_{DS} - \frac{1}{2}V_{DS}^2\right]$ en región lineal, donde el dato del enunciado $K \equiv k_n$.

\begin{ley}[Modos de Operación del Transistor NMOS ($V_{Tn} > 0$ en enriquecimiento)]
\begin{enumerate}[leftmargin=*]
    \item \textbf{1) Región de Corte (NMOS OFF):}
    \begin{equation}
        V_{GS} < V_{Tn} \implies I_D = 0
    \end{equation}
    El transistor equivale a un circuito abierto entre Drenador y Fuente ($I_D = I_S = 0$).

    \item \textbf{2) Región Lineal, Óhmica o Tríodo (NMOS ON con $V_{DS}$ pequeña):}
    \begin{equation}
        V_{GS} \ge V_{Tn} \quad \text{y} \quad V_{DS} < V_{GS} - V_{Tn}
    \end{equation}
    El canal está abierto de extremo a extremo y la corriente depende tanto de $V_{GS}$ como de $V_{DS}$:
    \begin{equation}
        I_{DS} = k_n \left[ (V_{GS} - V_{Tn}) V_{DS} - \frac{1}{2} V_{DS}^2 \right]
    \end{equation}

    \item \textbf{3) Región de Saturación (NMOS ON con $V_{DS}$ grande):}
    \begin{equation}
        V_{GS} \ge V_{Tn} \quad \text{y} \quad V_{DS} \ge V_{GS} - V_{Tn}
    \end{equation}
    El canal se estrangula junto al drenador (\textit{pinch-off}) y la corriente depende cuadráticamente solo de $V_{GS}$:
    \begin{equation}
        I_{DS} = \frac{k_n}{2} (V_{GS} - V_{Tn})^2
    \end{equation}
\end{enumerate}
\end{ley}

\begin{ley}[Modos de Operación del Transistor PMOS ($V_{Tp} < 0$ en enriquecimiento)]
En un transistor PMOS todas las desigualdades se invierten respecto al NMOS (o bien se trabaja con valores positivos invirtiendo los subíndices $V_{SG} = -V_{GS}$ y $V_{SD} = -V_{DS}$):
\begin{enumerate}[leftmargin=*]
    \item \textbf{1) Región de Corte (PMOS OFF):}
    \begin{equation}
        V_{GS} > V_{Tp} \quad (\text{es decir, } V_{SG} < |V_{Tp}|) \implies I_{SD} = 0
    \end{equation}
    \item \textbf{2) Región Lineal o Tríodo (PMOS ON):}
    \begin{equation}
        V_{GS} \le V_{Tp} \quad \text{y} \quad V_{DS} > V_{GS} - V_{Tp} \quad \left(\text{es decir, } V_{SD} < V_{SG} - |V_{Tp}|\right)
    \end{equation}
    \begin{equation}
        I_{SD} = k_p \left[ (V_{GS} - V_{Tp}) V_{DS} - \frac{1}{2} V_{DS}^2 \right] = k_p \left[ (V_{SG} - |V_{Tp}|) V_{SD} - \frac{1}{2} V_{SD}^2 \right]
    \end{equation}
    \item \textbf{3) Región de Saturación (PMOS ON):}
    \begin{equation}
        V_{GS} \le V_{Tp} \quad \text{y} \quad V_{DS} \le V_{GS} - V_{Tp} \quad \left(\text{es decir, } V_{SD} \ge V_{SG} - |V_{Tp}|\right)
    \end{equation}
    \begin{equation}
        I_{SD} = \frac{k_p}{2} (V_{GS} - V_{Tp})^2 = \frac{k_p}{2} (V_{SG} - |V_{Tp}|)^2
    \end{equation}
\end{enumerate}
\end{ley}

\subsection{Curvas Características y Cuadro Resumen de Resolución}
\begin{center}
\begin{tikzpicture}[scale=1.05]
    % Curva de transferencia ID vs VGS (izquierda)
    \draw[->, thick] (-5.8,0) -- (-1.8,0) node[right] {$V_{GS}$};
    \draw[->, thick] (-5.2,0) -- (-5.2,3.2) node[above] {$I_D$};
    \draw[very thick, fftRed] (-5.2,0) -- (-4.0,0);
    \draw[very thick, fftDarkBlue, domain=-4.0:-2.2, samples=40]
        plot (\x, {0.85*(\x + 4.0)*(\x + 4.0)});
    \node[below, font=\footnotesize] at (-4.0,0) {$V_T$};
    \node[above, font=\bfseries\footnotesize] at (-3.6,3.2) {Transferencia ($I_D - V_{GS}$)};

    % Curvas de salida ID vs VDS (derecha)
    \draw[->, thick] (0,0) -- (5.2,0) node[right] {$V_{DS}$};
    \draw[->, thick] (0,0) -- (0,3.2) node[above] {$I_D$};
    \node[above, font=\bfseries\footnotesize] at (2.6,3.2) {Salida ($I_D - V_{DS}$)};

    % Parábola límite Vds = Vgs - Vt
    \draw[dashed, fftAccent, thick, domain=0:2.5, samples=40]
        plot (\x, {0.42 * \x * \x});
    \node[above, fftAccent, font=\footnotesize] at (2.4,2.65) {$V_{DS} = V_{GS} - V_T$};

    % Curvas para varios Vgs
    \foreach \vgs/\col in {1.5/fftPrimary, 2.0/fftDarkBlue, 2.5/fftGreen} {
        \draw[thick, \col] (0,0) to[out=60, in=180] ({0.7*\vgs}, {0.35*\vgs*\vgs})
                          -- (4.6, {0.35*\vgs*\vgs}) node[right, font=\footnotesize] {$V_{GS\vgs}$};
    }
    \node[font=\bfseries\footnotesize, fftPrimary] at (0.8,1.8) {Lineal};
    \node[font=\bfseries\footnotesize, fftPrimary] at (3.4,1.1) {Saturación};
\end{tikzpicture}
\end{center}

\begin{aviso}[Cuadro Resumen de Resolución de Circuitos con Diodos y MOSFET]
\renewcommand{\arraystretch}{1.3}
\begin{center}
\resizebox{\linewidth}{!}{%
\begin{tabular}{>{\bfseries}l l l}
\toprule
Dispositivo & Estado / Región & Condición y Modelo Circuital \\
\midrule
Diodo (Modelo 1) & OFF / ON & $V_D < 0 \implies I_D = 0$ (abierto) \quad / \quad $I_D > 0 \implies V_D = 0$ (corto) \\
Diodo (Modelo 2) & OFF / ON & $V_D < V_\gamma \implies I_D = 0$ \quad / \quad $I_D > 0 \implies V_D = V_\gamma \approx 0{,}7\,\si{\volt}$ \\
Diodo (Modelo 3) & OFF / ON & $V_D < V_\gamma \implies I_D = 0$ \quad / \quad $I_D > 0 \implies V_D = V_\gamma + I_D r_d$ \\
Diodo Zener & Directo / OFF / Inverso & $V_D = 0{,}7\,\si{\volt}$ ($I_D>0$) \; / \; $-|V_Z|<V_D<0{,}7\,\si{\volt}$ \; / \; $V_D = -|V_Z|$ ($I_Z>0$) \\
\midrule
NMOS ($I_G = 0$) & 1) Corte (OFF) & $V_{GS} < V_{Tn} \implies I_D = 0$ \\
                 & 2) Lineal (Tríodo) & $V_{GS} > V_{Tn}$ y $V_{DS} < V_{GS}-V_{Tn} \implies I_D = k_n\left[(V_{GS}-V_{Tn})V_{DS} - \frac{1}{2}V_{DS}^2\right]$ \\
                 & 3) Saturación & $V_{GS} > V_{Tn}$ y $V_{DS} \ge V_{GS}-V_{Tn} \implies I_D = \frac{k_n}{2}(V_{GS}-V_{Tn})^2$ \\
\bottomrule
\end{tabular}%
}
\end{center}
\textbf{Regla de Oro para el MOSFET:}
\begin{enumerate}[leftmargin=*]
    \item Usar siempre que $I_G = 0$ para hallar $V_G$ (divisor de tensión en la puerta). Suponer inicialmente OFF ($I_D = 0$) para ver si $V_{GS} > V_T$.
    \item Cuando el transistor conduce ($V_{GS} > V_T$), \textbf{se suele suponer primero que está en SATURACIÓN}, se resuelven las 3 ecuaciones con 3 incógnitas ($I_D, V_{GS}, V_{DS}$) y se comprueba si $V_{DS} \ge V_{GS} - V_T$. Si no se cumple, está en la región \textbf{lineal} y se recalculan las ecuaciones con la fórmula de tríodo.
    \item \textbf{Caso especial (Puerta y Drenador cortocircuitados):} Si G y D están unidos (directamente o mediante una resistencia por la que no pasa corriente al ser $I_G = 0$), entonces $V_G = V_D \implies V_{DS} = V_{GS}$. Como $V_T > 0$, se cumple automáticamente que $V_{DS} > V_{GS} - V_T$, por lo que \textbf{si conduce, está garantizado en SATURACIÓN}.
\end{enumerate}
\end{aviso}

% ==============================================================================
% 4.4. COLECCIÓN COMPLETA DE EJERCICIOS RESUELTOS
% ==============================================================================
\section{Ejercicios Resueltos de Diodos y Transistores MOSFET}

% --- EJERCICIO 1 (Parcial 2 FFT) ---
\begin{ejercicio}[Circuito con Diodos de Dos Ramas --- Examen Parcial 2 FFT]
Determinar el estado de conducción (ON/OFF), las corrientes por cada rama y las caídas de tensión en el circuito de la figura empleando el segundo modelo lineal con $V_\gamma = 0{,}7\ \si{\volt}$.
\tcbline
\begin{center}
\begin{tikzpicture}[scale=0.92]
    \draw (0,0) to[V, v=$10\,\si{\volt}$] (0,2.5) -- (2,2.5) node[above] {$A$};
    \draw (2,2.5) to[D, l=$D_1$] (4,2.5)
          to[R, l=$R_1{=}1\,\si{\kilo\ohm}$, i=$I_{D1}$] (4,0) -- (0,0);
    \draw (2,2.5) -- (2,4) to[D, l=$D_2$] (4,4)
          to[V, l=$V_2{=}4\,\si{\volt}$] (6,4)
          to[R, l=$R_2{=}2\,\si{\kilo\ohm}$, i=$I_{D2}$] (6,0) -- (4,0);
    \draw (0,0) node[ground] {};
\end{tikzpicture}
\end{center}
\textbf{Solución Paso a Paso:}\\[4pt]
El potencial en el nudo $A$ común a ambos ánodos es $V_A = 10\ \si{\volt}$. Si suponemos ambos en OFF, $V_{K1} = 0\ \si{\volt}$ ($V_{D1} = 10\ \si{\volt} > 0{,}7\ \si{\volt}$) y $V_{K2} = 4\ \si{\volt}$ ($V_{D2} = 6\ \si{\volt} > 0{,}7\ \si{\volt}$). Por tanto, suponemos \textbf{$D_1$ ON y $D_2$ ON}:
\begin{align}
    V_{K1} &= V_A - V_\gamma = 10 - 0{,}7 = 9{,}3\ \si{\volt} \implies I_{D1} = \frac{9{,}3\ \si{\volt}}{1\ \si{\kilo\ohm}} = 9{,}3\ \si{\milli\ampere} > 0 \quad (\text{ON confirmado}) \\
    V_{K2} &= V_A - V_\gamma = 10 - 0{,}7 = 9{,}3\ \si{\volt} \implies I_{D2} = \frac{9{,}3 - 4}{2\ \si{\kilo\ohm}} = 2{,}65\ \si{\milli\ampere} > 0 \quad (\text{ON confirmado})
\end{align}
La corriente total suministrada por la fuente de $10\ \si{\volt}$ es $I_{\text{total}} = I_{D1} + I_{D2} = 9{,}3 + 2{,}65 = 11{,}95\ \si{\milli\ampere}$.
\end{ejercicio}

% --- EJERCICIO 2 (Ejercicio 4 de los apuntes) ---
\begin{ejercicio}[Característica de Transferencia de Limitador Simple --- Ejercicio 4 Apuntes]
Calcular y representar la característica de transferencia $V_o$ en función de $V_i$ si se toma la salida en bornes del diodo usando el 2.º modelo lineal ($V_\gamma = 0{,}7\ \si{\volt}$). Ambas resistencias valen $R = 1\ \si{\kilo\ohm}$.
\tcbline
\begin{center}
\begin{tikzpicture}[scale=0.9]
    \draw (0,0) node[ground] {} to[V, v=$V_i$] (0,2.2)
          to[R, l=$1\,\si{\kilo\ohm}$] (2.4,2.2) -- (4.2,2.2);
    \draw (2.4,2.2) to[R, l=$1\,\si{\kilo\ohm}$] (2.4,0) node[ground] {};
    \draw (4.2,2.2) to[D, l=$D$] (4.2,0) node[ground] {};
    \draw (4.2,2.2) to[short, -o] (5.0,2.2) node[right] {$V_o$};
\end{tikzpicture}
\end{center}
\textbf{Solución Paso a Paso:}\\[4pt]
Como el cátodo está a masa ($0\ \si{\volt}$), la tensión en el diodo es $V_D = V_{\text{Ánodo}} - V_{\text{Cátodo}} = V_o$.
\begin{itemize}[leftmargin=*]
    \item \textbf{Diodo OFF ($V_o < V_\gamma = 0{,}7\ \si{\volt}$):} El diodo es un circuito abierto y las dos resistencias de $1\ \si{\kilo\ohm}$ forman un divisor de tensión:
    \begin{equation}
        V_o = V_i \frac{R}{R + R} = \frac{V_i}{2}
    \end{equation}
    El diodo empieza a conducir cuando $V_o = 0{,}7\ \si{\volt} \implies V_i = 2 V_o = 2 \times 0{,}7 = 1{,}4\ \si{\volt}$.
    \item \textbf{Resumen de tramos:}
    \begin{itemize}
        \item Para $V_i < 1{,}4\ \si{\volt}$: \textbf{Diodo OFF} $\implies V_o = \dfrac{V_i}{2}$ (recta de pendiente $1/2$ que pasa por el origen).
        \item Para $V_i \ge 1{,}4\ \si{\volt}$: \textbf{Diodo ON} $\implies V_o = V_\gamma = 0{,}7\ \si{\volt}$ (constante fija sea cual sea $V_i$).
    \end{itemize}
\end{itemize}
\end{ejercicio}

% --- EJERCICIO 3 (Ejercicio 5 de los apuntes) ---
\begin{ejercicio}[Característica de Transferencia con Dos Ramas Antiparalelas --- Ejercicio 5 Apuntes]
Calcular y dibujar la característica de transferencia $V_o(V_i)$ del circuito de la figura usando el 2.º modelo lineal con $V_\gamma = 0{,}7\ \si{\volt}$ para $D_1$ y $D_2$. Todas las resistencias son de $R = 1\ \si{\kilo\ohm}$.
\tcbline
\begin{center}
\begin{tikzpicture}[scale=0.9]
    \draw (-0.5,2.4) node[left] {$V_i$} to[R, l=$1\,\si{\kilo\ohm}$, o-] (2.0,2.4) -- (4.2,2.4);
    \draw (4.2,2.4) to[short, -o] (5.0,2.4) node[right] {$V_o$};
    % Rama D1 hacia abajo
    \draw (2.0,2.4) to[D, l=$D_1$] (2.0,1.2) to[R, l=$1\,\si{\kilo\ohm}$] (2.0,0) node[ground] {};
    % Rama D2 hacia arriba
    \draw (3.8,1.2) to[D, l=$D_2$] (3.8,2.4);
    \draw (3.8,1.2) to[R, l=$1\,\si{\kilo\ohm}$] (3.8,0) node[ground] {};
\end{tikzpicture}
\end{center}
\textbf{Solución Paso a Paso:}\\[4pt]
\begin{enumerate}[label=\Alph*)]
    \item \textbf{Región A ($-V_\gamma \le V_i \le V_\gamma$, es decir, $-0{,}7\ \si{\volt} \le V_i \le 0{,}7\ \si{\volt}$):}
    Tanto $D_1$ como $D_2$ están en \textbf{OFF}. No circula corriente por ninguna rama ($I = 0$), luego no hay caída en la resistencia de entrada y:
    \begin{equation}
        V_o = V_i \quad (\text{recta de pendiente } 1)
    \end{equation}
    \item \textbf{Región B ($V_i > V_\gamma = 0{,}7\ \si{\volt}$):}
    Conduce \textbf{$D_1$ ON} y \textbf{$D_2$ OFF}. La corriente que circula por la malla única (resistencia de entrada, $D_1$ y su resistencia inferior $R$) es:
    \begin{equation}
        I = \frac{V_i - V_\gamma}{2R} \implies V_o = V_\gamma + I R = V_\gamma + \frac{V_i - V_\gamma}{2} = \frac{V_i + V_\gamma}{2} = \frac{V_i + 0{,}7}{2} \quad \left(\text{pendiente } \frac{1}{2}\right)
    \end{equation}
    \item \textbf{Región C ($V_i < -V_\gamma = -0{,}7\ \si{\volt}$):}
    Conduce \textbf{$D_2$ ON} y \textbf{$D_1$ OFF}. La corriente de malla es:
    \begin{equation}
        I = \frac{V_i + V_\gamma}{2R} \implies V_o = -V_\gamma + I R = \frac{V_i - V_\gamma}{2} = \frac{V_i - 0{,}7}{2} \quad \left(\text{pendiente } \frac{1}{2}\right)
    \end{equation}
\end{enumerate}
\end{ejercicio}

% --- EJERCICIO 4 (Ejercicio 6 de los apuntes) ---
\begin{ejercicio}[Red en Escalera con Dos Diodos Zener --- Ejercicio 6 Apuntes]
En el circuito de la figura, las tensiones inversas de ruptura de los diodos son $V_{Z1} = 10\ \si{\volt}$ y $V_{Z2} = 8\ \si{\volt}$, con $V_i = 20\ \si{\volt}$, $R_1 = 600\ \si{\ohm}$, $R_2 = 400\ \si{\ohm}$ y $R_3 = 300\ \si{\ohm}$. Calcular las corrientes que circulan a través de cada una de las resistencias.
\tcbline
\begin{center}
\begin{tikzpicture}[scale=0.9]
    \draw (0,0) node[ground] {} to[V, v=$20\,\si{\volt}$] (0,2.2)
          to[R, l=$R_1{=}600\,\si{\ohm}$] (2.5,2.2) node[above] {$A$}
          to[R, l=$R_2{=}400\,\si{\ohm}$] (5.0,2.2) node[above] {$B$}
          to[R, l=$R_3{=}300\,\si{\ohm}$] (5.0,0) -- (0,0);
    \draw (2.5,0) to[zzD, l=$D_1$] (2.5,2.2);
    \draw (3.8,2.2) -- (3.8,2.2);
    \draw (4.2,0) to[zzD, l=$D_2$] (4.2,2.2);
\end{tikzpicture}
\end{center}
\textbf{Solución Paso a Paso:}\\[4pt]
1) \textbf{Hipótesis inicial (ambos Zener en OFF):}
Calculamos la corriente y las tensiones en los nudos $A$ y $B$ sin diodos:
\begin{equation}
    I = \frac{20\ \si{\volt}}{R_1 + R_2 + R_3} = \frac{20\ \si{\volt}}{0{,}6 + 0{,}4 + 0{,}3\ \si{\kilo\ohm}} = \frac{20}{1{,}3} = 15{,}38\ \si{\milli\ampere}
\end{equation}
\begin{align}
    V_B &= I \cdot R_3 = 15{,}38 \times 0{,}3 = 4{,}61\ \si{\volt} < 8\ \si{\volt} \implies D_2\text{ está en OFF } (V_{D2} = -4{,}61\ \si{\volt} > -8\ \si{\volt}) \\
    V_A &= I (R_2 + R_3) = 15{,}38 \times 0{,}7 = 10{,}77\ \si{\volt} > 10\ \si{\volt} \implies D_1\text{ entra en ruptura inversa ON}
\end{align}
2) \textbf{Resolución con $D_1$ ON en Zener ($V_A = |V_{Z1}| = 10\ \si{\volt}$) y $D_2$ OFF:}
\begin{align}
    I_{R2} = I_{R3} &= \frac{V_A}{R_2 + R_3} = \frac{10\ \si{\volt}}{0{,}4 + 0{,}3\ \si{\kilo\ohm}} = \frac{10}{0{,}7} = 14{,}285\ \si{\milli\ampere} \\
    V_B &= I_{R3} \cdot R_3 = 14{,}285 \times 0{,}3 = 4{,}285\ \si{\volt} < 8\ \si{\volt} \implies D_2\text{ sigue en OFF (confirmado)} \\
    I_{R1} &= \frac{V_i - V_A}{R_1} = \frac{20 - 10}{0{,}6\ \si{\kilo\ohm}} = 16{,}67\ \si{\milli\ampere}
\end{align}
(La corriente inversa por el Zener $D_1$ es $I_{Z1} = I_{R1} - I_{R2} = 16{,}67 - 14{,}285 = 2{,}385\ \si{\milli\ampere} > 0$).
\end{ejercicio}

% --- EJERCICIO 5 (Ejercicio 8 de los apuntes) ---
\begin{ejercicio}[Comparación de los 3 Modelos Lineales del Diodo --- Ejercicio 8 Apuntes]
Calcular el valor de la tensión de salida $V_o$ en el circuito de la figura ($R_1 = R_2 = R_3 = 5\ \si{\kilo\ohm}$) sabiendo que el diodo $D_1$ cuando está en conducción se representa por:
a) un cortocircuito (modelo 1); \quad b) una fuente de $0{,}7\ \si{\volt}$ (modelo 2); \quad c) una fuente de $0{,}7\ \si{\volt}$ y resistencia de $20\ \si{\ohm}$ (modelo 3).
\tcbline
\begin{center}
\begin{tikzpicture}[scale=0.9]
    \draw (0,3.6) node[above] {$10\,\si{\volt}$} to[R, l=$R_1$] (0,1.8) node[left] {$V_x$}
          to[R, l=$R_2$] (0,0) node[ground] {};
    % Diodo D1 con cátodo en Vx y ánodo en Vo
    \draw (2.5,1.8) to[D, l_=$D_1$] (0,1.8);
    \draw (2.5,1.8) to[short, -o] (3.3,1.8) node[right] {$V_o$};
    \draw (2.5,1.8) to[R, l=$R_3$] (2.5,0.8) to[V, l=$6\,\si{\volt}$] (2.5,0) node[ground] {};
\end{tikzpicture}
\end{center}
\textbf{Solución Paso a Paso:}\\[4pt]
Comprobamos primero el estado de $D_1$ suponiendo que está en \textbf{OFF}:
\begin{equation}
    V_{\text{Ánodo}} = V_o = 6\ \si{\volt}, \qquad V_{\text{Cátodo}} = V_x = 10 \frac{R_2}{R_1 + R_2} = 5\ \si{\volt} \implies V_D = 6 - 5 = 1\ \si{\volt} > 0{,}7\ \si{\volt} \implies D_1\text{ ON}
\end{equation}
\begin{enumerate}[label=\alph*)]
    \item \textbf{Modelo 1 (cortocircuito, $V_\gamma = 0$):} $V_o = V_x$. Por nudos en $V_x$:
    \begin{equation}
        \frac{10 - V_x}{5} + \frac{6 - V_x}{5} = \frac{V_x}{5} \implies 16 = 3 V_x \implies V_x = V_o = 5{,}33\ \si{\volt}
    \end{equation}
    \item \textbf{Modelo 2 (fuente de $0{,}7\ \si{\volt}$):} $V_o = V_x + 0{,}7$. Aplicando Kirchhoff en el supernudo:
    \begin{equation}
        \frac{10 - V_x}{5} + \frac{6 - (V_x + 0{,}7)}{5} = \frac{V_x}{5} \implies 3 V_x = 15{,}3 \implies V_x = 5{,}1\ \si{\volt} \implies V_o = 5{,}1 + 0{,}7 = 5{,}8\ \si{\volt}
    \end{equation}
    \item \textbf{Modelo 3 ($V_\gamma = 0{,}7\ \si{\volt}$ y $r_d = 20\ \si{\ohm} = 0{,}02\ \si{\kilo\ohm}$):} La resistencia total de la rama derecha es $R_3 + r_d = 5{,}02\ \si{\kilo\ohm}$:
    \begin{equation}
        \frac{10 - V_x}{5} + \frac{6 - V_x - 0{,}7}{5{,}02} = \frac{V_x}{5} \implies 15{,}04 V_x = 76{,}7 \implies V_x \approx 5{,}1\ \si{\volt}
    \end{equation}
    La corriente por el diodo es $I_D = \frac{6 - 5{,}8}{5{,}02\ \si{\kilo\ohm}} \approx 40\ \si{\micro\ampere}$, por lo que $V_o = V_x + 0{,}7 + I_D r_d \approx 5{,}8\ \si{\volt}$ (al ser las resistencias de $5\ \si{\kilo\ohm}$ mucho mayores que $r_d = 20\ \si{\ohm}$, los modelos 2 y 3 dan prácticamente el mismo resultado).
\end{enumerate}
\end{ejercicio}

% --- EJERCICIO 6 (Ejercicio 9 de los apuntes) ---
\begin{ejercicio}[Circuito con Fuente de Corriente y Diodo --- Ejercicio 9 Apuntes]
Determinar el valor de la corriente $I$ y $I_2$ en el circuito de la figura suponiendo un diodo rectificador común ($V_\gamma = 0{,}7\ \si{\volt}$) con $I_1 = 8\ \si{\milli\ampere}$, $R_1 = 2{,}2\ \si{\kilo\ohm}$, $R_2 = 5{,}6\ \si{\kilo\ohm}$ y $R_3 = 3{,}3\ \si{\kilo\ohm}$.
\tcbline
\begin{center}
\begin{tikzpicture}[scale=0.9]
    \draw (0,0) node[ground] {} to[I, i=$I_1{=}8\,\si{\milli\ampere}$] (0,2.8)
          to[R, l=$R_1{=}2{,}2\,\si{\kilo\ohm}$] (2.4,2.8) node[above] {$A$}
          to[short, i=$I$] (4.4,2.8)
          to[R, l=$R_3{=}3{,}3\,\si{\kilo\ohm}$] (4.4,0) -- (0,0);
    \draw (2.4,2.8) to[D, l=$D_1$] (2.4,1.5) node[right] {$B$}
          to[R, l=$R_2{=}5{,}6\,\si{\kilo\ohm}$, i=$I_2$] (2.4,0);
\end{tikzpicture}
\end{center}
\textbf{Solución Paso a Paso:}\\[4pt]
Si suponemos $D_1\text{ OFF}$, toda la corriente $I_1 = 8\ \si{\milli\ampere}$ pasa por $R_3$, luego $V_A = 8\ \si{\milli\ampere} \times 3{,}3\ \si{\kilo\ohm} = 26{,}4\ \si{\volt} > 0{,}7\ \si{\volt}$, por lo que \textbf{$D_1$ conduce en ON}.
Sustituyendo $D_1$ por una fuente de $0{,}7\ \si{\volt}$ y aplicando LCK en el nudo $A$:
\begin{equation}
    8 = \frac{V_A - 0{,}7}{5{,}6} + \frac{V_A}{3{,}3} \implies 147{,}84 = 3{,}3 V_A - 2{,}31 + 5{,}6 V_A \implies 8{,}9 V_A = 150{,}15 \implies V_A = 16{,}87\ \si{\volt}
\end{equation}
Conocida la tensión $V_A$, hallamos las corrientes:
\begin{equation}
    I = \frac{V_A}{R_3} = \frac{16{,}87\ \si{\volt}}{3{,}3\ \si{\kilo\ohm}} = 5{,}11\ \si{\milli\ampere}, \qquad I_2 = \frac{V_A - 0{,}7}{R_2} = \frac{16{,}87 - 0{,}7}{5{,}6\ \si{\kilo\ohm}} = 2{,}89\ \si{\milli\ampere} > 0
\end{equation}
\end{ejercicio}

% --- EJERCICIO 7 (Ejercicio 10 de los apuntes) ---
\begin{ejercicio}[Análisis con el 3.er Modelo Lineal para Dos Diodos --- Ejercicio 10 Apuntes]
En el circuito de la figura, $D_1$ y $D_2$ se representan en conducción mediante una fuente de $0{,}7\ \si{\volt}$ en serie con $r_d = 20\ \si{\ohm} = 0{,}02\ \si{\kilo\ohm}$ ($R_1 = 5\ \si{\kilo\ohm}, R_2 = 2\ \si{\kilo\ohm}$). Determinar la tensión en el nudo $A$ si: a) $V_{in} = 10\ \si{\volt}$; \quad b) $V_{in} = -5\ \si{\volt}$.
\tcbline
\begin{center}
\begin{tikzpicture}[scale=0.9]
    \draw (0,0) node[ground] {} to[V, v=$V_{in}$] (0,2.0) -- (0.8,2.0);
    % Paralelo D1 y R1
    \draw (0.8,2.0) -- (0.8,2.8) to[D, l=$D_1$] (3.0,2.8) -- (3.0,2.0);
    \draw (0.8,2.0) -- (0.8,1.3) to[R, l=$R_1{=}5\,\si{\kilo\ohm}$] (3.0,1.3) -- (3.0,2.0);
    \draw (3.0,2.0) -- (4.0,2.0) node[above] {$A$} -- (5.2,2.0);
    % Paralelo R2 y D2 a masa
    \draw (3.8,2.0) to[R, l=$R_2{=}2\,\si{\kilo\ohm}$] (3.8,0) node[ground] {};
    \draw (5.2,0) node[ground] {} to[D, l_=$D_2$] (5.2,2.0);
\end{tikzpicture}
\end{center}
\textbf{Solución Paso a Paso:}\\[4pt]
Si ambos diodos están en OFF, el divisor de tensión da $V_A = \frac{R_2}{R_1 + R_2} V_{in} = \frac{2}{7} V_{in}$.
\begin{enumerate}[label=\alph*)]
    \item \textbf{Cuando $V_{in} = 10\ \si{\volt}$:} Con ambos OFF sería $V_A = \frac{20}{7} = 2{,}86\ \si{\volt}$, luego $V_{D1} = 10 - 2{,}86 = 7{,}14\ \si{\volt} > 0{,}7\ \si{\volt}$ y $V_{D2} = -2{,}86\ \si{\volt} < 0{,}7\ \si{\volt}$. Por tanto, \textbf{$D_1$ ON y $D_2$ OFF}. Aplicando Kirchhoff en el nudo $A$:
    \begin{equation}
        \frac{10 - 0{,}7 - V_A}{0{,}02} + \frac{10 - V_A}{5} = \frac{V_A}{2} \implies V_A = 9{,}211\ \si{\volt}
    \end{equation}
    \item \textbf{Cuando $V_{in} = -5\ \si{\volt}$:} Con ambos OFF sería $V_A = -\frac{10}{7} = -1{,}43\ \si{\volt}$, luego $V_{D2} = 0 - (-1{,}43) = 1{,}43\ \si{\volt} > 0{,}7\ \si{\volt}$ y $V_{D1} < 0$. Ahora \textbf{$D_2$ ON y $D_1$ OFF}:
    \begin{equation}
        \frac{-0{,}7 - V_A}{0{,}02} = \frac{V_A - (-5)}{5} + \frac{V_A}{2} \implies V_A = -0{,}71\ \si{\volt}
    \end{equation}
    con corrientes $I_{R1} = \frac{-0{,}71 - (-5)}{5} = 0{,}858\ \si{\milli\ampere}$, $I_{R2} = \frac{0{,}71}{2} = 0{,}355\ \si{\milli\ampere}$ e $I_{D2} = 0{,}858 - 0{,}355 = 0{,}503\ \si{\milli\ampere} > 0$.
\end{enumerate}
\end{ejercicio}

% --- EJERCICIO 8 (Polarización NMOS Examen UGR + Ejemplo orig-087) ---
\begin{ejercicio}[Punto de Trabajo $Q$ de Transistor NMOS con Divisor de Puerta --- Examen UGR]
Determinar el punto de reposo $Q(V_{GS}, I_D, V_{DS})$ y verificar la región de operación en el circuito de polarización de la figura, con $V_{DD} = 15\ \si{\volt}$, $V_T = 1{,}5\ \si{\volt}$, $k_n = 1\ \si{\milli\ampere\per\volt\squared}$, $R_1 = 10\ \si{\mega\ohm}$, $R_2 = 5\ \si{\mega\ohm}$, $R_D = 3\ \si{\kilo\ohm}$ y $R_S = 1\ \si{\kilo\ohm}$. (A continuación se incluye también la variante de clase con $R_1 = 150\ \si{\kilo\ohm}, R_2 = 100\ \si{\kilo\ohm}, R_D = 40\ \si{\kilo\ohm}, R_S = 5\ \si{\kilo\ohm}, k_n = 40\ \si{\micro\ampere\per\volt\squared}$ y $V_T = 2\ \si{\volt}$).
\tcbline
\begin{center}
\begin{tikzpicture}[scale=0.95]
    % Divisor de puerta a la izquierda
    \draw (0,4.2) node[above] {$V_{DD}=15\,\si{\volt}$}
          to[R, l={$R_1=10\,\si{\mega\ohm}$}] (0,2.0)
          to[R, l={$R_2=5\,\si{\mega\ohm}$}] (0,-0.5) node[ground] {};
    % Transistor NMOS en (2.5, 2.0)
    \draw (2.5, 2.0) node[nmos] (M) {};
    \draw (0,2.0) -- (M.gate);
    % RD vertical desde (M.drain) hacia arriba y RS desde (M.source) hacia abajo
    \draw (M.drain) to[R, l={$R_D=3\,\si{\kilo\ohm}$}] (2.5, 4.2) node[above] {$V_{DD}=15\,\si{\volt}$};
    \draw (M.source) to[R, l={$R_S=1\,\si{\kilo\ohm}$}] (2.5, -0.5) node[ground] {};
\end{tikzpicture}
\end{center}
\textbf{Solución Paso a Paso:}\\[4pt]
1) \textbf{Tensión en la Puerta ($I_G = 0$):}
\begin{equation}
    V_G = V_{DD} \frac{R_2}{R_1 + R_2} = 15 \times \frac{5}{10 + 5} = 5\ \si{\volt}
\end{equation}
Si el transistor estuviese en corte ($I_D = 0 \implies V_S = 0$), sería $V_{GS} = 5\ \si{\volt} > V_T = 1{,}5\ \si{\volt}$, luego \textbf{el NMOS conduce (ON)}.\\[3pt]
2) \textbf{Hipótesis de Saturación y Ecuaciones de Malla ($I_D = I_S$):}
\begin{equation}
    V_{GS} = V_G - I_D R_S = 5 - 1 \cdot I_D, \qquad I_D = \frac{k_n}{2}(V_{GS} - V_T)^2 = 0{,}5(3{,}5 - I_D)^2
\end{equation}
\begin{equation}
    I_D^2 - 9 I_D + 12{,}25 = 0 \implies I_D = \frac{9 - \sqrt{32}}{2} \approx 1{,}67\ \si{\milli\ampere}, \qquad V_{GS} = 5 - 1{,}67 = 3{,}33\ \si{\volt} > 1{,}5\ \si{\volt}
\end{equation}
(La otra raíz $I_D = 7{,}33\ \si{\milli\ampere}$ daría $V_{GS} = -2{,}33\ \si{\volt} < V_T$, imposible).\\[3pt]
3) \textbf{Cálculo de $V_{DS}$ y Comprobación de Saturación:}
\begin{equation}
    V_{DS} = V_{DD} - I_D(R_D + R_S) = 15 - 1{,}67 \times (3 + 1) = 8{,}32\ \si{\volt} \ge V_{GS} - V_T = 1{,}83\ \si{\volt} \quad (\text{Saturación OK})
\end{equation}
\textbf{Variante del Ejemplo de Clase (`orig-087.png`):} Con $R_1 = 150\ \si{\kilo\ohm}, R_2 = 100\ \si{\kilo\ohm}, R_D = 40\ \si{\kilo\ohm}, R_S = 5\ \si{\kilo\ohm}, k_n = 40 \times 10^{-6}\ \si{\ampere\per\volt\squared}$ y $V_T = 2\ \si{\volt}$:
\begin{equation}
    V_G = 15 \frac{100}{250} = 6\ \si{\volt}, \qquad \frac{6 - V_{GS}}{5\ \si{\kilo\ohm}} = 20 \times 10^{-6} (V_{GS} - 2)^2 \implies V_{GS}^2 + 6 V_{GS} - 56 = 0
\end{equation}
De donde $V_{GS} = 5{,}06\ \si{\volt} > 2\ \si{\volt}$, $I_D = \frac{6 - 5{,}06}{5\ \si{\kilo\ohm}} = 0{,}188\ \si{\milli\ampere}$ (o $0{,}1844\ \si{\milli\ampere}$ según redondeo) y $V_{DS} = 15 - I_D(45\ \si{\kilo\ohm}) = 6{,}54\ \si{\volt} > V_{GS} - V_T = 3{,}06\ \si{\volt}$ (\textbf{Saturación correcta}).
\end{ejercicio}

% --- EJERCICIO 9 (Ejercicio 12 de los apuntes) ---
\begin{ejercicio}[Punto de Trabajo en las 3 Regiones según $V_{GG}$ --- Ejercicio 12 Apuntes]
Hallar el punto de trabajo del MOSFET de canal N de la figura ($V_{SS} = -6\ \si{\volt}$, $R_1 = 5{,}6\ \si{\kilo\ohm}$ conectada entre masa $0\ \si{\volt}$ y drenador, $V_T = 2\ \si{\volt}$, $k_n = 2 \times 10^{-3}\ \si{\ampere\per\volt\squared} = 2\ \si{\milli\ampere\per\volt\squared}$) para:
a) $V_{GG} = -3{,}5\ \si{\volt}$; \quad b) $V_{GG} = -3\ \si{\volt}$; \quad c) $V_{GG} = -4\ \si{\volt}$.
\tcbline
\begin{center}
\begin{tikzpicture}[scale=0.9]
    \draw (2.5,1.8) node[nmos] (M) {};
    \draw (1.2,1.8) node[left] {$V_{GG}$} to[short, o-] (M.gate);
    \draw (M.drain) to[R, l=$R_1{=}5{,}6\,\si{\kilo\ohm}$] (2.5,3.8) node[ground] {};
    \draw (M.source) -- (2.5,0.6) node[below] {$V_{SS} = -6\,\si{\volt}$};
\end{tikzpicture}
\end{center}
\textbf{Solución Paso a Paso:}\\[4pt]
Como la fuente está fija a $V_S = V_{SS} = -6\ \si{\volt}$, tenemos siempre $V_{GS} = V_{GG} - (-6) = V_{GG} + 6$.
\begin{enumerate}[label=\alph*)]
    \item \textbf{Si $V_{GG} = -3{,}5\ \si{\volt}$:} $V_{GS} = -3{,}5 + 6 = 2{,}5\ \si{\volt} > V_T = 2\ \si{\volt} \implies \text{NMOS ON}$.
    Suponemos \textbf{saturación}:
    \begin{equation}
        I_D = \frac{k_n}{2}(V_{GS} - V_T)^2 = \frac{2}{2}(2{,}5 - 2)^2 = 0{,}25\ \si{\milli\ampere}
    \end{equation}
    Como $0 - V_D = I_D R_1 \implies V_D = -0{,}25 \times 5{,}6 = -1{,}4\ \si{\volt}$, resulta:
    \begin{equation}
        V_{DS} = V_D - V_{SS} = -1{,}4 - (-6) = 4{,}6\ \si{\volt} > V_{GS} - V_T = 0{,}5\ \si{\volt} \implies \textbf{Saturación OK}
    \end{equation}

    \item \textbf{Si $V_{GG} = -3\ \si{\volt}$:} $V_{GS} = -3 + 6 = 3\ \si{\volt} > V_T = 2\ \si{\volt} \implies \text{NMOS ON}$.
    Si suponemos saturación:
    \begin{equation}
        I_D = \frac{2}{2}(3 - 2)^2 = 1\ \si{\milli\ampere} \implies V_D = -1 \times 5{,}6 = -5{,}6\ \si{\volt} \implies V_{DS} = -5{,}6 - (-6) = 0{,}4\ \si{\volt}
    \end{equation}
    Como $V_{DS} = 0{,}4\ \si{\volt} < V_{GS} - V_T = 1\ \si{\volt}$, \textbf{NO está en saturación}, sino en la \textbf{región lineal (tríodo)}.
    Planteamos las ecuaciones en región lineal ($V_D = V_{SS} + V_{DS} = -6 + V_{DS}$):
    \begin{equation}
        I_D = \frac{-V_{SS} - V_{DS}}{R_1} = k_n \left[ (V_{GS} - V_T) V_{DS} - \frac{1}{2} V_{DS}^2 \right] \implies \frac{6 - V_{DS}}{5{,}6} = 2 \left( 1 \cdot V_{DS} - \frac{1}{2} V_{DS}^2 \right)
    \end{equation}
    Resolviendo la ecuación de segundo grado se obtienen dos raíces: $V_{DS} = 1{,}43\ \si{\volt}$ (se descarta porque es mayor que $V_{GS}-V_T = 1\ \si{\volt}$) y la solución válida:
    \begin{equation}
        V_{DS} = 0{,}75\ \si{\volt} < 1\ \si{\volt}, \qquad I_D = \frac{6 - 0{,}75}{5{,}6} = 0{,}93\ \si{\milli\ampere}, \qquad V_D = -6 + 0{,}75 = -5{,}25\ \si{\volt}
    \end{equation}

    \item \textbf{Si $V_{GG} = -4\ \si{\volt}$:} $V_{GS} = -4 - (-6) = 2\ \si{\volt} = V_T \implies I_D = 0\ \si{\milli\ampere}$ (el transistor está en el umbral de \textbf{OFF / Corte}, con $V_D = 0\ \si{\volt}$ y $V_{DS} = 6\ \si{\volt}$).
\end{enumerate}
\end{ejercicio}

% --- EJERCICIO 10 (Ejercicio 13 de los apuntes) ---
\begin{ejercicio}[MOSFET con Puerta y Drenador Cortocircuitados --- Ejercicio 13 Apuntes]
Determinar el valor de $I_D$, $V_{DS}$ y $V_{GS}$ en el circuito de la figura, sabiendo que $V_{DD} = 12\ \si{\volt}$, $V_{SS} = 0\ \si{\volt}$, $R_1 = 2\ \si{\kilo\ohm}$, $R_2 = 1\ \si{\mega\ohm}$, $V_T = 3\ \si{\volt}$ y $k_n = 0{,}48 \times 10^{-3}\ \si{\ampere\per\volt\squared} = 0{,}48\ \si{\milli\ampere\per\volt\squared}$.
\tcbline
\begin{center}
\begin{tikzpicture}[scale=0.9]
    \draw (2.5,1.8) node[nmos] (M) {};
    \draw (M.drain) -- (2.5,2.6) to[R, l_=$R_1{=}2\,\si{\kilo\ohm}$] (2.5,4.2) node[above] {$V_{DD}=12\,\si{\volt}$};
    \draw (M.source) -- (2.5,0.6) node[ground] {} node[right=4pt] {$V_{SS}=0\,\si{\volt}$};
    \draw (2.5,2.6) to[R, l_=$R_2{=}1\,\si{\mega\ohm}$] (0.5,2.6) -- (0.5,1.8) -- (M.gate);
\end{tikzpicture}
\end{center}
\textbf{Solución Paso a Paso:}\\[4pt]
1) \textbf{¿Conduce el transistor?} Como la corriente de puerta es siempre nula ($I_G = 0 \implies I_{R2} = 0$), no cae tensión en $R_2$, lo que obliga a que $V_G = V_D$ y por tanto \textbf{$V_{GS} = V_{DS}$}. Si supusiéramos el transistor en OFF ($I_D = 0$), sería $V_D = V_{DD} = 12\ \si{\volt} \implies V_{GS} = 12\ \si{\volt} > V_T = 3\ \si{\volt}$, contradicción: el transistor está \textbf{ON}.\\[3pt]
2) \textbf{¿En qué región conduce?} Al tener unidos drenador y puerta ($V_{DS} = V_{GS}$), se cumple siempre que $V_{DS} > V_{GS} - V_T$ (pues $V_T = 3\ \si{\volt} > 0$). Por tanto, \textbf{trabaja obligatoriamente en SATURACIÓN}.
Igualando la corriente del transistor con la que baja por $R_1$:
\begin{equation}
    I_D = \frac{k_n}{2}(V_{GS} - V_T)^2 = \frac{V_{DD} - V_{DS}}{R_1} \implies \frac{0{,}48}{2}(V_{GS} - 3)^2 = \frac{12 - V_{GS}}{2}
\end{equation}
Despejando la raíz positiva ($V_{GS} > 3\ \si{\volt}$):
\begin{equation}
    V_{GS} = V_{DS} = 6{,}43\ \si{\volt}, \qquad I_D = \frac{12 - 6{,}43}{2\ \si{\kilo\ohm}} = 2{,}785\ \si{\milli\ampere}
\end{equation}
\end{ejercicio}

% --- EJERCICIO 11 (Ejercicio 14 de los apuntes) ---
\begin{ejercicio}[Umbral de Conducción de un MOSFET --- Ejercicio 14 Apuntes]
En el circuito de la figura ($V_{DD} = 15\ \si{\volt}, V_{SS} = 5\ \si{\volt}, R_{G1} = 120\ \si{\ohm}, R_{G2} = 220\ \si{\ohm}, R_D = 4{,}7\ \si{\kilo\ohm}, V_T = 2\ \si{\volt}, k_n = 2\ \si{\milli\ampere\per\volt\squared}$):
a) Si $V_{GG} = 0\ \si{\volt}$, ¿cuál es el estado del transistor?
b) Si $V_{GG}$ aumenta desde $0$, ¿para qué tensión empieza a conducir?
c) En el instante en que entra en conducción, ¿en qué región de trabajo se encuentra?
\tcbline
\begin{center}
\begin{tikzpicture}[scale=0.9]
    \draw (0,0) node[ground] {} to[V, v=$V_{GG}$] (0,1.8)
          to[R, l=$R_{G1}{=}\SI{120}{\ohm}$] (2.0,1.8)
          to[R, l=$R_{G2}{=}\SI{220}{\ohm}$] (2.0,0) node[ground] {};
    \draw (3.5,1.8) node[nmos] (M) {};
    \draw (2.0,1.8) -- (M.gate);
    \draw (M.drain) to[R, l=$R_D{=}\SI{4.7}{\kilo\ohm}$] (3.5,3.8) node[above] {$V_{DD}=15\,\si{\volt}$};
    \draw (M.source) -- (3.5,0.6) node[below] {$V_{SS}=5\,\si{\volt}$};
\end{tikzpicture}
\end{center}
\textbf{Solución Paso a Paso:}\\[4pt]
\begin{enumerate}[label=\alph*)]
    \item \textbf{Para $V_{GG} = 0\ \si{\volt}$:} $V_G = 0\ \si{\volt}$ y $V_S = V_{SS} = 5\ \si{\volt} \implies V_{GS} = 0 - 5 = -5\ \si{\volt} < V_T = 2\ \si{\volt} \implies$ \textbf{NMOS OFF (en corte)}.
    \item \textbf{Tensión $V_{GG}$ para empezar a conducir:} El transistor empieza a conducir cuando $V_{GS} = V_T = 2\ \si{\volt}$, lo que exige en la puerta un potencial $V_G = V_S + V_T = 5 + 2 = 7\ \si{\volt}$. Como $I_G = 0$, aplicamos el divisor de tensión entre $R_{G1}$ y $R_{G2}$:
    \begin{equation}
        \frac{V_{GG} - 7}{120} = \frac{7}{220} \implies V_{GG} = \frac{7 \times 120}{220} + 7 = 10{,}82\ \si{\volt} \quad (\text{conduce para } V_{GG} > 10{,}82\ \si{\volt})
    \end{equation}
    \item \textbf{Región al inicio de la conducción:} Justo cuando empieza a conducir ($V_G = 7{,}01\ \si{\volt}$, $V_{GS} - V_T \to 0^+$), la corriente $I_D$ es pequeñísima ($I_D \approx 0$), por lo que la caída en $R_D$ es casi nula y $V_D \approx 15\ \si{\volt}$. Así, $V_{DS} \approx 15 - 5 = 10\ \si{\volt} \gg V_{GS} - V_T \approx 0\ \si{\volt}$, luego entra directamente en \textbf{SATURACIÓN}.
\end{enumerate}
\end{ejercicio}

% --- EJERCICIO 12 (Ejercicio 15 de los apuntes) ---
\begin{ejercicio}[Dos Transistores NMOS en Serie y Límite Lineal-Saturación --- Ejercicio 15 Apuntes]
En el circuito de la figura ($V_{DD} = 9\ \si{\volt}$; $M_1$: $V_{T1} = 1\ \si{\volt}, k_1 = 4\ \si{\milli\ampere\per\volt\squared}$; $M_2$: $V_{T2} = 2\ \si{\volt}, k_2 = 2\ \si{\milli\ampere\per\volt\squared}$):
a) Hallar el punto de trabajo y la potencia disipada en cada transistor si $V_{GG} = 3\ \si{\volt}$.
b) Calcular la tensión $V_{GG}$ máxima para que $M_1$ se mantenga en la región de saturación (límite lineal-saturación).
\tcbline
\begin{center}
\begin{tikzpicture}[scale=0.9]
    \draw (2.5,1.2) node[nmos] (M1) {};
    \draw (2.5,3.2) node[nmos] (M2) {};
    \draw (1.2,1.2) node[left] {$V_{GG}=3\,\si{\volt}$} to[short, o-] (M1.gate);
    \draw (M1.source) -- (2.5,0.2) node[ground] {};
    \draw (M1.drain) -- (M2.source) node[midway, right] {$D_1$};
    \draw (M2.drain) -- (2.5,4.2) node[above] {$V_{DD}=9\,\si{\volt}$};
    \draw (2.5,3.9) -- (1.4,3.9) -- (1.4,3.2) -- (M2.gate);
    \node[right=10pt] at (M1) {$M_1$};
    \node[right=10pt] at (M2) {$M_2$};
\end{tikzpicture}
\end{center}
\textbf{Solución Paso a Paso:}\\[4pt]
\begin{enumerate}[label=\alph*)]
    \item \textbf{Punto de trabajo y potencias para $V_{GG} = 3\ \si{\volt}$:}
    El transistor $M_2$ tiene su puerta conectada a su drenador ($V_{G2} = V_{D2} = 9\ \si{\volt}$), luego conduce en \textbf{saturación}. Para $M_1$, $V_{GS1} = 3\ \si{\volt} > V_{T1} = 1\ \si{\volt}$ (ON). Suponemos que $M_1$ también está en \textbf{saturación} e igualamos $I_{D1} = I_{D2}$ (teniendo en cuenta que $V_{S2} = V_{D1}$):
    \begin{equation}
        \frac{k_1}{2}(V_{GS1} - V_{T1})^2 = \frac{k_2}{2}(V_{GS2} - V_{T2})^2 \implies \frac{4}{2}(3 - 1)^2 = \frac{2}{2}(9 - V_{D1} - 2)^2
    \end{equation}
    \begin{equation}
        8 = (7 - V_{D1})^2 \implies V_{D1}^2 - 14 V_{D1} + 41 = 0 \implies V_{D1} = \frac{14 \pm \sqrt{32}}{2} = \begin{cases} 9{,}83\ \si{\volt} & (\text{NO, supera } 9\ \si{\volt}) \\ 4{,}17\ \si{\volt} & (\text{SÍ}) \end{cases}
    \end{equation}
    Comprobamos la saturación de $M_1$: $V_{DS1} = V_{D1} = 4{,}17\ \si{\volt} > V_{GS1} - V_{T1} = 3 - 1 = 2\ \si{\volt}$ (\textbf{Saturación OK}).
    La corriente común es $I_D = 8\ \si{\milli\ampere}$, con $V_{DS1} = 4{,}17\ \si{\volt}$ y $V_{DS2} = 9 - 4{,}17 = 4{,}83\ \si{\volt}$. Las potencias disipadas son:
    \begin{equation}
        P_{M1} = I_D \cdot V_{DS1} = 8 \times 4{,}17 = 33{,}36\ \si{\milli\watt}, \qquad P_{M2} = I_D \cdot V_{DS2} = 8 \times 4{,}83 = 38{,}64\ \si{\milli\watt}
    \end{equation}

    \item \textbf{Frontera entre saturación y región lineal de $M_1$:}
    El límite exacto ocurre cuando $V_{DS1} = V_{D1} = V_{GG} - V_{T1} = V_{GG} - 1$. Igualando de nuevo las corrientes $I_{D1} = I_{D2}$ en ese punto frontera:
    \begin{equation}
        \frac{4}{2}(V_{GG} - 1)^2 = \frac{2}{2}\left(9 - (V_{GG} - 1) - 2\right)^2 \implies 2(V_{GG} - 1)^2 = (8 - V_{GG})^2
    \end{equation}
    \begin{equation}
        V_{GG}^2 + 12 V_{GG} - 62 = 0 \implies V_{GG} = \frac{-12 + \sqrt{392}}{2} = 3{,}9\ \si{\volt}
    \end{equation}
    (Para $V_{GG} < 3{,}9\ \si{\volt}$, $M_1$ opera en saturación; para $V_{GG} > 3{,}9\ \si{\volt}$, $M_1$ opera en la región lineal).
\end{enumerate}
\end{ejercicio}

% --- EJERCICIO 13 (Ejercicio 17 de los apuntes) ---
\begin{ejercicio}[Diseño de Polarización de Dos NMOS Idénticos --- Ejercicio 17 Apuntes]
Los dos transistores NMOS del circuito de la figura son idénticos ($V_T = 0{,}6\ \si{\volt}$, $k_n = 4\ \si{\milli\ampere\per\volt\squared}$, $V_{DD} = 15\ \si{\volt}$, $R_2 = 1\ \si{\mega\ohm}$, $R_3 = 1{,}5\ \si{\kilo\ohm}$). Se desea que la corriente de drenador sea igual en ambos: $I_1 = I_2 = 2\ \si{\milli\ampere}$. Calcular $V_{GS}$ para $M_1$ y $M_2$, el valor de $R_1$ y justificar la zona de trabajo de ambos transistores.
\tcbline
\begin{center}
\begin{tikzpicture}[scale=0.9]
    % Divisor R1, R2 a la izquierda
    \draw (0,4.2) to[R, l_=$R_1$] (0,2.6) node[left] {$V_{G2}$}
          to[R, l_=$R_2{=}1\,\si{\mega\ohm}$] (0,1.0) node[ground] {};
    % M2 arriba a la izquierda, M1 a la derecha
    \draw (1.8,2.6) node[nmos] (M2) {};
    \draw (0,2.6) -- (M2.gate);
    \draw (M2.source) -- (1.8,1.5) to[I, i=$I_1{=}2\,\si{\milli\ampere}$] (1.8,0) node[ground] {};
    \draw (3.8,1.5) node[nmos] (M1) {};
    \draw (1.8,1.5) -- (M1.gate) node[midway, above, font=\footnotesize] {$V_{S2}=V_{G1}$};
    \draw (M1.source) to[R, l=$R_3{=}1{,}5\,\si{\kilo\ohm}$] (3.8,0) node[ground] {};
    % Alimentación común VDD = 15V
    \draw (0,4.2) -- (3.8,4.2) -- (M1.drain);
    \draw (M2.drain) -- (1.8,4.2);
    \node[above] at (2.0,4.2) {$V_{DD}=15\,\si{\volt}$};
    \node[right=10pt] at (M2) {$M_2$};
    \node[right=10pt] at (M1) {$M_1$};
\end{tikzpicture}
\end{center}
\textbf{Solución Paso a Paso:}\\[4pt]
1) \textbf{Análisis de $M_1$:} La tensión en la fuente de $M_1$ es $V_{S1} = I_2 R_3 = 2\ \si{\milli\ampere} \times 1{,}5\ \si{\kilo\ohm} = 3\ \si{\volt}$. Como su drenador está conectado directamente a $V_{DD} = 15\ \si{\volt}$, tenemos $V_{DS1} = 15 - 3 = 12\ \si{\volt}$, que es muy grande, luego \textbf{$M_1$ trabaja en saturación}:
\begin{equation}
    I_2 = \frac{k_n}{2}(V_{GS1} - V_T)^2 \implies 2 = \frac{4}{2}(V_{GS1} - 0{,}6)^2 \implies V_{GS1} - 0{,}6 = 1 \implies V_{GS1} = 1{,}6\ \si{\volt}
\end{equation}
(Comprobación: $V_{DS1} = 12\ \si{\volt} > V_{GS1} - V_T = 1\ \si{\volt}$, saturación verificada).
De aquí obtenemos la tensión en la puerta de $M_1$ (que es también la fuente de $M_2$):
\begin{equation}
    V_{G1} = V_{S2} = V_{S1} + V_{GS1} = 3 + 1{,}6 = 4{,}6\ \si{\volt}
\end{equation}
2) \textbf{Análisis de $M_2$ y cálculo de $R_1$:} Como ambos transistores son idénticos y queremos la misma corriente $I_1 = I_2 = 2\ \si{\milli\ampere}$ en saturación, debe cumplirse que:
\begin{equation}
    V_{GS2} = V_{GS1} = 1{,}6\ \si{\volt} \implies V_{G2} = V_{S2} + V_{GS2} = 4{,}6 + 1{,}6 = 6{,}2\ \si{\volt}
\end{equation}
Comprobamos la saturación de $M_2$: $V_{DS2} = V_{DD} - V_{S2} = 15 - 4{,}6 = 10{,}4\ \si{\volt} > V_{GS2} - V_T = 1\ \si{\volt}$ (\textbf{Saturación de $M_2$ OK}).
Finalmente, como $I_{G2} = 0$, la corriente por $R_1$ y $R_2$ es la misma:
\begin{equation}
    I_{R1} = I_{R2} = \frac{V_{G2}}{R_2} = \frac{6{,}2\ \si{\volt}}{1\ \si{\mega\ohm}} = 6{,}2\ \si{\micro\ampere} \implies R_1 = \frac{V_{DD} - V_{G2}}{I_{R1}} = \frac{15 - 6{,}2}{6{,}2\ \si{\micro\ampere}} = 1{,}42\ \si{\mega\ohm}
\end{equation}
\end{ejercicio}

\chapter{Fundamentos de Electrónica Digital y Lógica CMOS}
\label{cap:tema5}

\cuadrofuentes{
    \textbf{Asignatura:} Fundamentos Físicos y Tecnológicos (1.º Grado en Ingeniería Informática, UGR).\\[2pt]
    \textbf{Profesor:} Pedro Cartujo (Departamento de Electrónica y Tecnología de Computadores).\\[2pt]
    \textbf{Bibliografía de Referencia:} Rabaey, J. M., Chandrakasan, A., Nikolic, B., \textit{Digital Integrated Circuits: A Design Perspective}, Prentice Hall; Wakerly, J. F., \textit{Diseño Digital: Principios y Prácticas}, Pearson.
}

% ==============================================================================
% 5.1. PARÁMETROS CARACTERÍSTICOS DE UNA FAMILIA LÓGICA
% ==============================================================================
\section{Parámetros Característicos de una Familia Lógica}

Mientras que en la electrónica analógica las tensiones e intensidades pueden tomar cualquier valor continuo dentro de un rango, en la \textbf{electrónica digital} se trabaja con \textbf{lógica binaria}, donde las variables solo representan dos estados discretos: \textbf{'0' lógico} y \textbf{'1' lógico}.
\begin{itemize}[leftmargin=*]
    \item \textbf{Lógica Positiva (convenio estándar):} El estado '1' lógico se asocia al nivel alto de tensión ($V_{\text{alto}} \approx V_{DD}$) y el '0' lógico se asocia al nivel bajo ($V_{\text{bajo}} \approx 0\ \si{\volt}$).
    \item \textbf{Lógica Negativa:} El '1' lógico se asocia al nivel bajo de tensión y el '0' lógico al nivel alto.
\end{itemize}

\subsection{El Inversor, Niveles de Tensión y Márgenes de Ruido}
El circuito digital más elemental es el \textbf{inversor (puerta NOT)}, cuya tabla de verdad invierte el bit de entrada ($V_i = 0 \implies V_o = 1$; $V_i = 1 \implies V_o = 0$). La \textbf{función o característica de transferencia de tensión (VTC)} representa la tensión de salida $V_{out}$ en función de la tensión de entrada $V_{in}$:
\begin{itemize}[leftmargin=*]
    \item En un inversor ideal, la transición entre el nivel alto y el nivel bajo sería un escalón vertical instantáneo a mitad de rango.
    \item En una puerta lógica práctica real, la transición presenta una pendiente finita negativa; los puntos donde la pendiente vale $\frac{dV_{out}}{dV_{in}} = -1$ delimitan la región de transición e identifican los umbrales de entrada $V_{IL}$ y $V_{IH}$.
\end{itemize}

\begin{definicion}[Niveles de Tensión Normalizados de una Familia Lógica]
Para garantizar que una puerta lógica receptora interprete sin errores el estado lógico enviado por una puerta transmisora precedente, toda familia lógica especifica cuatro niveles límite de tensión:
\begin{itemize}[leftmargin=*]
    \item $V_{OH}$ (\textit{Output High}): Tensión \textbf{mínima} garantizada en la salida cuando la puerta entrega un '1' lógico.
    \item $V_{OL}$ (\textit{Output Low}): Tensión \textbf{máxima} garantizada en la salida cuando la puerta entrega un '0' lógico.
    \item $V_{IH}$ (\textit{Input High}): Tensión \textbf{mínima} que debe recibir una entrada para ser interpretada inequívocamente como un '1' lógico.
    \item $V_{IL}$ (\textit{Input Low}): Tensión \textbf{máxima} que puede recibir una entrada para ser interpretada inequívocamente como un '0' lógico.
\end{itemize}
La franja intermedia $[V_{IL}, V_{IH}]$ constituye la \textbf{zona prohibida o de indeterminación}.
\end{definicion}

\begin{ley}[Márgenes de Ruido en Nivel Alto y Nivel Bajo]
Los \textbf{márgenes de ruido} cuantifican la máxima amplitud de perturbación o ruido electromagnético que puede sumarse o restarse a la señal en el cable entre dos puertas sin provocar un error lógico en el receptor:
\begin{align}
    NM_H &= V_{OH} - V_{IH} \quad (\text{Margen de Ruido Alto, \textit{Noise Margin High}}) \\
    NM_L &= V_{IL} - V_{OL} \quad (\text{Margen de Ruido Bajo, \textit{Noise Margin Low}})
\end{align}
Para que una familia lógica sea funcional debe cumplirse estrictamente que $V_{OH} > V_{IH}$ ($NM_H > 0$) y $V_{IL} > V_{OL}$ ($NM_L > 0$). El margen de ruido global de la familia es el caso peor:
\begin{equation}
    NM = \min(NM_H, NM_L)
\end{equation}
\end{ley}

\begin{center}
\begin{tikzpicture}[scale=0.95]
    % Barra de salida (Transmisor) de 2.8 cm de ancho: (-4.5,0) a (-1.7,3.4)
    \node[font=\bfseries] at (-3.1,3.8) {Salida (Transmisor)};
    \draw[thick] (-4.5,0) rectangle (-1.7,3.4);
    \draw[fill=fftGreen!25] (-4.5,2.5) rectangle (-1.7,3.4) node[midway, font=\footnotesize] {'1' Lógico ($V_{OH}$)};
    \draw[fill=gray!20] (-4.5,0.9) rectangle (-1.7,2.5) node[midway, font=\footnotesize] {Zona Transición};
    \draw[fill=fftRed!25] (-4.5,0) rectangle (-1.7,0.9) node[midway, font=\footnotesize] {'0' Lógico ($V_{OL}$)};
    \node[left=3pt] at (-4.5,3.4) {$V_{DD}$};
    \node[left=3pt] at (-4.5,2.5) {$V_{OH}$};
    \node[left=3pt] at (-4.5,0.9) {$V_{OL}$};
    \node[left=3pt] at (-4.5,0) {$0\,\si{\volt}$};

    % Barra de entrada (Receptor) de 2.8 cm de ancho: (1.7,0) a (4.5,3.4)
    \node[font=\bfseries] at (3.1,3.8) {Entrada (Receptor)};
    \draw[thick] (1.7,0) rectangle (4.5,3.4);
    \draw[fill=fftGreen!25] (1.7,2.0) rectangle (4.5,3.4) node[midway, font=\footnotesize] {'1' Aceptado ($V_{IH}$)};
    \draw[fill=fftOrange!25] (1.7,1.3) rectangle (4.5,2.0) node[midway, font=\footnotesize] {Indeterminada};
    \draw[fill=fftRed!25] (1.7,0) rectangle (4.5,1.3) node[midway, font=\footnotesize] {'0' Aceptado ($V_{IL}$)};
    \node[right=3pt] at (4.5,3.4) {$V_{DD}$};
    \node[right=3pt] at (4.5,2.0) {$V_{IH}$};
    \node[right=3pt] at (4.5,1.3) {$V_{IL}$};
    \node[right=3pt] at (4.5,0) {$0\,\si{\volt}$};

    % Márgenes de ruido entre ambas barras
    \draw[<->, >=stealth, very thick, fftPrimary] (-1.5,2.5) -- (1.5,2.0) node[midway, above] {$NM_H$};
    \draw[<->, >=stealth, very thick, fftAccent] (-1.5,0.9) -- (1.5,1.3) node[midway, below] {$NM_L$};
\end{tikzpicture}
\end{center}

\subsection{Fan-In, Fan-Out, Tiempos de Retardo y Disipación de Potencia}
\begin{itemize}[leftmargin=*]
    \item \textbf{Fan-In (Abanico de Entrada):} Número de entradas que posee una misma puerta lógica (por ejemplo, una puerta NAND de 3 entradas tiene $\text{Fan-In} = 3$).
    \item \textbf{Fan-Out o Factor de Carga (Abanico de Salida):} Número máximo de entradas de puertas de la misma familia que se pueden conectar simultáneamente a la salida de una puerta sin que se estropee su funcionamiento ni los niveles lógicos salgan de los márgenes garantizados:
    \begin{equation}
        FO = \min\left( \left\lfloor \frac{|I_{OH}|}{|I_{IH}|} \right\rfloor, \left\lfloor \frac{|I_{OL}|}{|I_{IL}|} \right\rfloor \right)
    \end{equation}
    donde $I_{OH}, I_{OL}$ son las corrientes máximas que la salida puede suministrar o absorber en estado alto y bajo, e $I_{IH}, I_{IL}$ las corrientes requeridas por cada entrada conectada. En tecnologías MOS/CMOS, al ser $I_G = 0$, el \textit{fan-out} estático en continua es teóricamente infinito, quedando limitado en alta frecuencia por la suma de capacidades de entrada ($C_L = N \cdot C_{in}$) que ralentizan las conmutaciones.

    \item \textbf{Tiempos de Retardo de Propagación ($t_d$ o $t_p$):}
    Toda puerta lógica real presenta un retardo temporal entre el cambio en la entrada y la respuesta en la salida (medido entre los puntos del $50\%$ de excursión de ambas señales):
    \begin{itemize}
        \item $t_{pHL}$: tiempo de propagación cuando la salida pasa de nivel alto a bajo (\textit{High-to-Low}).
        \item $t_{pLH}$: tiempo de propagación cuando la salida pasa de nivel bajo a alto (\textit{Low-to-High}).
        \item \textbf{Retardo de propagación medio} (interesa que sea lo más pequeño posible):
        \begin{equation}
            t_d = t_p = \frac{t_{pHL} + t_{pLH}}{2}
        \end{equation}
    \end{itemize}

    \item \textbf{Disipación de Potencia Estática y Dinámica y Producto Retardo-Potencia:}
    \begin{enumerate}[label=\alph*)]
        \item \textbf{Potencia Estática ($P_{\text{est}}$):} Potencia consumida cuando la puerta permanece en un estado lógico fijo sin conmutar. En lógica NMOS con carga resistiva o activa es elevada (cuando la salida vale '0' circula corriente continua de $V_{DD}$ a masa), mientras que \textbf{en lógica CMOS estática es prácticamente nula} ($P_{\text{est}} = V_{DD} I_{\text{fuga}} \approx 0$).
        \item \textbf{Potencia Dinámica ($P_D$):} Potencia consumida al cargar y descargar la capacidad de carga equivalente $C_L$ a una frecuencia de conmutación $f$:
        \begin{equation}
            P_D = C_L \cdot V_{DD}^2 \cdot f
        \end{equation}
        \item \textbf{Producto Retardo-Potencia ($PDP$):} Figura de mérito energética que mide la energía consumida por cada transición lógica:
        \begin{equation}
            PDP = P \cdot t_p \quad [\si{\joule} \text{ o } \si{\pico\joule}]
        \end{equation}
    \end{enumerate}
\end{itemize}

% ==============================================================================
% 5.2. PUERTAS LÓGICAS CON TRANSISTORES MOSFET
% ==============================================================================
\section{Puertas Lógicas con Transistores MOSFET: Lógica NMOS y Lógica CMOS}

\subsection{El Transistor MOS como Interruptor Controlado}
En electrónica digital, los transistores MOSFET de enriquecimiento se utilizan como \textbf{interruptores controlados por la tensión de puerta}:
\begin{itemize}[leftmargin=*]
    \item \textbf{Transistor NMOS (conduce con '1', corta con '0'):}
    \begin{itemize}
        \item Si la entrada en la puerta es un \textbf{'0' lógico} ($V_G = 0\ \si{\volt} \implies V_{GS} < V_{Tn}$): el NMOS está en \textbf{CORTE (OFF)}, comportándose como un \textbf{interruptor abierto}.
        \item Si la entrada en la puerta es un \textbf{'1' lógico} ($V_G = V_{DD} \implies V_{GS} > V_{Tn}$): el NMOS \textbf{CONDUCE (ON)}, comportándose como un \textbf{interruptor cerrado} que descarga la salida a masa (transmite perfectamente el '0' lógico).
    \end{itemize}
    \item \textbf{Transistor PMOS (conduce con '0', corta con '1'):}
    \begin{itemize}
        \item Si la entrada en la puerta es un \textbf{'0' lógico} ($V_G = 0\ \si{\volt} \implies V_{SG} = V_{DD} > |V_{Tp}|$): el PMOS \textbf{CONDUCE (ON)}, comportándose como un \textbf{interruptor cerrado} que conecta la salida a $V_{DD}$ (transmite perfectamente el '1' lógico).
        \item Si la entrada en la puerta es un \textbf{'1' lógico} ($V_G = V_{DD} \implies V_{SG} = 0 < |V_{Tp}|$): el PMOS está en \textbf{CORTE (OFF)}, comportándose como un \textbf{interruptor abierto}.
    \end{itemize}
\end{itemize}

\subsection{Inversor y Puertas en Lógica NMOS (con Carga Pasiva o Activa)}
En la familia lógica \textbf{NMOS}, la red inferior de conmutación está formada por transistores NMOS de enriquecimiento y en la parte superior se coloca una \textbf{carga} conectada a $V_{DD}$. Como se estudia en los apuntes de clase, dicha carga de \textit{pull-up} puede ser de tres tipos:
\begin{enumerate}[leftmargin=*]
    \item \textbf{Una resistencia pasiva ($R_D$):} Fácil de analizar, pero ocupa mucha área en un chip integrado.
    \item \textbf{Un transistor NMOS de enriquecimiento como carga activa:} Con su puerta conectada a su propio drenador ($V_G = V_D = V_{DD}$), por lo que trabaja permanentemente en saturación cuando conduce.
    \item \textbf{Un transistor NMOS de deplexión como carga activa:} Con su puerta conectada a su fuente ($V_{GS} = 0 > V_{Td} < 0$), por lo que tiene canal permanentemente formado.
\end{enumerate}

\begin{center}
\begin{tikzpicture}[scale=0.92]
    % 1) Inversor NMOS con resistencia RD (x = -5.0)
    \node[font=\bfseries\footnotesize] at (-5.0,4.5) {1. Carga Resistiva ($R_D$)};
    \draw (-5.0,0.8) node[nmos, anchor=drain] (n1) {};
    \draw (n1.source) node[ground] {};
    \draw (-5.0,0.8) to[R, l=$R_D$] (-5.0,3.6) node[above] {$V_{DD}$};
    \draw (-5.0,1.0) to[short, *-o] (-3.8,1.0) node[right] {$V_o$};
    \draw (n1.gate) to[short, -o] (-6.3,0.1) node[left] {$V_i$};

    % 2) Inversor NMOS con carga NMOS enriquecimiento (x = 0)
    \node[font=\bfseries\footnotesize] at (0,4.5) {2. Carga NMOS Enriq.};
    \draw (0,0.8) node[nmos, anchor=drain] (n2) {};
    \draw (n2.source) node[ground] {};
    \draw (0,1.4) node[nmos, anchor=source] (nl2) {};
    \draw (0,0.8) -- (0,1.4);
    \draw (nl2.drain) -- (0,3.6) node[above] {$V_{DD}$};
    \draw (0,3.2) -- (-1.0,3.2) -- (-1.0,2.1) -- (nl2.gate);
    \draw (0,1.1) to[short, *-o] (1.2,1.1) node[right] {$V_o$};
    \draw (n2.gate) to[short, -o] (-1.3,0.1) node[left] {$V_i$};

    % 3) Inversor NMOS con carga NMOS deplexión (x = 5.0)
    \node[font=\bfseries\footnotesize] at (5.0,4.5) {3. Carga NMOS Deplex.};
    \draw (5.0,0.8) node[nmos, anchor=drain] (n3) {};
    \draw (n3.source) node[ground] {};
    \draw (5.0,1.4) node[nmos, anchor=source] (nl3) {};
    % Trazo grueso en canal para indicar deplexión
    \draw[ultra thick] ($(nl3.source)+(-0.18,0.22)$) -- ($(nl3.drain)+(-0.18,-0.22)$);
    \draw (5.0,0.8) -- (5.0,1.4);
    \draw (nl3.drain) -- (5.0,3.6) node[above] {$V_{DD}$};
    \draw (nl3.gate) -- (4.0,2.1) -- (4.0,1.2) -- (5.0,1.2);
    \draw (5.0,1.1) to[short, *-o] (6.2,1.1) node[right] {$V_o$};
    \draw (n3.gate) to[short, -o] (3.7,0.1) node[left] {$V_i$};
\end{tikzpicture}
\end{center}

\textbf{Cálculo de $V_{OH}$ y $V_{OL}$ en el Inversor NMOS con resistencia de pull-up $R_D$:}
\begin{itemize}[leftmargin=*]
    \item \textbf{Entrada a nivel bajo ($V_i = 0\ \si{\volt} \equiv \text{'0'}$):} El transistor NMOS inferior está en corte ($I_D = 0$). No cae tensión en la resistencia de carga $R_D$, por lo que la salida sube exactamente a:
    \begin{equation}
        V_{OH} = V_{DD} \equiv \text{'1' lógico}
    \end{equation}
    \item \textbf{Entrada a nivel alto ($V_i = V_{DD} \equiv \text{'1'}$):} El NMOS inferior conduce con $V_{GS} = V_{DD} > V_{Tn}$. Como la salida cae a un valor bajo $V_o = V_{OL} < V_{DD} - V_{Tn}$, el transistor trabaja en la \textbf{región óhmica (lineal)}. Igualando la corriente que baja por $R_D$ con la corriente del NMOS en zona lineal:
    \begin{equation}
        \frac{V_{DD} - V_{OL}}{R_D} = k_n \left[ (V_{DD} - V_{Tn}) V_{OL} - \frac{1}{2} V_{OL}^2 \right]
    \end{equation}
    Para que $V_{OL}$ sea cercano a $0\ \si{\volt}$ se requiere que $k_n R_D \gg 1$, pero siempre resulta $V_{OL} > 0\ \si{\volt}$ y existe un consumo constante de potencia estática $P = \frac{V_{DD}(V_{DD}-V_{OL})}{R_D}$ mientras la salida está en '0'.
\end{itemize}

\subsection{El Inversor Estático CMOS: Estructura, Tabla de Verdad y VTC}
La tecnología \textbf{CMOS} (\textit{Complementary MOS}) elimina el consumo de potencia estática uniendo un transistor \textbf{PMOS} en la parte superior (entre $V_{DD}$ y la salida) y un transistor \textbf{NMOS} en la parte inferior (entre la salida y masa), con ambas puertas conectadas a la entrada común $V_{in}$.

\begin{center}
\begin{tikzpicture}[scale=1.0]
    % Alimentación VDD arriba
    \draw (0,3.6) node[above] {$V_{DD}$} -- (0,3.2);
    % Transistor PMOS con anchor=source
    \draw (0,3.2) node[pmos, anchor=source] (p) {};
    % Transistor NMOS con anchor=drain debajo
    \draw (0,1.2) node[nmos, anchor=drain] (n) {};
    \draw (n.source) node[ground] {};
    % Unión de drenadores y salida Vout
    \draw (p.drain) -- (n.drain);
    \draw (0,1.5) to[short, *-o] (1.8,1.5) node[right] {$V_{out}$};
    % Unión de puertas y entrada Vin
    \draw (p.gate) -- (-1.3,2.5) -- (-1.3,0.5) -- (n.gate);
    \draw (-1.3,1.5) to[short, *-o] (-2.3,1.5) node[left] {$V_{in}$};
    % Etiquetas laterales limpias
    \node[right=12pt, fftPrimary, font=\footnotesize\bfseries] at (0.3,2.5) {PMOS (\textit{Pull-Up})};
    \node[right=12pt, fftDarkBlue, font=\footnotesize\bfseries] at (0.3,0.5) {NMOS (\textit{Pull-Down})};
\end{tikzpicture}
\end{center}

\textbf{Tabla de verdad física y funcionamiento del Inversor CMOS:}
\begin{center}
\renewcommand{\arraystretch}{1.25}
\begin{tabular}{c c c c c}
\toprule
\textbf{Entrada $V_{in}$} & \textbf{Estado PMOS} & \textbf{Estado NMOS} & \textbf{Salida $V_{out}$} & \textbf{Consumo Estático} \\
\midrule
$'0'\ (0\ \si{\volt})$ & \textbf{ON} (Lineal) & \textbf{OFF} (Corte) & $V_{OH} = V_{DD}\ ('1')$ & $I_{DD} = 0 \implies P_{\text{est}} = 0$ \\
$'1'\ (V_{DD})$ & \textbf{OFF} (Corte) & \textbf{ON} (Lineal) & $V_{OL} = 0\ \si{\volt}\ ('0')$ & $I_{DD} = 0 \implies P_{\text{est}} = 0$ \\
\bottomrule
\end{tabular}
\end{center}

\textbf{Curva de Transferencia de Tensión (VTC) y sus 5 Regiones:}
Al variar $V_{in}$ desde $0\ \si{\volt}$ hasta $V_{DD}$, el inversor CMOS recorre cinco regiones de operación:
\begin{enumerate}[leftmargin=*]
    \item \textbf{Región A ($0 \le V_{in} < V_{Tn}$):} NMOS en corte (OFF) y PMOS en región lineal $\implies V_{out} = V_{DD}$.
    \item \textbf{Región B ($V_{Tn} \le V_{in} < V_M$):} NMOS entra en saturación y PMOS sigue en región lineal.
    \item \textbf{Región C ($V_{in} = V_M \approx \frac{V_{DD}}{2}$):} \textbf{Ambos transistores están en saturación}, originando una caída casi vertical entre '1' y '0'.
    \item \textbf{Región D ($V_M < V_{in} \le V_{DD} - |V_{Tp}|$):} NMOS en región lineal y PMOS en saturación.
    \item \textbf{Región E ($V_{in} > V_{DD} - |V_{Tp}|$):} PMOS en corte (OFF) y NMOS en región lineal $\implies V_{out} = 0\ \si{\volt}$.
\end{enumerate}

\begin{center}
\begin{tikzpicture}[scale=1.0]
    \draw[->, thick] (0,0) -- (4.5,0) node[right] {$V_{in}$};
    \draw[->, thick] (0,0) -- (0,4.3) node[above] {$V_{out}$};
    \node[left] at (0,3.6) {$V_{DD}$};
    \node[below] at (3.6,0) {$V_{DD}$};
    \draw[very thick, fftAccent] (0,3.6) -- (1.0,3.6)
          to[out=0, in=100] (1.8,1.8)
          to[out=-80, in=180] (2.6,0) -- (3.6,0);
    \draw[dashed, gray] (1.8,0) node[below, text=black, font=\footnotesize] {$V_M{=}\frac{V_{DD}}{2}$} -- (1.8,1.8);
    \draw[dashed, gray] (0,1.8) node[left, text=black, font=\footnotesize] {$V_M$} -- (1.8,1.8);
    \filldraw[fftDarkBlue] (1.8,1.8) circle (2pt);
    \node[font=\footnotesize\bfseries, fftPrimary] at (0.5,3.2) {A};
    \node[font=\footnotesize\bfseries, fftPrimary] at (1.4,2.8) {B};
    \node[font=\footnotesize\bfseries, fftAccent] at (2.1,1.9) {C};
    \node[font=\footnotesize\bfseries, fftPrimary] at (2.3,1.0) {D};
    \node[font=\footnotesize\bfseries, fftPrimary] at (3.2,0.4) {E};
\end{tikzpicture}
\end{center}

% ==============================================================================
% 5.3. PUERTAS COMBINACIONALES Y METODOLOGÍA DE SÍNTESIS CMOS
% ==============================================================================
\section{Puertas CMOS Combinacionales: NOR, NAND y Síntesis Sistemática}

\begin{ley}[Regla Fundamental de Diseño de Puertas Lógicas CMOS]
Toda puerta CMOS estática consta de una \textbf{red inferior de Pull-Down (PDN, transistores NMOS)} entre la salida y masa, y una \textbf{red superior dual de Pull-Up (PUN, transistores PMOS)} entre $V_{DD}$ y la salida:
\begin{itemize}[leftmargin=*]
    \item \textbf{SUMA LÓGICA (OR, $+$):} Se implementa con \textbf{NMOS en PARALELO} abajo y \textbf{PMOS en SERIE} arriba.
    \item \textbf{MULTIPLICACIÓN LÓGICA (AND, $\cdot$):} Se implementa con \textbf{NMOS en SERIE} abajo y \textbf{PMOS en PARALELO} arriba.
    \item \textbf{Funciones NO negadas:} Como una etapa CMOS o NMOS simple invierte siempre la función de su red inferior ($V_o = \overline{F}$), si el problema pide una función positiva sin negar, p. ej. $Y = A \cdot B + C$, primero sintetizamos la función negada $\overline{Y} = \overline{A \cdot B + C}$ y a su salida conectamos un \textbf{inversor} para obtener $Y = \overline{\overline{Y}}$.
\end{itemize}
\end{ley}

\subsection{Puertas Básicas de 2 Entradas: NOR y NAND en CMOS}

\begin{minipage}[t]{0.48\textwidth}
\textbf{1. Puerta NOR CMOS ($V_o = \overline{A + B}$):}
\begin{itemize}[leftmargin=*]
    \item \textbf{PUN (arriba):} Dos PMOS ($M_1, M_2$) en \textbf{serie}.
    \item \textbf{PDN (abajo):} Dos NMOS ($M_3, M_4$) en \textbf{paralelo}.
\end{itemize}
\begin{center}
\footnotesize
\renewcommand{\arraystretch}{1.15}
\begin{tabular}{cc|cccc|c}
\toprule
$A$ & $B$ & $M_1$ & $M_2$ & $M_3$ & $M_4$ & $V_o = \overline{A+B}$ \\
\midrule
0 & 0 & ON & ON & OFF & OFF & \textbf{1} \\
0 & 1 & ON & OFF & OFF & ON & \textbf{0} \\
1 & 0 & OFF & ON & ON & OFF & \textbf{0} \\
1 & 1 & OFF & OFF & ON & ON & \textbf{0} \\
\bottomrule
\end{tabular}
\end{center}
\end{minipage}
\hfill
\begin{minipage}[t]{0.48\textwidth}
\textbf{2. Puerta NAND CMOS ($V_o = \overline{A \cdot B}$):}
\begin{itemize}[leftmargin=*]
    \item \textbf{PUN (arriba):} Dos PMOS ($M_1, M_2$) en \textbf{paralelo}.
    \item \textbf{PDN (abajo):} Dos NMOS ($M_3, M_4$) en \textbf{serie}.
\end{itemize}
\begin{center}
\footnotesize
\renewcommand{\arraystretch}{1.15}
\begin{tabular}{cc|cccc|c}
\toprule
$A$ & $B$ & $M_1$ & $M_2$ & $M_3$ & $M_4$ & $V_o = \overline{A\cdot B}$ \\
\midrule
0 & 0 & ON & ON & OFF & OFF & \textbf{1} \\
0 & 1 & ON & OFF & OFF & ON & \textbf{1} \\
1 & 0 & OFF & ON & ON & OFF & \textbf{1} \\
1 & 1 & OFF & OFF & ON & ON & \textbf{0} \\
\bottomrule
\end{tabular}
\end{center}
\end{minipage}

\vspace{6pt}
\begin{center}
\begin{tikzpicture}[scale=0.9]
    % ---------------- PUERTA NAND CMOS (izquierda, centrada en x = -3.4) ----------------
    \node[font=\bfseries] at (-3.4,5.0) {Puerta NAND ($V_o = \overline{A \cdot B}$)};
    % PUN: 2 PMOS en paralelo desde y = 4.0 separados 2.8 cm
    \draw (-4.8,4.0) node[pmos, anchor=source] (pa) {};
    \draw (-2.0,4.0) node[pmos, anchor=source] (pb) {};
    \draw (-4.8,4.0) -- (-2.0,4.0);
    \draw (-3.4,4.0) -- (-3.4,4.4) node[above] {$V_{DD}$};
    \draw (pa.drain) -- (pb.drain);
    % Nodo de salida y conexión a PDN
    \draw (-3.4,2.6) -- (-3.4,2.2);
    \draw (-3.4,2.4) to[short, *-o] (-1.2,2.4) node[right] {$V_o$};
    % PDN: 2 NMOS en serie hacia abajo desde (-3.4, 2.2)
    \draw (-3.4,2.2) node[nmos, anchor=drain] (na) {};
    \draw (na.source) node[nmos, anchor=drain] (nb) {};
    \draw (nb.source) node[ground] {};
    % Etiquetas de puerta limpias
    \node[left] at (pa.gate) {$A$};
    \node[left] at (pb.gate) {$B$};
    \node[left] at (na.gate) {$A$};
    \node[left] at (nb.gate) {$B$};

    % ---------------- PUERTA NOR CMOS (derecha, centrada en x = 3.4) ----------------
    \node[font=\bfseries] at (3.4,5.0) {Puerta NOR ($V_o = \overline{A + B}$)};
    % PUN: 2 PMOS en serie desde y = 4.4
    \draw (3.4,4.4) node[above] {$V_{DD}$} -- (3.4,4.1);
    \draw (3.4,4.1) node[pmos, anchor=source] (p1) {};
    \draw (p1.drain) node[pmos, anchor=source] (p2) {};
    % Nodo de salida y conexión a PDN en paralelo
    \draw (p2.drain) -- (3.4,0.9);
    \draw (3.4,1.1) to[short, *-o] (5.5,1.1) node[right] {$V_o$};
    % PDN: 2 NMOS en paralelo separados 2.8 cm (x = 2.0 y x = 4.8)
    \draw (2.0,0.9) -- (4.8,0.9);
    \draw (2.0,0.9) node[nmos, anchor=drain] (n1) {};
    \draw (4.8,0.9) node[nmos, anchor=drain] (n2) {};
    \draw (n1.source) -- (n2.source);
    \draw (3.4,-0.5) node[ground] {};
    % Etiquetas de puerta limpias
    \node[left] at (p1.gate) {$A$};
    \node[left] at (p2.gate) {$B$};
    \node[left] at (n1.gate) {$A$};
    \node[left] at (n2.gate) {$B$};
\end{tikzpicture}
\end{center}

% ==============================================================================
% 5.4. COLECCIÓN COMPLETA DE EJERCICIOS RESUELTOS
% ==============================================================================
\section{Ejercicios Resueltos de Exámenes de Electrónica Digital}

% --- EJERCICIO 1 (Ejemplo de clase orig-102 y orig-103: Y = A.B + C) ---
\begin{ejercicio}[Síntesis de $Y = A \cdot B + C$ en Lógica NMOS y en Lógica CMOS --- Ejemplo de Clase]
Diseñar a nivel de transistores la función booleana no negada:
\begin{equation}
    Y = A \cdot B + C
\end{equation}
a) Empleando lógica NMOS con carga NMOS de enriquecimiento. \quad b) Empleando lógica CMOS estática complementaria.
\tcbline
\textbf{Solución Paso a Paso:}\\[4pt]
Como la función pedida $Y = A \cdot B + C$ no tiene barra de negación superior, realizamos en una primera etapa la función invertida $\overline{V_o} = \overline{A \cdot B + C}$ y le encadenamos en cascada un \textbf{inversor} en la salida:
\begin{itemize}[leftmargin=*]
    \item \textbf{Red PDN (NMOS):} El producto $A \cdot B$ corresponde a dos NMOS ($N_A, N_B$) en \textbf{serie}, y la suma $+ C$ coloca dicho brazo en \textbf{paralelo} con el transistor $N_C$.
    \item \textbf{Red PUN en CMOS (PMOS):} Por dualidad, $P_A$ y $P_B$ van en \textbf{paralelo}, y el conjunto va en \textbf{serie} con $P_C$.
\end{itemize}

\begin{center}
\resizebox{\linewidth}{!}{%
\begin{tikzpicture}[scale=0.88]
    % --- a) LÓGICA NMOS (izquierda) ---
    \node[font=\bfseries\small] at (-4.2,5.9) {a) En Lógica NMOS con Carga Activa};
    % Carga NMOS enriquecimiento arriba
    \draw (-4.5,3.0) node[nmos, anchor=source] (nl) {};
    \draw (nl.drain) -- (-4.5,4.8) node[above] {$V_{DD}$};
    \draw (-4.5,4.5) -- (-5.5,4.5) -- (-5.5,3.7) -- (nl.gate);
    \draw (-4.5,3.0) -- (-4.5,2.4);
    \draw (-5.8,2.4) -- (-3.2,2.4);
    % Rama izquierda: A serie B
    \draw (-5.8,2.4) node[nmos, anchor=drain] (na1) {};
    \draw (na1.source) node[nmos, anchor=drain] (nb1) {};
    \node[left] at (na1.gate) {$A$};
    \node[left] at (nb1.gate) {$B$};
    \draw (nb1.source) -- (-5.8,-0.4);
    % Rama derecha: C
    \draw (-3.2,2.4) node[nmos, anchor=drain] (nc1) {};
    \node[left] at (nc1.gate) {$C$};
    \draw (nc1.source) -- (-3.2,-0.4);
    \draw (-5.8,-0.4) -- (-3.2,-0.4);
    \draw (-4.5,-0.4) node[ground] {};
    % Salida negada hacia el inversor NMOS
    \draw (-4.5,2.7) to[short, *-] (-1.8,2.7) node[midway, above] {$\overline{V_o}$} -- (-1.8,0.3);
    \draw (-0.8,1.0) node[nmos, anchor=drain] (ninv) {};
    \draw (-1.8,0.3) -- (ninv.gate);
    \draw (ninv.source) node[ground] {};
    \draw (-0.8,1.6) node[nmos, anchor=source] (nlinv) {};
    \draw (-0.8,1.0) -- (-0.8,1.6);
    \draw (nlinv.drain) -- (-0.8,3.4) node[above] {$V_{DD}$};
    \draw (-0.8,3.1) -- (-1.7,3.1) -- (-1.7,2.3) -- (nlinv.gate);
    \draw (-0.8,1.3) to[short, *-o] (0.1,1.3) node[right] {$Y$};

    % --- b) LÓGICA CMOS (derecha) ---
    \node[font=\bfseries\small] at (5.0,5.9) {b) En Lógica CMOS Complementaria};
    % PUN: PA || PB arriba (x=2.6 y x=5.4), en serie con PC abajo (x=4.0)
    \draw (2.6,4.8) -- (5.4,4.8);
    \draw (4.0,4.8) -- (4.0,5.1) node[above] {$V_{DD}$};
    \draw (2.6,4.8) node[pmos, anchor=source] (pa2) {};
    \draw (5.4,4.8) node[pmos, anchor=source] (pb2) {};
    \node[left] at (pa2.gate) {$A$};
    \node[left] at (pb2.gate) {$B$};
    \draw (pa2.drain) -- (pb2.drain);
    \draw (4.0,3.4) node[pmos, anchor=source] (pc2) {};
    \node[left] at (pc2.gate) {$C$};
    \draw (pc2.drain) -- (4.0,1.6);
    % PDN: (NA serie NB) en x=2.6 en paralelo con NC en x=5.4
    \draw (2.6,1.6) -- (5.4,1.6);
    \draw (2.6,1.6) node[nmos, anchor=drain] (na2) {};
    \draw (na2.source) node[nmos, anchor=drain] (nb2) {};
    \node[left] at (na2.gate) {$A$};
    \node[left] at (nb2.gate) {$B$};
    \draw (nb2.source) -- (2.6,-1.2);
    \draw (5.4,1.6) node[nmos, anchor=drain] (nc2) {};
    \node[left] at (nc2.gate) {$C$};
    \draw (nc2.source) -- (5.4,-1.2);
    \draw (2.6,-1.2) -- (5.4,-1.2);
    \draw (4.0,-1.2) node[ground] {};
    % Conexión al inversor CMOS en x=7.6
    \draw (4.0,1.8) to[short, *-] (6.5,1.8) node[midway, above] {$\overline{V_o}$};
    \draw (7.6,3.2) node[above] {$V_{DD}$} -- (7.6,2.9);
    \draw (7.6,2.9) node[pmos, anchor=source] (pinv) {};
    \draw (7.6,1.1) node[nmos, anchor=drain] (ninv2) {};
    \draw (ninv2.source) node[ground] {};
    \draw (pinv.drain) -- (ninv2.drain);
    \draw (pinv.gate) -- (6.5,2.2) -- (6.5,0.4) -- (ninv2.gate);
    \draw (7.6,1.3) to[short, *-o] (8.5,1.3) node[right] {$Y$};
\end{tikzpicture}%
}
\end{center}
\end{ejercicio}

% --- EJERCICIO 2 (Síntesis CMOS 8 transistores) ---
\begin{ejercicio}[Síntesis de Puerta CMOS de 8 Transistores: $Y = \overline{(A + B) \cdot C + D}$]
Diseñar a nivel de transistores la puerta estática complementaria CMOS que implementa directamente la función booleana negada:
\begin{equation}
    Y = \overline{(A + B) \cdot C + D}
\end{equation}
\tcbline
\begin{center}
\begin{tikzpicture}[scale=0.92]
    % Alimentación VDD arriba
    \draw (0,6.6) node[above] {$V_{DD}$} -- (0,6.2);

    % RED PUN: PD arriba en serie con [PC en paralelo con (PA en serie con PB)]
    \draw (0,6.2) node[pmos, anchor=source] (pd) {};
    \node[left] at (pd.gate) {$D$};
    \draw (-1.5,4.8) -- (1.5,4.8);
    % Rama izquierda PUN: PC (x = -1.5)
    \draw (-1.5,4.8) node[pmos, anchor=source] (pc) {};
    \node[left] at (pc.gate) {$C$};
    \draw (pc.drain) -- (-1.5,2.0);
    % Rama derecha PUN: PA en serie con PB (x = 1.5, separación 3.0 cm)
    \draw (1.5,4.8) node[pmos, anchor=source] (pa) {};
    \node[left] at (pa.gate) {$A$};
    \draw (pa.drain) node[pmos, anchor=source] (pb) {};
    \node[left] at (pb.gate) {$B$};
    \draw (-1.5,2.0) -- (1.5,2.0);

    % Nodo de salida Y
    \draw (0,2.0) -- (0,1.4);
    \draw (0,1.7) to[short, *-o] (3.0,1.7) node[right, font=\bfseries] {$Y = \overline{(A+B)\cdot C + D}$};

    % RED PDN: [NC en serie con (NA || NB)] en paralelo con ND
    \draw (-1.5,1.4) -- (2.0,1.4);
    % Rama izquierda PDN: NC en x = -1.5, debajo NA (x = -2.8) || NB (x = -0.2)
    \draw (-1.5,1.4) node[nmos, anchor=drain] (nc) {};
    \node[left] at (nc.gate) {$C$};
    \draw (-2.8,0.0) -- (-0.2,0.0);
    \draw (-2.8,0.0) node[nmos, anchor=drain] (na) {};
    \node[left] at (na.gate) {$A$};
    \draw (-0.2,0.0) node[nmos, anchor=drain] (nb) {};
    \node[left] at (nb.gate) {$B$};
    \draw (na.source) -- (nb.source);
    % Rama derecha PDN: ND en x = 2.5
    \draw (2.0,1.4) -- (2.5,1.4);
    \draw (2.5,1.4) node[nmos, anchor=drain] (nd) {};
    \node[left] at (nd.gate) {$D$};
    \draw (nd.source) -- (2.5,-1.4) -- (-1.5,-1.4);
    \draw (0,-1.4) node[ground] {};
\end{tikzpicture}
\end{center}
\textbf{Justificación Topológica:}\\[2pt]
- \textbf{Red PDN (4 NMOS):} $N_A \parallel N_B$ (suma $A+B$) en serie con $N_C$ (producto por $C$), y todo ello en paralelo con $N_D$ (suma $+D$).\\[2pt]
- \textbf{Red PUN (4 PMOS):} $P_A$ en serie con $P_B$, este bloque en paralelo con $P_C$, y el conjunto en serie con $P_D$ hacia $V_{DD}$. Total: \textbf{8 transistores}.
\end{ejercicio}

% --- EJERCICIO 3 (Examen Enero 2020 orig-104: Y = A.B + C.(D+E)) ---
\begin{ejercicio}[Síntesis de $Y = A \cdot B + C \cdot (D + E)$ --- Examen Enero 2020 FFT]
Realizar la función lógica $Y = A \cdot B + C \cdot (D + E)$:
a) Con lógica NMOS con carga NMOS de enriquecimiento. \quad b) Con lógica CMOS complementaria.
\tcbline
\textbf{Solución Paso a Paso:}\\[4pt]
La función tiene 5 entradas $\{A, B, C, D, E\}$ y no está negada, por lo que sintetizamos primero $\overline{Y} = \overline{A \cdot B + C \cdot (D + E)}$ y añadimos una etapa inversora a la salida:
\begin{itemize}[leftmargin=*]
    \item \textbf{Red de Pull-Down (NMOS, común a ambos apartados):}
    Rama 1 con $N_A$ y $N_B$ en \textbf{serie} ($A \cdot B$), conectada en \textbf{paralelo} con la Rama 2 formada por $N_C$ en \textbf{serie} con el paralelo $(N_D \parallel N_E)$ que realiza $C \cdot (D + E)$.
    \item \textbf{Red de Pull-Up en CMOS (PMOS):}
    Por dualidad de De Morgan, el bloque superior $(P_A \parallel P_B)$ se conecta en \textbf{serie} con el bloque inferior formado por $P_C$ en \textbf{paralelo} con la rama $(P_D \text{ en serie con } P_E)$.
\end{itemize}

\begin{center}
\begin{tikzpicture}[scale=0.88]
    \node[font=\bfseries] at (1.5,4.8) {a) Implementación en Lógica NMOS con Carga de Enriquecimiento};
    % Carga NMOS de enriquecimiento arriba en x = 0
    \draw (0,2.4) node[nmos, anchor=source] (nload) {};
    \draw (nload.drain) -- (0,4.1) node[above] {$V_{DD}$};
    \draw (0,3.8) -- (-1.0,3.8) -- (-1.0,3.1) -- (nload.gate);
    \draw (0,2.4) -- (0,1.8);
    % Dos ramas principales en paralelo: x = -2.2 (A serie B) y x = 1.6 (C serie (D || E))
    \draw (-2.2,1.8) -- (1.6,1.8);
    \draw (-2.2,1.8) node[nmos, anchor=drain] (na) {};
    \node[left] at (na.gate) {$A$};
    \draw (na.source) node[nmos, anchor=drain] (nb) {};
    \node[left] at (nb.gate) {$B$};
    \draw (1.6,1.8) node[nmos, anchor=drain] (nc) {};
    \node[left] at (nc.gate) {$C$};
    \draw (0.3,0.4) -- (2.9,0.4);
    \draw (0.3,0.4) node[nmos, anchor=drain] (nd) {};
    \node[left] at (nd.gate) {$D$};
    \draw (2.9,0.4) node[nmos, anchor=drain] (ne) {};
    \node[left] at (ne.gate) {$E$};
    \draw (nd.source) -- (ne.source);
    \draw (-2.2,-1.0) -- (1.6,-1.0);
    \draw (0,-1.0) node[ground] {};
    % Salida negada hacia el inversor NMOS a la derecha (x = 5.4)
    \draw (0,2.0) to[short, *-] (4.2,2.0) node[midway, above] {$\overline{Y}$} -- (4.2,-0.1);
    \draw (5.4,0.6) node[nmos, anchor=drain] (ninv) {};
    \draw (4.2,-0.1) -- (ninv.gate);
    \draw (ninv.source) node[ground] {};
    \draw (5.4,1.2) node[nmos, anchor=source] (nlinv) {};
    \draw (5.4,0.6) -- (5.4,1.2);
    \draw (nlinv.drain) -- (5.4,3.0) node[above] {$V_{DD}$};
    \draw (5.4,2.7) -- (4.4,2.7) -- (4.4,1.9) -- (nlinv.gate);
    \draw (5.4,0.9) to[short, *-o] (6.5,0.9) node[right, font=\bfseries] {$Y$};
\end{tikzpicture}
\end{center}

\vspace{4pt}
\begin{center}
\begin{tikzpicture}[scale=0.88]
    \node[font=\bfseries] at (1.8,7.2) {b) Implementación en Lógica CMOS Complementaria (12 Transistores)};
    % RED PUN: Arriba [PC || (PD serie PE)] y debajo en serie (PA || PB)
    \draw (0,6.4) -- (0,6.7) node[above] {$V_{DD}$};
    \draw (-1.5,6.4) -- (1.5,6.4);
    % PC a la izquierda (x = -1.5), PD serie PE a la derecha (x = 1.5)
    \draw (-1.5,6.4) node[pmos, anchor=source] (pc) {};
    \node[left] at (pc.gate) {$C$};
    \draw (pc.drain) -- (-1.5,3.6);
    \draw (1.5,6.4) node[pmos, anchor=source] (pd) {};
    \node[left] at (pd.gate) {$D$};
    \draw (pd.drain) node[pmos, anchor=source] (pe) {};
    \node[left] at (pe.gate) {$E$};
    \draw (-1.5,3.6) -- (1.5,3.6);
    % En serie debajo: PA (x = -1.5) || PB (x = 1.5)
    \draw (-1.5,3.6) node[pmos, anchor=source] (pa) {};
    \node[left] at (pa.gate) {$A$};
    \draw (1.5,3.6) node[pmos, anchor=source] (pb) {};
    \node[left] at (pb.gate) {$B$};
    \draw (pa.drain) -- (pb.drain);

    % Conexión intermedia entre PUN y PDN
    \draw (0,2.2) -- (0,1.6);

    % RED PDN: (NA serie NB en x = -2.2) || [NC en x = 1.6 serie (ND x = 0.3 || NE x = 2.9)]
    \draw (-2.2,1.6) -- (1.6,1.6);
    \draw (-2.2,1.6) node[nmos, anchor=drain] (na) {};
    \node[left] at (na.gate) {$A$};
    \draw (na.source) node[nmos, anchor=drain] (nb) {};
    \node[left] at (nb.gate) {$B$};
    \draw (1.6,1.6) node[nmos, anchor=drain] (nc) {};
    \node[left] at (nc.gate) {$C$};
    \draw (0.3,0.2) -- (2.9,0.2);
    \draw (0.3,0.2) node[nmos, anchor=drain] (nd) {};
    \node[left] at (nd.gate) {$D$};
    \draw (2.9,0.2) node[nmos, anchor=drain] (ne) {};
    \node[left] at (ne.gate) {$E$};
    \draw (nd.source) -- (ne.source);
    \draw (-2.2,-1.2) -- (1.6,-1.2);
    \draw (0,-1.2) node[ground] {};

    % Etapa inversora CMOS de salida (x = 5.6)
    \draw (0,1.9) to[short, *-] (4.4,1.9) node[midway, above] {$\overline{Y}$};
    \draw (5.6,3.3) node[above] {$V_{DD}$} -- (5.6,3.0);
    \draw (5.6,3.0) node[pmos, anchor=source] (pinv) {};
    \draw (5.6,1.2) node[nmos, anchor=drain] (ninv) {};
    \draw (ninv.source) node[ground] {};
    \draw (pinv.drain) -- (ninv.drain);
    \draw (pinv.gate) -- (4.4,2.3) -- (4.4,0.5) -- (ninv.gate);
    \draw (5.6,1.4) to[short, *-o] (6.6,1.4) node[right, font=\bfseries] {$Y$};
\end{tikzpicture}
\end{center}
\end{ejercicio}

% --- EJERCICIO 4 (Examen Enero 2019 orig-105: Y = A + B.(C+D)) ---
\begin{ejercicio}[Síntesis de $Y = A + B \cdot (C + D)$ --- Examen Enero 2019 FFT]
Realizar la función lógica $Y = A + B \cdot (C + D)$:
a) Con lógica NMOS con carga NMOS de deplexión. \quad b) Con lógica CMOS complementaria.
\tcbline
\textbf{Solución Paso a Paso:}\\[4pt]
Sintetizamos en la primera etapa $\overline{Y} = \overline{A + B \cdot (C + D)}$ y colocamos un inversor en la salida:
\begin{itemize}[leftmargin=*]
    \item \textbf{Red PDN (NMOS):} El transistor $N_A$ va en \textbf{paralelo} con la rama formada por $N_B$ en \textbf{serie} con el paralelo $(N_C \parallel N_D)$.
    \item \textbf{Red PUN en CMOS (PMOS):} El bloque superior formado por $P_B$ en \textbf{paralelo} con la rama $(P_C \text{ en serie con } P_D)$ se conecta en \textbf{serie} con $P_A$.
\end{itemize}

\begin{center}
\begin{tikzpicture}[scale=0.88]
    \node[font=\bfseries] at (1.5,4.8) {a) Implementación en Lógica NMOS con Carga de Deplexión};
    % Carga NMOS de deplexión arriba en x = 0 (puerta unida a fuente)
    \draw (0,2.4) node[nmos, anchor=source] (nload) {};
    \draw[ultra thick] ($(nload.source)+(-0.18,0.22)$) -- ($(nload.drain)+(-0.18,-0.22)$);
    \draw (nload.drain) -- (0,4.1) node[above] {$V_{DD}$};
    \draw (nload.gate) -- (-1.0,3.1) -- (-1.0,2.2) -- (0,2.2);
    \draw (0,2.4) -- (0,1.8);
    % Dos ramas en paralelo: NA en x = -2.2, y NB serie (NC || ND) en x = 1.6
    \draw (-2.2,1.8) -- (1.6,1.8);
    \draw (-2.2,1.8) node[nmos, anchor=drain] (na) {};
    \node[left] at (na.gate) {$A$};
    \draw (na.source) -- (-2.2,-1.0);
    \draw (1.6,1.8) node[nmos, anchor=drain] (nb) {};
    \node[left] at (nb.gate) {$B$};
    \draw (0.3,0.4) -- (2.9,0.4);
    \draw (0.3,0.4) node[nmos, anchor=drain] (nc) {};
    \node[left] at (nc.gate) {$C$};
    \draw (2.9,0.4) node[nmos, anchor=drain] (nd) {};
    \node[left] at (nd.gate) {$D$};
    \draw (nc.source) -- (nd.source);
    \draw (-2.2,-1.0) -- (1.6,-1.0);
    \draw (0,-1.0) node[ground] {};
    % Inversor con carga de deplexión a la derecha (x = 5.4)
    \draw (0,2.0) to[short, *-] (4.2,2.0) node[midway, above] {$\overline{Y}$} -- (4.2,-0.1);
    \draw (5.4,0.6) node[nmos, anchor=drain] (ninv) {};
    \draw (4.2,-0.1) -- (ninv.gate);
    \draw (ninv.source) node[ground] {};
    \draw (5.4,1.2) node[nmos, anchor=source] (nlinv) {};
    \draw[ultra thick] ($(nlinv.source)+(-0.18,0.22)$) -- ($(nlinv.drain)+(-0.18,-0.22)$);
    \draw (5.4,0.6) -- (5.4,1.2);
    \draw (nlinv.drain) -- (5.4,3.0) node[above] {$V_{DD}$};
    \draw (nlinv.gate) -- (4.4,1.9) -- (4.4,1.0) -- (5.4,1.0);
    \draw (5.4,0.9) to[short, *-o] (6.5,0.9) node[right, font=\bfseries] {$Y$};
\end{tikzpicture}
\end{center}

\vspace{4pt}
\begin{center}
\begin{tikzpicture}[scale=0.88]
    \node[font=\bfseries] at (1.8,7.2) {b) Implementación en Lógica CMOS Complementaria (10 Transistores)};
    % RED PUN: Arriba [PB || (PC serie PD)] y debajo en serie PA
    \draw (0,6.4) -- (0,6.7) node[above] {$V_{DD}$};
    \draw (-1.5,6.4) -- (1.5,6.4);
    \draw (-1.5,6.4) node[pmos, anchor=source] (pb) {};
    \node[left] at (pb.gate) {$B$};
    \draw (pb.drain) -- (-1.5,3.6);
    \draw (1.5,6.4) node[pmos, anchor=source] (pc) {};
    \node[left] at (pc.gate) {$C$};
    \draw (pc.drain) node[pmos, anchor=source] (pd) {};
    \node[left] at (pd.gate) {$D$};
    \draw (-1.5,3.6) -- (1.5,3.6);
    \draw (0,3.6) node[pmos, anchor=source] (pa) {};
    \node[left] at (pa.gate) {$A$};
    \draw (pa.drain) -- (0,1.6);

    % RED PDN: NA (x = -2.2) || [NB (x = 1.6) serie (NC x = 0.3 || ND x = 2.9)]
    \draw (-2.2,1.6) -- (1.6,1.6);
    \draw (-2.2,1.6) node[nmos, anchor=drain] (na) {};
    \node[left] at (na.gate) {$A$};
    \draw (na.source) -- (-2.2,-1.2);
    \draw (1.6,1.6) node[nmos, anchor=drain] (nb) {};
    \node[left] at (nb.gate) {$B$};
    \draw (0.3,0.2) -- (2.9,0.2);
    \draw (0.3,0.2) node[nmos, anchor=drain] (nc) {};
    \node[left] at (nc.gate) {$C$};
    \draw (2.9,0.2) node[nmos, anchor=drain] (nd) {};
    \node[left] at (nd.gate) {$D$};
    \draw (nc.source) -- (nd.source);
    \draw (-2.2,-1.2) -- (1.6,-1.2);
    \draw (0,-1.2) node[ground] {};

    % Etapa inversora CMOS de salida (x = 5.6)
    \draw (0,1.9) to[short, *-] (4.4,1.9) node[midway, above] {$\overline{Y}$};
    \draw (5.6,3.3) node[above] {$V_{DD}$} -- (5.6,3.0);
    \draw (5.6,3.0) node[pmos, anchor=source] (pinv) {};
    \draw (5.6,1.2) node[nmos, anchor=drain] (ninv) {};
    \draw (ninv.source) node[ground] {};
    \draw (pinv.drain) -- (ninv.drain);
    \draw (pinv.gate) -- (4.4,2.3) -- (4.4,0.5) -- (ninv.gate);
    \draw (5.6,1.4) to[short, *-o] (6.6,1.4) node[right, font=\bfseries] {$Y$};
\end{tikzpicture}
\end{center}
\end{ejercicio}

% --- EJERCICIO 5 (Márgenes de Ruido, Fan-Out y Potencia Dinámica) ---
\begin{ejercicio}[Cálculo de Márgenes de Ruido, Fan-Out y Disipación de Potencia]
Una familia lógica digital alimentada a $V_{DD} = 3{,}3\ \si{\volt}$ presenta los siguientes parámetros garantizados por el fabricante:
$V_{OH} = 3{,}0\ \si{\volt}$, $V_{OL} = 0{,}2\ \si{\volt}$, $V_{IH} = 2{,}0\ \si{\volt}$, $V_{IL} = 0{,}8\ \si{\volt}$; corrientes máximas de salida $|I_{OH}| = 400\ \si{\micro\ampere}$, $I_{OL} = 8\ \si{\milli\ampere}$ y corrientes de entrada $|I_{IH}| = 20\ \si{\micro\ampere}$, $|I_{IL}| = 0{,}4\ \si{\milli\ampere}$.
La capacidad de carga equivalente en la salida es $C_L = 20\ \si{\pico\farad}$, conmutando a una frecuencia $f = 50\ \si{\mega\hertz}$, con tiempos de propagación $t_{pHL} = 1{,}8\ \si{\nano\second}$ y $t_{pLH} = 2{,}2\ \si{\nano\second}$.
\begin{enumerate}[label=\alph*)]
    \item Calcular los márgenes de ruido en nivel alto y bajo ($NM_H, NM_L$) y el margen de ruido global.
    \item Determinar el factor de carga o abanico de salida (\textit{Fan-Out}).
    \item Calcular el tiempo de retardo medio $t_p$, la potencia dinámica disipada $P_D$ y el producto retardo-potencia $PDP$.
\end{enumerate}
\tcbline
\textbf{Solución Paso a Paso:}\\[4pt]
a) \textbf{Márgenes de ruido:}
\begin{align}
    NM_H &= V_{OH} - V_{IH} = 3{,}0\ \si{\volt} - 2{,}0\ \si{\volt} = 1{,}0\ \si{\volt} \\
    NM_L &= V_{IL} - V_{OL} = 0{,}8\ \si{\volt} - 0{,}2\ \si{\volt} = 0{,}6\ \si{\volt} \implies NM = \min(1{,}0; 0{,}6) = 0{,}6\ \si{\volt}
\end{align}
b) \textbf{Factor de carga (\textit{Fan-Out}):}
\begin{equation}
    FO_H = \left\lfloor \frac{|I_{OH}|}{|I_{IH}|} \right\rfloor = \frac{400\ \si{\micro\ampere}}{20\ \si{\micro\ampere}} = 20, \qquad FO_L = \left\lfloor \frac{I_{OL}}{|I_{IL}|} \right\rfloor = \frac{8\ \si{\milli\ampere}}{0{,}4\ \si{\milli\ampere}} = 20 \implies FO = \min(20, 20) = 20\text{ puertas}
\end{equation}
c) \textbf{Retardo de propagación, potencia dinámica y $PDP$:}
\begin{align}
    t_p &= \frac{t_{pHL} + t_{pLH}}{2} = \frac{1{,}8 + 2{,}2}{2} = 2{,}0\ \si{\nano\second} \\
    P_D &= C_L \cdot V_{DD}^2 \cdot f = (20 \times 10^{-12}\ \si{\farad}) \times (3{,}3\ \si{\volt})^2 \times (50 \times 10^6\ \si{\hertz}) = 10{,}89\ \si{\milli\watt} \\
    PDP &= P_D \cdot t_p = (10{,}89 \times 10^{-3}\ \si{\watt}) \times (2{,}0 \times 10^{-9}\ \si{\second}) = 21{,}78\ \si{\pico\joule}
\end{align}
\end{ejercicio}

\chapter{Amplificadores Operacionales y Conversión de Señales}
\label{cap:tema6}

\cuadrofuentes{
    \textbf{Asignatura:} Fundamentos Físicos y Tecnológicos (1.º Grado en Ingeniería Informática, UGR).\\[2pt]
    \textbf{Profesor:} Pedro Cartujo (Departamento de Electrónica y Tecnología de Computadores).\\[2pt]
    \textbf{Bibliografía de Referencia:} Franco, S., \textit{Design with Operational Amplifiers and Analog Integrated Circuits}, McGraw-Hill; Sedra, A. S., Smith, K. C., \textit{Circuitos Microelectrónicos}, Oxford University Press.\\[2pt]
    \textbf{Contenido:} El amplificador operacional real e ideal, modelo equivalente interno, curva de transferencia y límites de saturación, realimentación negativa (cortocircuito virtual) y positiva, catálogo completo de configuraciones lineales y no lineales deducidas por Kirchhoff, conversión A/D y D/A, y colección completa de ejercicios resueltos paso a paso.
}

% ==============================================================================
% 6.1. EL AMPLIFICADOR OPERACIONAL Y EL CORTOCIRCUITO VIRTUAL
% ==============================================================================
\section{El Amplificador Operacional y el Cortocircuito Virtual}

El \textbf{amplificador operacional} (denominado habitualmente \textit{Op-Amp} o AO) es un circuito integrado analógico de acoplamiento directo, constituido internamente por decenas de transistores, resistencias y fuentes de corriente integradas, que ofrece una ganancia de tensión diferencial extraordinariamente elevada. Su nombre histórico proviene de su concepción original en los computadores analógicos para ejecutar operaciones matemáticas sobre señales de tensión (suma, resta, multiplicación por una constante, integración y derivación).

\subsection{Terminales, Símbolo Circuital y Modelo Equivalente Interno}

El amplificador operacional posee cinco terminales fundamentales desde el punto de vista circuital externo:
\begin{itemize}[leftmargin=*]
    \item $v^+$ (\textbf{Entrada No Inversora}): Terminal marcado con el signo $+$. Una señal aplicada en este terminal produce a la salida una tensión en fase (del mismo signo).
    \item $v^-$ (\textbf{Entrada Inversora}): Terminal marcado con el signo $-$. Una señal aplicada en este terminal produce a la salida una tensión desfasada $180^\circ$ (de signo opuesto).
    \item $v_o$ (\textbf{Terminal de Salida}): Proporciona la tensión amplificada medida respecto a la referencia común (masa o tierra).
    \item $+V_{CC}$ (o $V_{pos}$) y $-V_{EE}$ (o $V_{neg}$) (\textbf{Terminales de Alimentación DC}): Suministran la energía eléctrica continua necesaria para polarizar los transistores internos del circuito integrado (típicamente $\pm 15\ \si{\volt}$ o $\pm 12\ \si{\volt}$). Estos terminales imponen los límites físicos de excursión de la salida (tensiones de saturación $+V_{sat}$ y $-V_{sat}$):
    \begin{equation}
        -V_{sat} \le v_o \le +V_{sat} \qquad (\text{donde } +V_{sat} \approx V_{CC} \text{ y } -V_{sat} \approx -V_{EE})
    \end{equation}
\end{itemize}

En su región lineal de operación, el comportamiento interno del amplificador operacional se modela mediante una \textbf{resistencia de entrada diferencial} $R_i$ conectada entre $v^+$ y $v^-$, y una \textbf{fuente de tensión controlada por tensión} de valor $A_v(v^+ - v^-)$ en serie con una \textbf{resistencia de salida} $R_o$:

\begin{center}
\begin{tikzpicture}[scale=0.95]
    % --- Símbolo Circuital (Izquierda) ---
    \node[font=\bfseries\sffamily, color=fftDarkBlue] at (-3.8, 2.5) {Símbolo Circuital Externo};
    \draw (-3.8,0) node[op amp] (op) {};
    \draw (op.-) -- ++(-0.6,0) to[short, -o] ++(-0.3,0) node[left] {$v^-$ (Inversora)};
    \draw (op.+) -- ++(-0.6,0) to[short, -o] ++(-0.3,0) node[left] {$v^+$ (No inversora)};
    \draw (op.out) to[short, -o] ++(0.7,0) node[right] {$v_o$};
    \draw (op.up) -- ++(0,0.6) to[short, -o] ++(0,0.2) node[above] {$+V_{CC}\ (+15\,\si{\volt})$};
    \draw (op.down) -- ++(0,-0.6) to[short, -o] ++(0,-0.2) node[below] {$-V_{EE}\ (-15\,\si{\volt})$};

    % --- Modelo Equivalente Interno (Derecha) ---
    \node[font=\bfseries\sffamily, color=fftDarkBlue] at (3.6, 2.6) {Modelo Equivalente Interno};
    % Triángulo envolvente punteado amplio
    \draw[dashed, thick, fftPrimary] (1.2, 2.3) -- (1.2, -2.3) -- (6.7, 0) -- cycle;
    % Entradas v- y v+
    \draw (0.3, 0.9) node[left] {$v^-$} to[short, o-, i=$i^-$] (1.8, 0.9);
    \draw (0.3, -0.9) node[left] {$v^+$} to[short, o-, i=$i^+$] (1.8, -0.9);
    \draw (1.8, 0.9) to[R, l=$R_i$] (1.8, -0.9);
    % Fuente dependiente y resistencia de salida
    \draw (3.2, -1.2) node[ground] {} to[V, l={\footnotesize $A_v(v^+ - v^-)$}] (3.2, 0)
          to[R, l=$R_o$] (5.6, 0) to[short, -o] (6.9, 0) node[right] {$v_o$};
\end{tikzpicture}
\end{center}

\subsection{Curva de Transferencia en Lazo Abierto y Aproximación Ideal}

Si representamos la tensión de salida $v_o$ en función de la tensión diferencial de entrada $v_d = v^+ - v^-$ cuando el operacional trabaja sin red exterior (en lazo abierto), se distinguen tres regiones:
\begin{enumerate}[leftmargin=*]
    \item \textbf{Región Lineal Estrecha:} Para diferencias $|v^+ - v^-| < \frac{V_{sat}}{A_v}$, la salida es directamente proporcional a la entrada diferencial:
    \begin{equation}
        v_o = A_v (v^+ - v^-)
    \end{equation}
    Dado que $A_v$ es del orden de $10^5$ a $10^6\ \si{\volt\per\volt}$, la pendiente en el origen es prácticamente vertical.
    \item \textbf{Saturación Positiva ($+V_{sat}$ o $V_{pos}$):} Cuando $v^+ - v^- > \frac{+V_{sat}}{A_v}$, la salida queda recortada en su techo de alimentación $+V_{sat}$.
    \item \textbf{Saturación Negativa ($-V_{sat}$ o $V_{neg}$):} Cuando $v^+ - v^- < \frac{-V_{sat}}{A_v}$, la salida queda recortada en su suelo de alimentación $-V_{sat}$.
\end{enumerate}

\begin{center}
\begin{tikzpicture}[scale=0.95]
    % Ejes coordenados
    \draw[->, >=stealth, thick, gray] (-4.0,0) -- (4.0,0) node[right, text=black] {$v_d = v^+ - v^-$};
    \draw[->, >=stealth, thick, gray] (0,-2.4) -- (0,2.4) node[above, text=black] {$v_o$};
    % Líneas asíntotas de saturación
    \draw[dashed, gray] (-3.6, 1.6) -- (3.6, 1.6);
    \draw[dashed, gray] (-3.6, -1.6) -- (3.6, -1.6);
    \node[left, color=fftDarkBlue] at (-3.6, 1.6) {$+V_{sat}\ (V_{pos})$};
    \node[right, color=fftDarkBlue] at (3.6, -1.6) {$-V_{sat}\ (V_{neg})$};
    % Curva de transferencia
    \draw[very thick, fftAccent] (-3.6, -1.6) -- (-0.35, -1.6) -- (0.35, 1.6) -- (3.6, 1.6);
    % Anotaciones de zonas
    \node[above, font=\footnotesize\bfseries, color=fftAccent] at (2.2, 1.65) {Saturación positiva};
    \node[below, font=\footnotesize\bfseries, color=fftAccent] at (-2.2, -1.65) {Saturación negativa};
    \node[right, font=\footnotesize, color=fftPrimary] at (0.4, 0.6) {Zona lineal: pendiente $A_v \to \infty$};
\end{tikzpicture}
\end{center}

\begin{definicion}[Características del Amplificador Operacional Ideal]
Para el análisis de circuitos electrónicos utilizamos el \textbf{modelo ideal del amplificador operacional}, definido por las siguientes condiciones límite:
\begin{enumerate}[leftmargin=*]
    \item \textbf{Ganancia de Tensión Diferencial Infinita:} $A_v \to \infty$ (en un operacional comercial real típico, $A_v \sim 10^5 - 10^6$).
    \item \textbf{Resistencia de Entrada Infinita:} $R_i \to \infty$ (en tecnología bipolar $R_i \sim \si{\mega\ohm}$ y en tecnología CMOS/JFET $R_i \sim 10^{12}\ \si{\ohm}$). Como consecuencia directa, \textbf{las corrientes que entran a los terminales inversor y no inversor son estrictamente nulas}:
    \begin{equation}
        i^+ = 0, \qquad i^- = 0
    \end{equation}
    \item \textbf{Resistencia de Salida Nula:} $R_o = 0\ \si{\ohm}$, de modo que la tensión de salida $v_o = A_v(v^+ - v^-)$ no cae al conectar cargas exteriores.
    \item \textbf{Ancho de Banda Infinito ($BW \to \infty$) y Rechazo al Modo Común Infinito ($CMRR \to \infty$):} Amplifica por igual cualquier frecuencia y responde exclusivamente a la diferencia $(v^+ - v^-)$.
\end{enumerate}
\end{definicion}

\subsection{Realimentación Negativa vs. Realimentación Positiva y Cortocircuito Virtual}

Se denomina \textbf{realimentación} (\textit{feedback}) a toda red que establece una conexión entre el terminal de salida $v_o$ y alguno de los terminales de entrada del operacional:

\begin{itemize}[leftmargin=*]
    \item \textbf{Realimentación Negativa (Estabilizadora --- Opera en Zona Lineal):}
    La salida $v_o$ se reconecta a la \textbf{entrada inversora} $v^-$ (por ejemplo, mediante un divisor resistivo $R_1 - R_2$). Con el paso del tiempo, el lazo corrige cualquier desequilibrio hasta forzar la igualdad de potenciales $v^- = v^+$:
    \begin{align*}
        \text{Si } v^- > v^+ &\implies (v^+ - v^-) < 0 \implies v_o \downarrow \implies v^- \downarrow \implies (v^+ - v^-) \uparrow \implies \text{equilibrio } v^+ = v^- \\
        \text{Si } v^- < v^+ &\implies (v^+ - v^-) > 0 \implies v_o \uparrow \implies v^- \uparrow \implies (v^+ - v^-) \downarrow \implies \text{equilibrio } v^+ = v^-
    \end{align*}
    \item \textbf{Realimentación Positiva (Regenerativa --- Lleva a Saturación):}
    La salida $v_o$ se reconecta a la \textbf{entrada no inversora} $v^+$. En este caso \textbf{no se igualan} $v^-$ y $v^+$; cualquier pequeña diferencia inicial se amplifica regenerativamente hasta saturar la salida en $+V_{CC}$ o $-V_{CC}$:
    \begin{align*}
        \text{Si } v^- > v^+ &\implies (v^+ - v^-) < 0 \implies v_o \downarrow \implies v^+ \downarrow \implies (v^+ - v^-) \downarrow \implies v_o \text{ se limita a } -V_{CC} \\
        \text{Si } v^- < v^+ &\implies (v^+ - v^-) > 0 \implies v_o \uparrow \implies v^+ \uparrow \implies (v^+ - v^-) \uparrow \implies v_o \text{ se limita a } +V_{CC}
    \end{align*}
\end{itemize}

\begin{ley}[Principio del Cortocircuito Virtual y Reglas de Oro]
Cuando un amplificador operacional ideal dispone de \textbf{realimentación negativa} y su salida se mantiene dentro de los límites lineales ($-V_{sat} < v_o < +V_{sat}$), el cociente finito entre la salida y la ganancia infinita obliga a que la tensión diferencial de entrada sea nula:
\begin{equation}
    v^+ - v^- = \frac{v_o}{A_v} \xrightarrow{A_v \to \infty} 0 \implies \mathbf{v^+ = v^-}
\end{equation}
Este fenómeno recibe el nombre de \textbf{cortocircuito virtual}: ambos terminales están exactamente al mismo potencial eléctrico ($v^+ = v^-$) pero sin que circule corriente entre ellos ($i^+ = i^- = 0$). Si además uno de los terminales está conectado directamente a tierra ($0\ \si{\volt}$), el otro terminal queda a $0\ \si{\volt}$ sin estar físicamente conectado a masa, lo que se denomina \textbf{masa virtual}.
\end{ley}

% ==============================================================================
% 6.2. CONFIGURACIONES LINEALES BÁSICAS CON OP-AMP
% ==============================================================================
\section{Configuraciones Lineales Básicas con Op-Amp}

A continuación se deducen paso a paso, aplicando las leyes de Kirchhoff y las dos reglas de oro del operacional ideal ($v^+ = v^-$ e $i^+ = i^- = 0$), todas las configuraciones lineales fundamentales del temario.

\subsection{1. Amplificador Inversor}

La tensión de entrada $V_i$ se aplica a través de la resistencia $R_1$ al terminal inversor ($v^-$), mientras que la resistencia de realimentación $R_2$ conecta la salida $V_o$ con $v^-$. El terminal no inversor se conecta directamente a masa ($V^+ = 0\ \si{\volt}$).

\begin{center}
\begin{tikzpicture}[scale=0.95]
    \draw (0,0) node[op amp] (op) {};
    % Entrada inversora y nodo de bifurcación
    \draw (op.-) -- ++(-0.6,0) coordinate (vn);
    \draw (vn) to[R, l_={$R_1$}, i<={$I_1$}] ++(-2.4,0) to[short, -o] ++(-0.2,0) node[left] {$V_i$};
    % Lazo de realimentación limpio por arriba
    \draw (vn) -- ++(0,1.35) to[R, l={$R_2$}, i={$I_2$}] ++(3.1,0) -| (op.out);
    % Salida
    \draw (op.out) to[short, -o] ++(0.9,0) node[right] {$V_o$};
    % Entrada no inversora a masa
    \draw (op.+) -- ++(-0.5,0) node[ground] {};
    \node[above right=1pt, font=\scriptsize, color=fftAccent] at (vn) {$V^-{=}0\,\si{\volt}$};
\end{tikzpicture}
\end{center}

\begin{itemize}[leftmargin=*]
    \item \textbf{Condiciones ideales (masa virtual):} Como hay realimentación negativa, $V^- = V^+ = 0\ \si{\volt}$, e $I^- = I^+ = 0$.
    \item \textbf{Ley de las Corrientes de Kirchhoff (LCK) en el nodo inversor:} Toda la corriente $I_1$ que llega por $R_1$ debe continuar por $R_2$ hacia la salida:
    \begin{equation}
        I_1 = I_2 \implies \frac{V_i - V^-}{R_1} = \frac{V^- - V_o}{R_2}
    \end{equation}
    \item Sustituyendo $V^- = 0\ \si{\volt}$:
    \begin{equation}
        \frac{V_i}{R_1} = -\frac{V_o}{R_2} \implies \boxed{V_o = -\frac{R_2}{R_1} V_i} \qquad \left(A_v = \frac{V_o}{V_i} = -\frac{R_2}{R_1}\right)
    \end{equation}
    \item \textbf{Resistencia de entrada del circuito:} Como el nodo $V^-$ está a masa virtual ($0\ \si{\volt}$), la resistencia equivalente que ve el generador $V_i$ es exactamente $R_{in} = \frac{V_i}{I_1} = R_1$.
\end{itemize}

\subsection{2. Amplificador No Inversor}

La señal de entrada $V_i$ se aplica directamente a la entrada no inversora ($V^+$). La red de realimentación negativa está formada por el divisor resistivo $R_1$ (entre $V^-$ y masa) y $R_2$ (entre $V^-$ y la salida $V_o$).

\begin{center}
\begin{tikzpicture}[scale=0.95]
    \draw (0,0) node[op amp] (op) {};
    % Entrada inversora y nodo de bifurcación
    \draw (op.-) -- ++(-0.6,0) coordinate (vn);
    \draw (vn) to[R, l_={$R_1$}] ++(-2.2,0) node[ground] {};
    % Lazo de realimentación limpio por arriba
    \draw (vn) -- ++(0,1.35) to[R, l={$R_2$}] ++(3.1,0) -| (op.out);
    % Entrada no inversora
    \draw (op.+) -- ++(-0.8,0) to[short, -o] ++(-0.3,0) node[left] {$V_i$};
    % Salida
    \draw (op.out) to[short, -o] ++(0.9,0) node[right] {$V_o$};
    \node[above right=1pt, font=\scriptsize, color=fftAccent] at (vn) {$V^-{=}V_i$};
\end{tikzpicture}
\end{center}

\begin{itemize}[leftmargin=*]
    \item \textbf{Condiciones ideales:} Por cortocircuito virtual, $V^- = V^+ = V_i$, con $I^- = I^+ = 0$.
    \item \textbf{LCK en el nodo inversor:} Aplicando el balance de corrientes entre la rama a masa ($R_1$) y la rama de realimentación ($R_2$):
    \begin{equation}
        \frac{0 - V^-}{R_1} = \frac{V^- - V_o}{R_2}
    \end{equation}
    \item Sustituyendo $V^- = V_i$ y despejando $V_o$:
    \begin{equation}
        -\frac{V_i}{R_1} = \frac{V_i}{R_2} - \frac{V_o}{R_2} \implies \frac{V_o}{R_2} = V_i \left(\frac{1}{R_2} + \frac{1}{R_1}\right) \implies \boxed{V_o = \left(1 + \frac{R_2}{R_1}\right) V_i}
    \end{equation}
    \item \textbf{Propiedades clave:} La ganancia $A_v = 1 + \frac{R_2}{R_1} \ge 1$ es siempre positiva (no invierte la fase) y la resistencia de entrada vista por $V_i$ es idealmente infinita ($R_{in} \to \infty$, puesto que $I^+ = 0$).
\end{itemize}

\subsection{3. Seguidor de Tensión (Buffer de Ganancia Unidad)}

Es el caso límite del amplificador no inversor cuando se hace $R_2 = 0\ \si{\ohm}$ (cortocircuito directo entre la salida y la entrada inversora) y $R_1 \to \infty$ (circuito abierto a masa):

\begin{center}
\begin{tikzpicture}[scale=0.95]
    \draw (0,0) node[op amp] (op) {};
    \draw (op.-) -- ++(-0.5,0) coordinate (vn);
    \draw (vn) -- ++(0,1.2) -- ++(2.9,0) -| (op.out);
    \draw (op.+) -- ++(-0.9,0) to[short, -o] ++(-0.3,0) node[left] {$V_i$};
    \draw (op.out) to[short, -o] ++(0.9,0) node[right] {$V_o = V_i$};
\end{tikzpicture}
\end{center}

\begin{equation}
    \boxed{V_o = V_i \qquad (A_v = 1)}
\end{equation}
\begin{aviso}[Utilidad del Seguidor de Tensión como Adaptador de Impedancias]
Aunque no amplifica en tensión ($A_v = 1$), el seguidor ofrece una impedancia de entrada infinita ($R_{in} \to \infty$, no consume corriente de la fuente previa) y una impedancia de salida nula ($R_{out} = 0$, capaz de entregar corriente a la carga sin caída de tensión). Se utiliza para \textbf{aislar etapas y eliminar el efecto de carga}.
\end{aviso}

\subsection{4. Sumador Inversor y Sumador No Inversor}

\subsubsection{a) Sumador Inversor}
Varias señales de entrada $V_1, V_2, \dots, V_n$ se aplican a través de sus respectivas resistencias $R_1, R_2, \dots, R_n$ al nodo común de la entrada inversora $V^-$, mientras que $V^+ = 0\ \si{\volt}$.

\begin{center}
\begin{tikzpicture}[scale=0.95]
    \draw (0,0) node[op amp] (op) {};
    \draw (op.-) -- ++(-0.6,0) coordinate (vn);
    % Entradas V1 y V2 hacia el nodo vn
    \draw (vn) -- ++(-0.4, 0.55) to[R, l_={$R_1$}] ++(-2.1,0) to[short, -o] ++(-0.2,0) node[left] {$V_1$};
    \draw (vn) -- ++(-0.4,-0.55) to[R, l={$R_2$}] ++(-2.1,0) to[short, -o] ++(-0.2,0) node[left] {$V_2$};
    % Realimentación RF
    \draw (vn) -- ++(0,1.45) to[R, l={$R_F$}] ++(3.1,0) -| (op.out);
    % Entrada + a masa y salida
    \draw (op.+) -- ++(-0.5,0) node[ground] {};
    \draw (op.out) to[short, -o] ++(0.9,0) node[right] {$V_o$};
\end{tikzpicture}
\end{center}

Aplicando LCK en el nodo inversor con la condición ideal de masa virtual ($V^- = V^+ = 0\ \si{\volt}$):
\begin{equation}
    \frac{V_1 - V^-}{R_1} + \frac{V_2 - V^-}{R_2} = \frac{V^- - V_o}{R_F} \implies \boxed{V_o = -R_F \left(\frac{V_1}{R_1} + \frac{V_2}{R_2}\right)}
\end{equation}
Para $n$ entradas: $V_o = -R_F \sum_{k=1}^n \frac{V_k}{R_k}$. En el caso particular en que todas las resistencias sean iguales ($R_1 = R_2 = R_F$), resulta $V_o = -(V_1 + V_2)$.

\subsubsection{b) Sumador No Inversor}
Las señales $V_1$ y $V_2$ se aplican mediante $R_1$ y $R_2$ al terminal no inversor ($V^+$), mientras que la red inversora incorpora el divisor $R_S - R_F$:

\begin{center}
\begin{tikzpicture}[scale=0.95]
    \draw (0,0) node[op amp] (op) {};
    % Nodo inversor
    \draw (op.-) -- ++(-0.6,0) coordinate (vn);
    \draw (vn) to[R, l_={$R_S$}] ++(-2.1,0) node[ground] {};
    \draw (vn) -- ++(0,1.35) to[R, l={$R_F$}] ++(3.1,0) -| (op.out);
    % Nodo no inversor con suma resistiva
    \draw (op.+) -- ++(-0.6,0) coordinate (vp);
    \draw (vp) -- ++(-0.4, 0.45) to[R, l_={$R_1$}] ++(-2.1,0) to[short, -o] ++(-0.2,0) node[left] {$V_1$};
    \draw (vp) -- ++(-0.4,-0.45) to[R, l={$R_2$}] ++(-2.1,0) to[short, -o] ++(-0.2,0) node[left] {$V_2$};
    % Salida
    \draw (op.out) to[short, -o] ++(0.9,0) node[right] {$V_o$};
\end{tikzpicture}
\end{center}

\begin{itemize}[leftmargin=*]
    \item En el terminal no inversor ($I^+ = 0$):
    \begin{equation}
        I_{R1} = I_{R2} \implies \frac{V_1 - V^+}{R_1} = \frac{V^+ - V_2}{R_2} \implies V^+ = \frac{R_2 V_1 + R_1 V_2}{R_1 + R_2} = V^-
    \end{equation}
    \item En el terminal inversor ($I^- = 0$):
    \begin{equation}
        \frac{0 - V^-}{R_S} = \frac{V^- - V_o}{R_F} \implies V_o = \left(1 + \frac{R_F}{R_S}\right) V^-
    \end{equation}
    \item Sustituyendo $V^-$ se obtiene la ecuación general del sumador no inversor:
    \begin{equation}
        \boxed{V_o = \left(1 + \frac{R_F}{R_S}\right) \left(\frac{R_2 V_1 + R_1 V_2}{R_1 + R_2}\right)}
    \end{equation}
    Si se eligen resistencias simétricas $R_1 = R_2$ y $R_F = R_S$, los factores $\left(1 + 1\right) = 2$ y $\frac{1}{2}$ se cancelan y resulta la suma pura sin inversión: $V_o = V_1 + V_2$.
\end{itemize}

\subsection{5. Amplificador Diferencial (Restador)}

Permite amplificar la diferencia entre dos señales de entrada $V_1$ y $V_2$:

\begin{center}
\begin{tikzpicture}[scale=0.95]
    \draw (0,0) node[op amp] (op) {};
    % Rama inversora (Nodo B)
    \draw (op.-) -- ++(-0.6,0) coordinate (vn);
    \node[above left, font=\footnotesize, color=fftPrimary] at (vn) {$B\ (V^-)$};
    \draw (vn) to[R, l_={$R_1$}] ++(-2.2,0) to[short, -o] ++(-0.2,0) node[left] {$V_1$};
    \draw (vn) -- ++(0,1.35) to[R, l={$R_2$}] ++(3.1,0) -| (op.out);
    % Rama no inversora (Nodo A)
    \draw (op.+) -- ++(-0.6,0) coordinate (vp);
    \node[below left, font=\footnotesize, color=fftPrimary] at (vp) {$A\ (V^+)$};
    \draw (vp) to[R, l={$R_3$}] ++(-2.2,0) to[short, -o] ++(-0.2,0) node[left] {$V_2$};
    \draw (vp) to[R, l={$R_4$}] ++(0,-1.8) node[ground] {};
    % Salida
    \draw (op.out) to[short, -o] ++(0.9,0) node[right] {$V_o$};
\end{tikzpicture}
\end{center}

\begin{itemize}[leftmargin=*]
    \item \textbf{Kirchhoff en el punto $A$ ($V^+$):} Como $I^+ = 0$, las resistencias $R_3$ y $R_4$ forman un divisor de tensión:
    \begin{equation}
        I_{R3} = I_{R4} \implies \frac{V_2 - V^+}{R_3} = \frac{V^+}{R_4} \implies V^+ = \left(\frac{R_4}{R_3 + R_4}\right) V_2 = V^-
    \end{equation}
    \item \textbf{Kirchhoff en el punto $B$ ($V^-$):} Como $I^- = 0$:
    \begin{equation}
        I_{R1} = I_{R2} \implies \frac{V_1 - V^-}{R_1} = \frac{V^- - V_o}{R_2} \implies V_o = \left(1 + \frac{R_2}{R_1}\right) V^- - \frac{R_2}{R_1} V_1
    \end{equation}
    \item Sustituyendo la expresión de $V^-$:
    \begin{equation}
        \boxed{V_o = \left(1 + \frac{R_2}{R_1}\right)\left(\frac{R_4}{R_3 + R_4}\right) V_2 - \frac{R_2}{R_1} V_1}
    \end{equation}
    \item \textbf{Caso de puente resistivo equilibrado ($\frac{R_2}{R_1} = \frac{R_4}{R_3}$, por ejemplo $R_1 = R_3$ y $R_2 = R_4$):}
    \begin{equation}
        \boxed{V_o = \frac{R_2}{R_1} (V_2 - V_1)}
    \end{equation}
\end{itemize}

\subsection{6. Integrador Inversor y Derivador Inversor}

Al sustituir una de las resistencias del amplificador inversor por un condensador $C$, se obtienen operadores dinámicos en el dominio del tiempo:

\subsubsection{a) Integrador Inversor}
La resistencia $R$ se sitúa en la entrada y el condensador $C$ en el lazo de realimentación negativa:

\begin{center}
\begin{tikzpicture}[scale=0.95]
    \draw (0,0) node[op amp] (op) {};
    \draw (op.-) -- ++(-0.5,0) coordinate (vn);
    \draw (vn) to[R, l_={$R$}, i<={$i_R(t)$}] ++(-2.3,0) to[short, -o] ++(-0.2,0) node[left] {$v_i(t)$};
    \draw (vn) -- ++(0,1.35) to[C, l={$C$}, i={$i_C(t)$}] ++(3.0,0) -| (op.out);
    \draw (op.+) -- ++(-0.5,0) node[ground] {};
    \draw (op.out) to[short, -o] ++(0.9,0) node[right] {$v_o(t)$};
\end{tikzpicture}
\end{center}

Con $V^- = V^+ = 0\ \si{\volt}$ e $i_R(t) = i_C(t)$:
\begin{equation}
    \frac{v_i(t) - 0}{R} = C \frac{d(0 - v_o(t))}{dt} = -C \frac{dv_o(t)}{dt} \implies \boxed{v_o(t) = -\frac{1}{R C} \int_0^t v_i(\tau)\, d\tau + v_o(0)}
\end{equation}
En régimen permanente sinusoidal ($\mathbf{Z}_C = \frac{1}{j\omega C}$), su función de transferencia fasorial es $H(j\omega) = -\frac{1}{j\omega R C}$ (comportamiento pasa-bajos con polo en el origen).

\subsubsection{b) Derivador Inversor}
El condensador $C$ se ubica en la rama de entrada y la resistencia $R$ en el lazo de realimentación:

\begin{center}
\begin{tikzpicture}[scale=0.95]
    \draw (0,0) node[op amp] (op) {};
    \draw (op.-) -- ++(-0.5,0) coordinate (vn);
    \draw (vn) to[C, l_={$C$}, i<={$i_C(t)$}] ++(-2.3,0) to[short, -o] ++(-0.2,0) node[left] {$v_i(t)$};
    \draw (vn) -- ++(0,1.35) to[R, l={$R$}, i={$i_R(t)$}] ++(3.0,0) -| (op.out);
    \draw (op.+) -- ++(-0.5,0) node[ground] {};
    \draw (op.out) to[short, -o] ++(0.9,0) node[right] {$v_o(t)$};
\end{tikzpicture}
\end{center}

Con $V^- = V^+ = 0\ \si{\volt}$ e $i_C(t) = i_R(t)$:
\begin{equation}
    i_C(t) = C \frac{dv_i(t)}{dt} = \frac{0 - v_o(t)}{R} \implies \boxed{v_o(t) = -R C \frac{dv_i(t)}{dt}}
\end{equation}

\subsection{7. Configuraciones No Lineales: Comparador y Disparador de Schmitt}

\subsubsection{a) Comparador Básico en Lazo Abierto}
Sin realimentación, el operacional compara instantáneamente las tensiones en sus entradas $V^+$ y $V^-$:

\begin{center}
\begin{tikzpicture}[scale=0.95]
    \draw (0,0) node[op amp] (op) {};
    \draw (op.-) -- ++(-0.8,0) to[short, -o] ++(-0.2,0) node[left] {$V_{ref}$};
    \draw (op.+) -- ++(-0.8,0) to[short, -o] ++(-0.2,0) node[left] {$V_{in}$};
    \draw (op.out) to[short, -o] ++(0.9,0) node[right] {$V_o \in \{+V_{pos}, -V_{neg}\}$};
\end{tikzpicture}
\end{center}

\begin{equation}
    V_o =
    \begin{cases}
        +V_{pos}\ (+V_{sat}) & \text{si } V_{in} > V_{ref} \iff (V_{in} - V_{ref}) > 0 \\
        -V_{neg}\ (-V_{sat}) & \text{si } V_{in} < V_{ref} \iff (V_{in} - V_{ref}) < 0
    \end{cases}
\end{equation}

\subsubsection{b) Disparador de Schmitt (\textit{Schmitt Trigger}) con Realimentación Positiva}
Para evitar conmutaciones erráticas producidas por el ruido cuando $V_i$ cruza el umbral de referencia, se introduce \textbf{realimentación positiva} hacia el terminal $V^+$ mediante las resistencias $R_1$ y $R_2$ (pudiendo acotar la salida con dos diodos Zener $D_1, D_2$ en antiserie):

\begin{center}
\begin{tikzpicture}[scale=0.95]
    \draw (0,0) node[op amp] (op) {};
    % Entrada inversora Vi
    \draw (op.-) -- ++(-0.8,0) to[short, -o] ++(-0.2,0) node[left] {$V_i$};
    % Entrada no inversora con realimentación positiva
    \draw (op.+) -- ++(-0.5,0) coordinate (vp);
    \draw (vp) -- ++(0,-0.9) coordinate (vdiv);
    \draw (vdiv) to[R, l={$R_2$}] ++(-2.1,0) to[short, -o] ++(-0.2,0) node[left] {$V_{ref}$};
    \draw (op.out) to[R, l={$R$}] ++(2.0,0) coordinate (vout);
    \draw (vdiv) to[R, l_={$R_1$}] ++(2.6,0) -| (op.out);
    % Rama Zener de recorte en la salida
    \draw (vout) to[short, -o] ++(1.0,0) node[right] {$V_o$};
    \draw (vout) to[zDo, l={$D_1$}] ++(0,-1.1) to[zDo, invert, l={$D_2$}] ++(0,-1.1) node[ground] {};
\end{tikzpicture}
\end{center}

El circuito presenta \textbf{dos tensiones umbral distintas} ($V_{TH}$ para subida y $V_{TL}$ para bajada), creando un ciclo de \textbf{histéresis} que inmuniza al comparador frente al ruido de alta frecuencia.

% ==============================================================================
% 6.3. CONVERSIÓN ANALÓGICA-DIGITAL (A/D) Y DIGITAL-ANALÓGICA (D/A)
% ==============================================================================
\section{Conversión Analógica-Digital (A/D) y Digital-Analógica (D/A)}

Los sistemas digitales y computadores procesan palabras binarias discretas en tiempo y amplitud, mientras que las magnitudes del mundo físico (temperatura, presión, sonido, tensión) son señales analógicas continuas. Los convertidores A/D (ADC) y D/A (DAC) constituyen la interfaz entre ambos dominios.

\subsection{Diagrama de Bloques y Etapas de la Digitalización (Conversión A/D)}

\begin{center}
\begin{tikzpicture}[
    scale=0.92,
    bloque/.style={
        rectangle, draw=fftPrimary, fill=fftGray, thick,
        rounded corners=3pt, minimum height=1.2cm, minimum width=2.4cm,
        align=center, font=\footnotesize\sffamily\bfseries
    }
]
    \node[align=center] (in) at (-1.8,0) {$v_a(t)$\\{\scriptsize (Analógica)}};
    \node[bloque] (sh) at (1.1,0) {1. Muestreo\\y Retención\\($f_s \ge 2f_{max}$)};
    \node[bloque] (quant) at (4.7,0) {2. Cuantificación\\($2^n$ niveles,\\$\Delta = V_{FS}/2^n$)};
    \node[bloque] (cod) at (8.3,0) {3. Codificación\\Binaria\\($n$ bits)};
    \node[align=center] (out) at (11.2,0) {$b_{n-1}\dots b_1 b_0$\\{\scriptsize (Palabra Digital)}};

    \draw[->, >=stealth, very thick, fftAccent] (in) -- (sh);
    \draw[->, >=stealth, very thick, fftAccent] (sh) -- (quant);
    \draw[->, >=stealth, very thick, fftAccent] (quant) -- (cod);
    \draw[->, >=stealth, very thick, fftAccent] (cod) -- (out);
\end{tikzpicture}
\end{center}

\begin{enumerate}[leftmargin=*]
    \item \textbf{Muestreo (\textit{Sampling}) y Teorema de Nyquist-Shannon:}
    Discretiza el eje temporal tomando muestras cada periodo $T_s = 1/f_s$: $x[k] = x(k T_s)$.
    \begin{ley}[Teorema del Muestreo de Nyquist-Shannon]
    Para poder reconstruir sin pérdida de información una señal analógica limitada en banda cuya componente de frecuencia máxima sea $f_{max}$, la frecuencia de muestreo $f_s$ debe cumplir:
    \begin{equation}
        \boxed{f_s \ge 2 f_{max}}
    \end{equation}
    A la frecuencia crítica $f_N = 2 f_{max}$ se la denomina \textbf{tasa de Nyquist}. Si $f_s < 2 f_{max}$, los espectros periódicos se solapan provocando distorsión irreversible por \textbf{aliasing} (solapamiento espectral).
    \end{ley}
    \item \textbf{Cuantificación (\textit{Quantization}):}
    Discretiza el eje de amplitudes asignando cada muestra analógica a uno de los $N_{niv} = 2^n$ escalones disponibles en un convertidor de $n$ bits con rango de fondo de escala $V_{FS}$ (\textit{Full Scale}):
    \begin{equation}
        \boxed{\Delta = \text{LSB} = \frac{V_{FS}}{2^n}} \qquad \left(\text{ o bien } \Delta = \frac{V_{FS}}{2^n - 1}\right)
    \end{equation}
    La diferencia entre el valor continuo real y el nivel discreto asignado introduce un \textbf{error de cuantificación} acotado por $|e_q| \le \frac{\Delta}{2}$.
    \item \textbf{Codificación Binaria:}
    Asigna a cada uno de los $2^n$ niveles cuantificados su representación única en binario natural o complemento a 2 de $n$ bits ($b_{n-1} b_{n-2} \dots b_1 b_0$).
\end{enumerate}

\subsection{Convertidores Digital-Analógicos (DAC) con Amplificador Operacional}

\subsubsection{1. DAC de Resistencias Ponderadas}
Consiste en un sumador inversor cuyas resistencias de entrada siguen una progresión geométrica de potencias de dos ($R, 2R, 4R, 8R, \dots, 2^{n-1}R$), gobernadas por conmutadores electrónicos activados por cada bit $b_k \in \{0, 1\}$ (donde $b_{n-1}$ es el MSB y $b_0$ es el LSB):

\begin{center}
\begin{tikzpicture}[scale=0.92]
    \draw (0,0) node[op amp] (op) {};
    \draw (op.-) -- ++(-0.6,0) coordinate (vn);
    % 4 ramas ponderadas R, 2R, 4R, 8R
    \draw (vn) -- ++(-0.4, 1.05) to[R, l_={$R$}] ++(-2.4,0) to[short, -o] ++(-0.2,0) node[left] {$b_3 V_{ref}$ (MSB)};
    \draw (vn) -- ++(-0.4, 0.35) to[R, l_={$2R$}] ++(-2.4,0) to[short, -o] ++(-0.2,0) node[left] {$b_2 V_{ref}$};
    \draw (vn) -- ++(-0.4,-0.35) to[R, l_={$4R$}] ++(-2.4,0) to[short, -o] ++(-0.2,0) node[left] {$b_1 V_{ref}$};
    \draw (vn) -- ++(-0.4,-1.05) to[R, l_={$8R$}] ++(-2.4,0) to[short, -o] ++(-0.2,0) node[left] {$b_0 V_{ref}$ (LSB)};
    % Realimentación Rf
    \draw (vn) -- ++(0,1.75) to[R, l={$R_f = R/2$}] ++(3.1,0) -| (op.out);
    % Entrada + a masa y salida
    \draw (op.+) -- ++(-0.5,0) node[ground] {};
    \draw (op.out) to[short, -o] ++(0.9,0) node[right] {$V_o$};
\end{tikzpicture}
\end{center}

\begin{equation}
    V_o = -R_f V_{ref} \left(\frac{b_3}{R} + \frac{b_2}{2R} + \frac{b_1}{4R} + \frac{b_0}{8R}\right) = -\frac{R_f}{R} V_{ref} \sum_{k=0}^{3} \frac{b_k}{2^{3-k}}
\end{equation}
\textbf{Limitación práctica:} Para resoluciones altas (p.\,ej. $n = 12$ o $16$ bits), la relación entre la resistencia del LSB ($2^{n-1}R$) y la del MSB ($R$) es enorme, haciendo inviable su fabricación precisa en circuitos integrados.

\subsubsection{2. DAC de Red en Escalera $R\text{--}2R$}
Resuelve el problema de tolerancias empleando \textbf{únicamente dos valores de resistencia} ($R$ y $2R$). En cada nudo de la escalera la resistencia equivalente vista hacia la izquierda o hacia el nudo inversor es siempre $2R$, dividiendo exactamente por $2$ la contribución de cada bit sucesivo:

\begin{center}
\begin{tikzpicture}[scale=0.90]
    \draw (0,0) node[op amp] (op) {};
    \draw (op.-) -- ++(-0.6,0) coordinate (n3);
    \draw (n3) to[R, l_={$R$}] ++(-2.1,0) coordinate (n2)
               to[R, l_={$R$}] ++(-2.1,0) coordinate (n1)
               to[R, l_={$R$}] ++(-2.1,0) coordinate (n0);
    % Resistencia de terminación 2R a masa
    \draw (n0) to[R, l_={$2R$}] ++(-1.9,0) node[ground] {};
    % Ramas verticales 2R de los bits
    \draw (n3) to[R, l={$2R$}] ++(0,-2.0) to[short, -o] ++(0,-0.2) node[below] {$b_3$ (MSB)};
    \draw (n2) to[R, l={$2R$}] ++(0,-2.0) to[short, -o] ++(0,-0.2) node[below] {$b_2$};
    \draw (n1) to[R, l={$2R$}] ++(0,-2.0) to[short, -o] ++(0,-0.2) node[below] {$b_1$};
    \draw (n0) to[R, l={$2R$}] ++(0,-2.0) to[short, -o] ++(0,-0.2) node[below] {$b_0$ (LSB)};
    % Realimentación
    \draw (n3) -- ++(0,1.4) to[R, l={$R_f = 2R$}] ++(3.1,0) -| (op.out);
    \draw (op.+) -- ++(-0.4,0) node[ground] {};
    \draw (op.out) to[short, -o] ++(0.9,0) node[right] {$V_o$};
\end{tikzpicture}
\end{center}

\begin{equation}
    \boxed{V_o = -V_{ref} \left(\frac{b_3}{2} + \frac{b_2}{4} + \frac{b_1}{8} + \frac{b_0}{16}\right) = -V_{ref} \sum_{k=0}^{n-1} b_k 2^{k - n}}
\end{equation}

% ==============================================================================
% 6.4. COLECCIÓN COMPLETA DE EJERCICIOS RESUELTOS DEL TEMA 6
% ==============================================================================
\section{Ejercicios Resueltos Paso a Paso de Amplificadores Operacionales}

\begin{ejercicio}[Ejercicio 1: Amplificador con Fuente Flotante]
Una fuente de tensión $V_f$ sin conexión directa a tierra se denomina \textbf{fuente flotante}. Este tipo de señal puede amplificarse a través del circuito de la figura, donde las dos resistencias conectadas en serie con la fuente valen $R_1$, la resistencia entre el nodo $A$ ($V^+$) y tierra vale $R_2$, y la resistencia de realimentación entre el nodo $B$ ($V^-$) y la salida $V_o$ también vale $R_2$. Determinar paso a paso la tensión de salida $V_o$ en función de $V_f$ y la ganancia del circuito.
\tcbline
\textbf{Esquema Circuital y Malla de Entrada:}\\[4pt]
\begin{center}
\begin{tikzpicture}[scale=0.95]
    \draw (0,0) node[op amp] (op) {};
    % Nodos B (V-) y A (V+)
    \draw (op.-) -- ++(-0.6,0) coordinate (vb);
    \draw (op.+) -- ++(-0.6,0) coordinate (va);
    \node[above, font=\footnotesize\bfseries, color=fftPrimary] at (vb) {$B\ (V_B)$};
    \node[below right, font=\footnotesize\bfseries, color=fftPrimary] at (va) {$A\ (V_A)$};
    % Resistencias R1 hacia la fuente flotante Vf
    \draw (vb) to[R, l_={$R_1$}, i_={$I$}] ++(-2.2,0) coordinate (topf);
    \draw (va) to[R, l={$R_1$}, i<={$I$}] ++(-2.2,0) coordinate (botf);
    \draw (botf) to[V, v={$V_f$}] (topf);
    % Resistencia R2 de A a tierra
    \draw (va) to[R, l={$R_2$}, i<={$I$}] ++(0,-1.8) node[ground] {};
    % Realimentación R2 de B a Vo
    \draw (vb) -- ++(0,1.35) to[R, l={$R_2$}, i<={$I$}] ++(3.1,0) -| (op.out);
    % Salida Vo
    \draw (op.out) to[short, -o] ++(0.9,0) node[right] {$V_o$};
\end{tikzpicture}
\end{center}

\textbf{Solución Paso a Paso:}\\[4pt]
1) \textbf{Cortocircuito virtual en la malla de la fuente flotante:}
Como el amplificador operacional tiene realimentación negativa ideal, $I^+ = I^- = 0$ y los potenciales de los nodos $A$ y $B$ son idénticos ($V_A = V_B$).
Por tanto, entre los nodos $A$ y $B$ la caída de tensión es nula ($V_B - V_A = 0$), cerrando una malla formada por la fuente $V_f$ y las dos resistencias $R_1$ en serie. La corriente $I$ que impulsa la fuente $V_f$ es:
\begin{equation}
    V_f = I \cdot R_1 + I \cdot R_1 = 2 R_1 I \implies I = \frac{V_f}{2 R_1}
\end{equation}
2) \textbf{Cálculo de la tensión $V_A = V_B$ aplicando la Ley de Ohm en la rama inferior:}
Como al terminal no inversor no entra corriente ($I^+ = 0$), la corriente $I$ que entra en la fuente flotante por su terminal negativo procede íntegramente desde tierra ($0\ \si{\volt}$) a través de la resistencia $R_2$ inferior:
\begin{equation}
    0 - V_A = I \cdot R_2 \implies V_A = -I \cdot R_2 = -\frac{R_2}{2 R_1} V_f = V_B
\end{equation}
3) \textbf{Cálculo de la tensión de salida $V_o$ en la rama superior:}
Como tampoco entra corriente al terminal inversor ($I^- = 0$), la misma corriente $I$ que sale de la fuente por la resistencia $R_1$ superior hacia el nodo $B$ continúa por la resistencia de realimentación $R_2$ hacia la salida $V_o$:
\begin{equation}
    V_B - V_o = I \cdot R_2 \implies V_o = V_B - I \cdot R_2 = -\frac{R_2}{2 R_1} V_f - \frac{R_2}{2 R_1} V_f = -\frac{R_2}{R_1} V_f
\end{equation}
\begin{equation}
    \boxed{V_o = -\frac{R_2}{R_1} V_f \qquad \implies \qquad A_v = \frac{V_o}{V_f} = -\frac{R_2}{R_1}}
\end{equation}
\end{ejercicio}

\begin{ejercicio}[Ejercicio 2: Dos Amplificadores Operacionales en Cascada (Inversor + Sumador Inversor)]
En el circuito de dos etapas en cascada de la figura, el primer amplificador operacional recibe una tensión de entrada de $-0{,}6\ \si{\volt}$ a través de una resistencia de $1\ \si{\kilo\ohm}$ y tiene una resistencia de realimentación de $3\ \si{\kilo\ohm}$. Su salida $V_1$ ataca a través de una resistencia de $2\ \si{\kilo\ohm}$ al nodo inversor de un segundo operacional, al cual también se conecta una tensión de $+0{,}5\ \si{\volt}$ mediante una resistencia de $1\ \si{\kilo\ohm}$. La resistencia de realimentación del segundo operacional es de $2\ \si{\kilo\ohm}$. Calcular paso a paso las tensiones $V_1$ y $V_2$.
\tcbline
\textbf{Esquema Circuital de las Dos Etapas en Cascada:}\\[4pt]
\begin{center}
\begin{tikzpicture}[scale=0.90]
    % Primer Op-Amp (Inversor)
    \draw (0,0) node[op amp] (op1) {};
    \draw (op1.-) -- ++(-0.5,0) coordinate (vn1);
    \draw (vn1) to[R, l_={$1\,\si{\kilo\ohm}$}] ++(-2.1,0) to[short, -o] ++(-0.2,0) node[left] {$-0{,}6\,\si{\volt}$};
    \draw (vn1) -- ++(0,1.35) to[R, l={$3\,\si{\kilo\ohm}$}] ++(2.9,0) -| (op1.out);
    \draw (op1.+) -- ++(-0.4,0) node[ground] {};
    \draw (op1.out) -- ++(0.5,0) coordinate (v1node);
    \node[below, font=\footnotesize\bfseries, color=fftPrimary] at (v1node) {$V_1$};

    % Segundo Op-Amp (Sumador Inversor)
    \draw (5.4,0) node[op amp] (op2) {};
    \draw (op2.-) -- ++(-0.5,0) coordinate (vn2);
    % Rama desde V1
    \draw (v1node) -- ++(0,0.49) to[R, l_={$2\,\si{\kilo\ohm}$}] (vn2);
    % Rama desde +0.5 V
    \draw (vn2) -- ++(-0.4, 0.75) to[R, l_={$1\,\si{\kilo\ohm}$}] ++(-1.9,0) to[short, -o] ++(-0.2,0) node[left] {$+0{,}5\,\si{\volt}$};
    % Realimentación 2 kOhm
    \draw (vn2) -- ++(0,1.5) to[R, l={$2\,\si{\kilo\ohm}$}] ++(2.9,0) -| (op2.out);
    \draw (op2.+) -- ++(-0.4,0) node[ground] {};
    \draw (op2.out) to[short, -o] ++(0.8,0) node[right] {$V_2$};
\end{tikzpicture}
\end{center}

\textbf{Solución Paso a Paso:}\\[4pt]
1) \textbf{Etapa 1 (Amplificador en configuración inversora):}
Aplicando la relación del amplificador inversor con $R_{in,1} = 1\ \si{\kilo\ohm}$ y $R_{f,1} = 3\ \si{\kilo\ohm}$:
\begin{equation}
    V_1 = -\frac{R_{f,1}}{R_{in,1}} V_i = -\left(\frac{3\ \si{\kilo\ohm}}{1\ \si{\kilo\ohm}}\right) \cdot (-0{,}6\ \si{\volt}) = +1{,}8\ \si{\volt}
\end{equation}
2) \textbf{Etapa 2 (Sumador inversor):}
El segundo operacional suma de forma inversora la tensión $+0{,}5\ \si{\volt}$ (por $1\ \si{\kilo\ohm}$) y la salida de la primera etapa $V_1 = +1{,}8\ \si{\volt}$ (por $2\ \si{\kilo\ohm}$), con resistencia de realimentación $R_F = 2\ \si{\kilo\ohm}$. Aplicando LCK en el nodo inversor del segundo operacional ($V^- = 0\ \si{\volt}$):
\begin{equation}
    \frac{0{,}5\ \si{\volt}}{1\ \si{\kilo\ohm}} + \frac{1{,}8\ \si{\volt}}{2\ \si{\kilo\ohm}} = -\frac{V_2}{2\ \si{\kilo\ohm}} \implies \frac{1{,}0 + 1{,}8}{2} = -\frac{V_2}{2}
\end{equation}
\begin{equation}
    \boxed{V_2 = -2\ \si{\kilo\ohm} \left(\frac{0{,}5}{1\ \si{\kilo\ohm}} + \frac{1{,}8}{2\ \si{\kilo\ohm}}\right) = -2 \cdot (0{,}5 + 0{,}9) = -2{,}8\ \si{\volt}}
\end{equation}
\end{ejercicio}

\begin{ejercicio}[Ejercicio 3: Dos Operacionales Inversores en Cascada con Realimentación Global $R_f$]
En el circuito de la figura, una fuente de tensión $v_s$ con resistencia serie $R_s = 1\ \si{\kilo\ohm}$ alimenta el nodo $B$ (de tensión $v_1$). A partir de $v_1$ se conectan en cascada dos etapas inversoras: la primera tiene resistencia de entrada $5\ \si{\kilo\ohm}$ y resistencia de realimentación $9\ \si{\kilo\ohm}$ (salida $v_2$), y la segunda tiene resistencia de entrada $1{,}2\ \si{\kilo\ohm}$ y resistencia de realimentación $6\ \si{\kilo\ohm}$ (salida final $v_o$). Además, existe una resistencia de realimentación global $R_f$ entre la salida final $v_o$ y el nodo $v_1$.
Hallar las tensiones $v_1, v_2, v_o$ y las corrientes $i_s, i_1, i_f$ en función de $v_s$ para:
\begin{enumerate}[label=\alph*)]
    \item $R_f = \infty$ (lazo global abierto).
    \item $R_f = 40\ \si{\kilo\ohm}$.
\end{enumerate}
\tcbline
\textbf{Esquema Circuital Completo:}\\[4pt]
\begin{center}
\begin{tikzpicture}[scale=0.88]
    % Nodo v1 y resistencia Rs
    \draw (-5.5, 0.49) to[V, v<={$v_s$}] (-5.5, -1.2) node[ground] {};
    \draw (-5.5, 0.49) to[R, l={$R_s{=}1\si{\kilo\ohm}$}, i={$i_s$}] (-3.2, 0.49) coordinate (v1);
    \node[below, font=\footnotesize\bfseries, color=fftPrimary] at (v1) {$B\ (v_1)$};

    % Primer Op-Amp
    \draw (0,0) node[op amp] (op1) {};
    \draw (op1.-) -- ++(-0.5,0) coordinate (vn1);
    \draw (v1) to[R, l_={$5\,\si{\kilo\ohm}$}, i={$i_1$}] (vn1);
    \draw (vn1) -- ++(0,1.25) to[R, l={$9\,\si{\kilo\ohm}$}] ++(2.9,0) -| (op1.out);
    \draw (op1.+) -- ++(-0.4,0) node[ground] {};
    \draw (op1.out) -- ++(0.4,0) coordinate (v2);
    \node[below, font=\footnotesize\bfseries, color=fftPrimary] at (v2) {$v_2$};

    % Segundo Op-Amp
    \draw (5.0,-0.49) node[op amp] (op2) {};
    \draw (op2.-) -- ++(-0.5,0) coordinate (vn2);
    \draw (v2) to[R, l_={$1{,}2\,\si{\kilo\ohm}$}] (vn2);
    \draw (vn2) -- ++(0,1.25) to[R, l={$6\,\si{\kilo\ohm}$}] ++(2.9,0) -| (op2.out);
    \draw (op2.+) -- ++(-0.4,0) node[ground] {};
    \draw (op2.out) -- ++(0.6,0) coordinate (vo);
    \draw (vo) to[short, -o] ++(0.6,0) node[right] {$v_o$};

    % Lazo superior Rf desde vo hasta v1
    \draw (vo) -- ++(0, 2.8) to[R, l_={$R_f$}, i_={$i_f$}] (-3.2, 2.31) -- (v1);
\end{tikzpicture}
\end{center}

\textbf{Solución Paso a Paso:}\\[4pt]
En primer lugar, como ambos operacionales están en configuración inversora (sus entradas inversoras están a masa virtual $0\ \si{\volt}$), las relaciones entre $v_1$, $v_2$ y $v_o$ son siempre:
\begin{align}
    v_2 &= -\frac{9\ \si{\kilo\ohm}}{5\ \si{\kilo\ohm}} v_1 = -1{,}8\, v_1 \\
    v_o &= -\frac{6\ \si{\kilo\ohm}}{1{,}2\ \si{\kilo\ohm}} v_2 = -5\, v_2 = -5 \left(-\frac{9}{5} v_1\right) = +9\, v_1
\end{align}
\textbf{Apartado a) Para $R_f = \infty$:}
Al estar $R_f$ en circuito abierto, la corriente de realimentación global es nula ($i_f = 0$). Como la entrada inversora del primer operacional está a masa virtual ($0\ \si{\volt}$), las resistencias $R_s = 1\ \si{\kilo\ohm}$ y $5\ \si{\kilo\ohm}$ quedan en serie hacia masa virtual, formando un divisor de tensión:
\begin{equation}
    i_1 = i_s = \frac{v_s}{1\ \si{\kilo\ohm} + 5\ \si{\kilo\ohm}} = \frac{v_s}{6\ \si{\kilo\ohm}}, \qquad v_1 = \frac{5\ \si{\kilo\ohm}}{1\ \si{\kilo\ohm} + 5\ \si{\kilo\ohm}} v_s = \frac{5}{6} v_s
\end{equation}
Sustituyendo $v_1$ en las ganancias de cada etapa:
\begin{equation}
    v_2 = -\frac{9}{5} \left(\frac{5}{6} v_s\right) = -\frac{3}{2} v_s = -1{,}5\, v_s, \qquad v_o = -\frac{6}{1{,}2} v_2 = -5 \left(-\frac{3}{2} v_s\right) = +\frac{15}{2} v_s = +7{,}5\, v_s
\end{equation}
\textit{(Nota: en el manuscrito original se arrastró el signo negativo al multiplicar $-\frac{6}{1{,}2} \cdot \left(-\frac{3}{2}\right)$ escribiendo $-7{,}5\,v_s$; algebraicamente el producto de dos etapas inversoras da ganancia positiva $+7{,}5\,v_s$).}\\[6pt]
\textbf{Apartado b) Para $R_f = 40\ \si{\kilo\ohm}$ ($i_f \ne 0$):}
Aplicamos LCK en el nodo $B$ ($v_1$), donde $i_s + i_f = i_1$:
\begin{equation}
    \frac{v_s - v_1}{1\ \si{\kilo\ohm}} + \frac{v_o - v_1}{40\ \si{\kilo\ohm}} = \frac{v_1 - 0}{5\ \si{\kilo\ohm}}
\end{equation}
Multiplicando toda la ecuación por $40\ \si{\kilo\ohm}$:
\begin{equation}
    40 (v_s - v_1) + (v_o - v_1) = 8 v_1 \implies 40 v_s - 40 v_1 + v_o - v_1 = 8 v_1
\end{equation}
Sustituyendo la relación de las dos etapas en cascada ($v_o = 9 v_1$):
\begin{equation}
    40 v_s - 40 v_1 + 9 v_1 - v_1 = 8 v_1 \implies 40 v_s = 40 v_1 \implies \boxed{v_1 = v_s}
\end{equation}
Por consiguiente, las demás tensiones e intensidades resultan ser:
\begin{equation}
    \boxed{v_2 = -\frac{9}{5} v_s, \qquad v_o = 9 v_s, \qquad i_s = \frac{v_s - v_1}{1\ \si{\kilo\ohm}} = 0, \qquad i_f = i_1 = \frac{v_1}{5\ \si{\kilo\ohm}} = \frac{v_s}{5\ \si{\kilo\ohm}}}
\end{equation}
\end{ejercicio}

\begin{ejercicio}[Ejercicio 4: Integrador Inversor con Entrada Sinusoidal]
En un circuito integrador inversor con Op-Amp ideal se tiene $R = 1\ \si{\kilo\ohm}$, $C = 1\ \si{\micro\farad}$ y una señal de entrada $v_1(t) = \sin(2000t)\ \si{\volt}$. Asumiendo que el condensador parte descargado en el instante inicial, $v_2(0) = 0\ \si{\volt}$, encontrar la expresión temporal de la tensión de salida $v_2(t)$ para $t > 0$.
\tcbline
\textbf{Solución Paso a Paso:}\\[4pt]
1) Por cortocircuito virtual en el integrador inversor ($V^+ = V^- = 0\ \si{\volt}$) e $I^+ = I^- = 0$:
\begin{equation}
    i_R(t) = i_C(t) \implies \frac{v_1(t)}{R} = C \frac{d(0 - v_2(t))}{dt} = -C \frac{dv_2(t)}{dt}
\end{equation}
2) Integrando entre $0$ y $t$ con $R C = (1 \times 10^3\ \si{\ohm}) \cdot (1 \times 10^{-6}\ \si{\farad}) = 10^{-3}\ \si{\second}$ y $\omega = 2000\ \si{\radian\per\second}$:
\begin{align}
    v_2(t) - v_2(0) &= -\frac{1}{R C} \int_0^t \sin(2000\tau)\, d\tau = -\frac{1}{10^{-3}} \left[-\frac{\cos(2000\tau)}{2000}\right]_0^t \\
    v_2(t) &= \frac{1000}{2000} \left[\cos(2000t) - \cos(0)\right] = \boxed{0{,}5 \left[\cos(2000t) - 1\right]\ \si{\volt}}
\end{align}
\end{ejercicio}

\begin{ejercicio}[Ejercicio 5: Comparador con Realimentación Positiva (Ciclo de Histéresis)]
Analizar el circuito comparador con realimentación positiva de la figura y dibujar su función de transferencia con histéresis. El operacional satura en $+10\ \si{\volt}$ y $-10\ \si{\volt}$, tiene su entrada inversora conectada a masa ($V^- = 0\ \si{\volt}$), la señal $V_i$ se aplica a $V^+$ a través de $R_1 = 24\ \si{\kilo\ohm}$, y la resistencia de realimentación entre $V^+$ y $V_o$ es $R_2 = 48\ \si{\kilo\ohm}$.
\tcbline
\textbf{Esquema Circuital y Lazo de Histéresis:}\\[4pt]
\begin{center}
\begin{tikzpicture}[scale=0.92]
    % --- Circuito a la izquierda ---
    \draw (-4.2,0) node[op amp, yscale=-1] (op) {};
    \draw (op.+) -- ++(-0.5,0) coordinate (vp);
    \draw (vp) to[R, l_={$R_1{=}24\si{\kilo\ohm}$}] ++(-2.2,0) to[short, -o] ++(-0.2,0) node[left] {$V_i$};
    \draw (vp) -- ++(0,1.35) to[R, l={$R_2{=}48\si{\kilo\ohm}$}] ++(3.0,0) -| (op.out);
    \draw (op.-) -- ++(-0.4,0) node[ground] {};
    \draw (op.out) to[short, -o] ++(0.8,0) node[right] {$V_o$};
    \node[above, font=\scriptsize] at (-4.2, 0.7) {$+10\si{\volt}$};
    \node[below, font=\scriptsize] at (-4.2,-0.7) {$-10\si{\volt}$};

    % --- Gráfica de Histéresis a la derecha ---
    \begin{scope}[xshift=3.4cm]
        \draw[->, >=stealth, gray] (-2.6,0) -- (2.6,0) node[right, text=black] {$V_i$};
        \draw[->, >=stealth, gray] (0,-1.8) -- (0,1.8) node[above, text=black] {$V_o$};
        % Lazo de histéresis
        \draw[very thick, fftAccent] (-2.3,-1.2) -- (1.0,-1.2) -- (1.0,1.2) -- (2.3,1.2);
        \draw[very thick, fftPrimary] (2.3,1.2) -- (-1.0,1.2) -- (-1.0,-1.2) -- (-2.3,-1.2);
        % Flechas de sentido del ciclo
        \draw[->, >=stealth, thick, fftAccent] (1.0,-0.4) -- (1.0,0.3);
        \draw[->, >=stealth, thick, fftPrimary] (-1.0,0.4) -- (-1.0,-0.3);
        % Etiquetas de umbrales y saturación
        \node[below right, font=\footnotesize] at (1.0,0) {$+5\si{\volt}$};
        \node[above left, font=\footnotesize] at (-1.0,0) {$-5\si{\volt}$};
        \node[above left, font=\footnotesize] at (0,1.2) {$+10\si{\volt}$};
        \node[below right, font=\footnotesize] at (0,-1.2) {$-10\si{\volt}$};
    \end{scope}
\end{tikzpicture}
\end{center}

\textbf{Solución Paso a Paso:}\\[4pt]
1) \textbf{Expresión del potencial $V^+$:} Como $I^+ = 0$, la corriente por $R_1$ es igual a la corriente por $R_2$:
\begin{equation}
    \frac{V_i - V^+}{24\ \si{\kilo\ohm}} = \frac{V^+ - V_o}{48\ \si{\kilo\ohm}} \implies 2(V_i - V^+) = V^+ - V_o \implies V^+ = \frac{2 V_i + V_o}{3}
\end{equation}
2) \textbf{Barrido ascendente (partiendo de $V_i \ll 0$, con $V_o = -10\ \si{\volt}$):}
El comparador conmutará a $+10\ \si{\volt}$ cuando $V^+$ alcance el potencial de $V^- = 0\ \si{\volt}$:
\begin{equation}
    V^+ = \frac{2 V_i - 10}{3} = 0 \implies V_{TH} = +5\ \si{\volt}
\end{equation}
Para $V_i > +5\ \si{\volt}$, la salida pasa a $V_o = +10\ \si{\volt}$.\\[4pt]
3) \textbf{Barrido descendente (regresando desde $V_i \gg 0$, con $V_o = +10\ \si{\volt}$):}
Ahora $V_o = +10\ \si{\volt}$, por lo que la conmutación de vuelta a $-10\ \si{\volt}$ ocurrirá cuando $V^+ = 0\ \si{\volt}$:
\begin{equation}
    V^+ = \frac{2 V_i + 10}{3} = 0 \implies V_{TL} = -5\ \si{\volt}
\end{equation}
\end{ejercicio}

\begin{ejercicio}[Ejercicio 6: Conversión de Onda Sinusoidal a Cuadrada y Triangular (Comparador + Integrador)]
En el circuito de la figura, una señal sinusoidal $v_i(t) = 1\sin(\omega t)\ \si{\volt}$ de frecuencia $f = 100\ \si{\kilo\hertz}$ ataca a un primer operacional conectado como comparador en lazo abierto (con límites de saturación $\pm 12\ \si{\volt}$ y $V^- = 0\ \si{\volt}$), cuya salida $V_A$ alimenta a un integrador inversor con resistencia $R$ y condensador $C = 10\ \si{\nano\farad}$. Sabiendo que en $t = 0$ es $v_o(0) = +1\ \si{\volt}$, calcular el valor de $R$ para que la señal triangular de salida $v_o(t)$ oscile simétricamente entre $+1\ \si{\volt}$ y $-1\ \si{\volt}$.
\tcbline
\textbf{Esquema Circuital (Comparador + Integrador):}\\[4pt]
\begin{center}
\begin{tikzpicture}[scale=0.92]
    % Etapa 1: Comparador
    \draw (0,0) node[op amp] (op1) {};
    \draw (op1.-) -- ++(-0.4,0) node[ground] {};
    \draw (op1.+) -- ++(-0.7,0) to[short, -o] ++(-0.2,0) node[left] {$v_i(t)$};
    \node[above, font=\scriptsize] at (op1.up) {$+12\si{\volt}$};
    \node[below, font=\scriptsize] at (op1.down) {$-12\si{\volt}$};
    \draw (op1.out) -- ++(0.5,0) coordinate (va);
    \node[below, font=\footnotesize\bfseries, color=fftPrimary] at (va) {$V_A$};

    % Etapa 2: Integrador
    \draw (5.0,-0.49) node[op amp] (op2) {};
    \draw (op2.-) -- ++(-0.5,0) coordinate (vn2);
    \draw (va) to[R, l={$R$}] (vn2);
    \draw (vn2) -- ++(0,1.35) to[C, l={$C{=}10\si{\nano\farad}$}] ++(2.9,0) -| (op2.out);
    \draw (op2.+) -- ++(-0.4,0) node[ground] {};
    \draw (op2.out) to[short, -o] ++(0.8,0) node[right] {$v_o(t)$};
\end{tikzpicture}
\end{center}

\textbf{Solución Paso a Paso:}\\[4pt]
1) \textbf{Periodo y salida del comparador ($V_A$):}
El periodo de la señal de entrada es $T = \frac{1}{f} = \frac{1}{100\ \si{\kilo\hertz}} = 10\ \si{\micro\second}$, luego cada semiperiodo dura $T/2 = 5\ \si{\micro\second}$.
Durante el primer semiperiodo $t \in (0, T/2)$, es $v_i(t) > 0\ \si{\volt}$, por lo que el comparador satura positivamente entregando una tensión constante $V_A = +12\ \si{\volt}$ (en el segundo semiperiodo entrega $V_A = -12\ \si{\volt}$, formando una onda cuadrada de $\pm 12\ \si{\volt}$).\\[4pt]
2) \textbf{Integración en el semiperiodo $[0, T/2]$:}
La ecuación del integrador inversor es:
\begin{equation}
    v_o(T/2) = -\frac{1}{R C} \int_0^{T/2} V_A\, dt + v_o(0)
\end{equation}
Para que $v_o(t)$ descienda linealmente desde $v_o(0) = +1\ \si{\volt}$ hasta $v_o(T/2) = -1\ \si{\volt}$ en el intervalo $T/2 = 5\ \si{\micro\second}$:
\begin{equation}
    -1\ \si{\volt} - 1\ \si{\volt} = -\frac{1}{R C} \cdot 12\ \si{\volt} \cdot \left(\frac{T}{2}\right) \implies -2 = -\frac{12 \cdot 5\ \si{\micro\second}}{R \cdot 10\ \si{\nano\farad}}
\end{equation}
\begin{equation}
    \boxed{R = \frac{12\ \si{\volt} \cdot 5 \times 10^{-6}\ \si{\second}}{2\ \si{\volt} \cdot 10 \times 10^{-9}\ \si{\farad}} = 3 \times 10^3\ \si{\ohm} = 3\ \si{\kilo\ohm}}
\end{equation}
\end{ejercicio}

\begin{ejercicio}[Ejercicio 7: Sumador-Inversor con Chequeo de Saturación --- Examen UGR]
Un amplificador operacional ideal está alimentado simétricamente con $\pm V_{CC} = \pm 15\ \si{\volt}$ ($V_{sat} = \pm 14\ \si{\volt}$). Dispone de una resistencia de realimentación $R_f = 20\ \si{\kilo\ohm}$ conectada a la entrada inversora, mientras la entrada no inversora está conectada a masa ($V^+ = 0\ \si{\volt}$). Al nudo inversor se aplican dos tensiones mediante sendas resistencias: $V_1 = +2\ \si{\volt}$ a través de $R_1 = 10\ \si{\kilo\ohm}$, y $V_2 = -3\ \si{\volt}$ a través de $R_2 = 5\ \si{\kilo\ohm}$.
\begin{enumerate}[label=\alph*)]
    \item Calcular la tensión de salida $V_o$.
    \item Si se modifica la señal 2 pasando a $V_2 = -6\ \si{\volt}$, determinar la nueva tensión de salida indicando si el operacional satura.
\end{enumerate}
\tcbline
\textbf{Esquema Circuital del Sumador Inversor:}\\[4pt]
\begin{center}
\begin{tikzpicture}[scale=0.95]
    \draw (0,0) node[op amp] (op) {};
    \draw (op.-) -- ++(-0.6,0) coordinate (vn);
    \draw (vn) -- ++(-0.4, 0.55) to[R, l_={$R_1{=}10\si{\kilo\ohm}$}] ++(-2.4,0) to[short, -o] ++(-0.2,0) node[left] {$V_1{=}+2\si{\volt}$};
    \draw (vn) -- ++(-0.4,-0.55) to[R, l={$R_2{=}5\si{\kilo\ohm}$}] ++(-2.4,0) to[short, -o] ++(-0.2,0) node[left] {$V_2$};
    \draw (vn) -- ++(0,1.45) to[R, l={$R_f{=}20\si{\kilo\ohm}$}] ++(3.1,0) -| (op.out);
    \draw (op.+) -- ++(-0.5,0) node[ground] {};
    \draw (op.out) to[short, -o] ++(0.9,0) node[right] {$V_o$};
    \node[above, font=\scriptsize] at (op.up) {$+15\si{\volt}$};
    \node[below, font=\scriptsize] at (op.down) {$-15\si{\volt}$};
\end{tikzpicture}
\end{center}

\textbf{Solución Paso a Paso:}\\[4pt]
a) Con $V_1 = +2\ \si{\volt}$ y $V_2 = -3\ \si{\volt}$:
\begin{equation}
    V_o = -R_f \left( \frac{V_1}{R_1} + \frac{V_2}{R_2} \right) = -20 \left( \frac{2}{10} + \frac{-3}{5} \right) = -20 (0{,}20 - 0{,}60) = -20 (-0{,}40) = +8{,}0\ \si{\volt}
\end{equation}
Como $|V_o| = 8{,}0\ \si{\volt} < V_{sat} = 14\ \si{\volt}$, trabaja en régimen lineal sin saturación.\\[4pt]
b) Con $V_2 = -6\ \si{\volt}$:
\begin{equation}
    V_{o,\text{calculado}} = -20 \left( \frac{2}{10} + \frac{-6}{5} \right) = -20 (0{,}20 - 1{,}20) = -20 (-1{,}00) = +20{,}0\ \si{\volt}
\end{equation}
Como $+20{,}0\ \si{\volt} > +V_{sat} = +14\ \si{\volt}$, el amplificador entra en \textbf{saturación positiva}:
\begin{equation}
    \boxed{V_o = +V_{sat} = +14{,}0\ \si{\volt}}
\end{equation}
\end{ejercicio}

\begin{ejercicio}[Ejercicio 8: Integrador Inversor ante Entrada Escalonada]
Un circuito integrador inversor con Op-Amp ideal dispone de $R = 50\ \si{\kilo\ohm}$ y $C = 2\ \si{\micro\farad}$, con tensiones de saturación $\pm 12\ \si{\volt}$. El condensador está inicialmente descargado ($V_o(0) = 0\ \si{\volt}$). En $t = 0$ se aplica una tensión escalón constante $V_i = +2\ \si{\volt}$.
\begin{enumerate}[label=\alph*)]
    \item Hallar la expresión temporal de $V_o(t)$ en régimen lineal y la pendiente de la rampa.
    \item Determinar el tiempo que tarda la salida en alcanzar la saturación negativa.
\end{enumerate}
\tcbline
\textbf{Solución Paso a Paso:}\\[4pt]
a) La constante de tiempo integradora es $R C = (50 \times 10^3\ \si{\ohm}) \times (2 \times 10^{-6}\ \si{\farad}) = 0{,}10\ \si{\second} = 100\ \si{\milli\second}$.
Para una entrada constante $V_i = +2\ \si{\volt}$:
\begin{equation}
    V_o(t) = -\frac{1}{R C} \int_0^t V_i\, d\tau + V_o(0) = -\frac{V_i}{R C} t = -\frac{2}{0{,}10} t = -20 t\ \si{\volt} \qquad \left(\frac{dV_o}{dt} = -20\ \si{\volt\per\second}\right)
\end{equation}
b) Alcanza la saturación negativa cuando $V_o(t_{sat}) = -V_{sat} = -12\ \si{\volt}$:
\begin{equation}
    -20 t_{sat} = -12 \implies \boxed{t_{sat} = \frac{12}{20} = 0{,}60\ \si{\second} = 600\ \si{\milli\second}}
\end{equation}
\end{ejercicio}

\chapter{Colección de Exámenes Oficiales Resueltos Paso a Paso}
\label{cap:tema7}

\cuadrofuentes{
    \textbf{Asignatura:} Fundamentos Físicos y Tecnológicos (1.º Grado en Ingeniería Informática, UGR).\\[2pt]
    \textbf{Profesor:} Pedro Cartujo (Departamento de Electrónica y Tecnología de Computadores).\\[2pt]
    \textbf{Contenido:} Problemas originales de los exámenes parciales y convocatorias oficiales de la UGR, resueltos de forma íntegra, sistemática y con todas las justificaciones algebraicas y numéricas.
}

% ==============================================================================
% 7.1. EXAMEN PARCIAL 1
% ==============================================================================
\section{Examen Parcial 1: Electromagnetismo y Circuitos DC/AC}

\begin{ejercicio}[Parcial 1 --- Problema 1: Campo Magnético y Fuerza entre Tres Hilos Paralelos]
Se tienen tres cables conductores rectilíneos, paralelos e indefinidos de longitud $L = 50\ \si{\centi\meter} = 0{,}5\ \si{\meter}$, situados coplanariamente en el vacío ($\mu_0 = 4\pi \times 10^{-7}\ \si{\tesla\meter\per\ampere}$).
Por los tres hilos circulan corrientes continuas en el mismo sentido ($+\hat{k}$ saliendo del plano):
\begin{equation}
    I_1 = 2\ \si{\milli\ampere}, \qquad I_2 = 4\ \si{\milli\ampere}, \qquad I_3 = 6\ \si{\milli\ampere}
\end{equation}
La distancia entre los cables 1 y 2 es $d_{12} = 1\ \si{\centi\meter} = 1 \times 10^{-2}\ \si{\meter}$, y entre los cables 2 y 3 es $d_{23} = 2\ \si{\centi\meter} = 2 \times 10^{-2}\ \si{\meter}$. Por consiguiente, la distancia entre el cable 1 y el cable 3 es $d_{13} = 3\ \si{\centi\meter} = 3 \times 10^{-2}\ \si{\meter}$.
\begin{enumerate}[label=\alph*)]
    \item Enunciar la Ley de Ampère y deducir la expresión del campo magnético creado por un conductor rectilíneo indefinido a una distancia $r$.
    \item Calcular los campos magnéticos producidos sobre cada uno de los cables por los otros dos, indicando magnitud, dirección y sentido.
    \item Calcular la fuerza magnética neta que experimenta cada cable indicando su dirección y sentido.
\end{enumerate}
\tcbline
\textbf{Solución:}\\[4pt]
a) \textbf{Deducción mediante la Ley de Ampère:}
Consideramos una curva circular amperiana de radio $r$ centrada coaxialmente en el hilo. Por simetría axial, el vector $\vec{B}$ es tangente en todos los puntos a la circunferencia y de módulo constante:
\begin{equation}
    \oint_C \vec{B} \cdot d\vec{\ell} = B \oint_C d\ell = B (2\pi r) = \mu_0 I_{enc} \implies B(r) = \frac{\mu_0 I}{2\pi r}
\end{equation}
b) \textbf{Campos Magnéticos en Cada Hilo:}\\[4pt]
1) \textbf{Sobre el Hilo 1} (en $x = 0$):
Los hilos 2 y 3 se encuentran a su derecha. Por la regla de la mano derecha, las corrientes salientes de 2 y 3 generan en la posición del hilo 1 un campo dirigido hacia arriba ($+\hat{j}$):
\begin{align}
    B_{2 \to 1} &= \frac{\mu_0 I_2}{2\pi d_{12}} = \frac{(4\pi \times 10^{-7}) \times (4 \times 10^{-3})}{2\pi \times (1 \times 10^{-2})} = 8 \times 10^{-8}\ \si{\tesla} \\
    B_{3 \to 1} &= \frac{\mu_0 I_3}{2\pi d_{13}} = \frac{(4\pi \times 10^{-7}) \times (6 \times 10^{-3})}{2\pi \times (3 \times 10^{-2})} = 4 \times 10^{-8}\ \si{\tesla}
\end{align}
\begin{equation}
    \vec{B}_1 = (B_{2 \to 1} + B_{3 \to 1})\hat{j} = (8 + 4) \times 10^{-8}\hat{j} = 12 \times 10^{-8}\ \si{\tesla}\ \hat{j} = 0{,}12\ \si{\micro\tesla}\ \hat{j}
\end{equation}
2) \textbf{Sobre el Hilo 2} (en $x = 1\ \si{\centi\meter}$):
El hilo 1 está a su izquierda $\implies$ genera campo hacia abajo ($-\hat{j}$). El hilo 3 está a su derecha $\implies$ genera campo hacia arriba ($+\hat{j}$):
\begin{align}
    B_{1 \to 2} &= \frac{\mu_0 I_1}{2\pi d_{12}} = \frac{(4\pi \times 10^{-7}) \times (2 \times 10^{-3})}{2\pi \times (1 \times 10^{-2})} = 4 \times 10^{-8}\ \si{\tesla} \quad (-\hat{j}) \\
    B_{3 \to 2} &= \frac{\mu_0 I_3}{2\pi d_{23}} = \frac{(4\pi \times 10^{-7}) \times (6 \times 10^{-3})}{2\pi \times (2 \times 10^{-2})} = 6 \times 10^{-8}\ \si{\tesla} \quad (+\hat{j})
\end{align}
\begin{equation}
    \vec{B}_2 = (-4 + 6) \times 10^{-8}\hat{j} = +2 \times 10^{-8}\ \si{\tesla}\ \hat{j} = 0{,}02\ \si{\micro\tesla}\ \hat{j}
\end{equation}
3) \textbf{Sobre el Hilo 3} (en $x = 3\ \si{\centi\meter}$):
Los hilos 1 y 2 están a su izquierda $\implies$ generan campo dirigido hacia abajo ($-\hat{j}$):
\begin{align}
    B_{1 \to 3} &= \frac{\mu_0 I_1}{2\pi d_{13}} = \frac{(4\pi \times 10^{-7}) \times (2 \times 10^{-3})}{2\pi \times (3 \times 10^{-2})} = 1{,}333 \times 10^{-8}\ \si{\tesla} \\
    B_{2 \to 3} &= \frac{\mu_0 I_2}{2\pi d_{23}} = \frac{(4\pi \times 10^{-7}) \times (4 \times 10^{-3})}{2\pi \times (2 \times 10^{-2})} = 4 \times 10^{-8}\ \si{\tesla}
\end{align}
\begin{equation}
    \vec{B}_3 = -(1{,}333 + 4) \times 10^{-8}\hat{j} = -5{,}333 \times 10^{-8}\ \si{\tesla}\ \hat{j}
\end{equation}
c) \textbf{Fuerzas Magnéticas Netas sobre Cada Cable ($\vec{F} = I \vec{L} \times \vec{B}$, con $\vec{L} = L \hat{k}$):}
Recordando que $\hat{k} \times \hat{j} = -\hat{i}$ y $\hat{k} \times (-\hat{j}) = +\hat{i}$:
\begin{align}
    \vec{F}_1 &= I_1 L (\hat{k} \times \vec{B}_1) = (2 \times 10^{-3}) \times 0{,}5 \times (12 \times 10^{-8}) (\hat{k} \times \hat{j}) \nonumber \\
              &= -1{,}2 \times 10^{-10}\ \si{\newton}\ \hat{i} \quad (\text{hacia la derecha, atracción}) \\
    \vec{F}_2 &= I_2 L (\hat{k} \times \vec{B}_2) = (4 \times 10^{-3}) \times 0{,}5 \times (2 \times 10^{-8}) (\hat{k} \times \hat{j}) \nonumber \\
              &= -0{,}4 \times 10^{-10}\ \si{\newton}\ \hat{i} \\
    \vec{F}_3 &= I_3 L (\hat{k} \times \vec{B}_3) = (6 \times 10^{-3}) \times 0{,}5 \times (-5{,}333 \times 10^{-8}) (\hat{k} \times \hat{j}) \nonumber \\
              &= +1{,}6 \times 10^{-10}\ \si{\newton}\ \hat{i}
\end{align}
Nótese que la suma vectorial de las fuerzas internas del sistema aislado es nula:
\begin{equation}
    \sum \vec{F} = (-1{,}2 - 0{,}4 + 1{,}6) \times 10^{-10}\ \hat{i} = \vec{0}
\end{equation}
(Verificando la Tercera Ley de Newton de acción y reacción).
\end{ejercicio}

\begin{ejercicio}[Parcial 1 --- Problema 2: Análisis por Nudos y Balance de Potencia en DC]
Dado un circuito de corriente continua compuesto por fuentes de tensión $V_1 = 12\ \si{\volt}$, $V_2 = 22\ \si{\volt}$ y resistencias $R_1 = 2\ \si{\ohm}$, $R_2 = 4\ \si{\ohm}$, $R_3 = 5\ \si{\ohm}$, $R_4 = 10\ \si{\ohm}$ y $R_5 = 8\ \si{\ohm}$, con una fuente de corriente $I_s = 1\ \si{\ampere}$.
Tras formular el sistema de ecuaciones por el método de los nudos, se obtienen las tensiones de nudo:
\begin{equation}
    V_b = 18\ \si{\volt}, \qquad V_c = 14\ \si{\volt}, \qquad V_d = -4\ \si{\volt}, \qquad V_e = 8\ \si{\volt}
\end{equation}
\begin{enumerate}[label=\alph*)]
    \item Calcular las intensidades que circulan por cada una de las ramas.
    \item Calcular la potencia disipada en cada una de las resistencias.
    \item Calcular las potencias absorbidas y entregadas por cada una de las fuentes.
    \item Comprobar que se satisface el balance de potencia global de Tellegen ($\sum P = 0$).
\end{enumerate}
\tcbline
\textbf{Solución:}\\[4pt]
a) \textbf{Corrientes de Rama:}\\[2pt]
- Rama 1 ($R_1 = 2\ \si{\ohm}$): $I_1 = \frac{V_b - V_c}{R_1} = \frac{18 - 14}{2} = \frac{4}{2} = 2\ \si{\ampere}$.\\[2pt]
- Rama 2 ($R_2 = 4\ \si{\ohm}$): $I_2 = \frac{V_c - V_e}{R_2} = \frac{14 - 8}{4} = 1{,}5\ \si{\ampere}$.\\[2pt]
- Rama 3 ($R_3 = 5\ \si{\ohm}$): $I_3 = \frac{V_b - V_e}{R_3} = \frac{18 - 8}{5} = \frac{10}{5} = 2\ \si{\ampere}$.\\[2pt]
- Rama 4 ($R_4 = 10\ \si{\ohm}$): $I_4 = \frac{V_d - 0}{R_4} = \frac{-4}{10} = -0{,}4\ \si{\ampere}$.\\[2pt]
- Rama 5 ($R_5 = 8\ \si{\ohm}$): $I_5 = \frac{V_e - 0}{R_5} = \frac{8}{8} = 1\ \si{\ampere}$.\\[6pt]
b) \textbf{Potencia Disipada en las Resistencias ($P_{R} = I^2 R$):}
\begin{align}
    P_{R1} &= 2^2 \times 2 = 8\ \si{\watt} \\
    P_{R2} &= 1{,}5^2 \times 4 = 9\ \si{\watt} \\
    P_{R3} &= 2^2 \times 5 = 20\ \si{\watt} \\
    P_{R4} &= (-0{,}4)^2 \times 10 = 1{,}6\ \si{\watt} \\
    P_{R5} &= 1^2 \times 8 = 8\ \si{\watt}
\end{align}
Potencia total disipada en elementos pasivos:
\begin{equation}
    \sum P_R = 8 + 9 + 20 + 1{,}6 + 8 = 46{,}6\ \si{\watt}
\end{equation}
c) \textbf{Potencias en las Fuentes (según convenio activo):}
\begin{itemize}[leftmargin=*]
    \item Fuente $V_2 = 22\ \si{\volt}$: La corriente que sale por su borne positivo es $I_{V2} = 2{,}5\ \si{\ampere}$.
    \begin{equation}
        P_{V2,entregada} = 22 \times 2{,}5 = 55\ \si{\watt}
    \end{equation}
    \item Fuente $V_1 = 12\ \si{\volt}$: La corriente entra por su borne positivo ($I_{V1} = -0{,}6\ \si{\ampere}$), por lo que absorbe potencia:
    \begin{equation}
        P_{V1,absorbida} = 12 \times 0{,}6 = 7{,}2\ \si{\watt}
    \end{equation}
    \item Fuente de corriente $I_s = 1\ \si{\ampere}$: Sometida a una diferencia de potencial que absorbe $1{,}2\ \si{\watt}$.
\end{itemize}
d) \textbf{Comprobación del Balance de Potencias:}
\begin{equation}
    P_{total,suministrada} = 55\ \si{\watt}
\end{equation}
\begin{equation}
    P_{total,absorbida} = \sum P_R + P_{fuentes,absorbidas} = 46{,}6 + 7{,}2 + 1{,}2 = 55{,}0\ \si{\watt}
\end{equation}
\begin{equation}
    \sum P = P_{suministrada} - P_{absorbida} = 55{,}0 - 55{,}0 = 0\ \si{\watt} \quad \text{[CORRECTO]}
\end{equation}
\end{ejercicio}

\begin{ejercicio}[Parcial 1 --- Problema 3: Circuito AC por Superposición y Transformación de Fuentes]
Determinar la tensión fasorial $\mathbf{V}_1$ y $\mathbf{V}_2$ en el circuito en régimen permanente sinusoidal compuesto por:
\begin{itemize}[leftmargin=*]
    \item Fuente de corriente 1: $\mathbf{I}_1 = 20\angle 0^\circ\ \si{\milli\ampere} = 0{,}02\ \si{\ampere}$.
    \item Fuente de corriente 2: $\mathbf{I}_2 = 50\angle -90^\circ\ \si{\milli\ampere} = -j 0{,}05\ \si{\ampere}$.
    \item Impedancias: $\mathbf{Z}_C = -j 20\ \si{\ohm}$, $\mathbf{Z}_L = j 40\ \si{\ohm}$ y resistencia $\mathbf{Z}_R = 25\ \si{\ohm}$.
\end{itemize}
Resolver el circuito aplicando el método de transformación de fuentes.
\tcbline
\textbf{Solución:}\\[4pt]
1) \textbf{Transformación de la Fuente $\mathbf{I}_1$ en Paralelo con $\mathbf{Z}_C$:}
La rama formada por la fuente de corriente $\mathbf{I}_1 = 0{,}02\ \si{\ampere}$ en paralelo con el condensador $\mathbf{Z}_C = -j 20\ \si{\ohm}$ se transforma en una fuente de tensión equivalente de Thévenin:
\begin{equation}
    \mathbf{V}_{s1} = \mathbf{I}_1 \cdot \mathbf{Z}_C = (0{,}02) \times (-j 20) = -j 0{,}4\ \si{\volt}
\end{equation}
conectada en serie con la impedancia $\mathbf{Z}_C = -j 20\ \si{\ohm}$.\\[4pt]
2) \textbf{Transformación de la Fuente $\mathbf{I}_2$ en Paralelo con $\mathbf{Z}_L$:}
La fuente $\mathbf{I}_2 = -j 0{,}05\ \si{\ampere}$ en paralelo con la bobina $\mathbf{Z}_L = j 40\ \si{\ohm}$ se transforma en:
\begin{equation}
    \mathbf{V}_{s2} = \mathbf{I}_2 \cdot \mathbf{Z}_L = (-j 0{,}05) \times (j 40) = -j^2 \times 2 = +2\ \si{\volt}
\end{equation}
conectada en serie con la impedancia $\mathbf{Z}_L = j 40\ \si{\ohm}$.\\[4pt]
3) \textbf{Circuito Serie Resultante:}
El circuito transformado consiste en un único lazo cerrado en serie formado por:
$\mathbf{V}_{s1}$, $\mathbf{Z}_C$, la resistencia $\mathbf{Z}_R = 25\ \si{\ohm}$, la bobina $\mathbf{Z}_L = j 40\ \si{\ohm}$ y la fuente $\mathbf{V}_{s2}$.
La impedancia total en serie es:
\begin{equation}
    \mathbf{Z}_{total} = \mathbf{Z}_C + \mathbf{Z}_R + \mathbf{Z}_L = -j 20 + 25 + j 40 = 25 + j 20\ \si{\ohm}
\end{equation}
La tensión total impulsora en el lazo es:
\begin{equation}
    \mathbf{V}_{neto} = \mathbf{V}_{s2} - \mathbf{V}_{s1} = 2 - (-j 0{,}4) = 2 + j 0{,}4\ \si{\volt}
\end{equation}
4) \textbf{Corriente de Lazo:}
\begin{equation}
    \mathbf{I}_{lazo} = \frac{\mathbf{V}_{neto}}{\mathbf{Z}_{total}} = \frac{2 + j 0{,}4}{25 + j 20}
\end{equation}
Multiplicamos por el conjugado del denominador ($25 - j 20$):
\begin{equation}
    \mathbf{I}_{lazo} = \frac{(2 + j 0{,}4)(25 - j 20)}{25^2 + 20^2} = \frac{50 - j 40 + j 10 + 8}{625 + 400} = \frac{58 - j 30}{1025} \approx 0{,}0566 - j 0{,}0293\ \si{\ampere}
\end{equation}
5) \textbf{Tensión en la Resistencia $\mathbf{V}_2$:}
\begin{equation}
    \mathbf{V}_2 = \mathbf{I}_{lazo} \cdot \mathbf{Z}_R = (0{,}0566 - j 0{,}0293) \times 25 = 1{,}415 - j 0{,}732\ \si{\volt} \approx 1{,}59 \angle -27{,}3^\circ\ \si{\volt}
\end{equation}
\end{ejercicio}

% ==============================================================================
% 7.2. EXAMEN PARCIAL 2
% ==============================================================================
\section{Examen Parcial 2: Dispositivos Electrónicos y Electrónica Digital}

\begin{ejercicio}[Parcial 2 --- Problema 1: Implementación de Función Booleana en Lógica NMOS y CMOS]
Diseñar a nivel físico de transistores la función lógica booleana:
\begin{equation}
    Y = \overline{(A + B) \cdot C + D}
\end{equation}
empleando:
\begin{enumerate}[label=\alph*)]
    \item Lógica NMOS con transistor de carga NMOS de deplexión.
    \item Tecnología CMOS estática complementaria completa.
\end{enumerate}
\tcbline
\textbf{Solución:}\\[4pt]
a) \textbf{Diseño con Lógica NMOS con Carga de Deplexión:}\\[2pt]
- La carga activa de Pull-Up está formada por un único transistor NMOS de deplexión con su puerta cortocircuitada a su fuente ($V_{GS} = 0$). Como $V_{th,dep} < 0$, $V_{GS} > V_{th,dep}$ y el transistor conduce permanentemente hacia $V_{DD}$.\\[2pt]
- La red de Pull-Down (PDN) descarga la salida a masa cuando la condición $(A+B)\cdot C + D$ es verdadera:
\begin{itemize}[leftmargin=*]
    \item Dos transistores NMOS de enriquecimiento $N_A$ y $N_B$ conectados en \textbf{paralelo} entre sí (operación OR: $A+B$).
    \item Este bloque paralelo se conecta en \textbf{serie} con el transistor $N_C$ (operación AND: $\cdot\, C$).
    \item Esta rama compuesta $[(N_A \parallel N_B) \text{ serie } N_C]$ se conecta en \textbf{paralelo} con el transistor $N_D$ (operación OR: $+ D$).
\end{itemize}
Total de transistores: $1$ NMOS de deplexión (carga) $+ 4$ NMOS de enriquecimiento (lógica) = \textbf{5 transistores}.\\[6pt]
b) \textbf{Diseño con Tecnología CMOS Complementaria:}\\[2pt]
- \textbf{Red de Pull-Down (PDN):} Idéntica a la red de conmutación del apartado anterior (4 transistores NMOS $N_A, N_B, N_C, N_D$). Conecta la salida $Y$ a masa cuando $F = 1$, dando $Y = 0$.\\[2pt]
- \textbf{Red de Pull-Up (PUN):} Aplica el grafo dual de De Morgan con transistores PMOS conectados a $V_{DD}$:
\begin{itemize}[leftmargin=*]
    \item El paralelo $(N_A \parallel N_B)$ pasa a ser una conexión en \textbf{serie} $(P_A \text{ serie } P_B)$.
    \item La serie con $C$ pasa a ser una conexión en \textbf{paralelo} $[(P_A \text{ serie } P_B) \parallel P_C]$.
    \item El paralelo con $D$ pasa a ser una conexión en \textbf{serie} de $P_D$ con el bloque anterior:
    \begin{equation}
        \text{Red PUN} = P_D \text{ en serie con } \left[ P_C \parallel (P_A \text{ en serie con } P_B) \right]
    \end{equation}
\end{itemize}
Total de transistores: $4$ NMOS $+ 4$ PMOS = \textbf{8 transistores}.\\[2pt]
Ventaja crítica de CMOS: consumo estático de potencia estrictamente nulo.
\end{ejercicio}

\begin{ejercicio}[Parcial 2 --- Problema 2: Discriminación de Estados ON/OFF en Red de Diodos]
En el circuito de la figura se tienen dos diodos $D_1$ y $D_2$ modelados con el segundo modelo lineal ($V_\gamma = 0{,}7\ \si{\volt}$).
Una fuente de $10\ \si{\volt}$ se conecta al ánodo de $D_1$ a través de una resistencia $R_1 = 4\ \si{\kilo\ohm}$. El cátodo de $D_1$ y el cátodo de $D_2$ confluyen en un nudo común $V_x$ conectado a masa mediante una resistencia $R_3 = 6\ \si{\kilo\ohm}$.
El ánodo de $D_2$ está alimentado por una fuente continua independiente de $3\ \si{\volt}$.
Determinar razonadamente qué diodos conducen (ON) y cuáles están bloqueados (OFF), calculando la tensión $V_x$ y las corrientes del circuito.
\tcbline
\textbf{Solución:}\\[4pt]
1) \textbf{Análisis Previo e Hipótesis Lógica:}
Si ambos diodos estuviesen apagados (OFF), por la resistencia $R_3$ no circularía corriente y la tensión sería $V_x = 0\ \si{\volt}$.
En ese caso:
\begin{align}
    V_{D1} &= 10 - 0 = 10\ \si{\volt} > 0{,}7\ \si{\volt} \quad (\text{tiende a encender}) \\
    V_{D2} &= 3 - 0 = 3\ \si{\volt} > 0{,}7\ \si{\volt} \quad (\text{tiende a encender})
\end{align}
Sin embargo, la fuente de $10\ \si{\volt}$ es significativamente mayor que la de $3\ \si{\volt}$. Al conducir $D_1$, elevará la tensión del cátodo común $V_x$, lo que podría polarizar en inversa a $D_2$.\\[4pt]
\textbf{Hipótesis 1: $D_1$ ON y $D_2$ OFF.}
Sustituimos $D_1$ por una batería en oposición de $V_\gamma = 0{,}7\ \si{\volt}$ y $D_2$ por un circuito abierto.
El circuito forma un divisor resistivo con $R_1 = 4\ \si{\kilo\ohm}$ y $R_3 = 6\ \si{\kilo\ohm}$:
\begin{equation}
    I_{D1} = \frac{10 - V_\gamma}{R_1 + R_3} = \frac{10 - 0{,}7}{4\ \si{\kilo\ohm} + 6\ \si{\kilo\ohm}} = \frac{9{,}3\ \si{\volt}}{10\ \si{\kilo\ohm}} = 0{,}93\ \si{\milli\ampere}
\end{equation}
La tensión en el nudo del cátodo común $V_x$ es:
\begin{equation}
    V_x = I_{D1} \cdot R_3 = 0{,}93\ \si{\milli\ampere} \times 6\ \si{\kilo\ohm} = 5{,}58\ \si{\volt}
\end{equation}
2) \textbf{Verificación Estricta de la Hipótesis 1:}
\begin{itemize}[leftmargin=*]
    \item Para $D_1$ (supuesto ON):
    \begin{equation}
        I_{D1} = 0{,}93\ \si{\milli\ampere} > 0 \quad \text{[COMPROBADO: conduce en directa]}
    \end{equation}
    \item Para $D_2$ (supuesto OFF):
    La tensión en su ánodo es $V_{A2} = 3\ \si{\volt}$ y en su cátodo es $V_x = 5{,}58\ \si{\volt}$.
    La caída ánodo-cátodo es:
    \begin{equation}
        V_{D2} = V_{A2} - V_x = 3 - 5{,}58 = -2{,}58\ \si{\volt} < V_\gamma = 0{,}7\ \si{\volt}
    \end{equation}
    (Comprobado: el diodo $D_2$ está fuertemente bloqueado en inversa).
\end{itemize}
3) \textbf{Demostración de que la Hipótesis Dual ($D_1$ ON y $D_2$ ON) es Físicamente Imposible:}
Si supusiéramos que ambos conducen simultáneamente:
$V_x = 3 - 0{,}7 = 2{,}3\ \si{\volt}$ (impuesto por $D_2$).
Entonces por $R_3$ circularía $I_{R3} = \frac{2{,}3}{6\ \si{\kilo\ohm}} = 0{,}383\ \si{\milli\ampere}$.
Pero por la rama 1 circularía $I_{D1} = \frac{10 - 0{,}7 - 2{,}3}{4\ \si{\kilo\ohm}} = \frac{7}{4} = 1{,}75\ \si{\milli\ampere}$.
Por LCK en el nudo $V_x$:
\begin{equation}
    I_{D2} = I_{R3} - I_{D1} = 0{,}383 - 1{,}75 = -1{,}367\ \si{\milli\ampere} < 0
\end{equation}
Corriente negativa contradice la conducción del diodo $\implies$ Hipótesis absurda.\\[4pt]
\textbf{Conclusión Final:}\\
\textbf{Estado:} $D_1$ está \textbf{ON} ($I_{D1} = 0{,}93\ \si{\milli\ampere}$); $D_2$ está \textbf{OFF} ($V_{D2} = -2{,}58\ \si{\volt}$).\\
\textbf{Tensión en el nudo:} $V_x = 5{,}58\ \si{\volt}$.
\end{ejercicio}

\begin{ejercicio}[Parcial 2 --- Problema 3: Circuito con Dos Amplificadores Operacionales Acoplados]
En el circuito de la figura se tienen dos amplificadores operacionales ideales $OA_1$ y $OA_2$ alimentados a $\pm 15\ \si{\volt}$.
\begin{itemize}[leftmargin=*]
    \item La primera etapa con $OA_1$ forma un sumador inversor donde se inyecta una corriente continua de control $I_{in} = -0{,}1\ \si{\milli\ampere}$ y la realimentación resistiva es $R_{f1} = 50\ \si{\kilo\ohm}$.
    \item La segunda etapa con $OA_2$ está configurada como un amplificador inversor con resistencia de entrada $R_1 = 10\ \si{\kilo\ohm}$ y realimentación $R_{f2} = 100\ \si{\kilo\ohm}$, excitado por una tensión de referencia $V_{ref} = 0{,}5\ \si{\volt}$.
\end{itemize}
Calcular las tensiones en los terminales de salida $V_A$ y $V_B$, así como la diferencia de potencial $V_{AB} = V_A - V_B$.
\tcbline
\textbf{Solución:}\\[4pt]
1) \textbf{Análisis de la Etapa 1 ($OA_1$):}
Por cortocircuito virtual en $OA_1$, como el terminal no inversor está a masa ($V^+ = 0$), el terminal inversor se encuentra a masa virtual:
\begin{equation}
    V^- = 0\ \si{\volt}
\end{equation}
Como no entra corriente al operacional ($I^- = 0$), toda la corriente de entrada debe circular por la resistencia de realimentación $R_{f1}$:
\begin{equation}
    I_{in} = \frac{V^- - V_A}{R_{f1}} \implies -0{,}1\ \si{\milli\ampere} = \frac{0 - V_A}{50\ \si{\kilo\ohm}} = -\frac{V_A}{50}
\end{equation}
Despejando la tensión de salida $V_A$:
\begin{equation}
    V_A = -50 \times (-0{,}1) = +5{,}0\ \si{\volt}
\end{equation}
2) \textbf{Análisis de la Etapa 2 ($OA_2$):}
La etapa 2 es un amplificador inversor estándar alimentado por $V_{ref} = 0{,}5\ \si{\volt}$:
\begin{equation}
    V_B = -\frac{R_{f2}}{R_1} V_{ref} = -\frac{100\ \si{\kilo\ohm}}{10\ \si{\kilo\ohm}} \times (0{,}5\ \si{\volt}) = -10 \times 0{,}5 = -5{,}0\ \si{\volt}
\end{equation}
3) \textbf{Diferencia de Potencial Diferencial $V_{AB}$:}
\begin{equation}
    V_{AB} = V_A - V_B = 5{,}0\ \si{\volt} - (-5{,}0\ \si{\volt}) = 5{,}0 + 5{,}0 = 10{,}0\ \si{\volt}
\end{equation}
Ambos operacionales operan holgadamente dentro de su rango lineal ($|V_A|, |V_B| = 5\ \si{\volt} < 14\ \si{\volt}$).
\end{ejercicio}


\end{document}
